\documentclass[aps,preprint]{revtex4}%
\usepackage{amsfonts}
\usepackage{amsmath}
\usepackage{amssymb}
\usepackage{graphicx}%
\setcounter{MaxMatrixCols}{30}
%TCIDATA{OutputFilter=latex2.dll}
%TCIDATA{Version=4.10.0.2347}
%TCIDATA{CSTFile=revtex4.cst}
%TCIDATA{Created=Sunday, December 08, 2002 14:36:59}
%TCIDATA{LastRevised=Saturday, January 25, 2003 11:09:34}
%TCIDATA{<META NAME="GraphicsSave" CONTENT="32">}
%TCIDATA{<META NAME="PrintViewPercent" CONTENT="100">}
%TCIDATA{<META NAME="DocumentShell" CONTENT="Articles\SW\REVTeX 4">}
%TCIDATA{Language=American English}
\newtheorem{theorem}{Theorem}
\newtheorem{acknowledgement}[theorem]{Acknowledgement}
\newtheorem{algorithm}[theorem]{Algorithm}
\newtheorem{axiom}[theorem]{Axiom}
\newtheorem{claim}[theorem]{Claim}
\newtheorem{conclusion}[theorem]{Conclusion}
\newtheorem{condition}[theorem]{Condition}
\newtheorem{conjecture}[theorem]{Conjecture}
\newtheorem{corollary}[theorem]{Corollary}
\newtheorem{criterion}[theorem]{Criterion}
\newtheorem{definition}[theorem]{Definition}
\newtheorem{example}[theorem]{Example}
\newtheorem{exercise}[theorem]{Exercise}
\newtheorem{lemma}[theorem]{Lemma}
\newtheorem{notation}[theorem]{Notation}
\newtheorem{problem}[theorem]{Problem}
\newtheorem{proposition}[theorem]{Proposition}
\newtheorem{remark}[theorem]{Remark}
\newtheorem{solution}[theorem]{Solution}
\newtheorem{summary}[theorem]{Summary}
\newenvironment{proof}[1][Proof]{\noindent\textbf{#1.} }{\ \rule{0.5em}{0.5em}}
\begin{document}
\preprint{DMFT}
\title[DMFT]{Density Matrix Functional Theory}
\author{Brian Weiner}
\affiliation{Pennsylvania State University}
\author{S. B. Trickey}
\affiliation{University of Florida}
\date{01/01/2003}
\startpage{1}
\maketitle
\tableofcontents


\section{ Introduction}

Density Functional Theory (DFT) \cite{DFTRev} has proved to be a powerful and
valuable computational tool in the study of molecular and extended quantum
mechanical systems. One of the most popular implementations of the density
functional approach is the Kohn-Sham (KS) method \cite{KS}. The basic
variables in the KS method are molecular orbitals, either restricted or
unrestricted depending on the nature of the spin symmetry of the system.
Correlation effects are included in the Exchange-Correlation (XC) functional
that together with the kinetic energy functional and external potential
determine the optimal orbitals. In order to extend this approach and allow
part of the correlation effects to be determined variationally, i.e. be
described by some of the active variables of the problem, one introduces
functionals whose values also depend on the occupation numbers of the
orbitals. In this treatment one is in effect considering functionals of the
First Order Reduced Density Operator (FORDO), and hence it is called Density
Matrix Functional Theory (DMFT) \cite{Staroverov02}.\cite{Yasuda02}%
\cite{Cioslowski02}\cite{Manby01}\cite{Yasuda01}\cite{Weitao}\cite{Holas99}%
\cite{Johnsen96}\cite{Morrison89}\cite{Wunsche80} In this method the kinetic
energy functional is now unambiguously defined and known in terms of natural
orbitals, occupation numbers and the matrix elements of the kinetic operator,
which is not the case in DFT, but the XC\ functional is still unknown. The
strategies normally used to produce functionals for computational purposes in
DFT, i.e. from heuristic principles based on a few necessary conditions or by
some form of semi empirical fitting procedure, are also the one most favored
in DMFT. Although many successful functionals have been developed, this way of
producing functionals has some major drawbacks - in particular it does not
readily lend itself to a systematic approximation theory leading to exact
functionals, nor does it develop them from basic quantum mechanical principles
in a transparent manner. This latter deficit makes it difficult to analyze and
correct inadequacies in an effective way. A major benefit of our approach is
that approximate functionals described by our constructive approach display
the ingredients that need to be included within a full rigorous quantum
mechanical context.

In DFT the Levy-Lieb \cite{LevyLieb} recipe provides a formal construction for
the total energy density functional, however, in order to numerically
implement this procedure one must carry out a constrained minimization over
\emph{pure }N-particle states keeping the density constant. An analogous
formal definition can be used in the context of density matrix functional
theory to define an energy functional of the density matrix, but again one
must implement a constrained minimization this time keeping the density matrix
constant. In this article we describe a systematic method for executing such a
constrained optimization bases on a transparent first principles framework
that allows one to outline concrete procedures for constructing density matrix
functionals to any desired degree of accuracy. Invariably in quantum theory
there is a gap between the formal development of a particular approach to the
solve the Schr\"{o}dinger equation and the practical implementation of that
approach, this article presents material that conforms to this pattern: we
develop in a formal manner a constructive procedure within a rigorous quantum
mechanical context the necessary and sufficient conditions that parameters
must obey in order to determine approximate DMFs and FORDOs, however in the
examples we present of practical implementations we only manage to formulate
necessary conditions that constrain the parameters, though still within a
rigorous constructive quantum mechanical context that could be further
extended to produce necessary and \emph{sufficient }conditions.

The approach we utilize is based on work presented in \cite{weiner 2002} that
describes how one can generate all states (pure and ensemble) by the action of
abelian semigroups \cite{semigroups}, of positive diagonal superoperator
transformations induced by diagonal operators, on special but easy to
construct \emph{cyclic} \cite{weiner2002}\emph{\ }reference states. If the
reference state is also a pure state then the subset of \emph{all pure states}
can be generated by the action of these abelian semigroups. In
\cite{weiner2002} it was shown that it is \emph{not sufficient }to consider
abelian groups of transformations as they only produce restricted classes of
states. The defining representations of these semigroups are built upon the
properties of generalized bases, in particular the set of generalized
eigenvectors of the space-spin position operator, and the theory of dual
vector spaces of operators, which are discussed in detail in \cite{weiner2002}.

We focus on abelian semigroups that are generated by operators that are
diagonal in the space-spin coordinate representation, which allow us to
express the expectation of the Hamiltonian with respect to (abbreviated wrt in
the proceeding) \emph{any} state in terms of the expectation value of an
effective Hamiltonian wrt a single \emph{fixed} reference independent particle
state. The effective Hamiltonian can be expanded in terms of multiparticle
local \cite{local} potentials and gauge transformations of the momentum
operator in a manner similar to Jastrow perturbation theory
\cite{Blazot&Ripka-Jastrow} which permits one to construct well defined
approximation procedures for DMFs and FORDOs. It should be noted that these
\emph{non unitary commutative }transformations preserve the diagonal form of
the potential terms in the Hamiltonian, producing less complicated expressions
for expectation values than would be obtained by using conventional unitary
transformations to generate new states, but at the cost of introducing higher
order multiparticle local potential operators. The use of diagonal operators
allows one to design approximations to the effective Hamiltonian that maintain
their form as one varies over all possible states consistent with this level
of approximation in contrast for example to the Neglect of Differential
Overlap (NDO) utilized in semi empirical molecular approximation theory
\cite{NDO}. On a more philosophical note the diagonal transformations
introduce a more classical flavor to the formulation without loosing
\emph{any} quantum mechanical rigor.

The map from N-particle states, represented as density operators, to FORDOs is
a bounded linear map showing that states that have a common FORDO differ from
each other by an operator that belongs to the kernel of this map, which is a
subspace of N-particle operator space. Hence all states with the same
FORDO\ have an \emph{identical }operator component\emph{\ }that lies in a
complimentary subspace to the kernel and each element in such a complimentary
space can be uniquely associated with a class of states that all have the same
FORDO. We show in this article how to construct sub semigroups of the abelian
semigroups described in \cite{weiner2002} that leave the kernel and its
complimentary subspaces invariant. In particular we describe maps that leave a
particular complimentary component invariant and thus transform the class of
states with the same FORDO to itself. This semigroup of maps allows one to
vary over the class of states with a common FORDO\ and thus perform the
constrained minimization producing the density matrix functional.

In order to produce concrete realizations of these semigroups we treat
operators as super vectors and the semigroup of transformations as super
operators acting on them. The formalism of dual vector spaces allows us to
express these objects in a generalized "bra-ket" notation and obtain
variational expressions involving parameters determining these semigroup
elements. One can consider nested sequences of sub semigroups, thus nested
sets of parameters corresponding to a certain level of truncation of the
effective Hamiltonian, to vary over and thus produce sequences of
approximations converging to the exact DMF.

To summarize, the approach we take in this article to construct functionals
from first principles to a prescribed degree of exactness is based on:

\begin{enumerate}
\item using the coordinate representation of states and observables

\item using the properties of the linear contraction map from N-particle
operator space to 1-particle operator space to decompose N-particle operator
space into a direct sum of the kernel of this map and its compliment

\item using this decomposition to classify pure and ensemble states according
to the FORDO they are associated with under this contraction map

\item using an abelian semigroup of positive super operators that generates
all states from a reference state

\item constructing a sub semigroup of this semigroup that generates all states
with a specified FORDO from a reference state (that also has the same FORDO)

\item expressing the expectation value of the Hamiltonian wrt any state in
terms of the expectation value of an effective Hamiltonian containing only
multiparticle potentials, kinetic energy and momentum operators wrt a fixed
independent particle reference state

\item using sequences of nested sub semigroups of the semigroup described in
point 5 above, that correspond to a certain level of truncation of the
effective Hamiltonian, to produce functionals of increasing exactness by
partial optimization of the energy functional over the parameters
characterizing these sub semigroups.
\end{enumerate}

We also describe semigroups that generate all complimentary components, thus
produce variation over all values of the arguments of these DMFs i.e over
different FORDOs. In a similar manner to the construction of the approximate
functionals one can use nested sequences of sub semigroups of super operators,
thus nested sets of parameters, to determine FORDOs optimized to prescribed
degrees of exactness for each of these approximate density matrix functionals
and thus obtain approximate solutions to the Schr\"{o}dinger equation.

In order to perform the optimizations to find DMFs and the optimized density
matrices in an efficient manner it is desirable to base the variation on a non
redundant parametrization of both the DMFs and their density matrix arguments.
This ensures that there is a one to one correspondence between the parameters
and the FORDOs and the values of the functionals. We describe how this can be
realized using specific representations of the abelian semigroups we introduce.

This article has the following structure: in section \ref{Cmap} we discuss the
contraction map from N-particle operator space to 1-particle operator space,
characterize its kernel and discuss the classification of states by their
FORDOs. In section \ref{Functionals} we describe the semigroup that maps
between states that have the same FORDO and thus the construction of exact
functionals. In section \ref{Domain} we discuss the semigroup that maps
between states with different FORDOs and thus the domains of the functionals,
and variation over these domains, while in section \ref{Standard} we relate
the constructive functionals and the ensuing variational equations based on
these functionals to the standard ones used in DMFT. In section \ref{approx}
we give an outline of how to construct approximate functionals and approximate
domains of these functionals by considering truncations of the effective
Hamiltonian and in section \ref{example} we present an example of such
constructions and finally we summarize the main points of the article in
section \ref{summary}.

\section{The Contraction Map\label{Cmap}}

This section deals with the mathematical properties of operators that
represent physical states of a system and the contraction map from N-particle
states to first order reduced one particle states. We use the properties of
this map to describe direct sum decompositions of the space of trace class
operators and hence classify states by the FORDO\ that they produce.

\subsection{States}

The states of an $N$-particle quantum mechanical system correspond to
positive, normalized trace class operators and form a convex subset,
$\mathcal{S}_{N}$, of the Banach space \cite{banach}, $\mathcal{B}_{1}\left(
\mathcal{H}^{N}\right)  $, of $N$-particle trace class operators, (in Appendix
\ref{ideals} we review the definition of the space of trace class operators
and some of its properties), defined as%
\begin{equation}
\mathcal{S}_{N}=\left\{  D^{N}|D^{N}\in\mathcal{B}_{1}\left(  \mathcal{H}%
^{N}\right)  ,D^{N}\geq0,TrD^{N}=1\right\}
\end{equation}
Pure states are idempotent operators i.e.
\begin{equation}
\left(  D^{N}\right)  ^{2}=D^{N}%
\end{equation}
and can be associated with, in a 1-1 fashion, with a ray $\left\{
\alpha\left\vert \Psi\right\rangle |\alpha\in\mathbb{C}\right\}  $ generated
by a vector $\left\vert \Psi\right\rangle $ in the N-particle Hilbert space
$\mathcal{H}^{N}$ by the relationship
\begin{equation}
D^{N}=\frac{\left\vert \Psi\right\rangle \left\langle \Psi\right\vert
}{\left\langle \Psi|\Psi\right\rangle }%
\end{equation}
States that are not idempotent are called mixed or ensemble states and
correspond to statistical mixtures of pure states.

In the analysis presented in this article we need to construct a procedure
that generates all states from a reference state. It is, however, difficult to
implement and construct maps that transform the set of \textit{normalized
}states $\mathcal{S}_{N}$ to itself, while we can explicitly characterize
linear maps transforming the cone of positive trace class operators,
$\mathcal{B}_{1+}\left(  \mathcal{H}^{1}\right)  $, of \textit{unnormalized
}states generated by $\mathcal{S}_{N}$, (see Appendix \ref{states}),
\begin{equation}
\mathcal{B}_{1+}\left(  \mathcal{H}^{1}\right)  =\left\{  \alpha D^{N}%
|D^{N}\in\mathcal{S}_{N};\alpha\geq0\right\}
\end{equation}
to itself. Thus we focus on unnormalized\textit{\ }states $\mathcal{D}^{N}%
\in\mathcal{B}_{1+}\left(  \mathcal{H}^{N}\right)  $ which produce normalized
states by the simple recipe
\begin{equation}
D^{N}=\frac{\mathcal{D}^{N}}{Tr\left\{  \mathcal{D}^{N}\right\}  }.
\end{equation}


\subsection{Definition of the Contraction Map}

In order to fully exploit the linearity of the contraction map it is
convenient to express its action on the full linear space $\mathcal{B}%
_{1}\left(  \mathcal{H}^{N}\right)  $, rather than just $\mathcal{B}%
_{1+}\left(  \mathcal{H}^{1}\right)  $ or $\mathcal{S}_{N}$, i.e. by
\begin{equation}
L_{N}^{1}:\mathcal{B}_{1}\left(  \mathcal{H}^{N}\right)  \rightarrow
\mathcal{B}_{1}\left(  \mathcal{H}^{1}\right)
\end{equation}
a map from trace class operators acting in $N$-particle space to trace class
operators acting in $1$-particle space defined in the space-spin coordinate
representation by
\begin{equation}
L_{N}^{1}\left(  X^{N}\right)  \left(  \mathsf{z}^{N},\mathsf{z}^{\prime
N}\right)  =X_{1}^{N}\left(  \mathsf{z}^{N},\mathsf{z}^{\prime N}\right)
=N\int X^{N}\left(  \mathsf{z}_{1}\mathsf{z}^{N-1},\mathsf{z}_{1}^{\prime
}\mathsf{z}^{N-1}\right)  d\mathsf{z}^{N-1} \label{contract}%
\end{equation}
where $\mathsf{z}^{N}=\mathsf{z}_{1}\cdots\mathsf{z}_{N},\mathsf{z}^{N-1}=$
$\mathsf{z}_{2}\cdots\mathsf{z}_{N}$, $d\mathsf{z}^{N-1}=d\mathsf{z}_{2}\cdots
d\mathsf{z}_{N}$ and $\left\{  \mathsf{z}_{j}\right\}  $ range over the
spectrum of the one particle space-spin coordinate operator. Although the
definition eq. $\left(  \ref{contract}\right)  $ utilizes the coordinate
representation the resulting contraction map is representation independent and
any basis or generalized basis of $\mathcal{H}^{N}$ \cite{Berezin1991} can be
used to define it.

\subsection{Direct Sum Decomposition of $\mathcal{B}_{1}\left(  \mathcal{H}%
^{N}\right)  \label{DirectSum}$}

The map $L_{N}^{1}$ is linear and bounded and thus its kernel is a closed
linear subspace, $\ker L_{N}^{1}$ of $\mathcal{B}_{1}\left(  \mathcal{H}%
^{N}\right)  $, defined by
\begin{equation}
\ker L_{N}^{1}=\left\{  X^{N}\left\vert \int X^{N}\left(  \mathsf{z}%
_{1}\mathsf{z}^{N-1},\mathsf{z}_{1}^{\prime}\mathsf{z}^{N-1}\right)
d\mathsf{z}^{N-1}\overset{a.e.}{=}0\quad\forall\mathsf{z}_{1},\mathsf{z}%
_{1}^{\prime}\right.  \right\}
\end{equation}
(See \cite{ae} for the definition of "$\overset{a.e.}{=}$" ). The space
$\mathcal{B}_{1}\left(  \mathcal{H}^{N}\right)  $ can be expressed as direct
sum of subspaces as
\begin{equation}
\mathcal{B}_{1}\left(  \mathcal{H}^{N}\right)  =\ker L_{N}^{1}\oplus\left[
\ker L_{N}^{1}\right]  ^{c} \label{dc}%
\end{equation}
where $\left[  \ker L_{N}^{1}\right]  ^{c}$ is the canonical complimentary
subspace of $\ker L_{N}^{1}$ in $\mathcal{B}_{1}\left(  \mathcal{H}%
^{N}\right)  $ (see Appendix \ref{ideals}). Hence any trace class operator,
$X$, can be decomposed uniquely as
\begin{equation}
X=\mathfrak{x}_{\ker}+\mathfrak{x}_{1} \label{orthog}%
\end{equation}
where $\mathfrak{x}_{\ker}\in\ker L_{N}^{1}$ and $\mathfrak{x}_{1}\in\left[
\ker L_{N}^{1}\right]  ^{c}$. It should, however, be noted that although the
map $L_{N}^{1}$ is a linear isomorphism from $\left[  \ker L_{N}^{1}\right]
^{c}$ to $\mathcal{B}_{1}\left(  \mathcal{H}^{1}\right)  \cite{stanline}$, it
\emph{does not }map the multiplicative structure of $\mathcal{B}_{1}\left(
\mathcal{H}^{N}\right)  $ i.e. $L_{N}^{1}\left(  AB\right)  \neq L_{N}%
^{1}\left(  A\right)  L_{N}^{1}\left(  B\right)  $.

As discussed in Appendix \ref{ideals} the direct sum decomposition of
$\mathcal{B}_{1}\left(  \mathcal{H}^{N}\right)  $, given in eq. $\left(
\ref{dc}\right)  $, with respect to $\ker L_{N}^{1}$ is not unique and we also
have expansions labelled by linear 1-1 maps (defined below)
\begin{equation}
\tilde{\Gamma}:\left[  \ker L_{N}^{1}\right]  ^{c}\rightarrow\mathcal{B}%
_{1}\left(  \mathcal{H}^{N}\right)
\end{equation}
given by
\begin{equation}
\mathcal{B}_{1}\left(  \mathcal{H}^{N}\right)  =\ker L_{N}^{1}\oplus
\mathcal{V}_{\tilde{\Gamma}} \label{nonorthog}%
\end{equation}
where $\mathcal{V}_{\tilde{\Gamma}}=\tilde{\Gamma}\left(  \left[  \ker
L_{N}^{1}\right]  ^{c}\right)  $. The subspaces $\mathcal{V}_{\tilde{\Gamma}}$
are isomorphic to $\left[  \ker L_{N}^{1}\right]  ^{c}$ and determined by the
maps $\tilde{\Gamma}$, which are defined in terms of \emph{arbitrary} linear
maps $\Gamma$ from $\left[  \ker L_{N}^{1}\right]  ^{c}$ to $\ker L_{N}^{1}$
by%
\begin{equation}
\tilde{\Gamma}\left(  \mathfrak{x}_{1}\right)  =\Gamma\left(  \mathfrak{x}%
_{1}\right)  +\mathfrak{x}_{1}.
\end{equation}
These non canonical decompositions lead to the operator expansions
\begin{equation}
X=\left[  \mathfrak{x}_{\ker}-\Gamma\left(  \mathfrak{x}_{1}\right)  \right]
+\left[  \mathfrak{x}_{1}+\Gamma\left(  \mathfrak{x}_{1}\right)  \right]
=\mathfrak{x}_{\Gamma}+\tilde{\Gamma}\left(  \mathfrak{x}_{1}\right)
\end{equation}
where $\tilde{\Gamma}\left(  \mathfrak{x}_{1}\right)  \in\mathcal{B}%
_{1}\left(  \mathcal{H}^{N}\right)  $ and $\mathfrak{x}_{\Gamma}%
=\mathfrak{x}_{\ker}-\Gamma\left(  \mathfrak{x}_{1}\right)  \in\ker L_{N}^{1}%
$. When the map $\Gamma$ is the zero map the decomposition eq. $\left(
\ref{nonorthog}\right)  $ becomes the decomposition eq. $\left(
\ref{orthog}\right)  $ and $\tilde{\Gamma}$ becomes the identity map $\left[
\ker L_{N}^{1}\right]  ^{c}\rightarrow\left[  \ker L_{N}^{1}\right]  ^{c}$. As
the subspaces $\left[  \ker L_{N}^{1}\right]  ^{c}$ and $\left\{
\mathcal{V}_{\tilde{\Gamma}}\right\}  $ are isomorphic the map $L_{N}^{1}$ is
also a linear isomorphism from the subspaces $\left\{  \mathcal{V}%
_{\tilde{\Gamma}}\right\}  $ to $\mathcal{B}_{1}\left(  \mathcal{H}%
^{1}\right)  $.

\subsection{Equivalence Class Decomposition of $\mathcal{B}_{1}\left(
\mathcal{H}^{N}\right)  $}

One can express $\mathcal{B}_{1}\left(  \mathcal{H}^{N}\right)  $ in terms of
equivalence classes of operators which are defined as%
\begin{equation}
\left[  X_{1}^{N}\right]  _{N}=\left\{  X|X\in\mathcal{B}_{1}\left(
\mathcal{H}^{N}\right)  ;L_{N}^{1}\left(  X\right)  =X_{1}^{N}\right\}  .
\end{equation}
All operators in the same equivalence class have the same component
$\mathfrak{x}_{1}$ belonging to $\left[  \ker L_{N}^{1}\right]  ^{c}$ in the
canonical decomposition of $\mathcal{B}_{1}\left(  \mathcal{H}^{N}\right)  $
wrt $\ker L_{N}^{1}$ and can be expressed as
\begin{equation}
X=\mathfrak{y}+\mathfrak{x}_{1}%
\end{equation}
where $\mathfrak{y}$ ranges over the elements of $\ker L_{N}^{1}$. One can
also use the non canonical decomposition of $\mathcal{B}_{1}\left(
\mathcal{H}^{N}\right)  $, which are much easier to implement and will be
utilized extensively in our construction of the functionals, to obtain the
alternative expressions
\begin{equation}
X=\mathfrak{\tilde{y}}+\tilde{\Gamma}\left(  \mathfrak{x}_{1}\right)
;\quad\tilde{\Gamma}\left(  \mathfrak{x}_{1}\right)  \in\mathcal{B}_{1}\left(
\mathcal{H}^{N}\right)  ;\,\mathfrak{\tilde{y}}\in\ker L_{N}^{1}%
\end{equation}
where $\mathfrak{\tilde{y}}=\mathfrak{y}-\Gamma\left(  \mathfrak{x}%
_{1}\right)  $. Hence the equivalence classes can be expressed either as
\begin{equation}
\left[  X_{1}^{N}\right]  _{N}=\left\{  X=\mathfrak{y}+\mathfrak{x}%
_{1}|\mathfrak{y}\in\ker L_{N}^{1};\,\mathfrak{x}_{1}\in\left[  \ker L_{N}%
^{1}\right]  ^{c}\right\}
\end{equation}
or alternatively as
\begin{equation}
\left[  X_{1}^{N}\right]  _{N}=\left\{  X=\mathfrak{\tilde{y}}+\tilde{\Gamma
}\left(  \mathfrak{x}_{1}\right)  |\mathfrak{\tilde{y}}\in\ker L_{N}%
^{1};\,\mathfrak{x}_{1}\in\left[  \ker L_{N}^{1}\right]  ^{c}\right\}
\end{equation}


The decomposition of $\mathcal{B}_{1}\left(  \mathcal{H}^{N}\right)  $ wrt
$\ker L_{N}^{1}$ can be formulated in a more pictorial fashion in Figure
\ref{Fig1}, which is useful in the analysis of the effect of the various
semigroups, as the vector bundle $\mathfrak{B}_{L_{N}^{1}}$ \begin{figure}[h]
\includegraphics{figa.eps}\caption{Vector Bundle showing fibers $F_{\tilde
{\Gamma}\left(  \mathfrak{x}_{1}\right)  }$ above a general representation of
the base $\mathcal{V}_{\Gamma}$.}%
\label{Fig1}%
\end{figure}whose fibers are $\ker L_{N}^{1}$ and whose base can be
represented by any of the subspaces $\mathcal{B}_{1}\left(  \mathcal{H}%
^{N}\right)  \diagup\ker L_{N}^{1}$ i.e. any $\mathcal{V}_{\Gamma}$. The
vector bundle $\mathfrak{B}_{L_{N}^{1}}$ is called "trivial" as%
\begin{equation}
\mathfrak{B}_{L_{N}^{1}}=\left[  \mathcal{B}_{1}\left(  \mathcal{H}%
^{N}\right)  \diagup\ker L_{N}^{1}\right]  \times\ker L_{N}^{1}\equiv
\mathcal{V}_{\Gamma}\times\ker L_{N}^{1}\ \label{VB}%
\end{equation}
i.e. all the fibers are isomorphic vector spaces. In order to produce concrete
representations of semigroups we will need to distinguish between these
isomorphic spaces and to this end we denote the fiber above the point
$\tilde{\Gamma}\left(  \mathfrak{x}_{1}\right)  \in\mathcal{V}_{\Gamma}$ as
$F_{\tilde{\Gamma}\left(  \mathfrak{x}_{1}\right)  }$.

If one uses the canonical representation, $\left[  \ker L_{N}^{1}\right]
^{c}$, of the base then the other representations of the base, $\left\{
\mathcal{V}_{\Gamma}\right\}  $, can be viewed as linear cross-sections
\cite{cross} $\left\{  \tilde{\Gamma}\left(  \mathfrak{x}_{1}\right)
|\mathfrak{x}_{1}\in\left[  \ker L_{N}^{1}\right]  ^{c}\right\}  $ of
$\mathfrak{B}_{L_{N}^{1}}$ as displayed in Figure \ref{Fig2} \begin{figure}[h]
\includegraphics{figb.eps}\caption{Vector Bundle showing the linear
cross-sections that represent the base $\mathcal{V}_{\Gamma}.$}%
\label{Fig2}%
\end{figure}

\subsection{The Action of the Contraction Map on States}

The contraction map restricted to $N$-particle density operators produces
First Order Reduced Density Operators (FORDOs). This set is the image of
$\mathcal{S}_{N}$ under $L_{N}^{1}$ and forms a convex subset, $\mathcal{S}%
_{N}^{1}=L_{N}^{1}\left(  \mathcal{S}_{N}\right)  $ of the cone of positive
operators, $\mathcal{B}_{1+}\left(  \mathcal{H}^{1}\right)  $, acting in
$\mathcal{H}^{1}$. The image $L_{N}^{1}\left(  \mathcal{B}_{1+}\left(
\mathcal{H}^{N}\right)  \right)  $, of the cone $\mathcal{B}_{1+}\left(
\mathcal{H}^{N}\right)  $ under the contraction map is the cone $\mathcal{C}%
_{N}^{1}\subset\mathcal{B}_{1+}\left(  \mathcal{H}^{1}\right)  $.

Properly normalized FORDO's are given in terms of the image of the cone
$\mathcal{B}_{1+}\left(  \mathcal{H}^{N}\right)  $ as
\begin{equation}
L_{N}^{1}\left(  D^{N}\right)  =\frac{L_{N}^{1}\left(  \alpha D^{N}\right)
}{Tr\left\{  \alpha D^{N}\right\}  }=\frac{L_{N}^{1}\left(  \mathcal{D}%
^{N}\right)  }{Tr\left\{  \mathcal{D}^{N}\right\}  }\equiv\frac{\mathcal{D}%
_{1}^{N}}{Tr\left\{  \mathcal{D}^{N}\right\}  }=\frac{\mathcal{D}_{1}^{N}%
}{Tr\left\{  \mathcal{D}_{1}^{N}\right\}  }=D_{1}^{N}%
\end{equation}
where $D^{N}\in\mathcal{S}_{N},\mathcal{D}^{N}\in\mathcal{B}_{1+}\left(
\mathcal{H}^{N}\right)  ,D_{1}^{N}\in\mathcal{S}_{N}^{1}$ and $\mathcal{D}%
_{1}^{N}\in\mathcal{C}_{N}^{1}$. In particular one can observe that
$\mathcal{C}_{N}^{1}$ determines and is determined by $\mathcal{S}_{N}^{1}$.

\subsection{Equivalence Classes of States}

In this subsection we describe how an equivalence class of states, that
produce a common FORDO, can be associated with a representative state, which
can be then be utilized as the independent variable of the density matrix
functionals that we describe subsequently.

Operators, $\mathcal{D}^{N}$, belonging to the cone $\mathcal{B}_{1+}\left(
\mathcal{H}^{N}\right)  $ can be decomposed canonically wrt $\ker L_{N}^{1}$
in a unique fashion as%
\begin{equation}
\mathcal{D}^{N}=\mathfrak{d}+\mathfrak{d}_{1}%
\end{equation}
and in terms of a non canonical compliment wrt $\ker L_{N}^{1}$ as%
\begin{equation}
\mathcal{D}^{N}=\mathfrak{d}_{\Gamma}+\tilde{\Gamma}\left(  \mathfrak{d}%
_{1}\right)  \label{decomp}%
\end{equation}
where $\mathfrak{d}$ and $\mathfrak{d}_{\Gamma}$ $\in\ker L_{N}^{1}$,
$\mathfrak{d}_{1}\in\left[  \ker L_{N}^{1}\right]  ^{c}$ and $\tilde{\Gamma
}\left(  \mathfrak{d}_{1}\right)  \in\mathcal{B}_{1}\left(  \mathcal{H}%
^{N}\right)  $ (see eqs. $\left(  \ref{orthog}\right)  $ and $\left(
\ref{nonorthog}\right)  $). The equivalence class of operators, $\left[
\mathcal{D}^{1}\right]  _{N}$, \emph{properly} contains the convex subset,
$\left\{  \mathcal{D}^{1}\right\}  _{N}$, of operators contained in the cone
$\mathcal{B}_{1+}\left(  \mathcal{H}^{N}\right)  $ of unnormalized states,
that contract to the same unnormalized FORDO $\mathcal{D}^{1}$. The fact that
there are non positive elements of $\left[  \mathcal{D}^{1}\right]  _{N}$ that
also contract to the unnormalized FORDO $\mathcal{D}^{1}$, is the origin of
Representability Problems \cite{coleman} in many body quantum mechanics.

All operators belonging to $\left\{  \mathcal{D}^{1}\right\}  _{N}$ have the
same fixed component $\mathfrak{d}_{1}$ belonging to $\left[  \ker L_{N}%
^{1}\right]  ^{c}$ and thus can be expressed as the set%
\begin{equation}
\left\{  \mathcal{D}^{1}\right\}  _{N}=\left\{  \mathcal{D}^{N}|;\mathcal{D}%
^{N}=\mathfrak{d}+\mathfrak{d}_{1}\right\}
\end{equation}
where $\mathfrak{d}$ ranges over $\ker L_{N}^{1}$ \emph{such that
}$\mathfrak{d}+\mathfrak{d}_{1}\geq0$ for a fixed $\mathfrak{d}_{1}$. We
denote the set of the $\left[  \ker L_{N}^{1}\right]  ^{c}$-components of
unnormalized states as $\mathfrak{C}_{1}$ i.e. the set of elements of $\left[
\ker L_{N}^{1}\right]  ^{c}$ that can be extended to elements of
$\mathcal{B}_{1+}\left(  \mathcal{H}^{N}\right)  $. This leads to the $1-1$
correspondence $\mathcal{D}^{1}\in\mathcal{B}_{1+}\left(  \mathcal{H}%
^{1}\right)  \leftrightarrow\mathfrak{d}_{1}\in\mathfrak{C}_{1}\subset
\mathcal{B}_{1}\left(  \mathcal{H}^{N}\right)  $. Hence we can denote each
class $\left\{  \mathcal{D}^{1}\right\}  _{N}$ equivalently as $\left\{
\mathfrak{d}_{1}\right\}  $ and each class $\left[  \mathcal{D}^{1}\right]
_{N}$ equivalently as $\left[  \mathfrak{d}_{1}\right]  $. One should note
that in general $\mathfrak{d}_{1}$ \emph{is not} an unnormalized state i.e.
\emph{does not belong to }$\mathcal{B}_{1+}\left(  \mathcal{H}^{N}\right)  $.
However we can define a set of $\Gamma^{\prime}s$ by
\begin{equation}
\mathfrak{P}=\left\{  \Gamma|\tilde{\Gamma}\left(  \mathfrak{d}_{1}\right)
\in\mathcal{B}_{1+}\left(  \mathcal{H}^{N}\right)  \,\forall\,\mathfrak{d}%
_{1}\in\mathfrak{C}_{1}\right\}
\end{equation}
and always choose a $\Gamma\in\mathfrak{P}$ (see Appendix \ref{noncanstate}),
and represent the base of the vector bundle $\mathfrak{B}_{L_{N}^{1}}$ by a
subspace $\mathcal{V}_{\Gamma}$ that contains the unnormalized states
\begin{equation}
\mathcal{V}_{\Gamma+}=\mathcal{V}_{\Gamma}\cap\mathcal{B}_{1+}\left(
\mathcal{H}^{N}\right)  \equiv\left\{  \tilde{\Gamma}\left(  \mathfrak{d}%
_{1}\right)  |\mathfrak{d}_{1}\in\ \mathfrak{C}_{1};\tilde{\Gamma}\left(
\mathfrak{d}_{1}\right)  \in\mathcal{B}_{1+}\left(  \mathcal{H}^{N}\right)
\right\}
\end{equation}
and the normalized states
\begin{equation}
\mathcal{N}_{\Gamma}=\left\{  \tilde{\Gamma}\left(  \mathfrak{d}_{1}\right)
|\mathfrak{d}_{1}\in\ \mathfrak{C}_{1};\tilde{\Gamma}\left(  \mathfrak{d}%
_{1}\right)  \in\mathcal{B}_{1+}\left(  \mathcal{H}^{N}\right)  ;Tr\left\{
\mathfrak{d}_{1}\right\}  =1\right\}  \label{NT}%
\end{equation}
and thus obtain the decomposition eq. $\left(  \ref{decomp}\right)  $. This is
a crucial point in our construction on two levels as

\begin{enumerate}
\item we can utilize an abelian semigroup of maps $\mathcal{B}_{1+}\left(
\mathcal{H}^{N}\right)  \rightarrow\mathcal{B}_{1+}\left(  \mathcal{H}%
^{N}\right)  $ to generate representatives of base elements of the fibers
$\left[  \mathcal{D}^{1}\right]  _{N}$ of $\mathfrak{B}_{L_{N}^{1}}$ from a
"global" reference state and

\item we can represent the base element, $\mathfrak{d}_{1}$, by a non unique
"local" reference \emph{state} $\tilde{\Gamma}\left(  \mathfrak{d}_{1}\right)
$ which is an object that has physical interpretable properties. (The use of
the word "local" implies that a particular reference state is associated with
just one fiber).
\end{enumerate}

In contrast to DFT the local reference states \emph{cannot all be chosen }to
be\emph{\ }Independent Particle States (IPS),\emph{\ }however the global
reference state \emph{can }always be chosen to be an IPS.

Clearly
\begin{equation}
D^{1}=\frac{L_{N}^{1}\left(  \mathfrak{d}_{1}\right)  }{Tr\left\{
\mathfrak{d}_{1}\right\}  }=\frac{L_{N}^{1}\left(  \tilde{\Gamma}\left(
\mathfrak{d}_{1}\right)  \right)  }{Tr\left\{  \tilde{\Gamma}\left(
\mathfrak{d}_{1}\right)  \right\}  } \label{D1}%
\end{equation}
and
\begin{equation}
Tr\left\{  \mathcal{D}^{N}\right\}  =Tr\left\{  \mathfrak{d}+\mathfrak{d}%
_{1}\right\}  =Tr\left\{  \mathfrak{d}_{\Gamma}+\tilde{\Gamma}\left(
\mathfrak{d}_{1}\right)  \right\}  =Tr\left\{  \tilde{\Gamma}\left(
\mathfrak{d}_{1}\right)  \right\}  =Tr\left\{  \mathfrak{d}_{1}\right\}
\end{equation}
so all operators belonging to $\left\{  \mathcal{D}^{1}\right\}  _{N}$ have
the same trace. Eq. $\left(  \ref{D1}\right)  $ also shows that although for
$\alpha\neq1$
\begin{equation}
\left\{  \mathcal{D}^{1}\right\}  _{N}\neq\left\{  \alpha\mathcal{D}%
^{1}\right\}  _{N}%
\end{equation}
the equivalence classes $\left\{  \alpha\mathcal{D}^{1}\right\}  _{N}$ for any
$\alpha>0$ lead to the same $D^{1}$.

\section{Construction of Functionals \label{Functionals}}

\subsection{Formal Definition}

The energy, $\mathcal{E}\left(  D^{N}\right)  $, of a state $D^{N}%
\in\mathcal{S}_{N}$ of a system described by a Hamiltonian $H$ is given by
\begin{equation}
\mathcal{E}\left(  D^{N}\right)  =Tr\left\{  HD^{N}\right\}  .
\end{equation}
As mentioned in the preceding it is convenient to define the energy functional
on the cone of positive operators $\mathcal{B}_{1+}\left(  \mathcal{H}%
^{N}\right)  $ by
\begin{equation}
\mathcal{E}\left(  \mathcal{D}^{N}\right)  =\frac{Tr\left\{  H\mathcal{D}%
^{N}\right\}  }{Tr\left\{  \mathcal{D}^{N}\right\}  } \label{EnFunc}%
\end{equation}
leading to a functional that is invariant to complex scaling i.e.%
\begin{equation}
\mathcal{E}\left(  \mathcal{D}^{N}\right)  =\mathcal{E}\left(  \alpha
\mathcal{D}^{N}\right)  \quad\forall\alpha\in\mathbb{C}%
\end{equation}
so in particular
\begin{equation}
\mathcal{E}\left(  D^{N}\right)  =\mathcal{E}\left(  \mathcal{D}^{N}\right)
\quad\text{when }D^{N}\in\mathcal{S}_{N}\text{ and }D^{N}=\left(
Tr\mathcal{D}^{N}\right)  ^{-1}\mathcal{D}^{N}%
\end{equation}


One can construct a functional on equivalence classes of unnormalized states
$\left\{  \mathcal{D}^{1}\right\}  _{N}$ $\equiv\left\{  \mathcal{D}%
^{N}=\mathfrak{d}_{\Gamma}+\tilde{\Gamma}\left(  \mathfrak{d}_{1}\right)
\right\}  \equiv\left\{  \mathfrak{d}_{1}\right\}  $ by an analogous procedure
to the Levy-Lieb one \cite{LevyLieb} used in DFT by defining
\begin{equation}
\mathcal{F}\left(  \mathfrak{d}_{1}\right)  =\mathcal{E}\left(  \mathfrak{d}%
_{\Gamma\#}+\tilde{\Gamma}\left(  \mathfrak{d}_{1}\right)  \right)
=\mathcal{E}\left(  \mathcal{D}_{\#\ }^{N}\right)  =\underset{\mathcal{D}%
^{N}\in\left\{  \mathcal{D}^{1}\right\}  _{N}}{Min}\mathcal{E}\left(
\mathcal{D}^{N}\right)  \label{func}%
\end{equation}
which is the value of the energy functional $\mathcal{E}$ when $\mathfrak{d}%
_{\Gamma}=\mathfrak{d}_{\Gamma\#};\mathfrak{d}_{\Gamma},\mathfrak{d}%
_{\Gamma\#}\in\ker L_{N}^{1}$ for a fixed $\tilde{\Gamma}\left(
\mathfrak{d}_{1}\right)  \in\mathcal{V}_{\Gamma}$. As $\mathcal{F}$ is
invariant to scaling of its argument this leads to a functional defined on
normalized states given by
\begin{equation}
\mathcal{F}\left(  d_{1}\right)  =\mathcal{F}\left(  \mathfrak{d}_{1}\right)
\quad\text{when }d_{1}=\left(  Tr\mathfrak{d}_{1}\right)  ^{-1}\mathfrak{d}%
_{1}%
\end{equation}
In order to obtain the functional defined in eq. $\left(  \ref{func}\right)  $
one must construct a procedure that allows one to vary over the elements of
$\left\{  \mathcal{D}^{1}\right\}  _{N}$ which are the positive elements of
$\left[  \mathcal{D}^{1}\right]  _{N}$. In the following subsection we
describe how to realize such a variation.

First we first describe an abelian $\ $semigroup, discussed in more detail in
\cite{Weiner2002}, that acts on certain ensemble reference elements of
$\left\{  \mathcal{D}^{1}\right\}  _{N}$, called cyclic elements, in a way
that generates all states in $\left\{  \mathcal{D}^{1}\right\}  _{N}$, then we
describe how if one chooses pure states as reference elements one can generate
all pure states, $ext\left\{  \mathcal{D}^{1}\right\}  _{N}$ in $\left\{
\mathcal{D}^{1}\right\}  _{N}$, finally we discuss how to accomplish this in
an non redundant fashion.

\subsection{Intra Equivalence Class Variation \label{Intra}}

Any $N-$body density operator, representing either a pure or an ensemble state
of a $N-$particle system, can be generated by the action of the representation
$\mathcal{\hat{T}}\left(  \mathfrak{A}_{pos}\right)  $ of the abelian
semigroup, $\mathfrak{A}_{pos}$, by positive \textit{super operators} $\hat
{T}$:$\mathcal{B}_{1}\left(  \mathcal{H}^{N}\right)  \rightarrow
\mathcal{B}_{1}\left(  \mathcal{H}^{N}\right)  $, on \textit{cyclic }reference
states $\mathcal{D}_{ref}^{N}\in\mathcal{B}_{1+}\left(  \mathcal{H}%
^{N}\right)  $ \cite{Weiner2002} according to
\begin{equation}
\mathcal{D}^{N}=\hat{T}\left(  \mathcal{D}_{ref}^{N}\right)  =T\mathcal{D}%
_{ref}^{N}T^{\dagger} \label{GP}%
\end{equation}
The defining representations, $\mathcal{\hat{T}}\left(  \mathfrak{A}%
_{pos}\right)  $ \cite{notation1}, of $\mathfrak{A}_{pos}$ is the realization
of this semigroup provided by the set of operators $T\in\mathcal{B}\left(
\mathcal{H}^{N}\right)  $ that are diagonal in the space-spin coordinate
representation so that (see Appendices \ref{symnot} and \ref{Groups} for a
detailed definition of the notation)
\begin{equation}
\mathcal{\hat{T}}\left(  \mathfrak{A}_{pos}\right)  \left(  \mathfrak{b}%
_{0},\cdots,\mathfrak{b}_{N}\right)  =\left(  \mathfrak{b}_{0},\cdots
,\mathfrak{b}_{N}\right)  \mathcal{R}_{\mathcal{B}_{1}\left(  \mathcal{H}%
^{N}\right)  }\left(  \mathfrak{A}_{pos}\right)  .
\end{equation}
where $\left(  \mathfrak{b}_{0},\cdots,\mathfrak{b}_{N}\right)  $ denotes an
array of sub bases elements $\left\{  \mathfrak{b}_{p}\right\}  $ constituted
by ket-bra's that have precisely $p$ values in common i.e.
\begin{equation}
\mathfrak{b}_{p}=\left\{  \left.  \left\vert \mathsf{z}^{N}\right\rangle
\left\langle \mathsf{z}^{\prime N}\right\vert \right\vert \mathsf{z}%
^{N}\overset{p}{=}\mathsf{z}^{\prime N}\right\}
\end{equation}
and $\mathcal{R}_{\mathcal{B}_{1}\left(  \mathcal{H}^{N}\right)  }\left(
s\right)  $ is given by the generalized matrix formed by integral kernels i.e.%
\begin{equation}
\mathcal{R}_{\mathcal{B}_{1}\left(  \mathcal{H}^{N}\right)  }\left(  s\right)
=\left\{  \hat{T}\left(  \mathsf{z}^{N},\mathsf{z}^{\prime N}\right)
\right\}  ;\quad s\in\mathfrak{A}_{pos}%
\end{equation}
These operators are defined by \emph{real} \emph{symmetric }integral kernels
$\nu\left(  \mathsf{z}^{N}\right)  $ (which is also positive) and
$\omega\left(  \mathsf{z}^{N}\right)  $ in the following manner
\begin{align}
&  T=\nu_{N}e^{i\omega_{N}}\in\mathcal{B}\left(  \mathcal{H}^{N}\right)
\nonumber\\
&  \nu_{N}=\int\nu\left(  \mathsf{z}^{N}\right)  \left\vert \mathsf{z}%
^{N}\right\rangle \left\langle \mathsf{z}^{N}\right\vert d\mathsf{z}^{N}%
\in\mathcal{B}\left(  \mathcal{H}^{N}\right) \nonumber\\
&  \omega_{N}=\int\omega\left(  \mathsf{z}^{N}\right)  \left\vert
\mathsf{z}^{N}\right\rangle \left\langle \mathsf{z}^{N}\right\vert
d\mathsf{z}^{N}\in\mathcal{B}\left(  \mathcal{H}^{N}\right) \nonumber\\
&  T=\int T\left(  \mathsf{z}^{N}\right)  \left\vert \mathsf{z}^{N}%
\right\rangle \left\langle \mathsf{z}^{N}\right\vert d\mathsf{z}^{N}=\int
\nu\left(  \mathsf{z}^{N}\right)  e^{i\omega\left(  \mathsf{z}^{N}\right)
}\left\vert \mathsf{z}^{N}\right\rangle \left\langle \mathsf{z}^{N}\right\vert
d\mathsf{z}^{N} \label{TDiag}%
\end{align}
where
\begin{equation}
\left\vert \mathsf{z}^{N}\right\rangle =\left\vert \mathsf{z}_{1}%
\cdots\mathsf{z}_{N}\right\rangle ;\quad1\leq i<j\leq N.
\end{equation}
are the generalized eigenvectors of the $N-$particle space-spin position
operator and
\begin{align}
\nu\left(  \mathsf{z}^{N}\right)   &  =\left\langle \mathsf{z}^{N}%
|\nu\right\rangle \nonumber\\
\omega\left(  \mathsf{z}^{N}\right)   &  =\left\langle \mathsf{z}^{N}%
|\omega\right\rangle
\end{align}
The super operators $\hat{T}\in\mathcal{\hat{T}}\left(  \mathfrak{A}%
_{pos}\right)  $ can decrease, but not increase the rank of $\mathcal{D}%
_{ref}^{N}$. One should note that the need and motivation for considering
semigroups and not groups is due to the fact that in order to produce all
possible integral kernels associated with states, we need to consider
$\nu\left(  \mathsf{z}^{N}\right)  $ that could be zero on a measurable set of
points that would produce singular operators $\nu_{N}$. \emph{All}
unnormalized density operators can hence be parameterized in terms of the
integral kernels $\left\{  \nu\left(  \mathsf{z}^{N}\right)  ,\omega\left(
\mathsf{z}^{N}\right)  \right\}  $ once a suitable cyclic reference state
\cite{weiner2002} has been selected.

The semigroup $\mathfrak{A}_{pos}$ contains a sub semigroup, $\mathfrak{A}%
_{\ker}$ of maps (notational note: if there is no possibility of
misunderstanding we will signify representations of $\mathfrak{A}$ by
$\mathfrak{A}$ itself), that leave the kernel $KerL_{N}^{1}$ invariant, which
are characterized by the following \textit{generalized} \textit{strong
orthogonality }condition on their integral kernels, $\hat{T}\left(
\mathsf{z}^{N},\mathsf{z}^{\prime N}\right)  =T\left(  \mathsf{z}^{N}\right)
T\left(  \mathsf{z}^{\prime N}\right)  ^{\ast}$, in the space-spin coordinate
representation
\begin{equation}
\int X\left(  \mathsf{z}_{1}\mathsf{z}^{N-1}\mathsf{,z}_{1}^{\prime}%
\mathsf{z}^{N-1}\right)  T\left(  \mathsf{z}_{1}\mathsf{z}^{N-1}\right)
T\left(  \mathsf{z}_{1}^{\prime}\mathsf{z}^{N-1}\right)  ^{\ast}%
d\mathsf{z}^{N-1}\overset{a.e.}{=}0\quad\forall X\in KerL_{N}^{1} \label{G}%
\end{equation}


The decomposition eq. $\left(  \ref{nonorthog}\right)  $ of $\mathcal{B}%
_{1}\left(  \mathcal{H}^{N}\right)  $ \emph{does} \emph{not }fully reduce the
representation of $\mathcal{R}_{\mathcal{B}_{1}\left(  \mathcal{H}^{N}\right)
}\left(  \mathfrak{A}_{\ker}\right)  $ carried by the basis $\left\{
\mathfrak{b}_{0},\cdots,\mathfrak{b}_{N}\right\}  $ of $\mathcal{B}_{1}\left(
\mathcal{H}^{N}\right)  $ \cite{weiner2002} i.e. $\mathcal{R}_{\mathcal{B}%
_{1}\left(  \mathcal{H}^{N}\right)  }\left(  \mathfrak{A}_{\ker}\right)  $
properly contains the direct sum of the representations $\mathcal{R}%
_{KerL_{N}^{1}}\left(  \mathfrak{A}_{\ker}\right)  $ and $\mathcal{R}%
_{\mathcal{V}_{\tilde{\Gamma}}}\left(  \mathfrak{A}_{\ker}\right)  $ so that%
\begin{equation}
\mathcal{R}_{\mathcal{B}_{1}\left(  \mathcal{H}^{N}\right)  }\left(
\mathfrak{A}_{\ker}\right)  \supset\mathcal{R}_{KerL_{N}^{1}}\left(
\mathfrak{A}_{\ker}\right)  \oplus\mathcal{R}_{\mathcal{V}_{\tilde{\Gamma}}%
}\left(  \mathfrak{A}_{\ker}\right)
\end{equation}
and hence $\mathcal{R}_{\mathcal{B}_{1}\left(  \mathcal{H}^{N}\right)
}\left(  \mathfrak{A}_{\ker}\right)  $ contains transformations that mix
$KerL_{N}^{1}$ and $\mathcal{V}_{\tilde{\Gamma}}$. This rather arcane
observation is significant in that transformations belonging to $\mathcal{R}%
_{KerL_{N}^{1}}\left(  \mathfrak{A}_{\ker}\right)  \oplus\mathcal{R}%
_{\mathcal{V}_{\tilde{\Gamma}}}\left(  \mathfrak{A}_{\ker}\right)  $ \emph{do
not }generate all states, while transformations belonging to $\mathcal{R}%
_{\mathcal{B}_{1}\left(  \mathcal{H}^{N}\right)  }\left(  \mathfrak{A}_{\ker
}\right)  $ \emph{do}.

We can construct representations, $\left\{  \mathcal{R}_{F_{\tilde{\Gamma
}\left(  \mathfrak{d}_{1}\right)  }}\left(  \mathfrak{A}_{\ker}\right)
\right\}  $, - ($F_{\tilde{\Gamma}\left(  \mathfrak{d}_{1}\right)  }$ stands
for the fiber above the element $\tilde{\Gamma}\left(  \mathfrak{d}%
_{1}\right)  $ of the base of the vector bundle $\mathfrak{B}_{L_{N}^{1}}$
defined in eq. $\left(  \ref{VB}\right)  $), of $\mathfrak{A}_{\ker}$ that are
all isomorphic to $\mathcal{R}_{KerL_{N}^{1}}\left(  \mathfrak{A}_{\ker
}\right)  $ (see Appendix \ref{uker}), constituted by super operators $\hat
{T}$ $\in\mathcal{\hat{T}}_{F_{\tilde{\Gamma}\left(  \mathfrak{d}_{1}\right)
}}\left(  \mathfrak{A}_{\ker}\right)  \subset\mathcal{\hat{T}}\left(
\mathfrak{A}_{\ker}\right)  $ that have the property
\begin{equation}
\hat{T}\left(  \mathcal{D}_{ref}^{N}\right)  =\hat{T}\left(  \mathfrak{d}%
_{ref}+\tilde{\Gamma}\left(  \mathfrak{d}_{1}\right)  \right)  =\hat{T}\left(
\mathfrak{d}_{ref}\right)  +\tilde{\Gamma}\left(  \mathfrak{d}_{1}\right)
\label{SF}%
\end{equation}
If the state $\mathcal{D}_{ref}^{N}$ is a cyclic reference state then this
representation generates all the unnormalized states belonging to $\left\{
\mathfrak{d}_{1}\right\}  $ by its action on $\mathcal{D}_{ref}^{N}$.

In generalized "bra-ket" form the super operator transformations belonging to
$\mathcal{\hat{T}}_{F_{\tilde{\Gamma}\left(  \mathfrak{d}_{1}\right)  }%
}\left(  \mathfrak{A}_{\ker}\right)  $ can be expanded as
\begin{align}
\hat{T}  &  =%
%TCIMACRO{\dint \limits_{\mathsf{z}^{N}\overset{2}{\neq}\mathsf{z}^{\prime N}%
%}}%
%BeginExpansion
{\displaystyle\int\limits_{\mathsf{z}^{N}\overset{2}{\neq}\mathsf{z}^{\prime
N}}}
%EndExpansion
T\left(  \mathsf{z}^{N}\right)  T\left(  \mathsf{z}^{\prime N}\right)  ^{\ast
}\left\vert \left\vert \mathsf{z}^{N}\right\rangle \left\langle \mathsf{z}%
^{\prime N}\right\vert \right)  \left(  \left\vert \mathsf{z}^{N}\right\rangle
\left\langle \mathsf{z}^{\prime N}\right\vert \right\vert d\mathsf{z}%
^{N}d\mathsf{z}^{\prime N}\nonumber\\
&  \qquad\qquad+%
%TCIMACRO{\dsum \limits_{1\leq\kappa,\lambda\leq\infty}}%
%BeginExpansion
{\displaystyle\sum\limits_{1\leq\kappa,\lambda\leq\infty}}
%EndExpansion
\hat{T}_{F\kappa\lambda}\left\vert k_{\kappa}\right)  \left(  k_{\lambda}%
^{d}\right\vert +%
%TCIMACRO{\dsum \limits_{1\leq\kappa,\lambda\leq\infty}}%
%BeginExpansion
{\displaystyle\sum\limits_{1\leq\kappa,\lambda\leq\infty}}
%EndExpansion
\hat{T}_{FB\kappa\lambda}\left\vert k_{\kappa}\right)  \left(  \mathfrak{x}%
_{\lambda}^{d}\right\vert \nonumber\\
&  \qquad\qquad\qquad\qquad\qquad+\left\vert \tilde{\Gamma}\left(
\mathfrak{d}_{1}\right)  \right)  \left(  \tilde{\Gamma}\left(  \mathfrak{d}%
_{1}\right)  ^{d}\right\vert +%
%TCIMACRO{\dsum \limits_{2\leq\kappa,\lambda\leq\infty}}%
%BeginExpansion
{\displaystyle\sum\limits_{2\leq\kappa,\lambda\leq\infty}}
%EndExpansion
\hat{T}_{B\kappa\lambda}\left\vert \mathfrak{x}_{\kappa}\right)  \left(
\mathfrak{x}_{\lambda}^{d}\right\vert \label{bk}%
\end{align}
where
\begin{equation}
\left\{  \left\{  \left\vert \left\vert \mathsf{z}^{N}\right\rangle
\left\langle \mathsf{z}^{\prime N}\right\vert \right)  ;\mathsf{z}^{N}%
\overset{2}{\neq}\mathsf{z}^{\prime N}\right\}  ,\left\{  \left\vert
k_{\lambda}\right)  ;1\leq\lambda\leq\infty\right\}  ,\left\{  \left\vert
\tilde{\Gamma}\left(  \mathfrak{d}_{1}\right)  \right)  ,\left\vert
\mathfrak{x}_{\kappa}\right)  ;2\leq\kappa\leq\infty\right\}  \right\}
\end{equation}
is a basis for $\mathcal{B}_{1}\left(  \mathcal{H}^{N}\right)  $ (see Appendix
\ref{bases} for details),
\begin{equation}
\left\{  \left\{  \left(  \left\vert \mathsf{z}^{N}\right\rangle \left\langle
\mathsf{z}^{\prime N}\right\vert \right\vert ;\mathsf{z}^{N}\overset{2}{\neq
}\mathsf{z}^{\prime N}\right\}  ,\left\{  \left(  k_{\lambda}^{d}\right\vert
;1\leq\lambda\leq\infty\right\}  ,\left(  \tilde{\Gamma}\left(  \mathfrak{d}%
_{1}\right)  \right\vert ,\left(  \mathfrak{x}_{\lambda}\right\vert \right\}
\end{equation}
an associated biorthogonal basis for $\mathcal{B}\left(  \mathcal{H}%
^{N}\right)  $ and the super matrix elements are given by
\begin{align}
\hat{T}_{F\kappa\lambda}  &  =\left(  k_{\kappa}^{d}\left\vert \hat{T}\right.
k_{\lambda}\right)  =\int k_{\kappa}^{d}\left(  \mathsf{z}^{N},\mathsf{z}%
^{\prime N}\right)  ^{\ast}T\left(  \mathsf{z}^{N}\right)  T\left(
\mathsf{z}^{\prime N}\right)  ^{\ast}k_{\lambda}\left(  \mathsf{z}^{\prime
N},\mathsf{z}^{N}\right)  d\mathsf{z}^{N}d\mathsf{z}^{\prime N}\nonumber\\
\hat{T}_{FB\kappa\lambda}  &  =\left(  k_{\kappa}^{d}\left\vert \hat
{T}\right.  \mathfrak{x}_{\lambda}\right)  =\int k_{\kappa}^{d}\left(
\mathsf{z}^{N},\mathsf{z}^{\prime N}\right)  ^{\ast}T\left(  \mathsf{z}%
^{N}\right)  T\left(  \mathsf{z}^{\prime N}\right)  ^{\ast}\mathfrak{x}%
_{\lambda}\left(  \mathsf{z}^{\prime N},\mathsf{z}^{N}\right)  d\mathsf{z}%
^{N}d\mathsf{z}^{\prime N}\nonumber\\
\hat{T}_{B\kappa\lambda}  &  =\left(  \mathfrak{x}_{\kappa}^{d}\left\vert
\hat{T}\right.  \mathfrak{x}_{\lambda}\right)  =\int\mathfrak{x}_{\kappa}%
^{d}\left(  \mathsf{z}^{N},\mathsf{z}^{\prime N}\right)  ^{\ast}T\left(
\mathsf{z}^{N}\right)  T\left(  \mathsf{z}^{\prime N}\right)  ^{\ast
}\mathfrak{x}_{\lambda}\left(  \mathsf{z}^{\prime N},\mathsf{z}^{N}\right)
d\mathsf{z}^{N}d\mathsf{z}^{\prime N} \label{matel}%
\end{align}


Hence \emph{All} unnormalized density operators with the same $\mathcal{D}%
^{1}$ can be parameterized in terms of integral kernels $\left\{  T\left(
\mathsf{z}^{N}\right)  \right\}  \equiv\left\{  \nu\left(  \mathsf{z}%
^{N}\right)  ,\omega\left(  \mathsf{z}^{N}\right)  \right\}  $ that generate
the semigroup representations, $\mathcal{\hat{T}}_{F_{\tilde{\Gamma}\left(
\mathfrak{d}_{1}\right)  }}\left(  \mathfrak{A}_{\ker}\right)  $, of super
operator transformations $\hat{T}$ that can be expressed as eq. $\left(
\ref{bk}\right)  $. The restrictions on $T\left(  \mathsf{z}^{N}\right)  $ and
thus on $\hat{T}$, so that these transformations belong to $\mathcal{\hat{T}%
}_{F_{\tilde{\Gamma}\left(  \mathfrak{d}_{1}\right)  }}\left(  \mathfrak{A}%
_{\ker}\right)  $, are produced by%
\begin{align}
\left(  \left\vert \mathsf{z}^{N}\right\rangle \left\langle \mathsf{z}^{\prime
N}\right\vert \left\vert \hat{T}\right.  \tilde{\Gamma}\left(  \mathfrak{d}%
_{1}\right)  \right)   &  =0;\,\mathsf{z}^{N}\overset{2}{\nsim}\mathsf{z}%
^{\prime N}\nonumber\\
\left(  \mathfrak{x}_{\lambda}^{d}\left\vert \hat{T}\right.  k_{\kappa
}\right)   &  =0;\,2\leq\lambda\leq\infty,1\leq\kappa\leq\infty\nonumber\\
\left(  \tilde{\Gamma}\left(  \mathfrak{d}_{1}\right)  ^{d}\left\vert \hat
{T}\right.  k_{\kappa}\right)   &  =0;\,1\leq\kappa\leq\infty\nonumber\\
\left(  \mathfrak{x}_{\lambda}^{d}\left\vert \hat{T}\right.  \tilde{\Gamma
}\left(  \mathfrak{d}_{1}\right)  \right)   &  =0;\,2\leq\lambda\leq
\infty\nonumber\\
\left(  \tilde{\Gamma}\left(  \mathfrak{d}_{1}\right)  ^{d}\left\vert \hat
{T}\right.  \mathfrak{x}_{\kappa}\right)   &  =0;\,2\leq\kappa\leq
\infty\label{IntraFib}%
\end{align}
where "$\overset{2}{\nsim"}$ denotes "not equal in at least two places" (see
Appendix $\ref{uker}$ for a more detailed discussion). Once a local reference
state for the fiber is generated by the construction described in the
following subsection one can act on it by transformations, $\hat{T}$,
satisfying eq. $\left(  \ref{IntraFib}\right)  $ by the algorithm described in
Appendix \ref{uker}. These transformations, however, do not generate states
belonging to $\left\{  \mathfrak{d}_{1}\right\}  $ in a non redundant fashion,
i.e. two or more different $T\left(  \mathsf{z}^{N}\right)  $'s may produce
the same state in $\left\{  \mathfrak{d}_{1}\right\}  $. One can, however,
consider a subset of the transformations described in eq. $\left(
\ref{bk}\right)  $ of the form
\begin{align}
\hat{T}  &  =%
%TCIMACRO{\dint \limits_{\mathsf{z}^{N}\overset{2}{\neq}\mathsf{z}^{\prime N}%
%}}%
%BeginExpansion
{\displaystyle\int\limits_{\mathsf{z}^{N}\overset{2}{\neq}\mathsf{z}^{\prime
N}}}
%EndExpansion
T\left(  \mathsf{z}^{N}\right)  T\left(  \mathsf{z}^{\prime N}\right)  ^{\ast
}\left\vert \left\vert \mathsf{z}^{N}\right\rangle \left\langle \mathsf{z}%
^{\prime N}\right\vert \right)  \left(  \left\vert \mathsf{z}^{N}\right\rangle
\left\langle \mathsf{z}^{\prime N}\right\vert \right\vert d\mathsf{z}%
^{N}d\mathsf{z}^{\prime N}\nonumber\\
&  +%
%TCIMACRO{\dsum \limits_{1\leq\kappa,\lambda\leq\infty}}%
%BeginExpansion
{\displaystyle\sum\limits_{1\leq\kappa,\lambda\leq\infty}}
%EndExpansion
\hat{T}_{F\kappa\lambda}\left\vert k_{\kappa}\right)  \left(  k_{\lambda}%
^{d}\right\vert +%
%TCIMACRO{\dsum \limits_{1\leq\kappa\leq\infty}}%
%BeginExpansion
{\displaystyle\sum\limits_{1\leq\kappa\leq\infty}}
%EndExpansion
\hat{T}_{FB\kappa\lambda}\left\vert k_{\kappa}\right)  \left(  \tilde{\Gamma
}\left(  \mathfrak{d}_{1}\right)  ^{d}\right\vert +\left\vert \tilde{\Gamma
}\left(  \mathfrak{d}_{1}\right)  \right)  \left(  \tilde{\Gamma}\left(
\mathfrak{d}_{1}\right)  ^{d}\right\vert \label{bknr}%
\end{align}
that produce a coset $\mathcal{C}_{F_{\tilde{\Gamma}\left(  \mathfrak{d}%
_{1}\right)  }}$ of $\mathcal{\hat{T}}_{F_{\tilde{\Gamma}\left(
\mathfrak{d}_{1}\right)  }}\left(  \mathfrak{A}_{\ker}\right)  $, (Appendix
\ref{uker}), based on a more restricted set of $T\left(  \mathsf{z}%
^{N}\right)  $, that forms a non redundant parameterization of the states
$\left\{  \mathfrak{d}_{1}\right\}  $. In order to construct transformations
belonging to this coset one must impose a further restriction on $T\left(
\mathsf{z}^{N}\right)  $ given by
\begin{equation}
\left(  \mathfrak{x}_{\lambda}^{d}\left\vert \hat{T}\right.  \mathfrak{x}%
_{\kappa}\right)  =0;\quad2\leq\lambda,\kappa\leq\infty
\end{equation}


Pure state functionals can be determined by the preceding by restricting
attention to the orbit of pure states $ext\left\{  \mathcal{D}^{1}\right\}
_{N}\equiv ext\left\{  \mathfrak{d}_{1}\right\}  $ contained in $\left\{
\mathfrak{d}_{1}\right\}  $. This orbit can be generated from a pure reference
state $\mathcal{D}_{ref}^{N}$ \cite{weiner2002} by transformation belonging to
$\mathcal{C}_{F_{\tilde{\Gamma}\left(  \mathfrak{d}_{1}\right)  }} $ of the
form eq. $\left(  \ref{bknr}\right)  $ where $\mathfrak{d}_{1}$ corresponds to
a pure state.

\subsection{Procedure for Generating Exact Functionals}

\subsubsection{Generalized Strong Orthogonality}

Starting from a given $\mathcal{D}_{0_{F}}^{N}\left(  \rho\right)  $ that is a
cyclic reference state for the $N$-particle states $\left\{  \rho\right\}  $
and has the expansion
\begin{equation}
\mathcal{D}_{0_{F}}^{N}\left(  \rho\right)  =\int\mathcal{D}_{0_{F}}%
^{N}\left[  \rho\right]  \left(  \mathsf{z}^{N},\mathsf{z}^{N\prime}\right)
\left\vert \left\vert \mathsf{z}^{N}\right\rangle \left\langle \mathsf{z}%
^{N\prime}\right\vert \right)  \left(  \left\vert \mathsf{z}^{N}\right\rangle
\left\langle \mathsf{z}^{N\prime}\right\vert \right\vert d\mathsf{z}%
^{N}d\mathsf{z}^{\prime N}%
\end{equation}
($\left[  \rho\right]  $ signifies that the argument $\rho$ is not the focus
of the current expression, but should not be ignored) we can obtain all of
$\left\{  \rho\right\}  $ in a nonredundent manner by $\hat{T}\in
\mathcal{C}_{F_{\tilde{\Gamma}\left(  \mathfrak{d}_{1}\right)  }}$ as $\hat
{T}\left(  \mathcal{D}_{0_{F}}^{N}\left(  \rho\right)  \right)  $ with
expansions
\begin{equation}
\hat{T}\left(  \mathcal{D}_{0_{F}}^{N}\left(  \rho\right)  \right)  =\int
T\left(  \mathsf{z}^{N}\right)  \mathcal{D}_{0_{F}}^{N}\left[  \rho\right]
\left(  \mathsf{z}^{N},\mathsf{z}^{N\prime}\right)  T\left(  \mathsf{z}%
^{\prime N}\right)  ^{\ast}\left\vert \left\vert \mathsf{z}^{N}\right\rangle
\left\langle \mathsf{z}^{\prime N}\right\vert \right)  \left(  \left\vert
\mathsf{z}^{N}\right\rangle \left\langle \mathsf{z}^{\prime N}\right\vert
\right\vert d\mathsf{z}^{N}d\mathsf{z}^{\prime N}%
\end{equation}
A \emph{necessary} condition that $\hat{T}\in$ $\mathcal{\hat{T}}\left(
\mathfrak{A}_{\ker}\right)  $ is that
\begin{equation}
L_{N}^{1}\left(  \hat{T}\left(  \mathcal{D}_{0_{F}}^{N}\left[  \rho\right]
\right)  \right)  =L_{N}^{1}\left(  \mathcal{D}_{0_{F}}^{N}\left[
\rho\right]  \right)
\end{equation}
If we express $\hat{T}$ as
\begin{equation}
\hat{T}=\hat{I}+\hat{\Omega}%
\end{equation}
we obtain that a \emph{necessary} condition for $\hat{T}\in$ $\mathcal{\hat
{T}}\left(  \mathfrak{A}_{\ker}\right)  $ is that
\begin{equation}
L_{N}^{1}\left(  \hat{\Omega}\left(  \mathcal{D}_{0_{F}}^{N}\left[
\rho\right]  \right)  \right)  =0
\end{equation}
which can be expressed in terms of integral kernels as
\begin{equation}
\int\Omega\left(  \mathsf{z}_{1},\mathsf{z}^{N-1},\mathsf{z}_{1}^{\prime
},\mathsf{z}^{N-1}\right)  \mathcal{D}_{0_{F}}^{N}\left[  \rho\right]  \left(
\mathsf{z}_{1},\mathsf{z}^{N-1},\mathsf{z}_{1}^{\prime},\mathsf{z}%
^{N-1}\right)  d\mathsf{z}^{N-1}\overset{a.e}{=}0 \label{fogso}%
\end{equation}
where the diagonal super operator $\hat{T}$ has an integral kernel given by
\begin{equation}
\hat{T}\left(  \mathsf{z}^{N},\mathsf{z}^{\prime N}\right)  =\delta\left(
\mathsf{z}^{N}\right)  \delta\left(  \mathsf{z}^{\prime N}\right)
+\Omega\left(  \mathsf{z}^{N},\mathsf{z}^{\prime N}\right)  \label{supker}%
\end{equation}
with
\begin{equation}
\delta\left(  \mathsf{z}^{N}\right)  =%
%TCIMACRO{\dprod \limits_{1\leq j\leq N}}%
%BeginExpansion
{\displaystyle\prod\limits_{1\leq j\leq N}}
%EndExpansion
\delta\left(  \mathsf{z}_{j}\right)  .
\end{equation}
We name the condition expressed in eq. $\left(  \ref{fogso}\right)  $ as the
\textit{first order generalized strong orthogonality }(fogso) wrt
$\mathcal{D}_{0_{F}}^{N}\left[  \rho\right]  $\textit{.}

In order to relate parts of the integral kernel of $\hat{T}$ to parts of the
integral kernel of $T$ it is convenient to use the expression
\begin{equation}
T\left(  \mathsf{z}^{N}\right)  =\nu\left(  \mathsf{z}^{N}\right)
e^{-i\omega\left(  \mathsf{z}^{N}\right)  }=\left(  \delta\left(
\mathsf{z}^{N}\right)  +\Phi\left(  \mathsf{z}^{N}\right)  \right)  ^{\frac
{1}{2}}e^{-i\omega\left(  \mathsf{z}^{N}\right)  } \label{T}%
\end{equation}
then consider the series expansion
\begin{align}
\nu\left(  \mathsf{z}^{N}\right)   &  =\left(  \delta\left(  \mathsf{z}%
^{N}\right)  +\Phi\left(  \mathsf{z}^{N}\right)  \right)  ^{\frac{1}{2}%
}=\delta\left(  \mathsf{z}^{N}\right)  -\frac{1}{2}\Phi\left(  \mathsf{z}%
^{N}\right)  -\frac{1}{8}\Phi\left(  \mathsf{z}^{N}\right)  ^{2}-\frac{1}%
{48}\Phi\left(  \mathsf{z}^{N}\right)  ^{3}-\cdots\nonumber\\
&  \equiv\delta\left(  \mathsf{z}^{N}\right)  +\Lambda\left(  \mathsf{z}%
^{N}\right)
\end{align}
where we define
\begin{equation}
\Lambda\left(  \mathsf{z}^{N}\right)  =-\left(  \frac{1}{2}\Phi\left(
\mathsf{z}^{N}\right)  +\frac{1}{8}\Phi\left(  \mathsf{z}^{N}\right)
^{2}+\frac{1}{48}\Phi\left(  \mathsf{z}^{N}\right)  ^{3}+\cdots\right)
\label{rootex}%
\end{equation}
which leads to
\begin{align}
&  T\left(  \mathsf{z}^{N}\right)  D_{0_{F}}^{N}\left[  \rho\right]  \left(
\mathsf{z}^{N},\mathsf{z}^{N\prime}\right)  T\left(  \mathsf{z}^{\prime
N}\right)  ^{\ast}\nonumber\\
&  =\left(  \delta\left(  \mathsf{z}^{N}\right)  +\Lambda\left(
\mathsf{z}^{N}\right)  \right)  D_{0_{F}}^{N}\left[  \rho\right]  \left(
\mathsf{z}^{N},\mathsf{z}^{N\prime}\right)  \left(  \delta\left(
\mathsf{z}^{\prime N}\right)  +\Lambda\left(  \mathsf{z}^{\prime N}\right)
\right)  e^{-i\left(  \omega\left(  \mathsf{z}_{1},\mathsf{z}^{N-1}\right)
-\omega\left(  \mathsf{z}_{1}^{\prime},\mathsf{z}^{\prime N-1}\right)
\right)  }\nonumber\\
&  =\left\{  \delta\left(  \mathsf{z}^{N}\right)  \delta\left(  \mathsf{z}%
^{\prime N}\right)  +\Lambda\left(  \mathsf{z}^{N}\right)  D_{0_{F}}%
^{N}\left[  \rho\right]  \left(  \mathsf{z}^{N},\mathsf{z}^{N\prime}\right)
+D_{0_{F}}^{N}\left[  \rho\right]  \left(  \mathsf{z}^{N},\mathsf{z}^{N\prime
}\right)  \Lambda\left(  \mathsf{z}^{\prime N}\right)  \right. \nonumber\\
&  \qquad\qquad+\Lambda\left(  \mathsf{z}^{N}\right)  D_{0_{F}}^{N}\left[
\rho\right]  \left(  \mathsf{z}^{N},\mathsf{z}^{N\prime}\right)
\Lambda\left(  \mathsf{z}^{\prime N}\right)  \left.  e^{-i\left(
\omega\left(  \mathsf{z}_{1},\mathsf{z}^{N-1}\right)  -\omega\left(
\mathsf{z}_{1}^{\prime},\mathsf{z}^{\prime N-1}\right)  \right)  }\right\}
\nonumber\\
&  =\delta\left(  \mathsf{z}^{N}\right)  \delta\left(  \mathsf{z}^{\prime
N}\right)  +\left\{  \Lambda\left(  \mathsf{z}^{N}\right)  D_{0_{F}}%
^{N}\left[  \rho\right]  \left(  \mathsf{z}^{N},\mathsf{z}^{N\prime}\right)
+D_{0_{F}}^{N}\left[  \rho\right]  \left(  \mathsf{z}^{N},\mathsf{z}^{N\prime
}\right)  \Lambda\left(  \mathsf{z}^{\prime N}\right)  \right. \nonumber\\
&  \left.  +\Lambda\left(  \mathsf{z}^{N}\right)  D_{0_{F}}^{N}\left[
\rho\right]  \left(  \mathsf{z}^{N},\mathsf{z}^{N\prime}\right)
\Lambda\left(  \mathsf{z}^{\prime N}\right)  \left\{  e^{-i\left(
\omega\left(  \mathsf{z}_{1},\mathsf{z}^{N-1}\right)  -\omega\left(
\mathsf{z}_{1}^{\prime},\mathsf{z}^{\prime N-1}\right)  \right)  }%
-\delta\left(  \mathsf{z}^{N}\right)  \delta\left(  \mathsf{z}^{\prime
N}\right)  \right\}  \right\}  \label{detker}%
\end{align}
Comparing eq. $\left(  \ref{detker}\right)  $ with eq. $\left(  \ref{supker}%
\right)  $ we deduce the correspondence
\begin{align}
&  \Omega\left(  \mathsf{z}^{N},\mathsf{z}^{\prime N}\right)  \mathcal{D}%
_{0_{F}}^{N}\left[  \rho\right]  \left(  \mathsf{z}^{N},\mathsf{z}^{\prime
N\prime}\right)  =\nonumber\\
&  \left\{  \Lambda\left(  \mathsf{z}^{N}\right)  D_{0_{F}}^{N}\left[
\rho\right]  \left(  \mathsf{z}^{N},\mathsf{z}^{N\prime}\right)  +D_{0_{F}%
}^{N}\left[  \rho\right]  \left(  \mathsf{z}^{N},\mathsf{z}^{N\prime}\right)
\Lambda\left(  \mathsf{z}^{\prime N}\right)  \right. \nonumber\\
&  \left.  +\Lambda\left(  \mathsf{z}^{N}\right)  D_{0_{F}}^{N}\left[
\rho\right]  \left(  \mathsf{z}^{N},\mathsf{z}^{N\prime}\right)
\Lambda\left(  \mathsf{z}^{\prime N}\right)  \left\{  e^{-i\left(
\omega\left(  \mathsf{z}_{1},\mathsf{z}^{N-1}\right)  -\omega\left(
\mathsf{z}_{1}^{\prime},\mathsf{z}^{\prime N-1}\right)  \right)  }%
-\delta\left(  \mathsf{z}^{N}\right)  \delta\left(  \mathsf{z}^{\prime
N}\right)  \right\}  \right\}
\end{align}
Thus $\hat{T}\in$ $\mathcal{\hat{T}}\left(  \mathfrak{A}_{\ker}\right)  $
implies that
\begin{align}
&  \int\left\{  \Lambda\left(  \mathsf{z}_{1},\mathsf{z}^{N-1}\right)
D_{0_{F}}^{N}\left[  \rho\right]  \left(  \mathsf{z}_{1},\mathsf{z}%
^{N-1},\mathsf{z}_{1}^{\prime},\mathsf{z}^{N-1}\right)  \Lambda\left(
\mathsf{z}_{1}^{\prime},\mathsf{z}^{N-1}\right)  +\right. \nonumber\\
&  \left.  \Lambda\left(  \mathsf{z}_{1},\mathsf{z}^{N-1}\right)  D_{0_{F}%
}^{N}\left[  \rho\right]  \left(  \mathsf{z}_{1},\mathsf{z}^{N-1}%
,\mathsf{z}_{1}^{\prime},\mathsf{z}^{N-1}\right)  +D_{0_{F}}^{N}\left[
\rho\right]  \left(  \mathsf{z}_{1},\mathsf{z}^{N-1},\mathsf{z}_{1}^{\prime
},\mathsf{z}^{N-1}\right)  \Lambda\left(  \mathsf{z}_{1}^{\prime}%
,\mathsf{z}^{N-1}\right)  \right\} \nonumber\\
&  \qquad\qquad\qquad\qquad\qquad\qquad\qquad\qquad\qquad\qquad e^{i\left(
\omega\left(  \mathsf{z}_{1},\mathsf{z}^{N-1}\right)  -\omega\left(
\mathsf{z}_{1}^{\prime},\mathsf{z}^{\prime N-1}\right)  \right)  }%
d\mathsf{z}^{N-1}\overset{a.e}{=}0
\end{align}
which further implies that
\begin{align}
&  \int\left\{  \Lambda\left(  \mathsf{z}_{1},\mathsf{z}^{N-1}\right)
D_{0_{F}}^{N}\left[  \rho\right]  \left(  \mathsf{z}_{1},\mathsf{z}%
^{N-1},\mathsf{z}_{1}^{\prime},\mathsf{z}^{\prime N-1}\right)  \Lambda\left(
\mathsf{z}_{1}^{\prime},\mathsf{z}^{\prime N-1}\right)  +\right. \nonumber\\
&  \left.  \Lambda\left(  \mathsf{z}_{1},\mathsf{z}^{N-1}\right)  D_{0_{F}%
}^{N}\left[  \rho\right]  \left(  \mathsf{z}_{1},\mathsf{z}^{N-1}%
,\mathsf{z}_{1}^{\prime},\mathsf{z}^{\prime N-1}\right)  +D_{0_{F}}^{N}\left[
\rho\right]  \left(  \mathsf{z}_{1},\mathsf{z}^{N-1},\mathsf{z}_{1}^{\prime
},\mathsf{z}^{\prime N-1}\right)  \Lambda\left(  \mathsf{z}_{1}^{\prime
},\mathsf{z}^{\prime N-1}\right)  \right\}  d\mathsf{z}^{N-1}\nonumber\\
&  \qquad\qquad\qquad\qquad\qquad\qquad\qquad\qquad\qquad\qquad\qquad
\qquad\qquad\qquad\qquad\qquad\qquad\qquad\overset{a.e}{=}0
\end{align}


Using the expression eq. $\left(  \ref{T}\right)  $ in the expression for the
super operator transformation $\hat{T}$ we observe that the fogso condition
implies that the integral kernels $\left\{  \Phi\left(  \mathsf{z}^{N}\right)
,e^{i\omega\left(  \mathsf{z}^{N}\right)  }\right\}  $ must satisfy
\begin{align}
&
%TCIMACRO{\dint }%
%BeginExpansion
{\displaystyle\int}
%EndExpansion
\left\{  \left(  \delta\left(  \mathsf{z}_{1},\mathsf{z}^{N-1}\right)
+\Phi\left(  \mathsf{z}_{1},\mathsf{z}^{N-1}\right)  \right)  ^{\frac{1}{2}%
}e^{-i\omega\left(  \mathsf{z}^{N}\right)  }D_{0_{F}}^{N}\left[  \rho\right]
\left(  \mathsf{z}_{1},\mathsf{z}^{N-1},\mathsf{z},\mathsf{z}^{N-1}\right)
\times\right. \label{invcon}\\
&  \quad\left(  \delta\left(  \mathsf{z}_{1}^{\prime},\mathsf{z}^{N-1}\right)
+\Phi\left(  \mathsf{z}_{1}^{\prime},\mathsf{z}^{N-1}\right)  ^{\ast}\right)
^{\frac{1}{2}}e^{i\omega\left(  \mathsf{z}^{N}\right)  }-\left.  D_{0_{F}}%
^{N}\left[  \rho\right]  \left(  \mathsf{z}_{1},\mathsf{z}^{N-1}%
,\mathsf{z}_{1}^{\prime},\mathsf{z}^{N-1}\right)  \right\}  d\mathsf{z}%
^{N-1}\overset{a.e}{=}0.\nonumber
\end{align}
By setting $\mathsf{z}_{1}^{\prime}=\mathsf{z}_{1}$ in eq.$\left(
\ref{invcon}\right)  $ we obtain the \textit{zeroth order generalized strong
orthogonality }(zogso)\emph{\ }condition\emph{\ }%
\begin{equation}%
%TCIMACRO{\dint }%
%BeginExpansion
{\displaystyle\int}
%EndExpansion
\Phi\left(  \mathsf{z}_{1},\mathsf{z}^{N-1}\right)  D_{0_{F}}^{N}\left[
\rho\right]  \left(  \mathsf{z}_{1},\mathsf{z}^{N-1},\mathsf{z}_{1}%
,\mathsf{z}^{N-1}\right)  d\mathsf{z}^{N-1}\overset{a.e}{=}0 \label{gso}%
\end{equation}
which is a necessary condition for $\hat{T}\in$ $\mathcal{\hat{T}}\left(
\mathfrak{A}_{\ker}\right)  $. If $\left\{  \Phi_{1}\left(  \mathsf{z}%
^{N}\right)  ,\Phi_{2}\left(  \mathsf{z}^{N}\right)  \right\}  $ satisfies the
zodso wrt $D_{0_{F}}^{N}\left[  \rho\right]  \left(  \mathsf{z}^{N}%
,\mathsf{z}^{N}\right)  $ then so does the product $\Phi_{1}\left(
\mathsf{z}^{N}\right)  \Phi_{2}\left(  \mathsf{z}^{N}\right)  $ and in fact
the set of generalized zeroth order strongly orthogonal $\ker$kernels wrt
$D_{0_{F}}^{N}\left[  \rho\right]  \left(  \mathsf{z}^{N},\mathsf{z}%
^{N}\right)  $ form an algebra.

If $D_{0_{F}}^{N}\left[  \rho\right]  $ is a pure state
\begin{equation}
D_{0_{F}}^{N}\left[  \rho\right]  \left(  \mathsf{z}^{N},\mathsf{z}^{\prime
N}\right)  =\Psi_{0_{F}}\left(  \mathsf{z}^{N}\right)  \Psi_{0_{F}}\left(
\mathsf{z}^{\prime N}\right)  ^{\ast}%
\end{equation}
we obtain
\begin{align}
&
%TCIMACRO{\dint }%
%BeginExpansion
{\displaystyle\int}
%EndExpansion
\left\{  \left(  \Phi\left(  \mathsf{z}_{1},\mathsf{z}^{N-1}\right)
\Phi\left(  \mathsf{z}_{1}^{\prime},\mathsf{z}^{N-1}\right)  ^{\ast}%
+\Phi\left(  \mathsf{z}_{1},\mathsf{z}^{N-1}\right)  +\Phi\left(
\mathsf{z}_{1}^{\prime},\mathsf{z}^{N-1}\right)  ^{\ast}\right)  \right.
\nonumber\\
&  \left.  \Psi_{0_{F}}\left(  \mathsf{z}_{1},\mathsf{z}^{N-1}\right)
\Psi_{0_{F}}^{\ast}\left(  \mathsf{z}_{1}^{\prime},\mathsf{z}^{N-1}\right)
d\mathsf{z}^{N-1}\right\}  \overset{a.e}{=}0\nonumber\\
&  \Rightarrow\int\Phi\left(  \mathsf{z}_{1},\mathsf{z}^{N-1}\right)
\Psi_{0_{F}}\left(  \mathsf{z}_{1},\mathsf{z}^{N-1}\right)  \Psi_{0_{F}}%
^{\ast}\left(  \mathsf{z}_{1}^{\prime},\mathsf{z}^{N-1}\right)  \Phi\left(
\mathsf{z}_{1}^{\prime},\mathsf{z}^{N-1}\right)  ^{\ast}d\mathsf{z}%
^{N-1}\nonumber\\
&  \left.  =\right.  -\int\left\{  \Phi\left(  \mathsf{z}_{1},\mathsf{z}%
^{N-1}\right)  \Psi_{0_{F}}\left(  \mathsf{z}_{1},\mathsf{z}^{N-1}\right)
\Psi_{0_{F}}^{\ast}\left(  \mathsf{z}_{1}^{\prime},\mathsf{z}^{N-1}\right)
\right. \nonumber\\
&  \left.  +\Psi_{0_{F}}\left(  \mathsf{z}_{1},\mathsf{z}^{N-1}\right)
\Psi_{0_{F}}^{\ast}\left(  \mathsf{z}_{1}^{\prime},\mathsf{z}^{N-1}\right)
\Phi\left(  \mathsf{z}_{1}^{\prime},\mathsf{z}^{N-1}\right)  ^{\ast}\right\}
d\mathsf{z}^{N-1}%
\end{align}


\subsubsection{Effective Hamiltonian}

The $N$-particle density operators $\mathcal{D}^{N}$ are parameterized by four
functions $\left\{  \nu_{B},\omega_{B},\nu_{F},\omega_{F}\right\}  $ of the
coordinates $\mathsf{z}^{N}$, the functions $\left\{  \nu_{B},\omega
_{B}\right\}  $ correspond to a particular FORDO and the functions $\left\{
\nu_{F},\omega_{F}\right\}  $ to different $N$-particle density operators that
have the same FORDO. The value of the energy functional, $\mathcal{F}$, at a
given FORDO $\mathcal{D}^{1}=\mathcal{D}^{1}\left(  \nu_{B},\omega_{B}\right)
$ is defined by
\begin{align}
\mathcal{F}\left(  \nu_{B},\omega_{B}\right)   &  =\mathcal{F}\left(  \nu
_{B},\omega_{B},\nu_{\#F},\omega_{\#F}\right) \nonumber\\
&  =\underset{\nu_{F},\omega_{F}}{Min}\left\{  \mathcal{E}\left(
\mathcal{D}^{N}\left(  \nu_{B},\omega_{B},\nu_{F},\omega_{F}\right)  \right)
\right\} \nonumber\\
&  =\underset{\nu_{F},\omega_{F}}{Min}\left\{  \frac{Tr\left\{  H\mathcal{D}%
^{N}\left(  \nu_{B},\omega_{B},\nu_{F},\omega_{F}\right)  \right\}
}{Tr\left\{  \mathcal{D}^{N}\left(  \nu_{B},\omega_{B},\nu_{F},\omega
_{F}\right)  \right\}  }\right\} \nonumber\\
&  =\frac{Tr\left\{  H\mathcal{D}^{N}\left(  \nu_{B},\omega_{B},\nu
_{F\#}\left(  \nu_{B},\omega_{B}\right)  ,\omega_{F\#}\left(  \nu_{B}%
,\omega_{B}\right)  \right)  \right\}  }{Tr\left\{  \mathcal{D}^{N}\left(
\nu_{B},\omega_{B},\nu_{F\#}\left(  \nu_{B},\omega_{B}\right)  ,\omega
_{F\#}\left(  \nu_{B},\omega_{B}\right)  \right)  \right\}  }%
\end{align}


where
\begin{align}
\mathcal{D}^{N}\left(  \nu_{B},\omega_{B},\nu_{F},\omega_{F}\right)   &
=T\left(  \nu_{F},\omega_{F}\right)  \mathcal{D}^{N}\left(  \nu_{B},\omega
_{B}\right)  T\left(  \nu_{F},\omega_{F}\right)  ^{\dagger}\nonumber\\
&  =T\left(  \nu_{F},\omega_{F}\right)  T\left(  \nu_{B},\omega_{B}\right)
\mathcal{D}_{ref}^{N}T\left(  \nu_{B},\omega_{B}\right)  ^{\dagger}T\left(
\nu_{F},\omega_{F}\right)  ^{\dagger}%
\end{align}
and $\mathcal{D}_{ref}^{N}$ is an IP reference state. The action of the
transformations $T\left(  \nu_{F},\omega_{F}\right)  $ and $T\left(  \nu
_{B},\omega_{B}\right)  $ can be incorporated into an effective Hamiltonian
$H\left(  \nu_{B},\omega_{B},\nu_{F},\omega_{F}\right)  $ defined by
\begin{equation}
H\left(  \nu_{B},\omega_{B},\nu_{F},\omega_{F}\right)  =T\left(  \nu
_{F},\omega_{F}\right)  ^{\dagger}T\left(  \nu_{B},\omega_{B}\right)
^{\dagger}HT\left(  \nu_{B},\omega_{B}\right)  T\left(  \nu_{F},\omega
_{F}\right)
\end{equation}
so that all expectation values can be taken wrt the IPS $\mathcal{D}_{ref}%
^{N}$.

The Hamiltonians, $H$, we consider are of the form%
\begin{equation}
H=H_{0}-\mathbf{P\bullet P}+V_{1}+V_{2}%
\end{equation}
where $H_{0}$ is a constant, $\mathbf{P}$ the total \emph{generalized
}momentum operator for the system, $V_{1}$ is the sum of all one particle
potential operators and $V_{2}$ is the sum of all two particle potential
operators. The \emph{generalized }momentum is defined in terms of one particle
\emph{generalized }momentum operators as
\begin{equation}
\mathbf{P=}%
%TCIMACRO{\dsum \limits_{1\leq j\leq N}}%
%BeginExpansion
{\displaystyle\sum\limits_{1\leq j\leq N}}
%EndExpansion
\mathbf{p}_{j}%
\end{equation}
where
\begin{equation}
\mathbf{p}_{j}=-i\mathbf{\nabla}_{j}-\mathbf{A}_{j}%
\end{equation}
$\mathbf{A}_{j}$ is the total external vector potential acting on particle $j$
and
\begin{equation}
\mathbf{\nabla}_{j}=\left(
\begin{array}
[c]{c}%
\frac{\partial}{\partial x_{j}}\\
\frac{\partial}{\partial y_{j}}\\
\frac{\partial}{\partial z_{j}}%
\end{array}
\right)  ;\quad\mathbf{A}_{j}=\left(
\begin{array}
[c]{c}%
A_{jx}\\
A_{jy}\\
A_{jz}%
\end{array}
\right)
\end{equation}
The potentials $\left\{  \mathbf{A}_{j}\right\}  $, $V_{1}$ and $V_{2}$ are
all diagonal in the space-spin representation. The operators $\left\{
T\left(  \nu_{F},\omega_{F}\right)  ,T\left(  \nu_{B},\omega_{B}\right)
\right\}  $ \emph{all} \emph{commute }with these potentials, which is of
course by design as we choose the representation $\mathcal{\hat{T}}\left(
\mathfrak{A}_{pos}\right)  $ to be by operators diagonal in the space-spin
representation and hence commute with the potential operators so that
\begin{align}
T\left(  \nu,\omega\right)  ^{\dagger}VT\left(  \nu,\omega\right)   &
=\nu_{N}e^{-i\omega_{N}}V\nu_{N}e^{i\omega_{N}}=\left(  \nu_{N}\right)
^{2}V\nonumber\\
T\left(  \nu,\omega\right)  ^{\dagger}\mathbf{A}_{j}T\left(  \nu
,\omega\right)   &  =\nu_{N}e^{-i\omega_{N}}\mathbf{A}_{j}\nu_{N}%
e^{i\omega_{N}}=\left(  \nu_{N}\right)  ^{2}\mathbf{A}_{j}%
\end{align}
where
\begin{align}
\left(  \nu_{N}\right)  ^{2}V &  =\int\nu\left(  \mathsf{z}^{N}\right)
^{2}V\left(  \mathsf{z}^{N}\right)  \left\vert \mathsf{z}^{N}\right\rangle
\left\langle \mathsf{z}^{N}\right\vert d\mathsf{z}^{N}\nonumber\\
\left(  \nu_{N}\right)  ^{2}A_{j\kappa} &  =\int\nu\left(  \mathsf{z}%
^{N}\right)  ^{2}A_{j\kappa}\left(  \mathsf{z}^{N}\right)  \left\vert
\mathsf{z}^{N}\right\rangle \left\langle \mathsf{z}^{N}\right\vert
d\mathsf{z}^{N};\qquad\kappa=x,y,z
\end{align}
i.e. the potentials are scaled by the integral kernels $\left\{  \nu\left(
\mathsf{z}^{N}\right)  \right\}  $. It is \emph{important }to note that while
the Hamiltonian $H$ only contains potentials describing up to two particle
interactions, the integral kernels $\left\{  \nu\left(  \mathsf{z}^{N}\right)
\right\}  $ are in general multiparticle functions and thus produce
multiparticle potentials $\left\{  \nu\left(  \mathsf{z}^{N}\right)
^{2}V\left(  \mathsf{z}^{N}\right)  \right\}  $ describing up to simultaneous
$N$-particle interactions. Furthermore the $N$-particle local scaling factors
$\nu\left(  \mathsf{z}^{N}\right)  $ contain phase information pertaining to
$1,2,\cdots,M<N$ particle interactions and the $N$-particle local phase
$\omega\left(  \mathsf{z}^{N}\right)  $ contains local scaling information
pertaining to $1,2,\cdots,M<N$ particle interactions as will discussed in the examples.

The effect of the unitary part $\left\{  e^{i\omega_{N}}\right\}  $ of the
transformations $\left\{  T\left(  \nu,\omega\right)  \right\}  $ on the
momentum operators $\left\{  -i\nabla_{j}\right\}  $ (for simplicity we
consider the case where there is no external vector potential) is to produce a
generalized gauge transformation i.e. producing the component operators
\begin{equation}
e^{-i\omega_{N}}\left(  -i\frac{\partial}{\partial x_{j}}\right)
e^{i\omega_{N}}=-i\frac{\partial}{\partial x_{j}}+\frac{\partial\omega_{N}%
}{\partial x_{j}}%
\end{equation}
and $1$-particle vector operators
\begin{equation}
e^{-i\omega_{N}}\left(  -i\nabla_{j}\right)  e^{i\omega_{N}}=-i\nabla
_{j}+\nabla_{j}\omega_{N}%
\end{equation}
(again note that $\omega_{N}$ are $N$-particle operators). The effect of the
scaling $\nu_{N}$ on this is
\begin{equation}
\left(  -i\nabla_{j}+\nabla_{j}\omega_{N}\right)  \nu_{N}=-i\nu_{N}\nabla
_{j}-i\left(  \nabla_{j}\nu_{N}\right)  +\nu_{N}\left(  \nabla_{j}\omega
_{N}\right)
\end{equation}
leading to (see Appendix \ref{KinEnTrans}) which is a sum of kinetic, momentum
and potential terms
\begin{align}
&  K\left(  \nu,\omega\right)  =T\left(  \nu,\omega\right)  ^{\dagger}\left(
\mathbf{P\bullet P}\right)  T\left(  \nu,\omega\right)  =\nonumber\\
&  \qquad-\nabla\bullet\left(  \nu^{2}\right)  \mathbb{I}\bullet\nabla
-\nabla\bullet\left(  \nu_{N}^{2}\nabla\omega_{N}-i\nu_{N}\nabla\nu_{N}\right)
\nonumber\\
&  \qquad\qquad-\left(  \nu_{N}^{2}\nabla\omega_{N}-i\nu_{N}\nabla\nu
_{N}\right)  \bullet\nabla+\left(  \nu_{N}\nabla\omega_{N}-i\nabla\nu
_{N}\right)  \bullet\mathbb{I}\bullet\left(  \nu_{N}\nabla\omega_{N}%
-i\nabla\nu_{N}\right)  \label{effkin}%
\end{align}
Thus the effective Hamiltonian is given by%
\begin{equation}
H_{eff}=H_{eff}\left(  \nu,\omega\right)  =K\left(  \nu,\omega\right)
+\left(  \nu_{N}\right)  ^{2}V \label{effham2}%
\end{equation}
where $V=V_{1}+V_{2}$ and the energy of the state $\mathcal{D}^{N}\left(
\nu,\omega\right)  =T\left(  \nu,\omega\right)  \mathcal{D}_{ref}^{N}T\left(
\nu,\omega\right)  ^{\dagger}$ by
\begin{equation}
\mathcal{E}\left(  \mathcal{D}^{N}\left(  \nu,\omega\right)  \right)
=\frac{Tr\left\{  H_{eff}\left(  \nu,\omega\right)  \mathcal{D}_{ref}%
^{N}\right\}  }{Tr\left\{  \mathcal{D}_{ref}^{N}\right\}  } \label{Enpar}%
\end{equation}


The variation over all states is divided into two steps; intra fiber
variations that leaves the FORDO constant, produced by transformations
$T\left(  \nu_{F},\omega_{F}\right)  $ and inter fiber ones that change the
FORDO, which is produced by transformations $T\left(  \nu_{B},\omega
_{B}\right)  $. The total transformation at any stage relative to a fixed
reference state $\mathcal{D}_{ref}^{N}$ is
\begin{equation}
T\left(  \nu_{FB},\omega_{FB}\right)  =T\left(  \nu_{B}+\nu_{F},\omega
_{B}+\omega_{F}\right)  =T\left(  \nu_{F},\omega_{F}\right)  T\left(  \nu
_{B},\omega_{B}\right)
\end{equation}
leading to an energy value $\mathcal{E}\left(  \mathcal{D}^{N}\left(  \nu
_{FB},\omega_{FB}\right)  \right)  $ which can be expressed using the
effective Hamiltonian defined in eqs. $\left(  \ref{effham1}\right)  $ and
$\left(  \ref{effham2}\right)  $ as shown in eq. $\left(  \ref{Enpar}\right)
$. The choice for the fixed "global" reference state $\mathcal{D}_{ref}^{N}$
that leads to the simplest computational schemes is in the pure state case a
cyclic pure IPS\ and in the ensemble case a cyclic state produced by a mixture
of pure IPS's. The expectation values in eq. $\left(  \ref{Enpar}\right)  $
then are of diagonal potentials and momentum and kinetic energy operators
diagonally weighted by these potentials wrt \emph{Independent Particle
States.} In general the potentials describe up to simultaneous $N$-particle
interactions. Diagonal potentials are often called "local" \cite{locpot} as
they describe physical interactions between spin-space points. Starting from
these expressions and partitioning these potentials into pure $0,1,\ldots,N$
parts one can develop approximations to both the intra and inter fiber
variations $\left(  \text{see section \ref{approx}}\right)  $ and thus to the
DMF's and the domains of these functionals leading to approximate solutions
for the energy of $N$-particle states and for FORDO associated with those
states $\left(  \text{see section \ref{Approximations} for an example}\right)
$.

It is important to note that

\begin{description}
\item[$\bullet$] by using space-spin diagonal operators to parameterize
\emph{all }states

\item[$\bullet$] we produce an effective Hamiltonian that contains
\emph{diagonal} multi particle potentials and \emph{diagonally} weighted
kinetic energy and momentum operators

\item[$\bullet$] whose expectation value with respect to a fixed Independent
Particle reference state gives the exact energy of the system.
\end{description}

The emphasis on the conservation of diagonal form of the potentials in the
Hamiltonian is crucial for reducing the complexity of the expressions involved
in the energy expectation value and for the practical implementation of any
numerical computational scheme.

\section{The Domain of the Functionals}

\label{Domain}

The domain, $\operatorname{Dom}\left(  \mathcal{F}\right)  $, of the
functional, $\mathcal{F}$, constructed in the previous section is the cone
$\mathcal{B}_{1+}\left(  \mathcal{H}^{1}\right)  $ of positive one particle
operators i.e. $\mathcal{F}=\mathcal{F}\left(  \mathcal{D}^{1}\right)  $. As
the functional defined in eq. $\left(  \ref{EnFunc}\right)  $ is invariant to
scaling of its argument it can be equivalently defined on $\mathcal{S}_{N}%
^{1}$ the set of FORDO's $\left.  \mathcal{F}\right\vert _{\mathcal{S}_{N}%
^{1}}$ as $\mathcal{F}=\mathcal{F}\left(  D^{1}\right)  $, which is of course
the aim of the construction.

By using the linear isomorphism
\begin{equation}
\left[  KerL_{N}^{1}\right]  ^{c}\cong\mathcal{B}_{1}\left(  \mathcal{H}%
^{N}\right)
\end{equation}
and the resulting cone isomorphism
\begin{equation}
\mathfrak{C}_{1}\cong\mathcal{C}_{N}^{1} \label{coneiso}%
\end{equation}
the functional can be equivalently defined as $\mathcal{F}\left(
\mathfrak{d}_{1}\right)  $ for $\mathfrak{d}_{1}\in\mathfrak{C}_{1}$. We can
extend this cone isomorphism by using the maps $\tilde{\Gamma}\in\mathfrak{P}$
introduced in subsection \ref{DirectSum} that have the property%
\begin{equation}
\tilde{\Gamma}:\mathfrak{C}_{1}\rightarrow\mathcal{V}_{\Gamma+}\subset
\mathcal{V}_{\Gamma}%
\end{equation}
(such $\tilde{\Gamma}$ are not in general unique) and equivalently define the
functional as $\mathcal{F}\left(  \mathcal{D}_{ref}^{N}\right)  =\mathcal{F}%
\left(  \tilde{\Gamma}\left(  \mathfrak{d}_{1}\right)  \right)  $, where
$\tilde{\Gamma}$ is associated with the choice of reference $\mathcal{D}%
_{ref}^{N}$. Finally using the convex set isomorphisms (see eq. $\left(
\ref{NT}\right)  $ for the definition of $\mathcal{N}_{\Gamma}$)
\begin{equation}
\mathcal{N}_{\Gamma}\cong\mathcal{S}_{N}^{1}%
\end{equation}
the functional can be defined on the normalized reference states as
$\mathcal{F}\left(  D_{ref}^{N}\right)  =\mathcal{F}\left(  \frac
{\mathcal{D}_{ref}^{N}}{Tr\mathcal{D}_{ref}^{N}}\right)  $.

\subsection{Inter Equivalence Class Variation (Variation over the Domain of
the Functionals)\label{Inter}}

In a analogous fashion to the intra equivalence class variation discussed in
section \ref{Functionals} we can construct irreducible representations,
$\left\{  \mathcal{R}_{B_{\tilde{\Gamma}\left(  \mathfrak{d}\right)  }}\left(
\mathfrak{A}_{\ker}\right)  \right\}  $ - ($B_{\tilde{\Gamma}\left(
\mathfrak{d}\right)  }$ stands for representations generated by the
representation $\mathcal{V}_{\tilde{\Gamma}}$ of the base and the element
$\mathfrak{d}\in\ker L_{N}^{1}$) of $\mathfrak{A}_{\ker}$, that are all
equivalent and isomorphic to $\mathcal{R}_{\ker L_{N}^{1}}\left(
\mathfrak{A}_{\ker}\right)  $, constituted by super operators $\hat{T}$
$\in\mathcal{R}_{\mathcal{B}_{1}\left(  \mathcal{H}^{N}\right)  }\left(
\mathfrak{A}_{\ker}\right)  $ that have the property (see Appendix
\ref{uker})
\begin{equation}
\hat{T}\left(  \mathcal{D}_{ref}^{N}\right)  =\hat{T}\left(  \mathfrak{d}%
+\tilde{\Gamma}\left(  \mathfrak{d}_{1}\right)  \right)  =\mathfrak{d}+\hat
{T}\left(  \tilde{\Gamma}\left(  \mathfrak{d}_{1}\right)  \right)  \label{SB}%
\end{equation}
where $\hat{T}\left(  \tilde{\Gamma}\left(  \mathfrak{d}_{1}\right)  \right)
\in\mathcal{V}_{\tilde{\Gamma}}$.

We can generate all of $\mathcal{V}_{\Gamma+}$ by the action of $\mathcal{R}%
_{B_{\mathfrak{d},\Gamma}}\left(  \mathfrak{A}_{\ker}\right)  $ on a fixed
cyclic reference state $\mathcal{D}_{ref}^{N}\in\mathcal{V}_{\Gamma+}$ by
super operators belonging to $\mathcal{R}_{B_{\mathfrak{d},\Gamma}}\left(
\mathfrak{A}_{\ker}\right)  $ expressed in generalized "bra-ket" form as
\begin{align}
\hat{T}  &  =%
%TCIMACRO{\dint \limits_{\mathsf{z}^{N}\,\overset{2}{\neq}\,\mathsf{z}^{\prime
%N}}}%
%BeginExpansion
{\displaystyle\int\limits_{\mathsf{z}^{N}\,\overset{2}{\neq}\,\mathsf{z}%
^{\prime N}}}
%EndExpansion
T\left(  \mathsf{z}^{N}\right)  T\left(  \mathsf{z}^{\prime N}\right)  ^{\ast
}\left\vert \left\vert \mathsf{z}^{N}\right\rangle \left\langle \mathsf{z}%
^{\prime N}\right\vert \right)  \left(  \left\vert \mathsf{z}^{N}\right\rangle
\left\langle \mathsf{z}^{\prime N}\right\vert \right\vert d\mathsf{z}%
^{N}d\mathsf{z}^{\prime N}\nonumber\\
&  \qquad\qquad+\left\vert \mathfrak{d}\right)  \left(  \mathfrak{d}%
^{d}\right\vert +%
%TCIMACRO{\dsum \limits_{2\leq\kappa,\lambda\leq\infty}}%
%BeginExpansion
{\displaystyle\sum\limits_{2\leq\kappa,\lambda\leq\infty}}
%EndExpansion
\hat{T}_{F\kappa\lambda}\left\vert k_{\kappa}\right)  \left(  k_{\lambda}%
^{d}\right\vert \nonumber\\
&  \qquad\qquad\qquad\qquad\qquad+%
%TCIMACRO{\dsum \limits_{1\leq\kappa,\lambda\leq\infty}}%
%BeginExpansion
{\displaystyle\sum\limits_{1\leq\kappa,\lambda\leq\infty}}
%EndExpansion
\hat{T}_{FB\kappa\lambda}\left\vert k_{\kappa}\right)  \left(  \mathfrak{x}%
_{\lambda}^{d}\right\vert +%
%TCIMACRO{\dsum \limits_{1\leq\kappa,\lambda\leq\infty}}%
%BeginExpansion
{\displaystyle\sum\limits_{1\leq\kappa,\lambda\leq\infty}}
%EndExpansion
\hat{T}_{B\kappa\lambda}\left\vert \mathfrak{x}_{\kappa}\right)  \left(
\mathfrak{x}_{\lambda}^{d}\right\vert
\end{align}
where $\left\vert \mathfrak{d}\right)  $ is chosen to be an element of the
basis of $\ker L_{N}^{1}$. The restrictions on $T\left(  \mathsf{z}%
^{N}\right)  $ and thus on $\hat{T}$ are produced by%
\begin{align}
\left(  \left\vert \mathsf{z}^{N}\right\rangle \left\langle \mathsf{z}^{\prime
N}\right\vert \left\vert \hat{T}\right.  \tilde{\Gamma}\left(  \mathfrak{d}%
_{1}\right)  \right)   &  =0;\,\mathsf{z}^{N}\overset{2}{\nsim}\mathsf{z}%
^{\prime N}\nonumber\\
\left(  \mathfrak{x}_{\lambda}^{d}\left\vert \hat{T}\right.  k_{\kappa
}\right)   &  =0;\,2\leq\lambda\leq\infty,1\leq\kappa\leq\infty\nonumber\\
\left(  \tilde{\Gamma}\left(  \mathfrak{d}_{1}\right)  ^{d}\left\vert \hat
{T}\right.  k_{\kappa}\right)   &  =0;\,1\leq\kappa\leq\infty\nonumber\\
\left(  \mathfrak{x}_{\lambda}^{d}\left\vert \hat{T}\right.  \tilde{\Gamma
}\left(  \mathfrak{d}_{1}\right)  \right)   &  =0;\,2\leq\lambda\leq
\infty\nonumber\\
\left(  \tilde{\Gamma}\left(  \mathfrak{d}_{1}\right)  ^{d}\left\vert \hat
{T}\right.  \mathfrak{x}_{\kappa}\right)   &  =0;\,2\leq\kappa\leq
\infty\label{InterFib}%
\end{align}
A nonredundent parameterization is provided by transformations of the form
\begin{align}
\hat{T}  &  =%
%TCIMACRO{\dint \limits_{\mathsf{z}^{N}\,\overset{2}{\neq}\,\mathsf{z}^{\prime
%N}}}%
%BeginExpansion
{\displaystyle\int\limits_{\mathsf{z}^{N}\,\overset{2}{\neq}\,\mathsf{z}%
^{\prime N}}}
%EndExpansion
T\left(  \mathsf{z}^{N}\right)  T\left(  \mathsf{z}^{\prime N}\right)  ^{\ast
}\left\vert \left\vert \mathsf{z}^{N}\right\rangle \left\langle \mathsf{z}%
^{\prime N}\right\vert \right)  \left(  \left\vert \mathsf{z}^{N}\right\rangle
\left\langle \mathsf{z}^{\prime N}\right\vert \right\vert d\mathsf{z}%
^{N}d\mathsf{z}^{\prime N}\nonumber\\
&  \qquad\qquad+\left\vert \mathfrak{d}\right)  \left(  \mathfrak{d}%
^{d}\right\vert +%
%TCIMACRO{\dsum \limits_{2\leq\kappa,\lambda\leq\infty}}%
%BeginExpansion
{\displaystyle\sum\limits_{2\leq\kappa,\lambda\leq\infty}}
%EndExpansion
\hat{T}_{F\kappa\lambda}\left\vert k_{\kappa}\right)  \left(  k_{\lambda}%
^{d}\right\vert \nonumber\\
&  \qquad\qquad\qquad\qquad\qquad+%
%TCIMACRO{\dsum \limits_{1\leq\kappa,\lambda\leq\infty}}%
%BeginExpansion
{\displaystyle\sum\limits_{1\leq\kappa,\lambda\leq\infty}}
%EndExpansion
\hat{T}_{FB\kappa\lambda}\left\vert k_{\kappa}\right)  \left(  \mathfrak{x}%
_{\lambda}^{d}\right\vert +%
%TCIMACRO{\dsum \limits_{1\leq\kappa,\lambda\leq\infty}}%
%BeginExpansion
{\displaystyle\sum\limits_{1\leq\kappa,\lambda\leq\infty}}
%EndExpansion
\hat{T}_{B\kappa\lambda}\left\vert \mathfrak{x}_{\kappa}\right)  \left(
\mathfrak{x}_{\lambda}^{d}\right\vert
\end{align}
which form the cosets $\mathfrak{C}_{B_{\tilde{\Gamma}\left(  \mathfrak{d}%
_{1}\right)  }}$.

In order to distinguish the use of a particular state as a fiber or base
reference state we will introduce the notation $\mathcal{D}_{0_{F}}^{N}$ and
$\mathcal{D}_{0_{B}}^{N}$. In a similar manner to the fiber case, when the
reference state $\mathcal{D}_{0_{B}}^{N}$ is pure, one generates all pure
states, $ext\left\{  \mathcal{V}_{\Gamma+}\right\}  $, contained in
$\mathcal{V}_{\Gamma+}$.

Any element of the cone $\mathcal{B}_{1+}\left(  \mathcal{H}^{N}\right)  $ can
thus be produced non redundantly in a two step process; starting from a base
cyclic reference element (a global reference state) $\mathcal{D}_{0_{B}}%
^{N}=\mathfrak{d}_{0_{B}\Gamma}+\tilde{\Gamma}\left(  \mathfrak{d}_{0_{B}%
1}\right)  $ one generates a fiber cyclic element (a local reference state)
$\hat{T}_{B}\left(  \mathcal{D}_{0_{B}}^{N}\right)  \in\left\{  \hat{T}%
_{B}\left(  \tilde{\Gamma}\left(  \mathfrak{d}_{0_{B}1}\right)  \right)
\right\}  $, where $\hat{T}_{B}\in\mathfrak{C}_{\mathcal{D}_{0_{B}}^{N}%
}\subset\mathcal{R}_{B_{\mathfrak{d},\tilde{\Gamma}}}\left(  \mathfrak{A}%
_{\ker}\right)  $, from which one can generate all of $\left\{  \hat{T}%
_{B}\left(  \tilde{\Gamma}\left(  \mathfrak{d}_{0_{B}1}\right)  \right)
\right\}  $ by the action of $\hat{T}_{F}\in\mathfrak{C}_{\hat{T}_{B}\left(
\mathcal{D}_{0_{B}}^{N}\right)  }\subset\mathcal{R}_{F_{\tilde{\Gamma}\left(
\mathfrak{d}_{0_{B}1}\right)  }}\left(  \mathfrak{A}_{\ker}\right)  $, i.e.
the final result is $\hat{T}_{F}\hat{T}_{B}\left(  \mathcal{D}_{0_{B}}%
^{N}\right)  $. The elements $\hat{T}_{F}$ and $\hat{T}_{B}$ are parameterized
by the integral kernels $\left\{  \nu_{F_{\tilde{\Gamma}\left(  \mathfrak{d}%
_{0_{B}1}\right)  }}\left(  \mathsf{z}^{N}\right)  ,\omega_{F_{\tilde{\Gamma
}\left(  \mathfrak{d}_{0_{B}1}\right)  }}\left(  \mathsf{z}^{N}\right)
\right\}  $ and $\left\{  \nu_{B_{\mathfrak{d},\tilde{\Gamma}}}\left(
\mathsf{z}^{N}\right)  ,\omega_{B_{\mathfrak{d},\tilde{\Gamma}}}\left(
\mathsf{z}^{N}\right)  \right\}  $ respectively which are characterized in
Appendix \ref{uker}

To summarize any element of $\mathcal{B}_{1+}\left(  \mathcal{H}^{N}\right)  $
can be nonredundently parameterized by integral kernels of super operators
belonging to the abelian semigroup cosets discussed in the previous sections
resulting in
\begin{align}
\mathcal{D}^{N}\left(  \nu_{B},\omega_{B},\nu_{F},\omega_{F}\right)   &
=\left[  \hat{T}\left(  \nu_{F},\omega_{F}\right)  \hat{T}\left(  \nu
_{B},\omega_{B}\right)  \right]  \left(  \mathcal{D}_{0_{B}}^{N}\right)
\nonumber\\
&  =T\left(  \nu_{F},\omega_{F}\right)  T\left(  \nu_{B},\omega_{B}\right)
\mathcal{D}_{0_{B}}^{N}T\left(  \nu_{B},\omega_{B}\right)  ^{\dagger}T\left(
\nu_{F},\omega_{F}\right)  ^{\dagger}%
\end{align}
where $\hat{T}\left(  \nu_{B},\omega_{B}\right)  \in\mathfrak{C}%
_{\mathcal{D}_{0_{B}}^{N}}$ and $\hat{T}\left(  \nu_{F},\omega_{F}\right)
\in\mathfrak{C}_{\hat{T}_{B}\left(  \mathcal{D}_{0_{B}}^{N}\right)  }$. This
parameterization of $N$-particle states explicitly partitions variables into
the sets $\left\{  \nu_{F},\omega_{F}\right\}  $ that preserve FORDOs and the
sets $\left\{  \nu_{B},\omega_{B}\right\}  $ that index $N$-particle
states$\ $that have different FORDOs. One can express the functionals defined
in section \ref{Functionals} as functionals of these integral kernels as
\begin{equation}
\mathcal{F}\left(  \nu_{B},\omega_{B}\right)  =\underset{\hat{T}\left(
\nu_{F},\omega_{F}\right)  \in\mathfrak{C}_{\hat{T}_{B}\left(  \mathcal{D}%
_{0_{B}}^{N}\right)  }}{Min}\mathcal{E}\left(  \mathcal{D}^{N}\left(  \nu
_{B},\omega_{B},\nu_{F},\omega_{F}\right)  \right)
\end{equation}
The integral kernels can themselves be expanded in terms of "basis" sets and
the variation reexpressed in terms of the coefficients of those expansions.

\section{Standard Functionals\label{Standard}}

One can associate contributions to the value of energy functional from various
parts of the Hamiltonian and obtain a standard decomposition in terms of
specific functionals. The Hamiltonian, $H$, for a system of $N$-electrons and
$M$-nuclei at fixed positions $\left\{  \mathbf{R}_{j}|1\leq j\leq M\right\}
$ expressed in atomic units is
\begin{equation}
H=\sum_{1\leq j<k\leq M}\frac{Z_{j}Z_{k}}{\left\Vert \mathbf{R}_{j}%
-\mathbf{R}_{k}\right\Vert }-\sum_{\substack{1\leq j\leq M\\1\leq k\leq
N}}\frac{Z_{j}}{\left\Vert \mathbf{R}_{j}-\mathbf{r}_{k}\right\Vert }%
-\sum_{1\leq j\leq N}^{N}\left\Vert \mathbf{p}_{j}\right\Vert ^{2}+\sum_{1\leq
j<k\leq N}\frac{1}{\left\Vert \mathbf{r}_{j}-\mathbf{r}_{k}\right\Vert
}+V_{ext} \label{hamilt}%
\end{equation}
where for future reference we denote the one and two particle potentials
$V_{1}$ and $V_{2}$ as
\begin{align}
V_{1}  &  =-\sum_{\substack{1\leq j\leq M\\1\leq k\leq N}}\frac{Z_{j}%
}{\left\Vert \mathbf{R}_{j}-\mathbf{r}_{k}\right\Vert }+V_{ext}\nonumber\\
V_{2}  &  =\sum_{1\leq j<k\leq N}\frac{1}{\left\Vert \mathbf{r}_{j}%
-\mathbf{r}_{k}\right\Vert }%
\end{align}
The Hamiltonian eq. $\left(  \ref{hamilt}\right)  $ can be expressed using the
obvious abbreviations as
\begin{equation}
V\equiv V_{n}+V_{en}+K_{e}+V_{ee}+V_{ext}%
\end{equation}
where nuclei $j$ has atomic number $Z_{j}$ and is at the position
$\mathbf{R}_{j}$, $\mathbf{r}_{j}$ is the position of electron $j$ that has
momentum $\mathbf{p}_{j}=\left(  \nabla_{\mathbf{r}_{j}}-\mathbf{A}%
_{ext\,j}\right)  ,$ where $\mathbf{A}_{ext\,j}$ is the sum of all external
vector potentials acting on particle $j$ and $V_{ext}$ is the sum of all
external scalar potentials. The use of the general definition of momentum
containing the vector potential is essential for the approximations introduced
in section \ref{approx}. The total energy functional is defined by
\begin{align}
\mathcal{F}_{tot}\left(  \nu_{B},\omega_{B}\right)   &  =\mathcal{F}%
_{tot}\left(  \nu_{B},\omega_{B},\nu_{\#F}\left(  \nu_{B},\omega_{B}\right)
,\omega_{\#F}\left(  \nu_{B},\omega_{B}\right)  \right) \nonumber\\
&  =\underset{\hat{T}\left(  \nu_{F},\omega_{F}\right)  \in\mathfrak{C}%
_{\hat{T}_{B}\left(  \mathcal{D}_{0_{B}}^{N}\right)  }}{Min}\mathcal{E}\left(
\mathcal{D}^{N}\left(  \nu_{B},\omega_{B},\nu_{F},\omega_{F}\right)  \right)
\nonumber\\
&  =\underset{\hat{T}\left(  \nu_{F},\omega_{F}\right)  \in\mathfrak{C}%
_{\hat{T}_{B}\left(  \mathcal{D}_{0_{B}}^{N}\right)  }}{Min}\frac{Tr\left\{
H\mathcal{D}^{N}\left(  \nu_{B},\omega_{B},\nu_{F},\omega_{F}\right)
\right\}  }{Tr\left\{  \mathcal{D}^{N}\left(  \nu_{B},\omega_{B},\nu
_{F},\omega_{F}\right)  \right\}  }%
\end{align}
and the individual functionals by
\begin{subequations}
\begin{align}
\mathcal{F}_{K_{e}}\left(  \nu_{B},\omega_{B}\right)   &  =\frac{Tr\left\{
K_{e}\mathcal{D}^{N}\left(  \nu_{B},\omega_{B},\nu_{\#F}\left(  \nu_{B}%
,\omega_{B}\right)  ,\omega_{\#F}\left(  \nu_{B},\omega_{B}\right)  \right)
\right\}  }{Tr\left\{  \mathcal{D}^{N}\left(  \nu_{B},\omega_{B},\nu
_{\#F}\left(  \nu_{B},\omega_{B}\right)  ,\omega_{\#F}\left(  \nu_{B}%
,\omega_{B}\right)  \right)  \right\}  }\\
\mathcal{F}_{V_{e}}\left(  \nu_{B},\omega_{B}\right)   &  =\frac{Tr\left\{
V_{e}\mathcal{D}^{N}\left(  \nu_{B},\omega_{B},\nu_{\#F}\left(  \nu_{B}%
,\omega_{B}\right)  ,\omega_{\#F}\left(  \nu_{B},\omega_{B}\right)  \right)
\right\}  }{Tr\left\{  \mathcal{D}^{N}\left(  \nu_{B},\omega_{B},\nu
_{\#F}\left(  \nu_{B},\omega_{B}\right)  ,\omega_{\#F}\left(  \nu_{B}%
,\omega_{B}\right)  \right)  \right\}  }\\
\mathcal{F}_{V_{en}}\left(  \nu_{B},\omega_{B}\right)   &  =\frac{Tr\left\{
V_{en}\mathcal{D}^{N}\left(  \nu_{B},\omega_{B},\nu_{\#F}\left(  \nu
_{B},\omega_{B}\right)  ,\omega_{\#F}\left(  \nu_{B},\omega_{B}\right)
\right)  \right\}  }{Tr\left\{  \mathcal{D}^{N}\left(  \nu_{B},\omega_{B}%
,\nu_{\#F}\left(  \nu_{B},\omega_{B}\right)  ,\omega_{\#F}\left(  \nu
_{B},\omega_{B}\right)  \right)  \right\}  }\\
\mathcal{F}_{V_{n}}\left(  \nu_{B},\omega_{B}\right)   &  =\frac{Tr\left\{
V_{n}\mathcal{D}^{N}\left(  \nu_{B},\omega_{B},\nu_{\#F}\left(  \nu_{B}%
,\omega_{B}\right)  ,\omega_{\#F}\left(  \nu_{B},\omega_{B}\right)  \right)
\right\}  }{Tr\left\{  \mathcal{D}^{N}\left(  \nu_{B},\omega_{B},\nu
_{\#F}\left(  \nu_{B},\omega_{B}\right)  ,\omega_{\#F}\left(  \nu_{B}%
,\omega_{B}\right)  \right)  \right\}  }\\
\mathcal{F}_{ext}\left(  \nu_{B},\omega_{B}\right)   &  =\frac{Tr\left\{
V_{ext}\mathcal{D}^{N}\left(  \nu_{B},\omega_{B},\nu_{\#F}\left(  \nu
_{B},\omega_{B}\right)  ,\omega_{\#F}\left(  \nu_{B},\omega_{B}\right)
\right)  \right\}  }{Tr\left\{  \mathcal{D}^{N}\left(  \nu_{B},\omega_{B}%
,\nu_{\#F}\left(  \nu_{B},\omega_{B}\right)  ,\omega_{\#F}\left(  \nu
_{B},\omega_{B}\right)  \right)  \right\}  }%
\end{align}
The electron-electron interaction functional $\mathcal{F}_{V_{e}}\left(
\nu_{B},\omega_{B}\right)  $ is conventionally decomposed into two parts one
describing the classical coulomb interaction between electrons and the other
exchange-correlation effects as
\end{subequations}
\begin{equation}
\mathcal{F}_{XC}\left(  \nu_{B},\omega_{B}\right)  =\mathcal{F}_{V_{e}}\left(
\nu_{B},\omega_{B}\right)  -\mathcal{F}_{coul}\left(  \nu_{B},\omega
_{B}\right)
\end{equation}
where the classical coulomb interaction energy is given by
\begin{align}
&  \mathcal{F}_{coul}\left(  \nu_{B},\omega_{B}\right)  =\frac{Tr\left\{
V_{coul}\mathcal{D}^{N}\left(  \nu_{B},\omega_{B},\nu_{\#F}\left(  \nu
_{B},\omega_{B}\right)  ,\omega_{\#F}\left(  \nu_{B},\omega_{B}\right)
\right)  \right\}  }{Tr\left\{  \mathcal{D}^{N}\left(  \nu_{B},\omega_{B}%
,\nu_{\#F}\left(  \nu_{B},\omega_{B}\right)  ,\omega_{\#F}\left(  \nu
_{B},\omega_{B}\right)  \right)  \right\}  }\nonumber\\
&  \left.  {}\right.  \qquad\qquad\qquad\qquad\qquad\qquad\qquad=\int
\frac{\gamma\left[  \nu_{B},\omega_{B}\right]  \left(  \mathsf{z}_{1}\right)
\gamma\left(  \mathsf{z}_{2}\right)  }{\left\Vert \mathbf{r}_{1}%
-\mathbf{r}_{2}\right\Vert }d\mathsf{z}_{1}d\mathsf{z}_{2}%
\end{align}
and the space-spin density $\gamma\left[  \nu_{B},\omega_{B}\right]  \left(
\mathsf{z}_{1}\right)  $ by%
\begin{equation}
\gamma\left[  \nu_{B},\omega_{B}\right]  \left(  \mathsf{z}_{1}\right)
=N\frac{\int\mathcal{D}^{N}\left[  \nu_{B},\omega_{B},\nu_{\#F}\left(  \nu
_{B},\omega_{B}\right)  ,\omega_{\#F}\left(  \nu_{B},\omega_{B}\right)
\right]  \left(  \mathsf{z}_{1}\mathsf{z}^{N-1},\mathsf{z}_{1}\mathsf{z}%
^{N-1}\right)  d\mathsf{z}^{N-1}}{\int\mathcal{D}^{N}\left[  \nu_{B}%
,\omega_{B},\nu_{\#F}\left(  \nu_{B},\omega_{B}\right)  ,\omega_{\#F}\left(
\nu_{B},\omega_{B}\right)  \right]  \left(  \mathsf{z}^{N},\mathsf{z}%
^{N}\right)  d\mathsf{z}^{N}}%
\end{equation}
Note that space-spin density $\gamma\left[  \nu_{B},\omega_{B}\right]  \left(
\mathsf{z}_{1}\right)  $ contains density contributions form the spin
components $\left\{  \rho_{\alpha},\rho_{\alpha\beta},\rho_{\beta}\right\}  $
which are needed when one considers spin effects and symmetry breaking
\cite{WeinerTrickey98}.

\section{Approximations\label{approx}}

The integral kernels of the diagonal operators $T$ are symmetric wrt
permutations of the coordinates $\left\{  \mathsf{z}_{1},\ldots,\mathsf{z}%
_{N}\right\}  $ thus they belong to a symmetric tensor product space, $F^{N}$,
based on the $1$-particle integral kernel space
\begin{equation}
F^{1}=\operatorname{linspan}\left\{  T\left(  \mathsf{z}_{1}\right)  \right\}
\end{equation}
that can be expressed as
\begin{equation}
F^{N}=%
%TCIMACRO{\dbigvee \nolimits^{N}}%
%BeginExpansion
{\displaystyle\bigvee\nolimits^{N}}
%EndExpansion
F^{1}%
\end{equation}
\cite{Greub,MM}. The symmetric tensor product is defined as
\begin{equation}
f_{1}\vee\cdots\vee f_{N}\left(  \mathsf{z}_{1},\ldots,\mathsf{z}_{N}\right)
=%
%TCIMACRO{\tsum \limits_{\sigma\in\mathcal{S}_{N}}}%
%BeginExpansion
{\textstyle\sum\limits_{\sigma\in\mathcal{S}_{N}}}
%EndExpansion
f_{1}\left(  \mathsf{z}_{\sigma\left(  1\right)  }\right)  \cdots f_{p}\left(
\mathsf{z}_{\sigma\left(  N\right)  }\right)  \equiv%
%TCIMACRO{\tsum \limits_{\sigma\in\mathcal{S}_{N}}}%
%BeginExpansion
{\textstyle\sum\limits_{\sigma\in\mathcal{S}_{N}}}
%EndExpansion
f_{\sigma\left(  1\right)  }\left(  \mathsf{z}_{1}\right)  \cdots
f_{\sigma\left(  N\right)  }\left(  \mathsf{z}_{p}\right)
\end{equation}
where $\mathcal{S}_{N}$ is the $N^{\text{th}}$-order symmetric group. The
space $F^{N}$ can also be expressed as the symmetric tensor product
\begin{equation}
F^{N}=F^{n_{1}}\vee\cdots\vee F^{n_{m}}%
\end{equation}
where $n_{j}\in\mathbb{N}$ for $1\leq j\leq m$ and $n_{1}+$ $\cdots+n_{m}=N$.

An element $T$ of $F^{n}$ is said to be indecomposable if it cannot be
expressed as a simple product of elements belonging to $\left\{  F^{m}|1\leq
m<n\right\}  $ i.e.
\begin{equation}
T\neq T_{n_{1}}\left(  \mathsf{z}_{1}^{n_{1}}\right)  \vee\cdots\vee T_{n_{m}%
}\left(  \mathsf{z}_{m}^{n_{m}}\right)
\end{equation}
where $T_{n_{j}}\in F^{n_{j}},$ $1\leq j\leq m$, $%
%TCIMACRO{\tsum \limits_{1\leq j\leq m}}%
%BeginExpansion
{\textstyle\sum\limits_{1\leq j\leq m}}
%EndExpansion
n_{j}=n$ and $\mathsf{z}_{j}^{n_{j}}=\mathsf{z}_{s_{j}+1}\cdots\mathsf{z}%
_{s_{j}+n_{j}}$ for $s_{j}=%
%TCIMACRO{\tsum \limits_{1\leq k\leq j-1}}%
%BeginExpansion
{\textstyle\sum\limits_{1\leq k\leq j-1}}
%EndExpansion
n_{k}$. By using this factorization property we can classify $F^{N}$ into
classes $\left\{  \mathcal{K}_{n_{1},\ldots,n_{m}}\right\}  $ of
indecomposable integral kernels indexed by $\left\{  n_{1},\ldots
,n_{m}\right\}  $ for $1\leq m\leq N$, $%
%TCIMACRO{\tsum \limits_{1\leq j\leq m}}%
%BeginExpansion
{\textstyle\sum\limits_{1\leq j\leq m}}
%EndExpansion
n_{j}=N$, $0\leq$ $n_{j}\leq N$ for $1\leq j\leq m$ of the form
\begin{equation}
T_{n_{1}}\left(  \mathsf{z}_{1}^{n_{1}}\right)  \vee\cdots\vee T_{n_{m}%
}\left(  \mathsf{z}_{m}^{n_{m}}\right)  \in F^{n_{1}}\vee\cdots\vee F^{n_{m}}
\label{intker}%
\end{equation}
If the space $F^{n_{j}}$ occurs more than once in the product eq. $\left(
\ref{intker}\right)  $ we distinguish kernels from different copies $F^{n_{j}}
$ by a further subscripts i.e. $T_{n_{j}k_{j}}\left(  \mathsf{z}_{j}^{n_{j}%
}\right)  $. Particular examples of such integral kernels are
\begin{equation}
T\left(  \mathsf{z}^{N}\right)  =T_{n_{1}}\left(  \mathsf{z}_{1}^{n_{1}%
}\right)  \vee\delta_{n_{2}}\left(  \mathsf{z}_{m}^{n_{2}}\right)  =T_{n_{1}%
}\left(  \mathsf{z}_{1},\ldots,\mathsf{z}_{n_{1}}\right)  \vee\delta\left(
\mathsf{z}_{n_{1}+1}\right)  \vee\cdots\vee\delta\left(  \mathsf{z}%
_{N}\right)
\end{equation}
where we note that $T\left(  \mathsf{z}^{N}\right)  $ is a product of
indecomposable integral kernels as each $\left\{  \delta\left(  \mathsf{z}%
_{j}\right)  |n_{1}+1\leq j\leq N\right\}  $ are indecomposable even though
$\delta_{n_{2}}\left(  \mathsf{z}_{m}^{n_{2}}\right)  $ is not, the product of
different one particle integral kernels
\begin{equation}
T\left(  \mathsf{z}^{N}\right)  =T_{11}\left(  \mathsf{z}_{1}\right)
\vee\cdots\vee T_{1N}\left(  \mathsf{z}_{N}\right)  \qquad\qquad\qquad
\qquad\qquad\qquad\qquad\qquad\qquad
\end{equation}
and the product of an identical one particle integral kernel%
\begin{equation}
T\left(  \mathsf{z}^{N}\right)  =T_{1}\left(  \mathsf{z}_{1}\right)
\vee\cdots\vee T_{1}\left(  \mathsf{z}_{N}\right)  \qquad\qquad\qquad
\qquad\qquad\qquad\qquad\qquad\qquad.
\end{equation}


The integral kernels in eq. $\left(  \ref{intker}\right)  $ correspond to
$N$-particle diagonal transformation $T_{n_{1}}\wedge\cdots\wedge T_{n_{m}}%
\in\mathcal{B}\left(  \mathcal{H}^{N}\right)  $ induced by diagonal
transformations $T_{n_{1}}\in\mathcal{B}\left(  \mathcal{H}^{n_{1}}\right)  $,
when one considers the antisymmetric tensor product of the spaces $\left\{
\mathcal{H}^{n_{j}}|1\leq j\leq m\right\}  $ that produce the factorization
\begin{equation}
\mathcal{H}^{N}=\mathcal{H}^{n_{1}}\wedge\cdots\wedge\mathcal{H}^{n_{m}},
\end{equation}
by the following relationship
\begin{equation}
T_{n_{1}}\wedge\cdots\wedge T_{n_{m}}\left\vert \Xi_{n_{1}}\cdots\Xi_{n_{m}%
}\right\rangle =\left\vert T_{n_{1}}\left(  \Xi_{n_{1}}\right)  \cdots
T_{n_{m}}\left(  \Xi_{n_{m}}\right)  \right\rangle \quad\forall\Xi_{n_{j}}%
\in\mathcal{H}^{n_{j}},1\leq j\leq m
\end{equation}
i.e.%
\begin{equation}
T_{n_{1}}\wedge\cdots\wedge T_{n_{m}}=\int T_{n_{1}}\left(  \mathsf{z}%
_{1}^{n_{1}}\right)  \vee\cdots\vee T_{n_{m}}\left(  \mathsf{z}_{m}^{n_{m}%
}\right)  \left\vert \mathsf{z}_{1}^{n_{1}}\cdots\mathsf{z}_{m}^{n_{m}%
}\right\rangle \left\langle \mathsf{z}_{1}^{n_{1}}\cdots\mathsf{z}_{m}^{n_{m}%
}\right\vert d\mathsf{z}_{1}^{n_{1}}\cdots d\mathsf{z}_{m}^{n_{m}}.
\end{equation}


We can expand a general element of $N$-particle integral kernel space $F^{N}$
as a sum over products of indecomposable elements as
\begin{equation}
T\left(  \mathsf{z}^{N}\right)  =%
%TCIMACRO{\dsum \limits_{1\leq m\leq N}}%
%BeginExpansion
{\displaystyle\sum\limits_{1\leq m\leq N}}
%EndExpansion%
%TCIMACRO{\dsum \limits_{n_{1},\ldots,n_{m}}}%
%BeginExpansion
{\displaystyle\sum\limits_{n_{1},\ldots,n_{m}}}
%EndExpansion
c_{n_{1}\ldots n_{m}}T_{n_{1}}\left(  \mathsf{z}_{1}^{n_{1}}\right)
\vee\cdots\vee T_{n_{m}}\left(  \mathsf{z}_{m}^{n_{m}}\right)
\label{kerexpan}%
\end{equation}
corresponding to the diagonal operator
\begin{align}
T  &  =\int\sum_{1\leq m\leq N}\nonumber\\
&
%TCIMACRO{\dsum \limits_{\substack{n_{1},\ldots,n_{m} \\n_{1}+\cdots+n_{m}=N
%\\0\leq n_{j}\leq N;1\leq j\leq m}}}%
%BeginExpansion
{\displaystyle\sum\limits_{\substack{n_{1},\ldots,n_{m} \\n_{1}+\cdots+n_{m}=N
\\0\leq n_{j}\leq N;1\leq j\leq m}}}
%EndExpansion
\sum_{\sigma\in\mathcal{S}_{N}}c_{n_{1}\ldots n_{m}}T_{n_{1}}\left(
\mathsf{z}_{\sigma\left(  1\right)  }\cdots\mathsf{z}_{\sigma\left(
n_{1}\right)  }\right)  \cdots T_{n_{m}}\left(  \mathsf{z}_{\sigma\left(
s_{m-1}\right)  }\cdots\mathsf{z}_{\sigma\left(  n_{m}\right)  }\right)
\left\vert \mathsf{z}^{N}\right\rangle \left\langle \mathsf{z}^{N}\right\vert
d\mathsf{z}^{N}.
\end{align}


Physically each indecomposable part describes simultaneous interactions within
subsets of particles. In order to catalogue these different types of
interactions we define the \textit{particle order} of the operator, (which
agrees with the conventional second quantized characterization of an
operator), by
\begin{equation}
\operatorname{order}\left(  T_{n_{1}}\wedge\cdots\wedge T_{n_{m}}\right)
=\sum_{\substack{\text{\textbf{distinct }}n_{j}\\1\leq j\leq m}%
}\operatorname{order}\left(  T_{n_{j}}\right)
\end{equation}
where $"\wedge"$ denotes antisymmetric tensor product of operators
\cite{antiten} and
\begin{equation}
\operatorname{order}\left(  T_{n_{j}}\right)  =\left\{
\begin{array}
[c]{c}%
0\quad\text{if }T_{n_{j}}\left(  \mathsf{z}\right)  =\delta\left(
\mathsf{z}\right) \\
n_{j}\quad\text{otherwise }%
\end{array}
\right.
\end{equation}
(the delta functions and scalar multiples of delta functions produce zero
particle operators), leading to
\begin{equation}
T=%
%TCIMACRO{\dsum \limits_{0\leq j\leq N}}%
%BeginExpansion
{\displaystyle\sum\limits_{0\leq j\leq N}}
%EndExpansion
T_{j}^{N}%
\end{equation}
where $T_{j}^{N}$ is the sum over all $j$-particle parts i.e. parts with
\textit{particle order }equal to $j$.

For example:

\begin{enumerate}
\item A one particle operator acting in $N$-particle space $\mathcal{H}^{N}$
produced by a product of identical one particle integral kernels has an
integral kernel
\begin{equation}
T\left(  \mathsf{z}^{N}\right)  =T_{1}^{N}\left(  \mathsf{z}^{N}\right)
=\left(  N!\right)  T_{1}\left(  \mathsf{z}_{1}\right)  \cdots T_{1}\left(
\mathsf{z}_{N}\right)  ;\operatorname{order}\left(  T_{1}\right)
=\operatorname{order}\left(  T\right)  =1
\end{equation}
and corresponds to a linear transformation of one particle space
$\mathcal{H}^{1}$ through%
\begin{equation}
T_{1}\wedge\cdots\wedge T_{1}\left\vert \varphi_{1}\cdots\varphi
_{N}\right\rangle =\left\vert T_{1}\varphi_{1}\cdots T_{1}\varphi
_{N}\right\rangle
\end{equation}
where $\left\vert \varphi_{j}\right\rangle \in\mathcal{H}^{1}$

\item A one particle operator acting in $N$-particle space $\mathcal{H}^{N}$
whose integral kernel is formed by a single one particle integral kernel
\begin{align}
&  T\left(  \mathsf{z}^{N}\right)  =T_{1}^{N}\left(  \mathsf{z}^{N}\right)
=\sum_{\sigma\in\mathcal{S}_{N}}T_{1}\left(  \mathsf{z}_{\sigma\left(
1\right)  }\right)  \delta\left(  \mathsf{z}_{\sigma\left(  2\right)
}\right)  \cdots\delta\left(  \mathsf{z}_{\sigma\left(  N\right)  }\right)
\nonumber\\
&  \qquad\qquad\qquad\qquad\qquad\qquad\qquad\qquad\qquad\qquad
\operatorname{order}\left(  T_{1}\right)  =\operatorname{order}\left(
T\right)  =1
\end{align}


\item A two particle operator acting in $N$-particle space $\mathcal{H}^{N}$
whose integral kernel is formed by a product of identical two particle
integral kernels (of course there must be an even number of particles in this
case)
\begin{align}
T\left(  \mathsf{z}^{N}\right)   &  =T_{2}^{\frac{N}{2}}\left(  \mathsf{z}%
^{N}\right)  =\sum_{\sigma\in\mathcal{S}_{N}}T_{2}\left(  \mathsf{z}%
_{\sigma\left(  1\right)  }\mathsf{z}_{\sigma\left(  2\right)  }\right)
\cdots T_{2}\left(  \mathsf{z}_{\sigma\left(  N-1\right)  }\mathsf{z}%
_{\sigma\left(  N\right)  }\right) \nonumber\\
&  \qquad\qquad\qquad\qquad\qquad\qquad\qquad\qquad\operatorname{order}\left(
T_{2}\right)  =\operatorname{order}\left(  T\right)  =2
\end{align}
which corresponds to a linear transformation of two particle space
$\mathcal{H}^{2}$ through%
\begin{equation}
T_{2}\wedge\cdots\wedge T_{2}\left\vert g_{1}\cdots g_{\frac{N}{2}%
}\right\rangle =\left\vert T_{2}g_{1}\cdots T_{2}g_{\frac{N}{2}}\right\rangle
\qquad\forall g_{j}\in\mathcal{H}^{2},1\leq j\leq\frac{N}{2}%
\end{equation}
where $\left\vert g_{i}\right\rangle \in\mathcal{H}^{2}$

\item A $p$-particle operator acting in $N$-particle space $\mathcal{H}^{N}$
whose integral kernel is formed in the following manner
\begin{align}
T\left(  \mathsf{z}^{N}\right)   &  =T_{p}^{N}\left(  \mathsf{z}^{N}\right)
=\sum_{\sigma\in\mathcal{S}_{N}}T_{1}\left(  \mathsf{z}_{\sigma\left(
1\right)  }\right)  ^{\alpha_{1}}\cdots T_{p}\left(  \mathsf{z}_{\sigma\left(
N\right)  }\right)  ^{\alpha_{p}}\nonumber\\
&  \operatorname{order}\left(  T_{j}^{\alpha_{j}}\right)  =1,1\leq j\leq
p;\alpha_{1}+\cdots+\alpha_{p}=N;\operatorname{order}\left(  T\right)  =p
\end{align}


\item A particular example of \#4 is a $N$-particle operator acting in
$N$-particle space $\mathcal{H}^{N}$ whose integral kernel is formed by a
product of different one particle integral kernels
\begin{align}
T\left(  \mathsf{z}^{N}\right)   &  =\sum_{\sigma\in\mathcal{S}_{N}}%
T_{1}\left(  \mathsf{z}_{\sigma\left(  1\right)  }\right)  \cdots T_{N}\left(
\mathsf{z}_{\sigma\left(  N\right)  }\right) \nonumber\\
&  \qquad\qquad\qquad\qquad\qquad\operatorname{order}\left(  T_{j}\right)
=1,1\leq j\leq N;\operatorname{order}\left(  T\right)  =N
\end{align}
This transformation corresponds to
\begin{equation}
T_{1}\wedge\cdots\wedge T_{N}\left\vert \varphi_{1}\cdots\varphi
_{N}\right\rangle =\left\vert T_{1}\varphi_{1}\cdots T_{N}\varphi
_{N}\right\rangle \qquad\forall\left\vert \varphi_{j}\right\rangle
\in\mathcal{H}^{N},1\leq j\leq N
\end{equation}
which\emph{\ is not induced }by a linear transformation of one particle space
$\mathcal{H}^{1}$ unless $T_{1}=\cdots=T_{N}$. The integral kernel of this
transformation can be developed using the polar decomposition of $T_{j}\left(
\mathsf{z}\right)  =\nu_{j}\left(  \mathsf{z}\right)  e^{i\omega_{j}\left(
\mathsf{z}\right)  }$ to give
\begin{align}
T\left(  \mathsf{z}^{N}\right)   &  =\sum_{\sigma\in\mathcal{S}_{N}}%
\nu_{\sigma\left(  1\right)  }\left(  \mathsf{z}_{1}\right)  e^{i\omega
_{\sigma\left(  1\right)  }\left(  \mathsf{z}_{1}\right)  }\cdots\nu
_{\sigma\left(  N\right)  }\left(  \mathsf{z}_{N}\right)  e^{i\omega
_{\sigma\left(  N\right)  }\left(  \mathsf{z}_{N}\right)  }\nonumber\\
&  =\sum_{\sigma\in\mathcal{S}_{N}}\nu_{\sigma\left(  1\right)  }\left(
\mathsf{z}_{1}\right)  \cdots\nu_{\sigma\left(  N\right)  }\left(
\mathsf{z}_{N}\right)  e^{i\left\{  \omega_{\sigma\left(  1\right)  }\left(
\mathsf{z}_{1}\right)  +\cdots+\omega_{\sigma\left(  N\right)  }\left(
\mathsf{z}_{N}\right)  \right\}  }%
\end{align}
One should be careful and note that in this expansion $\nu\left(
\mathsf{z}^{N}\right)  \neq\sum_{\sigma\in\mathcal{S}_{N}}\nu_{\sigma\left(
1\right)  }\left(  \mathsf{z}_{1}\right)  \cdots\nu_{\sigma\left(  N\right)
}\left(  \mathsf{z}_{N}\right)  ,$ and $\omega\left(  \mathsf{z}^{N}\right)
\neq\sum\limits_{1\leq j\leq N}\omega_{\sigma\left(  j\right)  }\left(
\mathsf{z}_{j}\right)  $ but each are in fact complicated functions of
$\left\{  \nu_{j},\omega_{j}|1\leq j\leq N\right\}  $. Thus although $T\left(
\mathsf{z}^{N}\right)  =\nu\left(  \mathsf{z}^{N}\right)  e^{i\omega\left(
\mathsf{z}^{N}\right)  }$ and the norm $\left\Vert \left\vert T_{1}\varphi
_{1}\cdots T_{N}\varphi_{N}\right\rangle \right\Vert $ \emph{is} independent
of\emph{ the }N-particle phase $\omega\left(  \mathsf{z}^{N}\right)  $ it is
\emph{not }independent of the individual particle phases $\left\{  \omega
_{j}\left(  \mathsf{z}_{j}\right)  \right\}  $. Another important point to
note from this example is that IPSs \emph{are not all} generated from a
reference IPS\ by one particle diagonal operators i.e. by transformations of
the form $T_{1}\wedge\cdots\wedge T_{1}$ but they \emph{are all }generated by
$N$-particle transformations of the form $T_{1}\wedge T_{2}\wedge\cdots\wedge
T_{N}$
\end{enumerate}

Using this classification scheme of $F^{N}$ one can decompose $T\left(
\mathsf{z}^{N}\right)  $ in terms of indecomposable parts indexed by $\left\{
n_{1},\ldots,n_{m}\right\}  $ e.g.
\begin{equation}
T\left(  \mathsf{z}^{N}\right)  =\sum_{1\leq m\leq N}%
%TCIMACRO{\tsum \limits_{n_{1},\ldots,n_{m}}}%
%BeginExpansion
{\textstyle\sum\limits_{n_{1},\ldots,n_{m}}}
%EndExpansion
T_{n_{1}\ldots n_{m}}\left(  \mathsf{z}^{N}\right)  \label{vwexpan}%
\end{equation}
and formulate approximations to the energy functional $\mathcal{E}\left(
\mathcal{D}^{N}\left(  \nu_{B},\omega_{B},\nu_{F},\omega_{F}\right)  \right)
$ by truncating the expansion eq. $\left(  \ref{vwexpan}\right)  $. For
instance one might choose the parts indexed by $\left\{  1,\cdots,1\right\}  $
which contains integral kernels of the form
\begin{equation}
f\left(  \mathsf{z}^{N}\right)  =f_{1}\left(  \mathsf{z}_{1}\right)
\vee\cdots\vee f_{N}\left(  \mathsf{z}_{N}\right)
\end{equation}
which includes $p$-particle operator integral kernels of the form%
\begin{equation}
f_{1}\left(  \mathsf{z}_{1}\right)  \vee\cdots\vee f_{p}\left(  \mathsf{z}%
_{p}\right)  \vee\delta\left(  \mathsf{z}_{p+1}\right)  \vee\cdots\vee
\delta\left(  \mathsf{z}_{N}\right)  ;f_{j}\neq f_{k},1\leq j<k\leq p;1\leq
p\leq N
\end{equation}
and $1$-particle operators integral kernels of the form
\begin{equation}
f_{1}\left(  \mathsf{z}_{1}\right)  \vee\cdots\vee f_{1}\left(  \mathsf{z}%
_{N}\right)  .
\end{equation}
This level of approximation would produce IP DMFs, i.e. DMFs determined by
IPSs whose domain are restricted to IPSs. The next level of approximation
would include both parts indexed by $\left\{  1,\cdots,1\right\}  $ and all or
a selection from integral kernels indexed by $\left\{  2,1,\cdots,1\right\}
,\left\{  2,2,1,\cdots,1\right\}  ,\ldots,\left\{  2,2,\cdots,2\right\}  $.

In this section we consider approximations generated first by the indices
$\left\{  1,\cdots,1\right\}  $, and then by the indices $\left\{  \left\{
1,\cdots,1\right\}  ,\left\{  2,1,\cdots,1\right\}  ,\left\{  2,2,\cdots
,2\right\}  \right\}  $ for \emph{pure }state functionals, $\mathcal{E}\left(
\mathcal{D}^{N}\left(  \nu,\omega\right)  \right)  $ for $\mathcal{D}%
^{N}\left(  \nu,\omega\right)  \in ext\left\{  \mathcal{B}_{1+}\left(
\mathcal{H}\right)  \right\}  $:

\subsection{Approximations indexed by $\left\{  1,\cdots,1\right\}  $}

The transformations that we consider in this case are
\begin{equation}
T_{1}\wedge\cdots\wedge T_{N}\left\vert \varphi_{1}\cdots\varphi
_{N}\right\rangle =\left\vert T_{1}\varphi_{1}\cdots T_{N}\varphi
_{N}\right\rangle
\end{equation}
That have integral kernels
\begin{subequations}
\begin{align}
T\left(  \mathsf{z}^{N}\right)   &  =T_{1}\left(  \mathsf{z}_{1}\right)
\vee\cdots\vee T_{N}\left(  \mathsf{z}_{N}\right)  =%
%TCIMACRO{\tsum \limits_{\sigma\in\mathcal{S}_{N}}}%
%BeginExpansion
{\textstyle\sum\limits_{\sigma\in\mathcal{S}_{N}}}
%EndExpansion
T_{1}\left(  \mathsf{z}_{\sigma\left(  1\right)  }\right)  \cdots T_{N}\left(
\mathsf{z}_{\sigma\left(  N\right)  }\right) \nonumber\\
&  =%
%TCIMACRO{\tsum \limits_{\sigma\in\mathcal{S}_{N}}}%
%BeginExpansion
{\textstyle\sum\limits_{\sigma\in\mathcal{S}_{N}}}
%EndExpansion
T_{\sigma\left(  1\right)  }\left(  \mathsf{z}_{1}\right)  \cdots
T_{\sigma\left(  N\right)  }\left(  \mathsf{z}_{N}\right) \nonumber\\
&  =\sum_{\sigma\in\mathcal{S}_{N}}\nu_{\sigma\left(  1\right)  }\left(
\mathsf{z}_{1}\right)  \cdots\nu_{\sigma\left(  N\right)  }\left(
\mathsf{z}_{N}\right)  e^{i\left\{  \omega_{\sigma\left(  1\right)  }\left(
\mathsf{z}_{1}\right)  +\cdots+\omega_{\sigma\left(  N\right)  }\left(
\mathsf{z}_{N}\right)  \right\}  }%
\end{align}
where
\end{subequations}
\begin{equation}
\omega_{\sigma\left(  j\right)  }\left(  \mathsf{z}_{j}\right)  =\delta\left(
\mathsf{z}_{1}\right)  \cdots\delta\left(  \mathsf{z}_{j-1}\right)
\omega_{\sigma\left(  j\right)  }\left(  \mathsf{z}_{j}\right)  \delta\left(
\mathsf{z}_{j+1}\right)  \cdots\delta\left(  \mathsf{z}_{N}\right)
\end{equation}
and the kernels $\left\{  \nu_{j}\left(  \mathsf{z}_{j}\right)  |1\leq j\leq
N\right\}  $ could be equal to each other or be delta functions as could the
kernels $\left\{  \omega_{j}\left(  \mathsf{z}_{j}\right)  |1\leq j\leq
N\right\}  $.

The first term we need to evaluate is the normalization
\begin{align}
&  Tr\left\{  \mathcal{D}^{N}\left(  \nu,\omega\right)  \right\}
=\left\langle T\left(  \nu,\omega\right)  \varphi_{1}\cdots\varphi
_{N}\right\vert \left.  T\left(  \nu,\omega\right)  \varphi_{1}\cdots
\varphi_{N}\right\rangle \nonumber\\
&  =\sum_{\sigma,\sigma^{\prime}\in\mathcal{S}_{N}}\left\langle T_{\sigma
\left(  1\right)  }\cdots T_{\sigma\left(  N\right)  }\varphi_{1}\cdots
\varphi_{N}\right.  \left\vert T_{\sigma^{\prime}\left(  1\right)  }\cdots
T_{\sigma^{\prime}\left(  N\right)  }\varphi_{1}\cdots\varphi_{N}\right\rangle
\nonumber\\
&  =\sum_{\sigma,\sigma^{\prime}\in\mathcal{S}_{N}}\left\langle T_{\sigma
\left(  1\right)  }\varphi_{1}\cdots T_{\sigma\left(  N\right)  }\varphi
_{N}\right.  \left\vert T_{\sigma^{\prime}\left(  1\right)  }\varphi_{1}\cdots
T_{\sigma^{\prime}\left(  N\right)  }\varphi_{N}\right\rangle \nonumber\\
&  =\left\langle \varphi_{1}\cdots\varphi_{N}\right\vert \left.
S_{eff}\left(  \nu,\omega\right)  \varphi_{1}\cdots\varphi_{N}\right\rangle
\equiv\mathcal{S}\left(  \nu,\omega\right)
\end{align}
and $\left\vert \varphi_{1}\cdots\varphi_{N}\right\rangle $ is a fixed IPS.
The effect of the multiparticle potential $T_{\sigma\left(  1\right)  }\cdots
T_{\sigma\left(  N\right)  }$ on an IPS\ is to produce a new IPS so that we
obtain the inner product between two different IPSs that gives
\begin{equation}
\mathcal{S}\left(  \nu,\omega\right)  =%
%TCIMACRO{\tsum \limits_{\sigma,\sigma^{\prime}\in\mathcal{S}_{N}}}%
%BeginExpansion
{\textstyle\sum\limits_{\sigma,\sigma^{\prime}\in\mathcal{S}_{N}}}
%EndExpansion
\det\left\{  \mathbf{S}_{\sigma\sigma^{\prime}}\right\}
\end{equation}


where
\begin{align}
\left(  \mathbf{S}_{\sigma\sigma^{\prime}}\right)  _{jk}  &  =\left\langle
\left.  \nu_{\sigma\left(  j\right)  }e^{i\omega_{\sigma\left(  j\right)  }%
}\varphi_{j}\right\vert \nu_{\sigma^{\prime}\left(  k\right)  }e^{i\omega
_{\sigma^{\prime}\left(  k\right)  }}\varphi_{k}\right\rangle =%
%TCIMACRO{\dint }%
%BeginExpansion
{\displaystyle\int}
%EndExpansion
\nu_{\sigma\left(  j\right)  }\left(  \mathsf{z}\right)  e^{-i\omega
_{\sigma\left(  j\right)  }}\varphi_{j}\left(  \mathsf{z}\right)  ^{\ast}%
\nu_{\sigma^{\prime}\left(  k\right)  }\left(  \mathsf{z}\right)
e^{i\omega_{\sigma^{\prime}\left(  k\right)  }}\varphi_{k}\left(
\mathsf{z}\right)  d\mathsf{z}\nonumber\\
&  =%
%TCIMACRO{\dint }%
%BeginExpansion
{\displaystyle\int}
%EndExpansion
\varphi_{j}\left(  \mathsf{z}\right)  ^{\ast}\varphi_{k}\left(  \mathsf{z}%
\right)  e^{i\left(  \omega_{\sigma^{\prime}\left(  k\right)  }\left(
\mathsf{z}\right)  -\omega_{\sigma\left(  j\right)  }\left(  \mathsf{z}%
\right)  \right)  }\nu_{\sigma\left(  j\right)  }\left(  \mathsf{z}\right)
\nu_{\sigma^{\prime}\left(  k\right)  }\left(  \mathsf{z}\right)  d\mathsf{z}%
\end{align}
The second term produced by the effective Hamiltonian is the expectation value
$\mathcal{K}\left(  \nu,\omega\right)  $ wrt the IPS\ $\left\vert \varphi
_{1}\cdots\varphi_{N}\right\rangle $
\begin{equation}
\mathcal{K}\left(  \nu,\omega\right)  =\mathcal{S}\left(  \nu,\omega\right)
^{-1}\left\langle \varphi_{1}\cdots\varphi_{N}\right\vert \left.  K\left(
\nu,\omega\right)  \varphi_{1}\cdots\varphi_{N}\right\rangle
\end{equation}
of the effective kinetic energy operator $K\left(  \nu,\omega\right)  $ given
by
\begin{align}
\mathcal{K}\left(  \nu,\omega\right)   &  =\mathcal{S}\left(  \nu
,\omega\right)  ^{-1}%
%TCIMACRO{\dsum \limits_{\sigma,\sigma^{\prime}\in\mathcal{S}_{N}}}%
%BeginExpansion
{\displaystyle\sum\limits_{\sigma,\sigma^{\prime}\in\mathcal{S}_{N}}}
%EndExpansion
\left\{  \left\langle \nu_{\sigma\left(  1\right)  }e^{i\omega_{\sigma\left(
1\right)  }}\varphi_{1}\cdots e^{i\omega_{\sigma\left(  N\right)  }}%
\nu_{\sigma\left(  N\right)  }\varphi_{N}\right.  \right. \nonumber\\
&  \qquad\qquad\left.  \left(
%TCIMACRO{\dsum \limits_{1\leq j\leq N}}%
%BeginExpansion
{\displaystyle\sum\limits_{1\leq j\leq N}}
%EndExpansion
\nabla_{j}^{2}\right)  \left.  e^{i\omega_{\sigma^{\prime}\left(  1\right)  }%
}\nu_{\sigma^{\prime}\left(  1\right)  }\varphi_{1}\cdots e^{i\omega
_{\sigma^{\prime}\left(  N\right)  }}\nu_{\sigma\prime\left(  N\right)
}\varphi_{N}\right\rangle \right\}
\end{align}
which can be evaluated by utilizing the expressions described in Appendix
\ref{Transamp} that have been developed by \cite{Verbeek1990} to give
\begin{equation}
\mathcal{K}\left(  \nu,\omega\right)  =%
%TCIMACRO{\tsum \limits_{\sigma,\sigma^{\prime}\in\mathcal{S}_{N}}}%
%BeginExpansion
{\textstyle\sum\limits_{\sigma,\sigma^{\prime}\in\mathcal{S}_{N}}}
%EndExpansion
\operatorname{Tr}\left\{  \mathbf{K}_{\sigma\sigma^{\prime}}\operatorname{adj}%
\left(  \mathbf{S}_{\sigma\sigma^{\prime}}\right)  \right\}
\end{equation}
where
\begin{equation}
\left(  \mathbf{K}_{\sigma\sigma^{\prime}}\right)  _{jk}=\int\left\{
\nabla\varphi_{j}\left(  \mathsf{z}\right)  ^{\ast}\bullet\nabla\varphi
_{k}\left(  \mathsf{z}\right)  \nu_{\sigma\left(  j\right)  }\left(
\mathsf{z}\right)  \nu_{\sigma^{\prime}\left(  k\right)  }\left(
\mathsf{z}\right)  e^{i\left\{  \omega_{\sigma^{\prime}\left(  k\right)
}\left(  \mathsf{z}\right)  -\,\omega_{\sigma\left(  j\right)  }\left(
\mathsf{z}\right)  \right\}  }\right\}  d\mathsf{z}%
\end{equation}


The expectation value of the external potential term
\begin{align}
&  \mathcal{V}_{1}\left(  \nu,\omega\right)  =\nonumber\\
&  \mathcal{S}\left(  \nu,\omega\right)  ^{-1}\left\langle \nu_{\sigma\left(
1\right)  }e^{i\omega_{\sigma\left(  1\right)  }}\varphi_{1}\cdots
e^{i\omega_{\sigma\left(  N\right)  }}\nu_{\sigma\left(  N\right)  }%
\varphi_{N}\left\vert \left(  V_{ext}+V_{en}\right)  \right.  e^{i\omega
_{\sigma^{\prime}\left(  1\right)  }}\nu_{\sigma^{\prime}\left(  1\right)
}\varphi_{1}\cdots e^{i\omega_{\sigma^{\prime}\left(  N\right)  }}\nu
_{\sigma\prime\left(  N\right)  }\varphi_{N}\right\rangle
\end{align}
can be evaluated to give
\begin{equation}
\mathcal{V}_{1}\left(  \nu,\omega\right)  =%
%TCIMACRO{\tsum \limits_{\sigma,\sigma^{\prime}\in\mathcal{S}_{N}}}%
%BeginExpansion
{\textstyle\sum\limits_{\sigma,\sigma^{\prime}\in\mathcal{S}_{N}}}
%EndExpansion
\operatorname{Tr}\left\{  \mathbf{V}_{1\sigma\sigma^{\prime}}%
\operatorname{adj}\left(  \mathbf{S}_{\sigma\sigma^{\prime}}\right)  \right\}
\end{equation}
where $V_{1}\left(  \mathsf{z}^{N}\right)  =V_{ext}\left(  \mathsf{z}%
^{N}\right)  +V_{en}\left(  \mathsf{z}^{N}\right)  =%
%TCIMACRO{\dsum \limits_{1\leq j\leq N}}%
%BeginExpansion
{\displaystyle\sum\limits_{1\leq j\leq N}}
%EndExpansion
\left\{  V_{ext}\left(  \mathsf{z}_{j}\right)  +V_{en}\left(  \mathsf{z}%
_{j}\right)  \right\}  =%
%TCIMACRO{\dsum \limits_{1\leq j\leq N}}%
%BeginExpansion
{\displaystyle\sum\limits_{1\leq j\leq N}}
%EndExpansion
\mathrm{v}\left(  \mathsf{z}_{j}\right)  .$and
\begin{equation}
\left(  \mathbf{V}_{1\sigma\sigma^{\prime}}\right)  _{jk}=\int\left\{
\mathrm{v}\left(  \mathsf{z}\right)  \varphi_{j}\left(  \mathsf{z}\right)
^{\ast}\varphi_{k}\left(  \mathsf{z}\right)  \nu_{\sigma\left(  j\right)
}\left(  \mathsf{z}\right)  \nu_{\sigma^{\prime}\left(  k\right)  }\left(
\mathsf{z}\right)  e^{i\left\{  \omega_{\sigma^{\prime}\left(  k\right)
}\left(  \mathsf{z}\right)  -\,\omega_{\sigma\left(  j\right)  }\left(
\mathsf{z}\right)  \right\}  }\right\}  d\mathsf{z}%
\end{equation}


The expectation value $\mathcal{V}_{2}\left(  \nu\right)  $ of the coulomb
potential term is given by
\begin{equation}
\mathcal{V}_{2}\left(  \nu,\omega\right)  =\mathcal{S}\left(  \nu
,\omega\right)  ^{-1}\left\langle \nu_{\sigma\left(  1\right)  }%
e^{i\omega_{\sigma\left(  1\right)  }}\varphi_{1}\cdots e^{i\omega
_{\sigma\left(  N\right)  }}\nu_{\sigma\left(  N\right)  }\varphi
_{N}\left\vert V_{2}\right.  e^{i\omega_{\sigma^{\prime}\left(  1\right)  }%
}\nu_{\sigma^{\prime}\left(  1\right)  }\varphi_{1}\cdots e^{i\omega
_{\sigma^{\prime}\left(  N\right)  }}\nu_{\sigma\prime\left(  N\right)
}\varphi_{N}\right\rangle
\end{equation}
which can be evaluated using the expressions in Appendix $\ref{Transamp}$ as
\begin{equation}
\mathcal{S}\left(  \nu,\omega\right)  ^{-1}%
%TCIMACRO{\tsum \limits_{\sigma,\sigma^{\prime}\in\mathcal{S}_{N}}}%
%BeginExpansion
{\textstyle\sum\limits_{\sigma,\sigma^{\prime}\in\mathcal{S}_{N}}}
%EndExpansion
Tr\left\{  \mathbf{G}_{\sigma,\sigma^{\prime}}\operatorname{adj}^{\left(
2\right)  }\left(  \mathbf{S}_{\sigma\sigma^{\prime}}\right)  \right\}
\end{equation}
where
\begin{align}
\left(  \mathbf{G}_{\sigma,\sigma^{\prime}}\right)  _{jklm}  &  =%
%TCIMACRO{\dint }%
%BeginExpansion
{\displaystyle\int}
%EndExpansion
\left\{  \nu_{\sigma\left(  j\right)  }\left(  \mathsf{z}_{1}\right)
\varphi_{j}\left(  \mathsf{z}_{1}\right)  ^{\ast}\nu_{\sigma\left(  k\right)
}\left(  \mathsf{z}_{1}\right)  \varphi_{k}\left(  \mathsf{z}_{1}\right)
^{\ast}\nu_{l}\left(  \mathsf{z}_{2}\right)  \varphi_{\sigma^{\prime}\left(
l\right)  }\left(  \mathsf{z}_{2}\right)  \nu_{m}\left(  \mathsf{z}%
_{2}\right)  \varphi_{\sigma^{\prime}\left(  m\right)  }\left(  \mathsf{z}%
_{2}\right)  \right. \nonumber\\
&  \qquad\qquad\qquad\qquad\qquad\times\frac{e^{i\left\{  \omega
_{\sigma^{\prime}\left(  l\right)  }\left(  \mathsf{z}_{2}\right)
+\omega_{\sigma^{\prime}\left(  m\right)  }\left(  \mathsf{z}_{2}\right)
-\,\omega_{\sigma\left(  j\right)  }\left(  \mathsf{z}_{1}\right)
-\omega_{\sigma\left(  k\right)  }\left(  \mathsf{z}_{1}\right)  \right\}  }%
}{\left\Vert \mathbf{r}_{1}-\mathbf{r}_{2}\right\Vert }\nonumber\\
&  -\nu_{\sigma\left(  j\right)  }\left(  \mathsf{z}_{2}\right)  \varphi
_{j}\left(  \mathsf{z}_{2}\right)  ^{\ast}\nu_{\sigma\left(  k\right)
}\left(  \mathsf{z}_{2}\right)  \varphi_{k}\left(  \mathsf{z}_{2}\right)
^{\ast}\nu_{l}\left(  \mathsf{z}_{1}\right)  \varphi_{\sigma^{\prime}\left(
l\right)  }\left(  \mathsf{z}_{1}\right)  \nu_{m}\left(  \mathsf{z}%
_{1}\right)  \varphi_{\sigma^{\prime}\left(  m\right)  }\left(  \mathsf{z}%
_{1}\right) \nonumber\\
&  \qquad\qquad\qquad\qquad\qquad\times\left.  \frac{e^{i\left\{
\omega_{\sigma^{\prime}\left(  l\right)  }\left(  \mathsf{z}_{1}\right)
+\omega_{\sigma^{\prime}\left(  m\right)  }\left(  \mathsf{z}_{1}\right)
-\,\omega_{\sigma\left(  j\right)  }\left(  \mathsf{z}_{2}\right)
-\omega_{\sigma\left(  k\right)  }\left(  \mathsf{z}_{2}\right)  \right\}  }%
}{\left\Vert \mathbf{r}_{1}-\mathbf{r}_{2}\right\Vert }\right\}
d\mathsf{z}_{1}d\mathsf{z}_{2}%
\end{align}
The value of the energy functional $\mathcal{E}\left(  \mathcal{D}^{N}\left(
\nu,\omega\right)  \right)  $ is then the sum
\begin{equation}
\mathcal{E}\left(  \mathcal{D}^{N}\left(  \nu_{1},\ldots,\nu_{N},\omega
_{1},\ldots,\omega_{N}\right)  \right)  =H_{0}+\mathcal{K}\left(  \nu
,\omega\right)  +\mathcal{V}_{1}\left(  \nu,\omega\right)  +\mathcal{V}%
_{2}\left(  \nu,\omega\right)
\end{equation}


In order to determine the value, $\mathcal{F}_{\left\{  1,\cdots,1\right\}
}\left(  \mathcal{D}_{ref}^{1}\right)  $, where $\mathcal{D}_{ref}^{1}%
=L_{N}^{1}\left(  \mathcal{D}_{ref}^{N}\right)  $, of the approximate DMF
$\mathcal{F}_{\left\{  1,\cdots,1\right\}  }$ at the reference IPS
\begin{equation}
\mathcal{D}_{ref}^{N}=\mathcal{D}^{N}\left(  \mathcal{D}_{ref}^{1}\right)
=\frac{\left\vert \varphi_{1}\cdots\varphi_{N}\right\rangle \left\langle
\varphi_{1}\cdots\varphi_{N}\right\vert }{\left\langle \varphi_{1}%
\cdots\varphi_{N}|\varphi_{1}\cdots\varphi_{N}\right\rangle }%
\end{equation}
we minimize $\mathcal{E}\left(  \mathcal{D}^{N}\left(  \nu,\omega\right)
\right)  $ over $\nu=\left(  \nu_{1},\cdots,\nu_{N}\right)  $ and
$\omega=\left(  \omega_{1},\cdots,\omega_{N}\right)  \mathsf{\ }$constraining
$\nu$ and $\omega$ to satisfy the \emph{necessary }generalized strong
orthogonality condition wrt $\left\vert \varphi_{1}\cdots\varphi
_{N}\right\rangle \left\langle \varphi_{1}\cdots\varphi_{N}\right\vert $ i.e.%
\begin{align}
&
%TCIMACRO{\tsum \limits_{\sigma,\sigma^{\prime},\sigma^{\prime\prime}%
%\in\mathcal{S}_{N}}}%
%BeginExpansion
{\textstyle\sum\limits_{\sigma,\sigma^{\prime},\sigma^{\prime\prime}%
\in\mathcal{S}_{N}}}
%EndExpansion%
%TCIMACRO{\dint }%
%BeginExpansion
{\displaystyle\int}
%EndExpansion
\left\{  \nu_{\sigma\left(  1\right)  }\left(  \mathsf{z}_{1}\right)
e^{i\omega_{\sigma\left(  1\right)  }\left(  \mathsf{z}_{1}\right)  }%
\varphi_{1}\left(  \mathsf{z}_{1}\right)  \cdots e^{i\omega_{\sigma\left(
N\right)  }\left(  \mathsf{z}_{N}\right)  }\nu_{\sigma\left(  N\right)
}\left(  \mathsf{z}_{N}\right)  \varphi_{N}\left(  \mathsf{z}_{N}\right)
\times\right. \nonumber\\
&  \nu_{\sigma^{\prime}\left(  1\right)  }\left(  \mathsf{z}_{1}^{\prime
}\right)  e^{-i\omega_{\sigma^{\prime}\left(  1\right)  }\left(
\mathsf{z}_{1}^{\prime}\right)  }\varphi_{1}\left(  \mathsf{z}_{1}^{\prime
}\right)  ^{\ast}\nu_{\sigma^{\prime}\left(  2\right)  }\left(  \mathsf{z}%
_{2}\right)  e^{-i\omega_{\sigma^{\prime}\left(  2\right)  }\left(
\mathsf{z}_{2}\right)  }\varphi_{2}\left(  \mathsf{z}_{2}\right)  ^{\ast
}\cdots e^{-i\omega_{\sigma\left(  N\right)  }\left(  \mathsf{z}_{N}\right)
}\nu_{\sigma\left(  N\right)  }\left(  \mathsf{z}_{N}\right)  \varphi
_{N}\left(  \mathsf{z}_{N}\right)  ^{\ast}\nonumber\\
&  \qquad\qquad\qquad\qquad\left.  -\varphi_{1}\left(  \mathsf{z}_{1}\right)
\cdots\varphi_{N}\left(  \mathsf{z}_{N}\right)  \varphi_{\sigma^{\prime\prime
}\left(  1\right)  }\left(  \mathsf{z}_{1}^{\prime}\right)  ^{\ast}%
\varphi_{\sigma^{\prime\prime}\left(  2\right)  }\left(  \mathsf{z}%
_{2}\right)  ^{\ast}\cdots\varphi_{\sigma^{\prime\prime}\left(  N\right)
}\left(  \mathsf{z}_{N}\right)  ^{\ast}\right\}  d\mathsf{z}^{N-1}%
\overset{a.e}{=}0
\end{align}


\subsection{Approximations indexed by $\left\{  2,\cdots,2\right\}  $}

The transformations that we consider in this case are
\begin{equation}
T_{1}\wedge\cdots\wedge T_{\frac{N}{2}}\left\vert \varphi_{1}\cdots\varphi
_{N}\right\rangle =%
%TCIMACRO{\dsum \limits_{\sigma\in\mathcal{S}_{N}}}%
%BeginExpansion
{\displaystyle\sum\limits_{\sigma\in\mathcal{S}_{N}}}
%EndExpansion
\left\vert T_{1}\left(  \varphi_{\sigma\left(  1\right)  }\varphi
_{\sigma\left(  2\right)  }\right)  \cdots T_{\frac{N}{2}}\left(
\varphi_{\sigma\left(  N-1\right)  }\varphi_{\sigma\left(  N\right)  }\right)
\right\rangle
\end{equation}
where
\begin{equation}
\left\langle \mathsf{z}_{2j-1},\mathsf{z}_{2j}\right.  \left\vert T_{j}\left(
\varphi_{\sigma\left(  2j-1\right)  }\varphi_{\sigma\left(  2j\right)
}\right)  \right\rangle =\int T_{j}\left(  \mathsf{z}_{2j-1},\mathsf{z}%
_{2j}\right)  \varphi_{\sigma\left(  2j-1\right)  }\left(  \mathsf{z}%
_{2j-1}\right)  \varphi_{\sigma\left(  2j\right)  }\left(  \mathsf{z}%
_{2j}\right)  d\mathsf{z}_{2j-1}d\mathsf{z}_{2j}%
\end{equation}
That have integral kernels
\begin{subequations}
\begin{align}
T\left(  \mathsf{z}^{N}\right)   &  =T_{1}\left(  \mathsf{z}_{1}%
,\mathsf{z}_{2}\right)  \vee\cdots\vee T_{\frac{N}{2}}\left(  \mathsf{z}%
_{N-1,}\mathsf{z}_{N}\right)  =%
%TCIMACRO{\tsum \limits_{\sigma\in\mathcal{S}_{N}}}%
%BeginExpansion
{\textstyle\sum\limits_{\sigma\in\mathcal{S}_{N}}}
%EndExpansion
T_{1}\left(  \mathsf{z}_{\sigma\left(  1\right)  },\mathsf{z}_{\sigma\left(
2\right)  }\right)  \cdots T_{\frac{N}{2}}\left(  \mathsf{z}_{\sigma\left(
N-1\right)  ,}\mathsf{z}_{\sigma\left(  N\right)  }\right) \nonumber\\
&  =\sum_{\sigma\in\mathcal{S}_{N}}\nu_{1}\left(  \mathsf{z}_{\sigma\left(
1\right)  },\mathsf{z}_{\sigma\left(  2\right)  }\right)  \cdots\nu_{\frac
{N}{2}}\left(  \mathsf{z}_{\sigma\left(  N-1\right)  },\mathsf{z}%
_{\sigma\left(  2N\right)  }\right)  e^{i\left\{  \omega_{1}\left(
\mathsf{z}_{\sigma\left(  1\right)  },\mathsf{z}_{\sigma\left(  2\right)
}\right)  +\cdots+\omega_{\frac{N}{2}}\left(  \mathsf{z}_{\sigma\left(
N-1\right)  ,}\mathsf{z}_{\sigma\left(  N\right)  }\right)  \right\}  }%
\end{align}
where
\end{subequations}
\begin{equation}
\omega_{j}\left(  \mathsf{z}_{\sigma\left(  2j-1\right)  },\mathsf{z}%
_{\sigma\left(  2j\right)  }\right)  =\delta\left(  \mathsf{z}_{1}\right)
\cdots\omega_{j}\left(  \mathsf{z}_{\sigma\left(  2j-1\right)  }%
,\mathsf{z}_{\sigma\left(  2j\right)  }\right)  \cdots\delta\left(
\mathsf{z}_{N}\right)
\end{equation}


In this case we must treat the effective overlap operator $S\left(  \nu
,\omega\right)  $ as a full $N$-particle operator
\begin{align}
S\left(  \nu,\omega\right)   &  =\int\left\{
%TCIMACRO{\tsum \limits_{\sigma,\sigma^{\prime}\in\mathcal{S}_{N}}}%
%BeginExpansion
{\textstyle\sum\limits_{\sigma,\sigma^{\prime}\in\mathcal{S}_{N}}}
%EndExpansion
T_{1}\left(  \mathsf{z}_{\sigma\left(  1\right)  },\mathsf{z}_{\sigma\left(
2\right)  }\right)  ^{\ast}\cdots T_{\frac{N}{2}}\left(  \mathsf{z}%
_{\sigma\left(  N-1\right)  ,}\mathsf{z}_{\sigma\left(  N\right)  }\right)
^{\ast}\right. \nonumber\\
&  \left.  T_{1}\left(  \mathsf{z}_{\sigma^{\prime}\left(  1\right)
},\mathsf{z}_{\sigma^{\prime}\left(  2\right)  }\right)  \cdots T_{\frac{N}%
{2}}\left(  \mathsf{z}_{\sigma^{\prime}\left(  N-1\right)  ,}\mathsf{z}%
_{\sigma^{\prime}\left(  N\right)  }\right)  \left\vert \mathsf{z}%
^{N}\right\rangle \left\langle \mathsf{z}^{N}\right\vert \right\}
d\mathsf{z}^{N}%
\end{align}
and take the expectation value of this operator wrt the IPS $\left\vert
\varphi_{1}\cdots\varphi_{N}\right\rangle $ producing
\begin{align}
&  \mathcal{S}\left(  \nu,\omega\right)  =\left\langle \varphi_{1}%
\cdots\varphi_{N}|S\left(  \nu,\omega\right)  \varphi_{1}\cdots\varphi
_{N}\right\rangle =\nonumber\\
&
%TCIMACRO{\tsum \limits_{\sigma,\sigma^{\prime},\sigma^{\prime\prime}%
%\in\mathcal{S}_{N}}}%
%BeginExpansion
{\textstyle\sum\limits_{\sigma,\sigma^{\prime},\sigma^{\prime\prime}%
\in\mathcal{S}_{N}}}
%EndExpansion%
%TCIMACRO{\dint }%
%BeginExpansion
{\displaystyle\int}
%EndExpansion
\left\{  \varphi_{_{1}}\left(  \mathsf{z}_{1}\right)  ^{\ast}\cdots\varphi
_{N}\left(  \mathsf{z}_{N}\right)  ^{\ast}\left(  -1\right)  ^{\pi\left(
\sigma^{\prime\prime}\right)  }\varphi_{1}\left(  \mathsf{z}_{\sigma
^{\prime\prime}\left(  1\right)  }\right)  \cdots\varphi_{N}\left(
\mathsf{z}_{\sigma^{\prime\prime}\left(  N\right)  }\right)  \right.
\times\nonumber\\
&  \qquad\qquad\qquad\qquad\qquad\qquad%
%TCIMACRO{\dprod \limits_{1\leq k\leq\frac{N}{2}}}%
%BeginExpansion
{\displaystyle\prod\limits_{1\leq k\leq\frac{N}{2}}}
%EndExpansion
\nu_{k}\left(  \mathsf{z}_{\sigma\left(  2k-1\right)  },\mathsf{z}%
_{\sigma\left(  2k\right)  }\right)  \nu_{k}\left(  \mathsf{z}_{\sigma
^{\prime}\left(  2k-1\right)  },\mathsf{z}_{\sigma^{\prime}\left(  2k\right)
}\right) \nonumber\\
&  \qquad\qquad\left.  e^{i\left\{  \omega_{1}\left(  \mathsf{z}%
_{\sigma^{\prime}\left(  1\right)  },\mathsf{z}_{\sigma^{\prime}\left(
2\right)  }\right)  -\omega_{1}\left(  \mathsf{z}_{\sigma\left(  1\right)
},\mathsf{z}_{\sigma\left(  2\right)  }\right)  +\cdots+\omega_{\frac{N}{2}%
}\left(  \mathsf{z}_{\sigma^{\prime}\left(  N-1\right)  ,}\mathsf{z}%
_{\sigma^{\prime}\left(  N\right)  }\right)  -\omega_{\frac{N}{2}}\left(
\mathsf{z}_{\sigma\left(  N-1\right)  ,}\mathsf{z}_{\sigma\left(  N\right)
}\right)  \right\}  }\right\}  d\mathsf{z}^{N}%
\end{align}
The effective kinetic energy is the expectation value $\mathcal{K}\left(
\nu,\omega\right)  $ wrt the IPS\ $\left\vert \varphi_{1}\cdots\varphi
_{N}\right\rangle $
\begin{equation}
\mathcal{K}\left(  \nu,\omega\right)  =\mathcal{S}\left(  \nu,\omega\right)
^{-1}\left\langle \varphi_{1}\cdots\varphi_{N}\right\vert \left.  K\left(
\nu,\omega\right)  \varphi_{1}\cdots\varphi_{N}\right\rangle
\end{equation}
of the effective kinetic energy operator $K\left(  \nu,\omega\right)  $ given
by%
\begin{align}
K\left(  \nu,\omega\right)   &  =\,\underset{I}{\underbrace{-\nabla
\bullet\left(  \nu^{2}\right)  \mathbb{I}\bullet\nabla}}-\underset
{II}{\underbrace{\nabla\bullet\left(  \nu^{2}\nabla\omega-i\nu\nabla
\nu\right)  }}\nonumber\\
&  -\underset{III}{\underbrace{\left(  \nu^{2}\nabla\omega-i\nu\nabla
\nu\right)  \bullet\nabla}}+\underset{IV}{\underbrace{\left(  \nu\nabla
\omega-i\nabla\nu\right)  \bullet\mathbb{I}\bullet\left(  \nu\nabla
\omega-i\nabla\nu\right)  }} \label{KE2}%
\end{align}
where the $N-$particle integral kernels $\nu_{N},\omega_{N}$ are functions of
the two particle integral kernels $\left\{  \nu_{1},\ldots,\nu_{\frac{N}{2}%
},\omega_{1},\ldots,\omega_{\frac{N}{2}}\right\}  $ i.e.%
\begin{align}
\nu &  =\nu\left(  \nu_{1},\ldots,\nu_{\frac{N}{2}},\omega_{1},\ldots
,\omega_{\frac{N}{2}}\right) \nonumber\\
\omega &  =\omega\left(  \nu_{1},\ldots,\nu_{\frac{N}{2}},\omega_{1}%
,\ldots,\omega_{\frac{N}{2}}\right)
\end{align}
the various parts of the effective kinetic energy operator eq. $\left(
\ref{KE2}\right)  $ can be expanded as

\begin{description}
\item[I]
\begin{align}
&  -\nabla\bullet\left(  \nu^{2}\right)  \mathbb{I}\bullet\nabla\left.
=\right. \nonumber\\
&  -%
%TCIMACRO{\tsum \limits_{\sigma,\sigma^{\prime}\in\mathcal{S}_{N}}}%
%BeginExpansion
{\textstyle\sum\limits_{\sigma,\sigma^{\prime}\in\mathcal{S}_{N}}}
%EndExpansion%
%TCIMACRO{\dsum \limits_{1\leq j\leq N}}%
%BeginExpansion
{\displaystyle\sum\limits_{1\leq j\leq N}}
%EndExpansion
\nabla_{j}\bullet\int\left\{
%TCIMACRO{\dprod \limits_{1\leq k\leq\frac{N}{2}}}%
%BeginExpansion
{\displaystyle\prod\limits_{1\leq k\leq\frac{N}{2}}}
%EndExpansion
\nu_{k}\left(  \mathsf{z}_{\sigma\left(  2k-1\right)  },\mathsf{z}%
_{\sigma\left(  2k\right)  }\right)  \nu_{k}\left(  \mathsf{z}_{\sigma
^{\prime}\left(  2k-1\right)  },\mathsf{z}_{\sigma^{\prime}\left(  2k\right)
}\right)  \right. \nonumber\\
&  \left.  \exp\left\{  i%
%TCIMACRO{\dsum \limits_{1\leq k\leq\frac{N}{2}}}%
%BeginExpansion
{\displaystyle\sum\limits_{1\leq k\leq\frac{N}{2}}}
%EndExpansion
\left\{  \omega_{k}\left(  \mathsf{z}_{\sigma^{\prime}\left(  2k-1\right)
},\mathsf{z}_{\sigma^{\prime}\left(  2k\right)  }\right)  -\omega_{k}\left(
\mathsf{z}_{\sigma\left(  2k-1\right)  },\mathsf{z}_{\sigma\left(  2k\right)
}\right)  \right\}  \right\}  \right\}  \left\vert \mathsf{z}^{N}\right\rangle
\left\langle \mathsf{z}^{N}\right\vert d\mathsf{z}^{N}\mathbb{I}\bullet
\nabla_{j}%
\end{align}


\item[II]
\begin{align}
&  -\nabla\bullet\left(  \nu^{2}\nabla\omega-i\nu\nabla\nu\right)  \left.
=\right. \nonumber\\
&  -%
%TCIMACRO{\tsum \limits_{\sigma,\sigma^{\prime}\in\mathcal{S}_{N}}}%
%BeginExpansion
{\textstyle\sum\limits_{\sigma,\sigma^{\prime}\in\mathcal{S}_{N}}}
%EndExpansion%
%TCIMACRO{\dsum \limits_{1\leq j\leq N}}%
%BeginExpansion
{\displaystyle\sum\limits_{1\leq j\leq N}}
%EndExpansion
\nabla_{j}\bullet\int\left\{  \nabla_{j}\left\{
%TCIMACRO{\dsum \limits_{1\leq k\leq\frac{N}{2}}}%
%BeginExpansion
{\displaystyle\sum\limits_{1\leq k\leq\frac{N}{2}}}
%EndExpansion
\left\{  \omega_{k}\left(  \mathsf{z}_{\sigma^{\prime}\left(  2k-1\right)
},\mathsf{z}_{\sigma^{\prime}\left(  2k\right)  }\right)  \right\}  \right\}
\right.  \times\nonumber\\
&  \qquad\qquad\qquad\qquad%
%TCIMACRO{\dprod \limits_{1\leq k\leq\frac{N}{2}}}%
%BeginExpansion
{\displaystyle\prod\limits_{1\leq k\leq\frac{N}{2}}}
%EndExpansion
\nu_{k}\left(  \mathsf{z}_{\sigma\left(  2k-1\right)  },\mathsf{z}%
_{\sigma\left(  2k\right)  }\right)  \nu_{k}\left(  \mathsf{z}_{\sigma
^{\prime}\left(  2k-1\right)  },\mathsf{z}_{\sigma^{\prime}\left(  2k\right)
}\right) \nonumber\\
&  \left.  -i\nabla_{j}\left\{
%TCIMACRO{\dprod \limits_{1\leq k\leq\frac{N}{2}}}%
%BeginExpansion
{\displaystyle\prod\limits_{1\leq k\leq\frac{N}{2}}}
%EndExpansion
\nu_{k}\left(  \mathsf{z}_{\sigma^{\prime}\left(  2k-1\right)  }%
,\mathsf{z}_{\sigma^{\prime}\left(  2k\right)  }\right)  \right\}
%TCIMACRO{\dprod \limits_{1\leq k\leq\frac{N}{2}}}%
%BeginExpansion
{\displaystyle\prod\limits_{1\leq k\leq\frac{N}{2}}}
%EndExpansion
\nu_{k}\left(  \mathsf{z}_{\sigma\left(  2k-1\right)  },\mathsf{z}%
_{\sigma\left(  2k\right)  }\right)  \right\}  \left\vert \mathsf{z}%
^{N}\right\rangle \left\langle \mathsf{z}^{N}\right\vert d\mathsf{z}^{N}%
\end{align}


\item[III]
\begin{align}
&  -\left(  \nu^{2}\nabla\omega-i\nu\nabla\nu\right)  \bullet\nabla\left.
=\right. \nonumber\\
&  -%
%TCIMACRO{\tsum \limits_{\sigma,\sigma^{\prime}\in\mathcal{S}_{N}}}%
%BeginExpansion
{\textstyle\sum\limits_{\sigma,\sigma^{\prime}\in\mathcal{S}_{N}}}
%EndExpansion%
%TCIMACRO{\dsum \limits_{1\leq j\leq N}}%
%BeginExpansion
{\displaystyle\sum\limits_{1\leq j\leq N}}
%EndExpansion
\int\left\{
%TCIMACRO{\dprod \limits_{1\leq k\leq\frac{N}{2}}}%
%BeginExpansion
{\displaystyle\prod\limits_{1\leq k\leq\frac{N}{2}}}
%EndExpansion
\nu_{k}\left(  \mathsf{z}_{\sigma\left(  2k-1\right)  },\mathsf{z}%
_{\sigma\left(  2k\right)  }\right)  \nu_{j}\left(  \mathsf{z}_{\sigma
^{\prime}\left(  2k-1\right)  },\mathsf{z}_{\sigma^{\prime}\left(  2k\right)
}\right)  \right.  \times\nonumber\\
&  \qquad\qquad\qquad\qquad\nabla_{j}\left\{
%TCIMACRO{\dsum \limits_{1\leq k\leq\frac{N}{2}}}%
%BeginExpansion
{\displaystyle\sum\limits_{1\leq k\leq\frac{N}{2}}}
%EndExpansion
\left\{  \omega_{k}\left(  \mathsf{z}_{\sigma^{\prime}\left(  2k-1\right)
},\mathsf{z}_{\sigma^{\prime}\left(  2k\right)  }\right)  \right\}  \right\}
\nonumber\\
&  \left.  -i%
%TCIMACRO{\dprod \limits_{1\leq k\leq\frac{N}{2}}}%
%BeginExpansion
{\displaystyle\prod\limits_{1\leq k\leq\frac{N}{2}}}
%EndExpansion
\nu_{k}\left(  \mathsf{z}_{\sigma\left(  2k-1\right)  },\mathsf{z}%
_{\sigma\left(  2k\right)  }\right)  \nabla_{j}\left\{
%TCIMACRO{\dprod \limits_{1\leq k\leq\frac{N}{2}}}%
%BeginExpansion
{\displaystyle\prod\limits_{1\leq k\leq\frac{N}{2}}}
%EndExpansion
\nu_{k}\left(  \mathsf{z}_{\sigma^{\prime}\left(  2k-1\right)  }%
,\mathsf{z}_{\sigma^{\prime}\left(  2k\right)  }\right)  \right\}  \right\}
\left\vert \mathsf{z}^{N}\right\rangle \left\langle \mathsf{z}^{\prime
N}\right\vert d\mathsf{z}^{N}\bullet\nabla_{j}%
\end{align}


\item[IV]
\begin{align}
&  \left(  \nu\nabla\omega-i\nabla\nu\right)  \bullet\mathbb{I}\bullet\left(
\nu\nabla\omega-i\nabla\nu\right)  \left.  =\right. \nonumber\\
&
%TCIMACRO{\tsum \limits_{\sigma,\sigma^{\prime}\in\mathcal{S}_{N}}}%
%BeginExpansion
{\textstyle\sum\limits_{\sigma,\sigma^{\prime}\in\mathcal{S}_{N}}}
%EndExpansion%
%TCIMACRO{\dsum \limits_{1\leq j\leq N}}%
%BeginExpansion
{\displaystyle\sum\limits_{1\leq j\leq N}}
%EndExpansion
\int%
%TCIMACRO{\dprod \limits_{1\leq k\leq\frac{N}{2}}}%
%BeginExpansion
{\displaystyle\prod\limits_{1\leq k\leq\frac{N}{2}}}
%EndExpansion
\nu_{k}\left(  \mathsf{z}_{\sigma\left(  2k-1\right)  },\mathsf{z}%
_{\sigma\left(  2k\right)  }\right)  \times\nonumber\\
&  \left\{  \nabla_{j}\left\{
%TCIMACRO{\dsum \limits_{1\leq k\leq\frac{N}{2}}}%
%BeginExpansion
{\displaystyle\sum\limits_{1\leq k\leq\frac{N}{2}}}
%EndExpansion
\left\{  \omega_{k}\left(  \mathsf{z}_{\sigma\left(  2k-1\right)  }%
,\mathsf{z}_{\sigma\left(  2k\right)  }\right)  \right\}  \right\}
-i\nabla_{j}\left\{
%TCIMACRO{\dprod \limits_{1\leq k\leq\frac{N}{2}}}%
%BeginExpansion
{\displaystyle\prod\limits_{1\leq k\leq\frac{N}{2}}}
%EndExpansion
\nu_{k}\left(  \mathsf{z}_{\sigma\left(  2k-1\right)  },\mathsf{z}%
_{\sigma\left(  2k\right)  }\right)  \right\}  \right\}  \left\vert
\mathsf{z}^{N}\right\rangle \left\langle \mathsf{z}^{N}\right\vert
d\mathsf{z}^{N}\nonumber\\
&  \bullet\mathbb{I}\bullet%
%TCIMACRO{\tsum \limits_{\sigma,\sigma^{\prime}\in\mathcal{S}_{N}}}%
%BeginExpansion
{\textstyle\sum\limits_{\sigma,\sigma^{\prime}\in\mathcal{S}_{N}}}
%EndExpansion%
%TCIMACRO{\dsum \limits_{1\leq j\leq N}}%
%BeginExpansion
{\displaystyle\sum\limits_{1\leq j\leq N}}
%EndExpansion
\int\left\{
%TCIMACRO{\dprod \limits_{1\leq k\leq\frac{N}{2}}}%
%BeginExpansion
{\displaystyle\prod\limits_{1\leq k\leq\frac{N}{2}}}
%EndExpansion
\nu_{k}\left(  \mathsf{z}_{\sigma^{\prime}\left(  2k-1\right)  }%
,\mathsf{z}_{\sigma^{\prime}\left(  2k\right)  }\right)  \right.  \times\\
&  \nabla_{j}\left\{
%TCIMACRO{\dsum \limits_{1\leq k\leq\frac{N}{2}}}%
%BeginExpansion
{\displaystyle\sum\limits_{1\leq k\leq\frac{N}{2}}}
%EndExpansion
\left\{  \omega_{k}\left(  \mathsf{z}_{\sigma^{\prime}\left(  2k-1\right)
},\mathsf{z}_{\sigma^{\prime}\left(  2k\right)  }\right)  \right\}  \right\}
\left.  -i\nabla_{j}\left\{
%TCIMACRO{\dprod \limits_{1\leq k\leq\frac{N}{2}}}%
%BeginExpansion
{\displaystyle\prod\limits_{1\leq k\leq\frac{N}{2}}}
%EndExpansion
\nu_{k}\left(  \mathsf{z}_{\sigma^{\prime}\left(  2k-1\right)  }%
,\mathsf{z}_{\sigma^{\prime}\left(  2k\right)  }\right)  \right\}  \right\}
\left\vert \mathsf{z}^{N}\right\rangle \left\langle \mathsf{z}^{N}\right\vert
d\mathsf{z}^{N}%
\end{align}

\end{description}

leading to the expectation values

\begin{description}
\item[I]
\begin{align}
&  \mathcal{K}\left(  \nu,\omega\right)  =\left\langle \varphi_{1}%
\cdots\varphi_{N}|K\left(  \nu,\omega\right)  \varphi_{1}\cdots\varphi
_{N}\right\rangle =%
%TCIMACRO{\tsum \limits_{\sigma,\sigma^{\prime},\sigma^{\prime\prime}%
%\in\mathcal{S}_{N}}}%
%BeginExpansion
{\textstyle\sum\limits_{\sigma,\sigma^{\prime},\sigma^{\prime\prime}%
\in\mathcal{S}_{N}}}
%EndExpansion%
%TCIMACRO{\dint }%
%BeginExpansion
{\displaystyle\int}
%EndExpansion
\nonumber\\
&  \left\{  \left\{  \varphi_{_{1}}\left(  \mathsf{z}_{1}\right)  ^{\ast
}\cdots\nabla_{j}\varphi_{j}\left(  \mathsf{z}_{j}\right)  ^{\ast}%
\cdots\varphi_{N}\left(  \mathsf{z}_{N}\right)  ^{\ast}\right\}  \bullet
\nabla_{j}\left\{  \left(  -1\right)  ^{\pi\left(  \sigma^{\prime\prime
}\right)  }\varphi_{1}\left(  \mathsf{z}_{\sigma^{\prime\prime}\left(
1\right)  }\right)  \cdots\varphi_{N}\left(  \mathsf{z}_{\sigma^{\prime\prime
}\left(  N\right)  }\right)  \right\}  \right. \nonumber\\
&  \left.
%TCIMACRO{\dprod \limits_{1\leq k\leq\frac{N}{2}}}%
%BeginExpansion
{\displaystyle\prod\limits_{1\leq k\leq\frac{N}{2}}}
%EndExpansion
\nu_{k}\left(  \mathsf{z}_{\sigma\left(  2k-1\right)  },\mathsf{z}%
_{\sigma\left(  2k\right)  }\right)  \nu_{k}\left(  \mathsf{z}_{\sigma
^{\prime}\left(  2k-1\right)  },\mathsf{z}_{\sigma^{\prime}\left(  2k\right)
}\right)  e^{i%
%TCIMACRO{\dsum \limits_{1\leq k\leq\frac{N}{2}}}%
%BeginExpansion
{\displaystyle\sum\limits_{1\leq k\leq\frac{N}{2}}}
%EndExpansion
\left\{  \omega_{k}\left(  \mathsf{z}_{\sigma^{\prime}\left(  2k-1\right)
},\mathsf{z}_{\sigma^{\prime}\left(  2k\right)  }\right)  -\omega_{k}\left(
\mathsf{z}_{\sigma\left(  2k-1\right)  },\mathsf{z}_{\sigma\left(  2k\right)
}\right)  \right\}  }\right\}  d\mathsf{z}^{N}%
\end{align}


\item[II]
\begin{align}
&  -%
%TCIMACRO{\tsum \limits_{\sigma,\sigma^{\prime},\sigma^{\prime\prime}%
%\in\mathcal{S}_{N}}}%
%BeginExpansion
{\textstyle\sum\limits_{\sigma,\sigma^{\prime},\sigma^{\prime\prime}%
\in\mathcal{S}_{N}}}
%EndExpansion%
%TCIMACRO{\dsum \limits_{1\leq j\leq N}}%
%BeginExpansion
{\displaystyle\sum\limits_{1\leq j\leq N}}
%EndExpansion
\int\left\{  \left\{  \left(  -1\right)  ^{\pi\left(  \sigma^{\prime\prime
}\right)  }\varphi_{1}\left(  \mathsf{z}_{\sigma^{\prime\prime}\left(
1\right)  }\right)  \cdots\varphi_{N}\left(  \mathsf{z}_{\sigma^{\prime\prime
}\left(  N\right)  }\right)  \right\}  \right. \nonumber\\
&  \varphi_{_{1}}\left(  \mathsf{z}_{1}\right)  ^{\ast}\cdots\nabla_{j}%
\varphi_{j}\left(  \mathsf{z}_{j}\right)  ^{\ast}\cdots\varphi_{N}\left(
\mathsf{z}_{N}\right)  ^{\ast}\bullet\left\{  \nabla_{j}\left\{
%TCIMACRO{\dsum \limits_{1\leq k\leq\frac{N}{2}}}%
%BeginExpansion
{\displaystyle\sum\limits_{1\leq k\leq\frac{N}{2}}}
%EndExpansion
\left\{  \omega_{k}\left(  \mathsf{z}_{\sigma^{\prime}\left(  2k-1\right)
},\mathsf{z}_{\sigma^{\prime}\left(  2k\right)  }\right)  \right\}  \right\}
\right. \nonumber\\
&  \qquad\qquad\left.  \qquad\right.  \qquad%
%TCIMACRO{\dprod \limits_{1\leq k\leq\frac{N}{2}}}%
%BeginExpansion
{\displaystyle\prod\limits_{1\leq k\leq\frac{N}{2}}}
%EndExpansion
\nu_{k}\left(  \mathsf{z}_{\sigma\left(  2k-1\right)  },\mathsf{z}%
_{\sigma\left(  2k\right)  }\right)  \nu_{k}\left(  \mathsf{z}_{\sigma
^{\prime}\left(  2k-1\right)  },\mathsf{z}_{\sigma^{\prime}\left(  2k\right)
}\right) \nonumber\\
&  \qquad\qquad\left.  -\,i\nabla_{j}\left\{
%TCIMACRO{\dprod \limits_{1\leq k\leq\frac{N}{2}}}%
%BeginExpansion
{\displaystyle\prod\limits_{1\leq k\leq\frac{N}{2}}}
%EndExpansion
\nu_{k}\left(  \mathsf{z}_{\sigma^{\prime}\left(  2k-1\right)  }%
,\mathsf{z}_{\sigma^{\prime}\left(  2k\right)  }\right)  \right\}
%TCIMACRO{\dprod \limits_{1\leq k\leq\frac{N}{2}}}%
%BeginExpansion
{\displaystyle\prod\limits_{1\leq k\leq\frac{N}{2}}}
%EndExpansion
\nu_{k}\left(  \mathsf{z}_{\sigma\left(  2k-1\right)  },\mathsf{z}%
_{\sigma\left(  2k\right)  }\right)  \right\}  d\mathsf{z}^{N}%
\end{align}


\item[III]
\begin{align}
&  -%
%TCIMACRO{\tsum \limits_{\sigma,\sigma^{\prime},\sigma^{\prime\prime}%
%\in\mathcal{S}_{N}}}%
%BeginExpansion
{\textstyle\sum\limits_{\sigma,\sigma^{\prime},\sigma^{\prime\prime}%
\in\mathcal{S}_{N}}}
%EndExpansion%
%TCIMACRO{\dsum \limits_{1\leq j\leq N}}%
%BeginExpansion
{\displaystyle\sum\limits_{1\leq j\leq N}}
%EndExpansion
\int\varphi_{_{1}}\left(  \mathsf{z}_{1}\right)  ^{\ast}\cdots\varphi
_{N}\left(  \mathsf{z}_{N}\right)  ^{\ast}\left\{
%TCIMACRO{\dprod \limits_{1\leq k\leq\frac{N}{2}}}%
%BeginExpansion
{\displaystyle\prod\limits_{1\leq k\leq\frac{N}{2}}}
%EndExpansion
\nu_{k}\left(  \mathsf{z}_{\sigma\left(  2k-1\right)  },\mathsf{z}%
_{\sigma\left(  2k\right)  }\right)  \nu_{k}\left(  \mathsf{z}_{\sigma
^{\prime}\left(  2k-1\right)  },\mathsf{z}_{\sigma^{\prime}\left(  2k\right)
}\right)  \right.  \times\nonumber\\
&  \qquad\qquad\qquad%
%TCIMACRO{\dsum \limits_{1\leq k\leq\frac{N}{2}}}%
%BeginExpansion
{\displaystyle\sum\limits_{1\leq k\leq\frac{N}{2}}}
%EndExpansion
\nabla_{j}\left\{  \omega_{k}\left(  \mathsf{z}_{\sigma^{\prime}\left(
2k-1\right)  },\mathsf{z}_{\sigma^{\prime}\left(  2k\right)  }\right)
\right\}  -i%
%TCIMACRO{\dprod \limits_{1\leq k\leq\frac{N}{2}}}%
%BeginExpansion
{\displaystyle\prod\limits_{1\leq k\leq\frac{N}{2}}}
%EndExpansion
\nu_{k}\left(  \mathsf{z}_{\sigma\left(  2k-1\right)  },\mathsf{z}%
_{\sigma\left(  2k\right)  }\right)  \times\nonumber\\
&  \nabla_{j}\left\{
%TCIMACRO{\dprod \limits_{1\leq k\leq\frac{N}{2}}}%
%BeginExpansion
{\displaystyle\prod\limits_{1\leq k\leq\frac{N}{2}}}
%EndExpansion
\nu_{k}\left(  \mathsf{z}_{\sigma^{\prime}\left(  2k-1\right)  }%
,\mathsf{z}_{\sigma^{\prime}\left(  2k\right)  }\right)  \right\}  \left.
\bullet\nabla_{j}\left\{  \left(  -1\right)  ^{\pi\left(  \sigma^{\prime
\prime}\right)  }\varphi_{1}\left(  \mathsf{z}_{\sigma^{\prime\prime}\left(
1\right)  }\right)  \cdots\varphi_{N}\left(  \mathsf{z}_{\sigma^{\prime\prime
}\left(  N\right)  }\right)  \right\}  \right\}  d\mathsf{z}^{N}%
\end{align}


\item[IV]
\begin{align}
&  +%
%TCIMACRO{\tsum \limits_{\sigma,\sigma^{\prime},\sigma^{\prime\prime}%
%\in\mathcal{S}_{N}}}%
%BeginExpansion
{\textstyle\sum\limits_{\sigma,\sigma^{\prime},\sigma^{\prime\prime}%
\in\mathcal{S}_{N}}}
%EndExpansion%
%TCIMACRO{\dsum \limits_{1\leq j\leq N}}%
%BeginExpansion
{\displaystyle\sum\limits_{1\leq j\leq N}}
%EndExpansion
\int\varphi_{_{1}}\left(  \mathsf{z}_{1}\right)  ^{\ast}\cdots\varphi
_{N}\left(  \mathsf{z}_{N}\right)  ^{\ast}\left(  -1\right)  ^{\pi\left(
\sigma^{\prime\prime}\right)  }\varphi_{1}\left(  \mathsf{z}_{\sigma
^{\prime\prime}\left(  1\right)  }\right)  \cdots\varphi_{N}\left(
\mathsf{z}_{\sigma^{\prime\prime}\left(  N\right)  }\right)  \times\nonumber\\
&  \left\{
%TCIMACRO{\dprod \limits_{1\leq k\leq\frac{N}{2}}}%
%BeginExpansion
{\displaystyle\prod\limits_{1\leq k\leq\frac{N}{2}}}
%EndExpansion
\nu_{k}\left(  \mathsf{z}_{\sigma\left(  2k-1\right)  },\mathsf{z}%
_{\sigma\left(  2k\right)  }\right)  \nabla_{j}\left\{
%TCIMACRO{\dsum \limits_{1\leq k\leq\frac{N}{2}}}%
%BeginExpansion
{\displaystyle\sum\limits_{1\leq k\leq\frac{N}{2}}}
%EndExpansion
\left\{  \omega_{k}\left(  \mathsf{z}_{\sigma\left(  2k-1\right)  }%
,\mathsf{z}_{\sigma\left(  2k\right)  }\right)  \right\}  \right\}  \right.
\nonumber\\
&  -i\nabla_{j}\left\{
%TCIMACRO{\dprod \limits_{1\leq k\leq\frac{N}{2}}}%
%BeginExpansion
{\displaystyle\prod\limits_{1\leq k\leq\frac{N}{2}}}
%EndExpansion
\nu_{k}\left(  \mathsf{z}_{\sigma\left(  2k-1\right)  },\mathsf{z}%
_{\sigma\left(  2k\right)  }\right)  \right\}  \bullet\left\{  \nabla
_{j}\left\{
%TCIMACRO{\dsum \limits_{1\leq k\leq\frac{N}{2}}}%
%BeginExpansion
{\displaystyle\sum\limits_{1\leq k\leq\frac{N}{2}}}
%EndExpansion
\left\{  \omega_{k}\left(  \mathsf{z}_{\sigma^{\prime}\left(  2k-1\right)
},\mathsf{z}_{\sigma^{\prime}\left(  2k\right)  }\right)  \right\}  \right\}
\right.  \times\nonumber\\
&  \qquad\qquad%
%TCIMACRO{\dprod \limits_{1\leq k\leq\frac{N}{2}}}%
%BeginExpansion
{\displaystyle\prod\limits_{1\leq k\leq\frac{N}{2}}}
%EndExpansion
\nu_{k}\left(  \mathsf{z}_{\sigma^{\prime}\left(  2k-1\right)  }%
,\mathsf{z}_{\sigma^{\prime}\left(  2k\right)  }\right)  \left.  -i\nabla
_{j}\left\{
%TCIMACRO{\dprod \limits_{1\leq k\leq\frac{N}{2}}}%
%BeginExpansion
{\displaystyle\prod\limits_{1\leq k\leq\frac{N}{2}}}
%EndExpansion
\nu_{k}\left(  \mathsf{z}_{\sigma^{\prime}\left(  2k-1\right)  }%
,\mathsf{z}_{\sigma^{\prime}\left(  2k\right)  }\right)  \right\}  \right\}
d\mathsf{z}^{N}%
\end{align}
The expectation value of the external potential term
\begin{equation}
\mathcal{V}_{1}\left(  \nu,\omega\right)  =\mathcal{S}\left(  \nu
,\omega\right)  ^{-1}\left\langle \varphi_{1}\cdots\varphi_{N}\left\vert
\left(  V_{ext}+V_{en}\right)  \right.  \varphi_{1}\cdots\varphi
_{N}\right\rangle
\end{equation}
is given in terms of the effective operators%
\begin{align}
&
%TCIMACRO{\tsum \limits_{\sigma,\sigma^{\prime}\in\mathcal{S}_{N}}}%
%BeginExpansion
{\textstyle\sum\limits_{\sigma,\sigma^{\prime}\in\mathcal{S}_{N}}}
%EndExpansion%
%TCIMACRO{\dsum \limits_{1\leq j\leq N}}%
%BeginExpansion
{\displaystyle\sum\limits_{1\leq j\leq N}}
%EndExpansion
\int\left\{  \left\{  V_{ext}\left(  \mathsf{z}_{j}\right)  +V_{en}\left(
\mathsf{z}_{j}\right)  \right\}  T_{1}\left(  \mathsf{z}_{\sigma\left(
1\right)  },\mathsf{z}_{\sigma\left(  2\right)  }\right)  ^{\ast}\cdots
T_{\frac{N}{2}}\left(  \mathsf{z}_{\sigma\left(  N-1\right)  ,}\mathsf{z}%
_{\sigma\left(  N\right)  }\right)  ^{\ast}\right. \nonumber\\
&  \qquad\qquad\qquad\qquad\left.  T_{1}\left(  \mathsf{z}_{\sigma^{\prime
}\left(  1\right)  },\mathsf{z}_{\sigma^{\prime}\left(  2\right)  }\right)
\cdots T_{\frac{N}{2}}\left(  \mathsf{z}_{\sigma^{\prime}\left(  N-1\right)
,}\mathsf{z}_{\sigma^{\prime}\left(  N\right)  }\right)  \left\vert
\mathsf{z}^{N}\right\rangle \left\langle \mathsf{z}^{N}\right\vert \right\}
d\mathsf{z}^{N}%
\end{align}
and can be evaluated to give
\begin{align}
\mathcal{V}_{1}\left(  \nu,\omega\right)   &  =%
%TCIMACRO{\tsum \limits_{\sigma,\sigma^{\prime},\sigma^{\prime\prime}%
%\in\mathcal{S}_{N}}}%
%BeginExpansion
{\textstyle\sum\limits_{\sigma,\sigma^{\prime},\sigma^{\prime\prime}%
\in\mathcal{S}_{N}}}
%EndExpansion%
%TCIMACRO{\dsum \limits_{1\leq j\leq N}}%
%BeginExpansion
{\displaystyle\sum\limits_{1\leq j\leq N}}
%EndExpansion
\int\left\{  \varphi_{_{1}}\left(  \mathsf{z}_{1}\right)  ^{\ast}\cdots
\varphi_{N}\left(  \mathsf{z}_{N}\right)  ^{\ast}\left(  -1\right)
^{\pi\left(  \sigma^{\prime\prime}\right)  }\varphi_{1}\left(  \mathsf{z}%
_{\sigma^{\prime\prime}\left(  1\right)  }\right)  \cdots\varphi_{N}\left(
\mathsf{z}_{\sigma^{\prime\prime}\left(  N\right)  }\right)  \right.
\times\nonumber\\
&  \left\{  V_{ext}\left(  \mathsf{z}_{j}\right)  +V_{en}\left(
\mathsf{z}_{j}\right)  \right\}
%TCIMACRO{\dprod \limits_{1\leq k\leq\frac{N}{2}}}%
%BeginExpansion
{\displaystyle\prod\limits_{1\leq k\leq\frac{N}{2}}}
%EndExpansion
\nu_{k}\left(  \mathsf{z}_{\sigma\left(  2k-1\right)  },\mathsf{z}%
_{\sigma\left(  2k\right)  }\right)  \nu_{k}\left(  \mathsf{z}_{\sigma
^{\prime}\left(  2k-1\right)  },\mathsf{z}_{\sigma^{\prime}\left(  2k\right)
}\right) \nonumber\\
&  \left.  e^{i\left\{  \omega_{1}\left(  \mathsf{z}_{\sigma^{\prime}\left(
1\right)  },\mathsf{z}_{\sigma^{\prime}\left(  2\right)  }\right)  -\omega
_{1}\left(  \mathsf{z}_{\sigma\left(  1\right)  },\mathsf{z}_{\sigma\left(
2\right)  }\right)  +\cdots+\omega_{\frac{N}{2}}\left(  \mathsf{z}%
_{\sigma^{\prime}\left(  N-1\right)  ,}\mathsf{z}_{\sigma^{\prime}\left(
N\right)  }\right)  -\omega_{\frac{N}{2}}\left(  \mathsf{z}_{\sigma\left(
N-1\right)  ,}\mathsf{z}_{\sigma\left(  N\right)  }\right)  \right\}
}\right\}  d\mathsf{z}^{N}%
\end{align}

\end{description}

The expectation value $\mathcal{V}_{2}\left(  \nu\right)  $ of the coulomb
potential term is given by
\begin{equation}
\mathcal{V}_{2}\left(  \nu,\omega\right)  =\mathcal{S}\left(  \nu
,\omega\right)  ^{-1}\left\langle \varphi_{1}\cdots\varphi_{N}\left\vert
V_{2}\right.  \varphi_{1}\cdots\varphi_{N}\right\rangle
\end{equation}
which can be evaluated in terms of the effective coulomb operator
\begin{align}
&
%TCIMACRO{\tsum \limits_{\sigma,\sigma^{\prime}\in\mathcal{S}_{N}}}%
%BeginExpansion
{\textstyle\sum\limits_{\sigma,\sigma^{\prime}\in\mathcal{S}_{N}}}
%EndExpansion%
%TCIMACRO{\dsum \limits_{1\leq j<k\leq N}}%
%BeginExpansion
{\displaystyle\sum\limits_{1\leq j<k\leq N}}
%EndExpansion
\int\left\{  V_{2}\left(  \mathsf{z}_{j,}\mathsf{z}_{k}\right)  T_{1}\left(
\mathsf{z}_{\sigma\left(  1\right)  },\mathsf{z}_{\sigma\left(  2\right)
}\right)  ^{\ast}\cdots T_{\frac{N}{2}}\left(  \mathsf{z}_{\sigma\left(
N-1\right)  ,}\mathsf{z}_{\sigma\left(  N\right)  }\right)  ^{\ast}\right.
\nonumber\\
&  \qquad\qquad\qquad\qquad\left.  T_{1}\left(  \mathsf{z}_{\sigma^{\prime
}\left(  1\right)  },\mathsf{z}_{\sigma^{\prime}\left(  2\right)  }\right)
\cdots T_{\frac{N}{2}}\left(  \mathsf{z}_{\sigma^{\prime}\left(  N-1\right)
,}\mathsf{z}_{\sigma^{\prime}\left(  N\right)  }\right)  \left\vert
\mathsf{z}^{N}\right\rangle \left\langle \mathsf{z}^{N}\right\vert \right\}
d\mathsf{z}^{N}%
\end{align}
to give%
\begin{align}
\mathcal{V}_{2}\left(  \nu,\omega\right)   &  =%
%TCIMACRO{\tsum \limits_{\sigma,\sigma^{\prime},\sigma^{\prime\prime}%
%\in\mathcal{S}_{N}}}%
%BeginExpansion
{\textstyle\sum\limits_{\sigma,\sigma^{\prime},\sigma^{\prime\prime}%
\in\mathcal{S}_{N}}}
%EndExpansion%
%TCIMACRO{\dsum \limits_{1\leq j<k\leq N}}%
%BeginExpansion
{\displaystyle\sum\limits_{1\leq j<k\leq N}}
%EndExpansion
\int\left\{  \varphi_{_{1}}\left(  \mathsf{z}_{1}\right)  ^{\ast}\cdots
\varphi_{N}\left(  \mathsf{z}_{N}\right)  ^{\ast}\left(  -1\right)
^{\pi\left(  \sigma^{\prime\prime}\right)  }\varphi_{1}\left(  \mathsf{z}%
_{\sigma^{\prime\prime}\left(  1\right)  }\right)  \cdots\varphi_{N}\left(
\mathsf{z}_{\sigma^{\prime\prime}\left(  N\right)  }\right)  \right.
\times\nonumber\\
&  V_{2}\left(  \mathsf{z}_{j},\mathsf{z}_{k}\right)
%TCIMACRO{\dprod \limits_{1\leq k\leq\frac{N}{2}}}%
%BeginExpansion
{\displaystyle\prod\limits_{1\leq k\leq\frac{N}{2}}}
%EndExpansion
\nu_{k}\left(  \mathsf{z}_{\sigma\left(  2k-1\right)  },\mathsf{z}%
_{\sigma\left(  2k\right)  }\right)  \nu_{k}\left(  \mathsf{z}_{\sigma
^{\prime}\left(  2k-1\right)  },\mathsf{z}_{\sigma^{\prime}\left(  2k\right)
}\right) \nonumber\\
&  \left.  e^{i\left\{  \omega_{1}\left(  \mathsf{z}_{\sigma^{\prime}\left(
1\right)  },\mathsf{z}_{\sigma^{\prime}\left(  2\right)  }\right)  -\omega
_{1}\left(  \mathsf{z}_{\sigma\left(  1\right)  },\mathsf{z}_{\sigma\left(
2\right)  }\right)  +\cdots+\omega_{\frac{N}{2}}\left(  \mathsf{z}%
_{\sigma^{\prime}\left(  N-1\right)  ,}\mathsf{z}_{\sigma^{\prime}\left(
N\right)  }\right)  -\omega_{\frac{N}{2}}\left(  \mathsf{z}_{\sigma\left(
N-1\right)  ,}\mathsf{z}_{\sigma\left(  N\right)  }\right)  \right\}
}\right\}  d\mathsf{z}^{N}%
\end{align}


where%
\begin{equation}
V_{2}\left(  \mathsf{z}_{j},\mathsf{z}_{k}\right)  =\frac{1}{\left\Vert
\mathbf{r}_{j}-\mathbf{r}_{k}\right\Vert }\delta\left(  \xi_{j}-\xi
_{k}\right)
\end{equation}
for $\mathsf{z}_{j}\equiv\left(  \mathbf{r}_{j},\xi_{j}\right)  $ and
$\mathsf{z}_{k}\equiv\left(  \mathbf{r}_{k},\xi_{k}\right)  .$

The value of the energy functional $\mathcal{E}\left(  \mathcal{D}^{N}\left(
\nu,\omega\right)  \right)  $ is then the sum
\begin{equation}
\mathcal{E}\left(  \mathcal{D}^{N}\left(  \nu_{1},\ldots,\nu_{\frac{N}{2}%
},\omega_{1},\ldots,\omega_{\frac{N}{2}}\right)  \right)  =H_{0}%
+\mathcal{K}\left(  \nu,\omega\right)  +\mathcal{V}_{1}\left(  \nu
,\omega\right)  +\mathcal{V}_{2}\left(  \nu,\omega\right)
\end{equation}


and the value of $\mathcal{F}_{\left\{  2,\cdots,2\right\}  }\left(
\mathcal{D}_{ref}^{1}\right)  $ at the reference IPS
\begin{equation}
\mathcal{D}_{ref}^{N}=\mathcal{D}^{N}\left(  \mathcal{D}_{ref}^{1}\right)
=\frac{\left\vert \varphi_{1}\cdots\varphi_{N}\right\rangle \left\langle
\varphi_{1}\cdots\varphi_{N}\right\vert }{\left\langle \varphi_{1}%
\cdots\varphi_{N}|\varphi_{1}\cdots\varphi_{N}\right\rangle }%
\end{equation}
is given by minimizing $\mathcal{E}\left(  \mathcal{D}^{N}\left(  \nu
,\omega\right)  \right)  $ over $\nu=\left(  \nu_{1},\cdots,\nu\frac{_{N}}%
{2}\right)  $ and $\omega=\left(  \omega_{1},\cdots,\omega\frac{_{N}}%
{2}\right)  \mathsf{\ }$constraining $\nu$ and $\omega$ to satisfy the
\emph{necessary }generalized strong orthogonality condition wrt $\left\vert
\varphi_{1}\cdots\varphi_{N}\right\rangle \left\langle \varphi_{1}%
\cdots\varphi_{N}\right\vert $ i.e.
\begin{align}
&
%TCIMACRO{\tsum \limits_{\sigma,\sigma^{\prime},\sigma^{\prime\prime}%
%\in\mathcal{S}_{N}}}%
%BeginExpansion
{\textstyle\sum\limits_{\sigma,\sigma^{\prime},\sigma^{\prime\prime}%
\in\mathcal{S}_{N}}}
%EndExpansion%
%TCIMACRO{\dint }%
%BeginExpansion
{\displaystyle\int}
%EndExpansion
\left\{  \left(  -1\right)  ^{\pi\left(  \sigma^{\prime\prime}\right)
}\varphi_{_{1}}\left(  \mathsf{z}_{1}^{\prime}\right)  ^{\ast}\varphi_{_{2}%
}\left(  \mathsf{z}_{2}\right)  ^{\ast}\cdots\varphi_{N}\left(  \mathsf{z}%
_{N}\right)  ^{\ast}\left(  -1\right)  ^{\pi\left(  \sigma^{\prime\prime
}\right)  }\varphi_{1}\left(  \mathsf{z}_{\sigma^{\prime\prime}\left(
1\right)  }\right)  \cdots\varphi_{N}\left(  \mathsf{z}_{\sigma^{\prime\prime
}\left(  N\right)  }\right)  \right.  \times\nonumber\\
&  \qquad\qquad\qquad%
%TCIMACRO{\dprod \limits_{1\leq k\leq\frac{N}{2}}}%
%BeginExpansion
{\displaystyle\prod\limits_{1\leq k\leq\frac{N}{2}}}
%EndExpansion
\nu_{k}\left(  \mathsf{z}_{\sigma\left(  2k-1\right)  },\mathsf{z}%
_{\sigma\left(  2k\right)  }\right)  \nu_{k}\left(  \mathsf{z}_{\sigma
^{\prime}\left(  2k-1\right)  },\mathsf{z}_{\sigma^{\prime}\left(  2k\right)
}\right) \nonumber\\
&  \qquad\qquad\qquad\qquad\qquad e^{i\left\{  \omega_{1}\left(
\mathsf{z}_{\sigma^{\prime}\left(  1\right)  },\mathsf{z}_{\sigma^{\prime
}\left(  2\right)  }\right)  -\omega_{1}\left(  \mathsf{z}_{\sigma\left(
1\right)  },\mathsf{z}_{\sigma\left(  2\right)  }\right)  +\cdots
+\omega_{\frac{N}{2}}\left(  \mathsf{z}_{\sigma^{\prime}\left(  N-1\right)
,}\mathsf{z}_{\sigma^{\prime}\left(  N\right)  }\right)  -\omega_{\frac{N}{2}%
}\left(  \mathsf{z}_{\sigma\left(  N-1\right)  ,}\mathsf{z}_{\sigma\left(
N\right)  }\right)  \right\}  }\nonumber\\
&  \qquad\qquad\qquad\qquad\left.  -\varphi_{1}\left(  \mathsf{z}_{1}\right)
\cdots\varphi_{N}\left(  \mathsf{z}_{N}\right)  \varphi_{\sigma^{\prime\prime
}\left(  1\right)  }\left(  \mathsf{z}_{1}^{\prime}\right)  ^{\ast}%
\varphi_{\sigma^{\prime\prime}\left(  2\right)  }\left(  \mathsf{z}%
_{2}\right)  ^{\ast}\cdots\varphi_{\sigma^{\prime\prime}\left(  N\right)
}\left(  \mathsf{z}_{N}\right)  ^{\ast}\right\}  d\mathsf{z}^{N-1}%
\overset{a.e}{=}0
\end{align}


\section{Summary\label{summary}}

\appendix

\section{Symbolic Notation\label{symnot}}

In this article we make extensive utilization of representations of super
vectors and super operators with respect to different bases and the
transformations between these representations. In order to succinctly display
these representations we use a straightforward generalization of the notation
introduced by L\"{o}wdin \cite{Lowdin} for Hilbert to spaces to dual spaces of
super vectors. In this appendix we briefly review L\"{o}wdins symbolic notation.

In this formalism a given basis of a Hilbert space $\mathcal{H}$ is written as
a \emph{row }vector of the vectors constituting the basis i.e. $\left\{
\left.  \left\vert v_{j}\right\rangle \right\vert 1\leq j\leq n\right\}  $ and
is displayed as
\begin{equation}
\left(  \left\vert v_{1}\right\rangle ,\cdots,\left\vert v_{n}\right\rangle
\right)  \equiv\left(  \underline{\left\vert v\right\rangle }\right)  _{R}.
\end{equation}
The basis of the dual vector space $\mathcal{H}^{d}$ (the vector space of
bounded linear functionals on $\mathcal{H}$) $\left\{  \left.  \left\vert
w_{j}^{d}\right\rangle \right\vert 1\leq j\leq n\right\}  $ is also displayed
as a row vector $\left(  \underline{\left\vert w^{d}\right\rangle }\right)
_{R}$, but their action on $\mathcal{H}$ as linear functionals is expressed in
terms of the column vector%
\begin{equation}
\left(
\begin{array}
[c]{c}%
\left\langle w_{1}^{d}\right\vert \\
\vdots\\
\left\langle w_{n}^{d}\right\vert
\end{array}
\right)  \equiv\left(  \underline{\left\vert w\right\rangle }\right)  _{C}%
\end{equation}
Hilbert spaces are self dual i.e. $\mathcal{H}^{d}=\mathcal{H}$ giving the
standard relationships
\begin{subequations}
\begin{align}
\left(
\begin{array}
[c]{c}%
\left\langle w_{1}^{d}\right\vert \\
\vdots\\
\left\langle w_{n}^{d}\right\vert
\end{array}
\right)   &  =\left(  \left\vert w_{1}^{d}\right\rangle ,\cdots,\left\vert
w_{n}^{d}\right\rangle \right)  ^{\dagger}\\
\left(  \left\vert v_{1}\right\rangle ,\cdots,\left\vert v_{n}\right\rangle
\right)   &  =\left(
\begin{array}
[c]{c}%
\left\langle v_{1}\right\vert \\
\vdots\\
\left\langle v_{n}\right\vert
\end{array}
\right)  ^{\dagger}%
\end{align}
where all the basis elements belong to $\mathcal{H}$. It is possible to change
from a row vector to column vector without changing a bra to ket using the
transpose operation i.e.
\end{subequations}
\begin{subequations}
\begin{align}
\left(
\begin{array}
[c]{c}%
\left\vert v_{1}\right\rangle \\
\vdots\\
\left\vert v_{n}\right\rangle
\end{array}
\right)   &  =\left(  \left\vert v_{1}\right\rangle ,\cdots,\left\vert
v_{n}\right\rangle \right)  ^{t}\\
\left(
\begin{array}
[c]{c}%
\left\langle v_{1}\right\vert \\
\vdots\\
\left\langle v_{n}\right\vert
\end{array}
\right)   &  =\left(  \left\langle v_{1}\right\vert ,\cdots,\left\langle
v_{n}\right\vert \right)  ^{t}%
\end{align}
When we deal with non self dual spaces the superscript "$d$" denotes an
element of the dual space.

Any vector $\left\vert a\right\rangle \in\mathcal{H}$ can then be represented
as $\mathbf{a}$\textbf{\ }wrt the basis $\left\{  \left.  \left\vert
v_{j}\right\rangle \right\vert 1\leq j\leq n\right\}  $ as
\end{subequations}
\begin{equation}
\left\vert a\right\rangle =\left(  \left\vert v_{1}\right\rangle
,\cdots,\left\vert v_{n}\right\rangle \right)  \left(
\begin{array}
[c]{c}%
a_{1}\\
\vdots\\
a_{n}%
\end{array}
\right)  \equiv\left(  \underline{\left\vert v\right\rangle }\right)
_{R}\mathbf{a}_{v}%
\end{equation}
and any operator $T$ as
\begin{equation}
T\left(  \underline{\left\vert v\right\rangle }\right)  _{R}=\left(
\underline{\left\vert Tv\right\rangle }\right)  _{R}=\left(  \underline
{\left\vert v\right\rangle }\right)  _{R}\mathbb{T}%
\end{equation}
where $\mathbb{T}$ is the matrix representing $T$ wrt the basis $\left\{
\left\vert v\right\rangle \right\}  $ and
\begin{equation}
T\left\vert v_{j}\right\rangle =%
%TCIMACRO{\dsum \limits_{1\leq k\leq n}}%
%BeginExpansion
{\displaystyle\sum\limits_{1\leq k\leq n}}
%EndExpansion
\left\vert v_{k}\right\rangle T_{kj};\quad1\leq j\leq n.
\end{equation}
Using this notation allows one to obtain the transpose and hermitian conjugate
of matrices as
\begin{equation}
\left(  \left(  \underline{\left\vert Tv\right\rangle }\right)  \right)
_{R}^{t}=\left(  \left(  \underline{\left\vert v\right\rangle }\right)
\mathbb{T}\right)  _{R}^{t}=\mathbb{T}^{t}\left(  \underline{\left\vert
v\right\rangle }\right)  _{R}^{t}\equiv\mathbb{T}^{t}\left(  \underline
{\left\vert v\right\rangle }\right)  _{C}%
\end{equation}
and
\begin{equation}
\left(  \underline{\left\vert Tv\right\rangle }\right)  _{R}^{\dagger
}=\mathbb{T}^{\dagger}\left(  \left\vert \underline{v}\right\rangle \right)
_{R}^{\dagger}\equiv\mathbb{T}^{\dagger}\left(  \underline{\left\langle
v\right\vert }\right)  _{C}%
\end{equation}


The \emph{metric matrix} of the basis $\left\{  \left\vert v\right\rangle
\right\}  $ can be expressed as
\begin{equation}
\left(
\begin{array}
[c]{c}%
\left\langle v_{1}\right\vert \\
\vdots\\
\left\langle v_{n}\right\vert
\end{array}
\right)  \left(  \left\vert v_{1}\right\rangle ,\cdots,\left\vert
v_{n}\right\rangle \right)  =\left(  \underline{\left\langle v\right\vert
}\right)  _{R}^{t}\left(  \underline{\left\vert v\right\rangle }\right)
_{R}=\left(
\begin{array}
[c]{ccc}%
\left\langle v_{1}|v_{1}\right\rangle  & \cdots & \left\langle v_{1}%
|v_{n}\right\rangle \\
\vdots & \vdots & \vdots\\
\left\langle v_{n}|v_{1}\right\rangle  & \cdots & \left\langle v_{n}%
|v_{1}\right\rangle
\end{array}
\right)
\end{equation}
and a \emph{resolution of the identity} as
\begin{equation}
\left(  \left\vert v_{1}\right\rangle ,\cdots,\left\vert v_{n}\right\rangle
\right)  \left(
\begin{array}
[c]{c}%
\left\langle v_{1}^{d}\right\vert \\
\vdots\\
\left\langle v_{n}^{d}\right\vert
\end{array}
\right)  =\left(  \underline{\left\vert v\right\rangle }\right)  _{R}\left(
\underline{\left\langle v^{d}\right\vert }\right)  _{R}^{t}=%
%TCIMACRO{\dsum \limits_{1\leq j\leq n}}%
%BeginExpansion
{\displaystyle\sum\limits_{1\leq j\leq n}}
%EndExpansion
\left\vert v_{j}\right\rangle \left\langle v_{j}^{d}\right\vert
\end{equation}
where $\left\{  \underline{\left\vert v^{d}\right\rangle }\right\}  $ is a
biorthogonal basis to $\left\{  \underline{\left\vert v\right\rangle
}\right\}  $ i.e.
\begin{equation}
\left(
\begin{array}
[c]{ccc}%
\left\langle v_{1}^{d}|v_{1}\right\rangle  & \cdots & \left\langle v_{1}%
^{d}|v_{n}\right\rangle \\
\vdots & \ddots & \vdots\\
\left\langle v_{n}^{d}|v_{1}\right\rangle  & \cdots & \left\langle v_{n}%
^{d}|v_{n}\right\rangle
\end{array}
\right)  =\left(
\begin{array}
[c]{ccc}%
1 & \cdots & 0\\
0 & \ddots & 0\\
0 & \cdots & 1
\end{array}
\right)  \equiv\mathbb{I}_{n}%
\end{equation}


Transformations between any two bases $\left\{  \underline{\left\vert
v\right\rangle }\right\}  $ and $\left\{  \underline{\left\vert w\right\rangle
}\right\}  $ is simply displayed using the resolution of the identity as
\begin{equation}
\left(  \underline{\left\vert w\right\rangle }\right)  _{R}=\left(
\underline{\left\vert w\right\rangle }\right)  _{R}\left(  \underline
{\left\langle w^{d}\right\vert }\right)  _{R}^{t}\left(  \underline{\left\vert
v\right\rangle }\right)  _{R}%
\end{equation}
so that $\left(  \underline{\left\langle w^{d}\right\vert }\right)
^{t}\left(  \underline{\left\vert v\right\rangle }\right)  $ is the
transformation matrix from the basis $\left\{  \underline{\left\vert
v\right\rangle }\right\}  $ to the basis, $\left\{  \underline{\left\vert
w\right\rangle }\right\}  $ which is explicitly given by
\begin{equation}
\left(  \underline{\left\langle w^{d}\right\vert }\right)  _{R}^{t}\left(
\underline{\left\vert v\right\rangle }\right)  _{R}=\left(
\begin{array}
[c]{ccc}%
\left\langle w_{1}^{d}|v_{1}\right\rangle  & \cdots & \left\langle w_{1}%
^{d}|v_{n}\right\rangle \\
\vdots & \ddots & \vdots\\
\left\langle w_{n}^{d}|v_{1}\right\rangle  & \cdots & \left\langle w_{n}%
^{d}|v_{n}\right\rangle
\end{array}
\right)
\end{equation}
The transformation between the representation $\mathbf{a}_{v}$ of vector
$\left\vert a\right\rangle $ on the $\left\{  \underline{\left\vert
v\right\rangle }\right\}  $ basis to representation $\mathbf{a}_{w}$ on the
$\left\{  \underline{\left\vert w\right\rangle }\right\}  $ basis is given by%
\begin{equation}
\left(  \underline{\left\vert w\right\rangle }\right)  _{R}\mathbf{a}%
_{w}=\left(  \underline{\left\vert w\right\rangle }\right)  _{R}\left(
\underline{\left\langle w^{d}\right\vert }\right)  _{R}^{t}\left(
\underline{\left\vert v\right\rangle }\right)  _{R}\mathbf{a}_{v}%
\end{equation}
showing that
\begin{equation}
\mathbf{a}_{w}=\left(  \underline{\left\langle w^{d}\right\vert }\right)
_{R}^{t}\left(  \underline{\left\vert v\right\rangle }\right)  _{R}%
\mathbf{a}_{v}%
\end{equation}


\section{States, Cones and Convex Sets\label{states}}

A convex set is the convex closure of its extreme points in the relevant
topology i.e.%
\begin{equation}
S=\overline{ext}\left\{  S\right\}
\end{equation}
$x$ is an extreme point when it cannot be expressed as a convex sum of two
other elements of $S$ i.e. when
\begin{align}
x  &  =\left(  1-\alpha\right)  x_{1}+\alpha x_{2};\quad\alpha\geq0;\quad
x_{1},x_{2}\in S\nonumber\\
&  \Longrightarrow\alpha=0
\end{align}
Any element of the convex set can be expanded as
\begin{align}
x  &  =\sum\alpha_{i}x_{i}\\
\sum\alpha_{i}  &  =1;\quad\alpha_{i}\geq0\\
x_{i}  &  \in ext\left\{  S\right\}
\end{align}
The coefficients $\left\{  \alpha_{i}\right\}  $ are not unique. The cone $C$
generated by $S$ is given by
\begin{equation}
C=\left\{  \alpha x|\forall\alpha\geq0,\quad x\in S\right\}
\end{equation}
Thus $y\in C$ if
\begin{align}
y  &  =\alpha\sum\alpha_{i}x_{i}=\sum\alpha\alpha_{i}x_{i}=\sum\beta_{i}%
x_{i}\\
\sum\beta_{i}  &  =\alpha\geq0,\quad\beta_{i}\geq0
\end{align}
Hence the extreme points of $S$ generate $C$. The convex set $S$ is called the
base of the cone $C$.

\section{Subspaces of Operators\label{ideals}}

\subsection{Definitions}

The set of trace class operators is defined as
\begin{equation}
\mathcal{B}_{1}\left(  \mathcal{H}\right)  =\left\{  X|X\in\mathcal{B}\left(
\mathcal{H}\right)  ;Tr\left\{  \left(  X^{\dagger}X\right)  ^{\frac{1}{2}%
}\right\}  <\infty\right\}
\end{equation}
where $\mathcal{B}\left(  \mathcal{H}\right)  $ is the space of bounded
operators acting in the Hilbert space $\mathcal{H}$. $\mathcal{B}_{1}\left(
\mathcal{H}\right)  $ is a Banach space equipped with the norm
\begin{equation}
\left\Vert X\right\Vert _{1}=Tr\left\{  \left(  X^{\dagger}X\right)
^{\frac{1}{2}}\right\}
\end{equation}
The space $\mathcal{B}\left(  \mathcal{H}\right)  $ is also a Banach space
with a norm defined as
\begin{equation}
\left\Vert X\right\Vert =\sup_{\left\vert \Psi\right\rangle \in\mathcal{H}%
}\left\{  \frac{\left\langle \Psi|X\Psi\right\rangle }{\left\langle \Psi
|\Psi\right\rangle }\right\}
\end{equation}
and is the dual space of $\mathcal{B}_{1}\left(  \mathcal{H}\right)  $ i.e.
the set of continuous (wrt the topology given by $\left\Vert \ \right\Vert
_{1}$) linear functionals defined on $\mathcal{B}_{1}\left(  \mathcal{H}%
\right)  $. The closure of $\mathcal{B}_{1}\left(  \mathcal{H}\right)  $ wrt
$\left\Vert {}\right\Vert $ produces $\mathcal{B}\left(  \mathcal{H}\right)
.$ The set of trace class operators is also a subset of the set of
Hilbert-Schmidt operators, which is defined as
\begin{equation}
\mathcal{B}_{2}\left(  \mathcal{H}\right)  =\left\{  X|X\in\mathcal{B}\left(
\mathcal{H}\right)  ;Tr\left\{  \left(  X^{\dagger}X\right)  \right\}
<\infty\right\}
\end{equation}
and is also a Banach space equipped with the norm
\begin{equation}
\left\Vert X\right\Vert _{2}=\left\{  Tr\left(  X^{\dagger}X\right)  \right\}
^{\frac{1}{2}}.
\end{equation}
and in fact a Hilbert space equipped with the inner product
\begin{equation}
\left(  X|Y\right)  _{HS}=Tr\left\{  \left(  X^{\dagger}Y\right)  \right\}  .
\end{equation}
The closure of $\mathcal{B}_{1}\left(  \mathcal{H}\right)  $ wrt $\left\Vert
{}\right\Vert _{2}$ produces $\mathcal{B}_{2}\left(  \mathcal{H}\right)  $.
The relationship between these spaces is
\begin{equation}
\mathcal{B}\left(  \mathcal{H}\right)  \supseteq\mathcal{B}_{2}\left(
\mathcal{H}\right)  \supseteq\mathcal{B}_{1}\left(  \mathcal{H}\right)
\end{equation}
and their norms is
\begin{equation}
\left\Vert X\right\Vert \leq\left\Vert X\right\Vert _{2}\leq\left\Vert
X\right\Vert _{1}%
\end{equation}
It can be shown that the product of two Hilbert-Schmidt operators is always a
trace class operator i.e.
\begin{equation}
AB\in\mathcal{B}_{1}\left(  \mathcal{H}\right)  \text{ if }A\in\mathcal{B}%
_{2}\left(  \mathcal{H}\right)  \text{ and }B\in\mathcal{B}_{2}\left(
\mathcal{H}\right)  .
\end{equation}


\subsection{A Note on Closures}

The topological closure of a set in a Banach space is the union of the set and
the set of all the limit points of all Cauchy sequences \cite{cauchy}
generated by the set. The algebraic closure of a set in a Banach space is the
totality of all linear combinations of the elements of the set. A subspace in
a Banach space\ is the set of linear combinations, that have finite norm,
generated by the subset, if this does not include all the limit points of the
Cauchy sequences generated by the set then the subspace is open. Thus the set,
$\mathcal{T},$ of operators that have finite trace class norm is a subset of
$\mathcal{B}\left(  \mathcal{H}\right)  $ and its topological closures wrt the
different norms $\left\Vert {}\right\Vert ,\left\Vert {}\right\Vert _{2},$ and
$\left\Vert {}\right\Vert _{1}$ produce
\begin{subequations}
\begin{align}
\mathcal{B}\left(  \mathcal{H}\right)   &  =\left\Vert {}\right\Vert
-closure\left\{  \mathcal{T}\right\} \\
&  \left\vert \cup\right. \nonumber\\
\mathcal{B}_{2}\left(  \mathcal{H}\right)   &  =\left\Vert {}\right\Vert
_{2}-closure\left\{  \mathcal{T}\right\} \\
&  \left\vert \cup\right. \nonumber\\
\mathcal{B}_{1}\left(  \mathcal{H}\right)   &  =\left\Vert {}\right\Vert
_{1}-closure\left\{  \mathcal{T}\right\}
\end{align}
the equalities occur when space $\mathcal{H}$ is finite dimensional.

\subsection{Complimentary Subspaces}

As $\mathcal{B}_{1}\left(  \mathcal{H}^{N}\right)  $ \emph{is not }a Hilbert
space there is no property of orthogonality and hence no orthogonal complement
to a given subspace. Recall that a subspace $\mathcal{W}^{c}$ is complementary
to $\mathcal{W}$ in a space $\mathcal{V}$ if
\end{subequations}
\begin{equation}
\mathcal{W}\cap\mathcal{W}^{c}=\left\{  0\right\}
\end{equation}
and $\mathcal{W}\cup\mathcal{W}^{c}$ linearly generates $\mathcal{V}$. In
general there is a set, $\mathfrak{V}=\left\{  \mathcal{W}_{\Gamma}%
^{c}\right\}  $, of subspaces%
\begin{equation}
\mathcal{W}_{\Gamma}^{c}=\tilde{\Gamma}\left(  \mathcal{W}^{c}\right)
=\Gamma\left(  \mathcal{W}\right)  \oplus\mathcal{W}^{c}%
\end{equation}
that are complementary to a subspace $\mathcal{W}$ \cite{Geroch}, where
$\mathcal{W}^{c}$ is any complimentary subspace and
\begin{align}
\tilde{\Gamma}  &  :\mathcal{W}^{c}\rightarrow\mathcal{B}_{1}\left(
\mathcal{H}^{N}\right) \nonumber\\
\tilde{\Gamma}\left(  X\right)   &  =X+\Gamma\left(  X\right)  ;\,X\in
\mathcal{W}^{c}%
\end{align}
and $\Gamma$ is an \emph{arbitrary} linear map%
\begin{equation}
\Gamma:\mathcal{W}^{c}\rightarrow\mathcal{W}%
\end{equation}
One can always choose $\mathcal{W}^{c}$ to be a canonical representative
subspace, $\mathcal{W}_{0}^{c}$, from this set that is constituted of vectors,
$X$, that have no components belonging to $\mathcal{W}$ i.e.
\begin{equation}
\mathcal{P}_{\mathcal{W}}\left(  X\right)  =0
\end{equation}
where $\mathcal{P}_{\mathcal{W}}$ is the projector onto $\mathcal{W}$ and
hence%
\begin{equation}
\mathcal{P}_{\mathcal{W}}\mathcal{P}_{\mathcal{W}_{0}^{c}}=0 \label{PP}%
\end{equation}
(As $\mathcal{B}_{1}\left(  \mathcal{H}^{N}\right)  $ is not an inner product
space the adjoint operation is not defined on super operators acting on
$\mathcal{B}_{1}\left(  \mathcal{H}^{N}\right)  $ so one cannot consider the
property of orthogonality as this property is expressed through the condition
eq. $\left(  \ref{PP}\right)  $ \emph{combined} \emph{with} the self
adjointness of projectors). All other complimentary subspaces $\left\{
\mathcal{W}_{\Gamma}^{c}\right\}  $ can be generated from the canonical one by
the map $\tilde{\Gamma}$ as
\begin{equation}
\mathcal{W}_{\Gamma}^{c}=\left\{  \tilde{\Gamma}\left(  X\right)  \equiv
X+\Gamma\left(  X\right)  |X\in\mathcal{W}_{0}^{c};\Gamma\left(  X\right)
\in\mathcal{W}\right\}
\end{equation}
The canonical complimentary space $\mathcal{W}_{0}^{c}$ is produced when
$\Gamma\equiv0$. The space $\mathcal{B}_{1}\left(  \mathcal{H}^{N}\right)  $
can be decomposed as the direct sum
\begin{equation}
\mathcal{B}_{1}\left(  \mathcal{H}^{N}\right)  =\mathcal{W\oplus W}_{\Gamma
}^{c}%
\end{equation}
where $\mathcal{W}$ is fixed and $\mathcal{W}_{\Gamma}^{c}$ is any of the
complimentary subspaces belonging to $\mathfrak{V}$. Elements $X\in
\mathcal{B}_{1}\left(  \mathcal{H}^{N}\right)  $ can be decomposed canonically
as
\begin{equation}
X=X_{\mathcal{W}}+X_{\mathcal{W}_{0}^{c}} \label{cana}%
\end{equation}
where $X_{\mathcal{W}}\in\mathcal{W}$ and $X_{\mathcal{W}_{0}^{c}}%
\in\mathcal{W}_{0}^{c}$ and decomposed in general wrt to $\mathcal{W}$ as
\begin{equation}
X=X_{\mathcal{W}_{\Gamma}}+X_{\mathcal{W}_{\Gamma}^{c}} \label{noncana}%
\end{equation}
with $X_{\mathcal{W}_{\Gamma}}\in\mathcal{W}$ and $X_{\mathcal{W}_{\Gamma}%
^{c}}\in\mathcal{W}_{\Gamma}^{c}\subset\mathcal{B}_{1}\left(  \mathcal{H}%
^{N}\right)  $.

These decompositions can be visualized more clearly in terms of the bases
$\left\{  \left\vert \mathbf{w}\right)  ,\left\vert \mathbf{w}_{0}^{c}\right)
\right\}  $ and $\left\{  \left\vert \mathbf{w}\right)  ,\left\vert
\mathbf{w}_{\Gamma}^{c}\right)  \right\}  $ of $\mathcal{B}_{1}\left(
\mathcal{H}^{N}\right)  $ (see subsection \ref{bases} for a detailed
discussion of super operator bases), where the bold symbols denote \emph{row}
vectors of basis elements i.e.%
\begin{align}
\left\vert \mathbf{w}\right)   &  \equiv\left\vert w_{1}\right)
,\ldots,\left\vert w_{d}\right) \nonumber\\
\left\vert \mathbf{w}_{0}^{c}\right)   &  \equiv\left\vert w_{1}^{c}\right)
,\ldots,\left\vert w_{d^{c}}^{c}\right) \nonumber\\
\left\vert \mathbf{w}_{\Gamma}^{c}\right)   &  \equiv\left\vert w_{1}%
^{c}\right)  ,\ldots,\left\vert w_{d^{c}}^{c}\right)  . \label{basis}%
\end{align}
and $d=\dim\left\{  \mathcal{W}\right\}  ;d^{c}=\dim\left\{  \mathcal{W}%
^{c}\right\}  =\dim\left\{  \mathcal{W}_{\Gamma}^{c}\right\}  $. The map
$\Gamma:\mathcal{W}^{c}\rightarrow\mathcal{W}$ can be expressed as
\begin{equation}
\Gamma\left(  \left\vert \mathbf{w}_{0}^{c}\right)  \right)  =\left(
\left\vert \mathbf{w}_{0}\right)  ,\left\vert \mathbf{w}_{0}^{c}\right)
\right)  \mathbf{\Gamma}=\left(  \left\vert \mathbf{w}_{0}\right)  ,\left\vert
\mathbf{w}_{0}^{c}\right)  \right)  \left(
\begin{array}
[c]{c}%
\mathbf{\Gamma}\\
\mathbb{O}%
\end{array}
\right)
\end{equation}
where $\mathbf{\Gamma}$ is a $d\times d^{c}$ matrix. The map, $\tilde{\Gamma}%
$, that defines the complimentary subspace $\mathcal{W}^{c}$ through the basis
$\left\{  \left\vert \mathbf{w}_{\Gamma}^{c}\right)  \right\}  $, is defined
as
\begin{equation}
\left(  \left\vert \mathbf{w}_{\Gamma}^{c}\right)  \right)  =\tilde{\Gamma
}\left(  \left\vert \mathbf{w}_{0}^{c}\right)  \right)  =\left(  \left\vert
\mathbf{w}\right)  ,\left\vert \mathbf{w}_{0}^{c}\right)  \right)
\mathbf{\tilde{\Gamma}}=\left(  \left\vert \mathbf{w}\right)  ,\left\vert
\mathbf{w}_{0}^{c}\right)  \right)  \left(
\begin{array}
[c]{c}%
\mathbf{\Gamma}\\
\mathbb{I}_{d^{c}}%
\end{array}
\right)
\end{equation}
The canonical expansion eq. $\left(  \ref{cana}\right)  $ is based on the
basis $\left\{  \left\vert \mathbf{w}\right)  ,\left\vert \mathbf{w}_{0}%
^{c}\right)  \right\}  $ while the non canonical expansion eq. $\left(
\ref{noncana}\right)  $ on the basis $\left\{  \left\vert \mathbf{w}_{\Gamma
}\right)  ,\left\vert \mathbf{w}_{\Gamma}^{c}\right)  \right\}  $ so that%
\begin{equation}
X=\left(  \left\vert \mathbf{w}\right)  ,\left\vert \mathbf{w}_{0}^{c}\right)
\right)  \left(
\begin{array}
[c]{c}%
\mathbb{X}_{W}\\
\mathbb{X}_{W_{0}^{c}}%
\end{array}
\right)  =\left(  \left\vert \mathbf{w}\right)  ,\left\vert \mathbf{w}%
_{\Gamma}^{c}\right)  \right)  \left(
\begin{array}
[c]{c}%
\mathbb{X}_{W_{\Gamma}}\\
\mathbb{X}_{W_{\Gamma}^{c}}%
\end{array}
\right)
\end{equation}
where $\mathbb{X}$ denotes vectors of expansion coefficients with respect to
the various bases. We thus have that
\begin{subequations}
\begin{align}
X  &  =\left(  \left\vert \mathbf{w}\right)  \right)  \mathbb{X}%
_{W}+\left\vert \mathbf{w}_{0}^{c}\right)  \mathbb{X}_{W_{0}^{c}}\\
X  &  =\left(  \left\vert \mathbf{w}\right)  \right)  \mathbb{X}_{W_{\Gamma}%
}+\left(  \left\vert \mathbf{w}_{\Gamma}^{c}\right)  \right)  \mathbb{X}%
_{W_{\Gamma}^{c}}\nonumber\\
&  =\left(  \left\vert \mathbf{w}\right)  \right)  \mathbb{X}_{W_{\Gamma}%
}+\left(  \left\vert \mathbf{w}\right)  \right)  \mathbf{\Gamma}%
\mathbb{X}_{W_{\Gamma}^{c}}+\left(  \left\vert \mathbf{w}_{0}^{c}\right)
\right)  \mathbb{X}_{W_{\Gamma}^{c}}%
\end{align}
comparing expansion coefficients of like bases elements one concludes that
\end{subequations}
\begin{equation}
\mathbb{X}_{W_{\Gamma}^{c}}=\mathbb{X}_{W_{0}^{c}}%
\end{equation}
and%
\begin{equation}
\mathbb{X}_{W_{\Gamma}}+\mathbf{\Gamma}\mathbb{X}_{W_{0}^{c}}=\mathbb{X}_{W}%
\end{equation}
hence%
\begin{equation}
\mathbb{X}_{W_{\Gamma}}=\mathbb{X}_{W}-\mathbf{\Gamma}\mathbb{X}_{W_{0}^{c}}%
\end{equation}
These expansion coefficient changes can be expressed in a basis free manner as%
\begin{align}
X_{\mathcal{W}_{\Gamma}}  &  \equiv\Pi_{\Gamma}\left(  X\right)
=X_{\mathcal{W}}-\Gamma\left(  X_{\mathcal{W}_{0}^{c}}\right) \nonumber\\
X_{\mathcal{W}_{\Gamma}^{c}}  &  \equiv\Pi_{\Gamma}^{c}\left(  X\right)
=X_{\mathcal{W}_{0}^{c}}+\Gamma\left(  X_{\mathcal{W}_{0}^{c}}\right)
\label{Csubspace}%
\end{align}
where we have introduced the projection maps $\Pi_{\Gamma}:\mathcal{B}%
_{1}\left(  \mathcal{H}^{N}\right)  \rightarrow\mathcal{W}$ and $\Pi_{\Gamma
}^{c}:\mathcal{B}_{1}\left(  \mathcal{H}^{N}\right)  \rightarrow
\mathcal{W}_{\Gamma}^{c}$. The matrix representation of these maps wrt the
basis eq. $\left(  \ref{basis}\right)  $ are
\begin{subequations}
\begin{align}
\Pi_{\Gamma}\left(  \left\vert \mathbf{w}_{0}\right)  ,\left\vert
\mathbf{w}_{0}^{c}\right)  \right)   &  =\left(  \left\vert \mathbf{w}%
_{0}\right)  ,\left\vert \mathbf{w}_{0}^{c}\right)  \right)  \left(
\begin{array}
[c]{cc}%
\mathbb{O} & -\mathbf{\Gamma}\\
\mathbb{O} & \mathbb{I}%
\end{array}
\right) \\
\Pi_{\Gamma}^{c}\left(  \left\vert \mathbf{w}_{0}\right)  ,\left\vert
\mathbf{w}_{0}^{c}\right)  \right)   &  =\left(  \left\vert \mathbf{w}%
_{0}\right)  ,\left\vert \mathbf{w}_{0}^{c}\right)  \right)  \left(
\begin{array}
[c]{cc}%
\mathbb{O} & \mathbf{\Gamma}\\
\mathbb{O} & \mathbb{I}%
\end{array}
\right)
\end{align}
and they form a resolution of the identity
\end{subequations}
\begin{equation}
I_{\mathcal{B}_{1}\left(  \mathcal{H}^{N}\right)  }=\Pi_{\Gamma}+\Pi_{\Gamma
}^{c}%
\end{equation}
Clearly we obtain the canonical decomposition when $\Gamma=0$ $\Leftrightarrow
$ $\Pi_{\Gamma}^{c}=P_{\mathcal{W}_{0}^{c}}$ - the projector onto
$\mathcal{W}_{0}^{c}$.

\subsection{Super Operator Bases \label{bases}}

\subsubsection{Proper Basis}

The subspaces $\mathcal{W}$ and $\mathcal{W}_{0}^{c}$ are contained in the
Hilbert-Schmidt Hilbert space $\mathcal{B}_{2}\left(  \mathcal{H}^{N}\right)
$ as well as the Banach space of Trace Class operators $\mathcal{B}_{1}\left(
\mathcal{H}^{N}\right)  $ and are orthogonal wrt the Hilbert-Schmidt inner
product i.e.
\begin{equation}
\left(  X|Y\right)  _{HS}=Tr\left\{  X^{\dagger}Y\right\}  =0\quad\forall
X\in\mathcal{W}\text{ and }\forall Y\in\mathcal{W}_{0}^{c}%
\end{equation}
and hence $\mathcal{W}_{0}^{c}\subseteq\mathcal{W}^{\bot}$ where
$\mathcal{W}^{\bot}$ is the Hilbert-Schmidt closed orthogonal compliment of
$\mathcal{W}$. The Hilbert-Schmidt closure of $\mathcal{W}_{0}^{c}$ is
$\mathcal{W}^{\bot}$ i.e. $\mathcal{W}_{0}^{c}$ is dense in $\mathcal{W}%
^{\bot}$ and, except for the finite dimensional case, contains elements that
\emph{are not }trace class operators. The closures of the subspaces
$\mathcal{W}_{\Gamma}^{c}$ in the $\left\Vert {}\right\Vert _{2}$- topology
produce non orthogonal compliments to $\mathcal{W}$ in $\mathcal{B}_{2}\left(
\mathcal{H}^{N}\right)  $.

The space $\mathcal{B}_{1}\left(  \mathcal{H}^{N}\right)  $ and thus the
spaces $\mathcal{B}\left(  \mathcal{H}^{N}\right)  $ and $\mathcal{B}%
_{2}\left(  \mathcal{H}^{N}\right)  $ contain the space, $\mathcal{F}$, of
finite rank operators defined by any complete orthonormal bases, $\left\{
\left\vert \Phi_{I}\right\rangle \right\}  $, of $\mathcal{H}^{N}$ as the set
of finite linear combinations of the ket-bra operators $\left\{  \left\vert
\Phi_{I}\right\rangle \left\langle \Phi_{J}\right\vert \right\}  $. The spaces
$\mathcal{B}\left(  \mathcal{H}^{N}\right)  $, $\mathcal{B}_{2}\left(
\mathcal{H}^{N}\right)  $ and $\mathcal{B}_{1}\left(  \mathcal{H}^{N}\right)
$ are generated as the $\left\Vert {}\right\Vert $, $\left\Vert {}\right\Vert
_{2}$, and $\left\Vert {}\right\Vert _{1}$ closures of $\mathcal{F}$. It can
be shown that the set of operators $\left\{  \left\vert \Phi_{I}\right\rangle
\left\langle \Phi_{J}\right\vert \right\}  $ form a basis for each of these
spaces, i.e. all operators in these three spaces have a unique expansion in
terms of these basis elements. One can transform basis of this type to ones
adapted to the decomposition
\begin{equation}
\mathcal{B}_{1}\left(  \mathcal{H}^{N}\right)  =\mathcal{W\oplus W}_{0}^{c}%
\end{equation}
(However one should note that the $\left\Vert {}\right\Vert _{2}-$closures of
these bases of $\mathcal{W}$ and $\mathcal{W}_{0}^{c}$ contain elements that
have infinite trace).

\subsubsection{Generalized Bases for $\mathcal{B}\left(  \mathcal{H}%
^{N}\right)  $ and $\mathcal{B}_{1}\left(  \mathcal{H}^{N}\right)  $}

The duality between $\mathcal{B}_{1}\left(  \mathcal{H}^{N}\right)  $ and
$\mathcal{B}\left(  \mathcal{H}^{N}\right)  $ can be expressed in terms of the
scalar product
\begin{equation}
\left(  X|Y\right)  =Tr\left\{  X^{\dagger}Y\right\}  ;\quad X\in
\mathcal{B}\left(  \mathcal{H}^{N}\right)  ,Y\in\mathcal{B}_{1}\left(
\mathcal{H}^{N}\right)
\end{equation}
Following the Dirac convention we can denote operators, $Y\in\mathcal{B}%
_{1}\left(  \mathcal{H}^{N}\right)  $, as $\left\vert Y\right)  $ and
operators, $X\in\mathcal{B}\left(  \mathcal{H}^{N}\right)  $, as $\left(
X\right\vert $ thus producing the super operator%
\begin{equation}
\left\vert Y\right)  \left(  X\right\vert :\mathcal{B}_{1}\left(
\mathcal{H}^{N}\right)  \rightarrow\mathcal{B}_{1}\left(  \mathcal{H}%
^{N}\right)
\end{equation}
whose action is given by
\begin{equation}
\left\vert Y\right)  \left(  X\right\vert \,Z=Tr\left\{  X^{\dagger}Z\right\}
\left\vert Y\right)  \equiv Tr\left\{  X^{\dagger}Z\right\}  Y
\end{equation}


We can introduce a generalized basis, $\left\{  \left\vert \mathsf{z}%
^{N}\right\rangle \left\langle \mathsf{z}^{\prime N}\right\vert \right\}  $,
\emph{for both} of the dual spaces $\mathcal{B}\left(  \mathcal{H}^{N}\right)
$ and $\mathcal{B}_{1}\left(  \mathcal{H}^{N}\right)  $, composed of ket-bra
operators that correspond to the generalized eigenvectors of the $N$-particle
space-spin coordinate operator. None of the basis elements $\left\vert
\mathsf{z}^{N}\right\rangle \left\langle \mathsf{z}^{\prime N}\right\vert $
belong to $\mathcal{B}\left(  \mathcal{H}^{N}\right)  $ as the action of any
individual $\left\vert \mathsf{z}^{N}\right\rangle \left\langle \mathsf{z}%
^{\prime N}\right\vert $ on any vector belonging to $\mathcal{H}^{N}$ produces
a vector $\notin\mathcal{H}^{N}$. The notation $\left\vert \left\vert
\mathsf{z}^{N}\right\rangle \left\langle \mathsf{z}^{\prime N}\right\vert
\right)  $ signifies that $\left\vert \mathsf{z}^{N}\right\rangle \left\langle
\mathsf{z}^{\prime N}\right\vert $ is to be considered within the
$\mathcal{B}_{1}\left(  \mathcal{H}^{N}\right)  $-topology, while the notation
$\left(  \left\vert \mathsf{z}^{N}\right\rangle \left\langle \mathsf{z}%
^{\prime N}\right\vert \right\vert $ signifies that $\left\vert \mathsf{z}%
^{N}\right\rangle \left\langle \mathsf{z}^{\prime N}\right\vert $ is to be
considered within the $\mathcal{B}\left(  \mathcal{H}^{N}\right)  $-topology.

This basis allows one to define a super operator resolution of the identity on
$\mathcal{B}_{1}\left(  \mathcal{H}^{N}\right)  $ by
\begin{equation}
\hat{I}_{\mathcal{B}_{1}\left(  \mathcal{H}^{N}\right)  }=%
%TCIMACRO{\dint }%
%BeginExpansion
{\displaystyle\int}
%EndExpansion
\left\vert \left\vert \mathsf{z}^{N}\right\rangle \left\langle \mathsf{z}%
^{\prime N}\right\vert \right)  \left(  \left\vert \mathsf{z}^{N}\right\rangle
\left\langle \mathsf{z}^{\prime N}\right\vert \right\vert d\mathsf{z}%
^{N}d\mathsf{z}^{\prime N},
\end{equation}
which can be expressed symbolically as
\begin{equation}
\hat{I}_{\mathcal{B}_{1}\left(  \mathcal{H}^{N}\right)  }=\left(  \left(
\underline{\left\vert \mathsf{z}^{N}\right\rangle \left\langle \mathsf{z}%
^{\prime N}\right\vert }\right\vert \right)  \left(  \left\vert \underline
{\left\vert \mathsf{z}^{N}\right\rangle \left\langle \mathsf{z}^{\prime
N}\right\vert }\right)  \right)  ^{t},
\end{equation}
producing the generalized expansion of an element $X\in\mathcal{B}_{1}\left(
\mathcal{H}^{N}\right)  $ as%
\begin{align}
X  &  =%
%TCIMACRO{\dint }%
%BeginExpansion
{\displaystyle\int}
%EndExpansion
\left\vert \left\vert \mathsf{z}^{N}\right\rangle \left\langle \mathsf{z}%
^{\prime N}\right\vert \right)  \left(  \left\vert \mathsf{z}^{N}\right\rangle
\left\langle \mathsf{z}^{\prime N}\right\vert |X\right)  d\mathsf{z}%
^{N}d\mathsf{z}^{\prime N}\nonumber\\
&  =%
%TCIMACRO{\dint }%
%BeginExpansion
{\displaystyle\int}
%EndExpansion
\left\vert \left\vert \mathsf{z}^{N}\right\rangle \left\langle \mathsf{z}%
^{\prime N}\right\vert \right)  Tr\left\{  \left\vert \mathsf{z}%
^{N}\right\rangle \left\langle \mathsf{z}^{\prime N}\right\vert ^{\dagger
}X\right\}  d\mathsf{z}^{N}d\mathsf{z}^{\prime N}\nonumber\\
&  =%
%TCIMACRO{\dint }%
%BeginExpansion
{\displaystyle\int}
%EndExpansion
\left\vert \left\vert \mathsf{z}^{N}\right\rangle \left\langle \mathsf{z}%
^{\prime N}\right\vert \right)  \left\langle \mathsf{z}^{N}|X\mathsf{z}%
^{\prime N}\right\rangle d\mathsf{z}^{N}d\mathsf{z}^{\prime N}\nonumber\\
&  =%
%TCIMACRO{\dint }%
%BeginExpansion
{\displaystyle\int}
%EndExpansion
X\left(  \mathsf{z}^{N},\mathsf{z}^{\prime N}\right)  \left\vert
\mathsf{z}^{N}\right\rangle \left\langle \mathsf{z}^{\prime N}\right\vert
d\mathsf{z}^{N}d\mathsf{z}^{\prime N}%
\end{align}
where $X\left(  \mathsf{z}^{N},\mathsf{z}^{\prime N}\right)  $ is the integral
kernel corresponding to $X$.

The basis $\left\{  \left\vert \mathsf{z}^{N}\right\rangle \left\langle
\mathsf{z}^{N\prime}\right\vert \right\}  $ can be graded into sub bases
$\left\{  \mathfrak{b}_{m}\right\}  $ which are defined as
\begin{equation}
\mathfrak{b}_{m}=\left\{  \left\vert \mathsf{z}^{N}\right\rangle \left\langle
\mathsf{z}^{N\prime}\right\vert |\mathsf{z}^{N}\overset{m}{=}\mathsf{z}%
^{N\prime}\right\}
\end{equation}
where $\mathsf{z}^{N}\overset{m}{=}\mathsf{z}^{N\prime}$ denotes that the
coordinates $\mathsf{z}^{N}$ and $\mathsf{z}^{N\prime}$ have exactly $m$
values in common, which could be equivalently written as $\mathsf{z}%
^{N}\overset{N-m}{\neq}\mathsf{z}^{N\prime}$ i.e. $\mathsf{z}^{N}$ and
$\mathsf{z}^{N\prime}$ have exactly $N-m$ values different. These bases
produce a graded decomposition of $\mathcal{B}_{1}\left(  \mathcal{H}%
^{N}\right)  $ as the direct sum
\begin{equation}
\mathcal{B}_{1}\left(  \mathcal{H}^{N}\right)  =%
%TCIMACRO{\dbigoplus \limits_{0\leq m\leq N}}%
%BeginExpansion
{\displaystyle\bigoplus\limits_{0\leq m\leq N}}
%EndExpansion
\mathfrak{B}_{m}%
\end{equation}
where $\mathfrak{B}_{m}$ is defined as the closed linear span of the basis
$\left\{  \mathfrak{b}_{m}\right\}  $ i.e.
\begin{equation}
\mathfrak{B}_{m}=\overline{\operatorname{linspan}}\left\{  \left.  \underset
{}{\left\vert \mathsf{z}^{N}\right\rangle \left\langle \mathsf{z}^{N\prime
}\right\vert }\right\vert \left\vert \mathsf{z}^{N}\right\rangle \left\langle
\mathsf{z}^{N\prime}\right\vert \in\mathfrak{b}_{m}\right\}
\end{equation}
and the closure is taken wrt the trace norm topology. In order to analyze the
kernels of the contraction maps $\left\{  L_{N}^{p}|0\leq p\leq N\right\}  $
and in particular $L_{N}^{1}$ it is useful to consider the sequence of
subspaces $\left\{  \mathfrak{V}_{n}\right\}  $, that are generated by
$\left\{  \left\vert \mathsf{z}^{N}\right\rangle \left\langle \mathsf{z}%
^{N\prime}\right\vert \right\}  $ that have at least $n$ differences between
$\mathsf{z}^{N}$ and $\mathsf{z}^{\prime N}$ (denoted by $\mathsf{z}%
^{N}\overset{n}{\nsim}$ $\mathsf{z}^{\prime N}$) or equivalently as generated
by $\left\{  \left\vert \mathsf{z}^{N}\right\rangle \left\langle
\mathsf{z}^{N\prime}\right\vert \right\}  $ that have at least $N-n$ values in
common between $\mathsf{z}^{N}$ and $\mathsf{z}^{\prime N}$ (denoted by
$\mathsf{z}^{N}\overset{N-n}{\sim}$ $\mathsf{z}^{\prime N})$, which are
defined as
\begin{equation}
\mathfrak{V}_{n}=%
%TCIMACRO{\dbigoplus \limits_{0\leq m\leq N-n}}%
%BeginExpansion
{\displaystyle\bigoplus\limits_{0\leq m\leq N-n}}
%EndExpansion
\mathfrak{B}_{m}%
\end{equation}
so that
\begin{equation}
\mathfrak{V}_{n}=\overline{\operatorname{linspan}}\left\{  \left.  \left\vert
\mathsf{z}^{N}\right\rangle \left\langle \mathsf{z}^{N\prime}\right\vert
\overset{}{}\right\vert \mathsf{z}^{N}\overset{n}{\nsim}\mathsf{z}^{N\prime
}\right\}  =\overline{\operatorname{linspan}}\left\{  \mathfrak{v}%
_{n}\right\}
\end{equation}
where $\mathfrak{v}_{n}\equiv\left\{  \left.  \underset{}{\left\vert
\mathsf{z}^{N}\right\rangle \left\langle \mathsf{z}^{N\prime}\right\vert
}\right\vert \mathsf{z}^{N}\overset{n}{\nsim}\mathsf{z}^{N\prime}\right\}  $
$\equiv\left\{  \left.  \underset{}{\left\vert \mathsf{z}^{N}\right\rangle
\left\langle \mathsf{z}^{N\prime}\right\vert }\right\vert \mathsf{z}%
^{N}\overset{N-n}{=}\mathsf{z}^{N\prime}\right\}  $ and hence $\mathfrak{V}%
_{0}=\mathcal{B}_{1}\left(  \mathcal{H}^{N}\right)  $. These subspaces
partially characterize the kernel of the maps $\left\{  L_{N}^{p}\right\}  $
as they have the property
\begin{equation}
\mathfrak{V}_{0}\supset\ker L_{N}^{p}\supset\mathfrak{V}_{p}.
\end{equation}


In particular one can use the basis $\mathfrak{v}_{2}$ of $\mathfrak{V}_{2}$
to construct a new mixed generalized-proper basis of $\mathcal{B}_{1}\left(
\mathcal{H}^{N}\right)  $ that is adapted to the decomposition eq. $\left(
\ref{orthog}\right)  $. There are three types of elements in this adapted basis:

\begin{enumerate}
\item The generalized operators $\mathfrak{v}_{2}=\left\{  \left\vert
\mathsf{z}^{N}\right\rangle \left\langle \mathsf{z}^{\prime N}\right\vert
\left\vert \mathsf{z}^{N}\overset{2}{\nsim}\mathsf{z}^{\prime N}\right.
\right\}  $. These span the space $\mathfrak{V}_{2}$ but which we will denote
as $\mathcal{K}_{OD}$ in this context. This basis could be changed into a
proper basis by defining trace class "wavepackets"
\[
l_{\mu}=%
%TCIMACRO{\dint }%
%BeginExpansion
{\displaystyle\int}
%EndExpansion
l_{\mu}\left(  \mathsf{z}^{N};\mathsf{z}^{\prime N}\right)  \left\vert
\mathsf{z}^{N}\right\rangle \left\langle \mathsf{z}^{\prime N}\right\vert
d\mathsf{z}^{N}d\mathsf{z}^{\prime N};\qquad\mathsf{z}^{N}\overset{2}{\nsim
}\mathsf{z}^{\prime N}%
\]
but for the purposes of this article we find it more convenient to stay with
the generalized basis.

\item The proper trace class operators that are elements of $\ker L_{N}^{1}$
\begin{align*}
&  \left\{  k_{\kappa}=%
%TCIMACRO{\dint }%
%BeginExpansion
{\displaystyle\int}
%EndExpansion
k_{\kappa}\left(  \mathsf{z}_{1}\mathsf{z}^{N-1};\mathsf{z}_{1}^{\prime
}\mathsf{z}^{N-1}\right)  \left\vert \mathsf{z}_{1}\mathsf{z}^{N-1}%
\right\rangle \left\langle \mathsf{z}_{1}^{\prime}\mathsf{z}^{N-1}\right\vert
d\mathsf{z}_{1}d\mathsf{z}_{1}^{\prime}d\mathsf{z}^{N-1}\right. \\
&  \qquad\qquad\left.  \left\vert
%TCIMACRO{\dint }%
%BeginExpansion
{\displaystyle\int}
%EndExpansion
k_{\kappa}\left(  \mathsf{z}_{1}\mathsf{z}^{N-1};\mathsf{z}_{1}^{\prime
}\mathsf{z}^{N-1}\right)  d\mathsf{z}^{N-1}\right.  \overset{a.e}{=}%
0\quad\forall\mathsf{z}_{1}\mathsf{z}_{1}^{\prime};1\leq\kappa\leq
\infty\overset{}{}\right\}
\end{align*}


which span a space which we will denote $\mathcal{K}_{D}$.

\item The proper trace class operators that are elements of $\left[  \ker
L_{N}^{1}\right]  ^{c}$%
\begin{align*}
&  \left\{  \mathfrak{x}_{\lambda}=%
%TCIMACRO{\dint }%
%BeginExpansion
{\displaystyle\int}
%EndExpansion
\mathfrak{x}_{\lambda}\left(  \mathsf{z}_{1}\mathsf{z}^{N-1};\mathsf{z}%
_{1}^{\prime}\mathsf{z}^{N-1}\right)  \left\vert \mathsf{z}_{1}\mathsf{z}%
^{N-1}\right\rangle \left\langle \mathsf{z}_{1}^{\prime}\mathsf{z}%
^{N-1}\right\vert d\mathsf{z}_{1}d\mathsf{z}_{1}^{\prime}d\mathsf{z}%
^{N-1}|\right. \\
&  \qquad\qquad\qquad\left.  \int\mathfrak{x}_{\lambda}\left(  \mathsf{z}%
_{1}\mathsf{z}^{N-1};\mathsf{z}_{1}^{\prime}\mathsf{z}^{N-1}\right)
d\mathsf{z}^{N-1}\neq0\quad\forall\mathsf{z}_{1}\mathsf{z}_{1}^{\prime}%
\in\Delta;1\leq\lambda\leq\infty\overset{}{}\right\}
\end{align*}


where $\Delta$ is a set of non zero measure. These span the space $\left[
\ker L_{N}^{1}\right]  ^{c}$.
\end{enumerate}

The basis elements of type 1 and type 2 together clearly span $\ker L_{N}^{1}
$, which can be expressed as the direct sum%

\begin{equation}
\ker L_{N}^{1}=\mathcal{K}_{OD}\oplus\mathcal{K}_{D}.
\end{equation}
$\mathcal{K}_{OD}$ can be further decomposed as direct a integral
\cite{dirint} of subspaces, $\mathcal{K}_{\left\vert \mathsf{z}_{1}%
\mathsf{z}^{N-1}\right\rangle \left\langle \mathsf{z}_{1}^{\prime}%
\mathsf{z}^{\prime N-1}\right\vert }$, indexed by the generalized basis
elements $\left\{  \left\vert \mathsf{z}^{N}\right\rangle \left\langle
\mathsf{z}^{\prime N}\right\vert \right\}  $, where $\mathsf{z}^{N}\overset
{2}{\nsim}\mathsf{z}^{\prime N}$ i.e.
\begin{equation}
\mathcal{K}_{OD}\equiv\mathfrak{V}_{2}=\int_{\oplus}\mathcal{K}_{\left\vert
\mathsf{z}^{N}\right\rangle \left\langle \mathsf{z}^{\prime N}\right\vert
}d\mathsf{z}^{N}d\mathsf{z}^{\prime N};\quad\mathsf{z}^{N}\overset{2}{\nsim
}\mathsf{z}^{\prime N}.
\end{equation}


One can define a basis of the dual space $\mathcal{B}\left(  \mathcal{H}%
^{N}\right)  $ that is biorthogonal to the basis of $\mathcal{B}_{1}\left(
\mathcal{H}^{N}\right)  $ just described by%
\begin{equation}
\left\{  \left\{  \left\vert \mathsf{z}^{N}\right\rangle \left\langle
\mathsf{z}^{\prime N}\right\vert \left\vert \mathsf{z}^{N}\overset{2}{\nsim
}\mathsf{z}^{\prime N}\right.  \right\}  ,\left\{  k_{\kappa}^{d}\right\}
,\left\{  \mathfrak{x}_{\lambda}^{d}\right\}  \right\}
\end{equation}
where
\begin{align}
\left(  k_{\kappa}^{d}|k_{\lambda}\right)   &  =\delta_{\kappa\lambda
}\nonumber\\
\left(  \mathfrak{x}_{\kappa}^{d}|\mathfrak{x}_{\lambda}\right)   &
=\delta_{\kappa\lambda}%
\end{align}
and
\begin{equation}
\left(  k_{\kappa}^{d}|b\right)  =\left(  \mathfrak{x}_{\kappa}^{d}|b\right)
=0
\end{equation}
for all other basis elements $\left\{  \left\vert b\right)  \right\}  $ of
$\mathcal{B}_{1}\left(  \mathcal{H}^{N}\right)  .$ The dual basis elements of
$\mathcal{B}\left(  \mathcal{H}^{N}\right)  $ are explicitly expressed as:

\begin{enumerate}
\item The same generalized operators $\left\{  \left\vert \mathsf{z}%
^{N}\right\rangle \left\langle \mathsf{z}^{\prime N}\right\vert \left\vert
\mathsf{z}^{N}\overset{2}{\nsim}\mathsf{z}^{\prime N}\right.  \right\}  $.
These span a space which we will denote as $\mathcal{K}_{OD}^{d}$. This basis
could be changed into a proper basis by defining bounded operator
"wavepackets"
\begin{equation}
l_{\mu}^{d}=%
%TCIMACRO{\dint }%
%BeginExpansion
{\displaystyle\int}
%EndExpansion
l_{\mu}^{d}\left(  \mathsf{z}^{N};\mathsf{z}^{\prime N}\right)  \left\vert
\mathsf{z}^{N}\right\rangle \left\langle \mathsf{z}^{\prime N}\right\vert
d\mathsf{z}^{N}d\mathsf{z}^{\prime N};\qquad\mathsf{z}^{N}\overset{2}{\nsim
}\mathsf{z}^{\prime N}%
\end{equation}
but again for the purposes of this article we find it more convenient to stay
with the generalized basis.

\item The proper bounded operators%
\begin{align}
&  \left\{  k_{\kappa}^{d}=%
%TCIMACRO{\dint }%
%BeginExpansion
{\displaystyle\int}
%EndExpansion
k_{\kappa}^{d}\left(  \mathsf{z}_{1}\mathsf{z}^{N-1};\mathsf{z}_{1}^{\prime
}\mathsf{z}^{N-1}\right)  \left\vert \mathsf{z}_{1}\mathsf{z}^{N-1}%
\right\rangle \left\langle \mathsf{z}_{1}^{\prime}\mathsf{z}^{N-1}\right\vert
d\mathsf{z}_{1}d\mathsf{z}_{1}^{\prime}d\mathsf{z}^{N-1}\right. \nonumber\\
&  \qquad\qquad\qquad\left.  \left\vert
%TCIMACRO{\dint }%
%BeginExpansion
{\displaystyle\int}
%EndExpansion
k_{\kappa}^{d}\left(  \mathsf{z}_{1}\mathsf{z}^{N-1};\mathsf{z}_{1}^{\prime
}\mathsf{z}^{N-1}\right)  d\mathsf{z}^{N-1}\right.  \overset{a.e}{=}%
0\quad\forall\mathsf{z}_{1}\mathsf{z}_{1}^{\prime};1\leq\kappa\leq
\infty\overset{}{}\right\}
\end{align}


which span a space which we will denote $\mathcal{K}_{D}^{d}$.

\item The proper bounded operators%
\begin{align}
&  \left\{  \mathfrak{x}_{\lambda}^{d}=%
%TCIMACRO{\dint }%
%BeginExpansion
{\displaystyle\int}
%EndExpansion
\mathfrak{x}_{\lambda}^{d}\left(  \mathsf{z}_{1}\mathsf{z}^{N-1}%
;\mathsf{z}_{1}^{\prime}\mathsf{z}^{N-1}\right)  \left\vert \mathsf{z}%
_{1}\mathsf{z}^{N-1}\right\rangle \left\langle \mathsf{z}_{1}^{\prime
}\mathsf{z}^{N-1}\right\vert d\mathsf{z}_{1}d\mathsf{z}_{1}^{\prime
}d\mathsf{z}^{N-1}|\right. \nonumber\\
&  \qquad\qquad\qquad\left.  \int\mathfrak{x}_{\lambda}^{d}\left(
\mathsf{z}_{1}\mathsf{z}^{N-1};\mathsf{z}_{1}^{\prime}\mathsf{z}^{N-1}\right)
d\mathsf{z}^{N-1}\neq0\quad\forall\mathsf{z}_{1}\mathsf{z}_{1}^{\prime}%
\in\Delta;1\leq\lambda\leq\infty\overset{}{}\right\}
\end{align}


where $\Delta$ is a set of non zero measure. These span the space $\left[
\ker L_{N}^{1}\right]  ^{cd}$. It \emph{should} be noted that in general%
\begin{equation}
k_{\kappa}^{d}\left(  \mathsf{z}_{1}\mathsf{z}^{N-1};\mathsf{z}_{1}^{\prime
}\mathsf{z}^{N-1}\right)  \neq k_{\kappa}\left(  \mathsf{z}_{1}\mathsf{z}%
^{N-1};\mathsf{z}_{1}^{\prime}\mathsf{z}^{N-1}\right)  ^{\ast}%
\end{equation}
and
\begin{equation}
\mathfrak{x}_{\lambda}^{d}\left(  \mathsf{z}_{1}\mathsf{z}^{N-1}%
;\mathsf{z}_{1}^{\prime}\mathsf{z}^{N-1}\right)  \neq\mathfrak{x}_{\lambda
}\left(  \mathsf{z}_{1}\mathsf{z}^{N-1};\mathsf{z}_{1}^{\prime}\mathsf{z}%
^{N-1}\right)  ^{\ast}.
\end{equation}

\end{enumerate}

The biorthogonality of the preceding bases can be displayed symbolically as%
\begin{equation}
\left(
\begin{array}
[c]{c}%
\left(  \underline{\left\vert \mathsf{z}^{N}\right\rangle _{2}\left\langle
\mathsf{z}^{\prime N}\right\vert }\right\vert \\
\left(  \underline{k^{d}}\right\vert \\
\left(  \underline{\mathfrak{x}^{d}}\right\vert
\end{array}
\right)  \left(  \left\vert \underline{\left\vert \mathsf{z}^{N}\right\rangle
_{2}\left\langle \mathsf{z}^{\prime N}\right\vert }\right)  ,\underline
{\left\vert k\right)  },\underline{\left\vert \mathfrak{x}^{d}\right)
}\right)  =\left(
\begin{array}
[c]{ccc}%
\mathbb{I}_{\mathcal{K}_{OD}} & \mathbb{O} & \mathbb{O}\\
\mathbb{O} & \mathbb{I}_{\mathcal{K}_{D}} & \mathbb{O}\\
\mathbb{O} & \mathbb{O} & \mathbb{I}_{\left[  \ker L_{N}^{1}\right]  ^{c}}%
\end{array}
\right)
\end{equation}
This basis also provides us with a resolution of the identity symbolically
expressed as%
\begin{equation}
\hat{I}_{\mathcal{B}_{1}\left(  \mathcal{H}^{N}\right)  }=\left(
\underline{\left\vert \left\vert \mathsf{z}^{N}\right\rangle _{2}\left\langle
\mathsf{z}^{\prime N}\right\vert \right)  },\underline{\left\vert k\right)
},\underline{\left\vert \mathfrak{x}\right)  }\right)  \left(
\begin{array}
[c]{c}%
\underline{\left(  \left\vert \mathsf{z}^{N}\right\rangle _{2}\left\langle
\mathsf{z}^{\prime N}\right\vert \right\vert }\\
\underline{\left(  k^{d}\right\vert }\\
\underline{\left(  \mathfrak{x}^{d}\right\vert }%
\end{array}
\right)  \label{rescan}%
\end{equation}
where $\left\vert \left\vert \mathsf{z}^{N}\right\rangle _{2}\left\langle
\mathsf{z}^{\prime N}\right\vert \right)  $ denotes the basis $\left\{
\left\vert \mathsf{z}^{N}\right\rangle \left\langle \mathsf{z}^{\prime
N}\right\vert \left\vert \mathsf{z}^{N}\overset{2}{\nsim}\mathsf{z}^{\prime
N}\right.  \right\}  $.

Any element $X\in\mathcal{B}_{1}\left(  \mathcal{H}^{N}\right)  $ can thus be
expanded as
\begin{equation}
X=X_{\mathcal{K}_{OD}}+X_{\mathcal{K}_{D}}+X_{\left[  \ker L_{N}^{1}\right]
^{c}} \label{kerdec}%
\end{equation}
where
\begin{align}
X_{\mathcal{K}_{OD}}  &  =\int\left(  \left\vert \mathsf{z}^{N}\right\rangle
\left\langle \mathsf{z}^{\prime N}\right\vert |X\right)  \left\vert
\mathsf{z}^{N}\right\rangle \left\langle \mathsf{z}^{\prime N}\right\vert
d\mathsf{z}^{N}d\mathsf{z}^{\prime N}\nonumber\\
&  \int Tr\left\{  \left\vert \mathsf{z}^{N}\right\rangle \left\langle
\mathsf{z}^{\prime N}\right\vert ^{\dagger}X\right\}  \left\vert
\mathsf{z}^{N}\right\rangle \left\langle \mathsf{z}^{\prime N}\right\vert
d\mathsf{z}^{N}d\mathsf{z}^{\prime N}\nonumber\\
&  =\int\left\langle \mathsf{z}^{N}|X\mathsf{z}^{\prime N}\right\rangle
\left\vert \mathsf{z}^{N}\right\rangle \left\langle \mathsf{z}^{\prime
N}\right\vert d\mathsf{z}^{N}d\mathsf{z}^{\prime N}\nonumber\\
&  =\int X\left(  \mathsf{z}^{N},\mathsf{z}^{\prime N}\right)  \left\vert
\mathsf{z}^{N}\right\rangle \left\langle \mathsf{z}^{\prime N}\right\vert
d\mathsf{z}^{N}d\mathsf{z}^{\prime N};\quad\mathsf{z}^{N}\overset{2}{\nsim
}\mathsf{z}^{\prime N}%
\end{align}%
\begin{equation}
\qquad X_{\mathcal{K}_{D}}=%
%TCIMACRO{\dsum \limits_{\kappa}}%
%BeginExpansion
{\displaystyle\sum\limits_{\kappa}}
%EndExpansion
\left(  k_{\kappa}^{d}|X\right)  k_{\kappa}=%
%TCIMACRO{\dsum \limits_{\kappa}}%
%BeginExpansion
{\displaystyle\sum\limits_{\kappa}}
%EndExpansion
Tr\left\{  k_{\kappa}^{d\dagger}X\right\}  k_{\kappa}\qquad\qquad\qquad
\end{equation}
and
\begin{equation}
X_{\left[  \ker L_{N}^{1}\right]  ^{c}}=%
%TCIMACRO{\dsum \limits_{\kappa}}%
%BeginExpansion
{\displaystyle\sum\limits_{\kappa}}
%EndExpansion
\left(  \mathfrak{x}_{\lambda}^{d}|X\right)  \mathfrak{x}_{\lambda}=%
%TCIMACRO{\dsum \limits_{\kappa}}%
%BeginExpansion
{\displaystyle\sum\limits_{\kappa}}
%EndExpansion
Tr\left\{  \mathfrak{x}_{\lambda}^{d\dagger}X\right\}  \mathfrak{x}_{\lambda
}\qquad\qquad\qquad
\end{equation}
When analyzing the properties of an operator $X\in\mathcal{B}_{1}\left(
\mathcal{H}^{N}\right)  $ it behooves us to both consider its expansion wrt
the subspace decomposition $\mathcal{B}_{1}\left(  \mathcal{H}^{N}\right)  $
induced by $\ker L_{N}^{1}$ as in eq. $\left(  \ref{kerdec}\right)  $ and its
expansion wrt to subspace decomposition of $\mathcal{B}_{1}\left(
\mathcal{H}^{N}\right)  $ given by $\left\{  \mathfrak{B}_{m}|0\leq m\leq
N\right\}  $ i.e.
\begin{equation}
X=\sum_{0\leq m\leq N}X_{\mathfrak{B}_{m}}%
\end{equation}
where
\begin{equation}
X_{\mathfrak{B}_{m}}=%
%TCIMACRO{\dint \limits_{\mathsf{z}^{\prime N}\,\overset{m}{=}\,\mathsf{z}^{N}%
%}}%
%BeginExpansion
{\displaystyle\int\limits_{\mathsf{z}^{\prime N}\,\overset{m}{=}%
\,\mathsf{z}^{N}}}
%EndExpansion
X_{m}\left(  \mathsf{z}^{\prime N},\mathsf{z}^{N}\right)  \left\vert
\mathsf{z}^{\prime N}\right\rangle \left\langle \mathsf{z}^{N}\right\vert
d\mathsf{z}^{N}d\mathsf{z}^{\prime N}%
\end{equation}
Explicitly using the resolution of the identity for the basis $\left\{
\underline{\left\vert \left\vert \mathsf{z}^{N}\right\rangle _{2}\left\langle
\mathsf{z}^{\prime N}\right\vert \right)  },\underline{\left\vert k\right)
},\underline{\left\vert \mathfrak{x}\right)  }\right\}  $ we obtain%
\begin{align}
X  &  =\left(  \underline{\left\vert \left\vert \mathsf{z}^{N}\right\rangle
_{2}\left\langle \mathsf{z}^{\prime N}\right\vert \right)  },\underline
{\left\vert k\right)  },\underline{\left\vert \mathfrak{x}\right)  }\right)
\left(
\begin{array}
[c]{c}%
\underline{\left(  \left\vert \mathsf{z}^{N}\right\rangle _{2}\left\langle
\mathsf{z}^{\prime N}\right\vert \right\vert }\\
\underline{\left(  k^{d}\right\vert }\\
\underline{\left(  \mathfrak{x}^{d}\right\vert }%
\end{array}
\right)  \left\vert X\right) \nonumber\\
&  =\left(  \underline{\left\vert \left\vert \mathsf{z}^{N}\right\rangle
_{2}\left\langle \mathsf{z}^{\prime N}\right\vert \right)  },\underline
{\left\vert k\right)  },\underline{\left\vert \mathfrak{x}\right)  }\right)
\mathbb{X}_{can}%
\end{align}
and for the basis $\left\{  \mathfrak{b}_{0},\cdots,\mathfrak{b}_{N}\right\}
$%
\begin{equation}
X=\left(  \mathfrak{b}_{0},\cdots,\mathfrak{b}_{N}\right)  \left(
\begin{array}
[c]{c}%
\mathfrak{b}_{0}\\
\vdots\\
\mathfrak{b}_{N}%
\end{array}
\right)  \left\vert X\right)  =\left(  \mathfrak{b}_{0},\cdots,\mathfrak{b}%
_{N}\right)  \mathbb{X}_{\mathcal{B}_{1}\left(  \mathcal{H}^{N}\right)  }%
\end{equation}
leading to the transformation from the representation $\mathbb{X}_{can}$ to
$\mathbb{X}_{\mathcal{B}_{1}\left(  \mathcal{H}^{N}\right)  }$ given as
\begin{align}
\mathbb{X}_{can}  &  =\left(
\begin{array}
[c]{c}%
\underline{\left(  \left\vert \mathsf{z}^{N}\right\rangle _{2}\left\langle
\mathsf{z}^{\prime N}\right\vert \right\vert }\\
\underline{\left(  k^{d}\right\vert }\\
\underline{\left(  \mathfrak{x}^{d}\right\vert }%
\end{array}
\right)  \left(  \mathfrak{b}_{0},\cdots,\mathfrak{b}_{N}\right)
\mathbb{X}_{\mathcal{B}_{1}\left(  \mathcal{H}^{N}\right)  }\nonumber\\
&  =\left(
\begin{array}
[c]{ccc}%
\mathbb{I} & \mathbb{O} & \mathbb{O}\\
\mathbb{O} & \left(  \left.  \underline{\left(  k^{d}\right\vert }\right\vert
\mathfrak{b}_{N-1}\right)  & \left(  \left.  \underline{\left(  k^{d}%
\right\vert }\right\vert \mathfrak{b}_{N}\right) \\
\mathbb{O} & \left(  \left.  \underline{\left(  \mathfrak{x}^{d}\right\vert
}\right\vert \mathfrak{b}_{N-1}\right)  & \left(  \left.  \underline{\left(
\mathfrak{x}^{d}\right\vert }\right\vert \mathfrak{b}_{N}\right)
\end{array}
\right)  \mathbb{X}_{\mathcal{B}_{1}\left(  \mathcal{H}^{N}\right)  }%
\end{align}
and the transformation from the representation $\mathbb{X}_{\mathcal{B}%
_{1}\left(  \mathcal{H}^{N}\right)  }$ to $\mathbb{X}_{can}$ as
\begin{align}
\mathbb{X}_{\mathcal{B}_{1}\left(  \mathcal{H}^{N}\right)  }  &  =\left(
\begin{array}
[c]{c}%
\mathfrak{b}_{0}\\
\vdots\\
\mathfrak{b}_{N}%
\end{array}
\right)  \left(  \underline{\left\vert \left\vert \mathsf{z}^{N}\right\rangle
_{2}\left\langle \mathsf{z}^{\prime N}\right\vert \right)  },\underline
{\left\vert k\right)  },\underline{\left\vert \mathfrak{x}\right)  }\right)
\mathbb{X}_{can}\nonumber\\
&  =\left(
\begin{array}
[c]{ccc}%
\mathbb{I} & \mathbb{O} & \mathbb{O}\\
\mathbb{O} & \left(  \mathfrak{b}_{N-1}\left\vert \underline{\left\vert
k\right)  }\right.  \right)  & \left(  \mathfrak{b}_{N-1}\left\vert
\underline{\left\vert \mathfrak{x}\right)  }\right.  \right) \\
\mathbb{O} & \left(  \mathfrak{b}_{N}\left\vert \underline{\left\vert
k\right)  }\right.  \right)  & \left(  \mathfrak{b}_{N}\left\vert
\underline{\left\vert \mathfrak{x}\right)  }\right.  \right)
\end{array}
\right)  \mathbb{X}_{can}%
\end{align}
Explicitly the expansion coefficients that change are given by
\begin{subequations}
\begin{align}
X_{k_{j}}  &  =%
%TCIMACRO{\dint \limits_{\mathsf{z}_{1}^{N}\,\overset{_{1}}{=}\,\mathsf{z}%
%_{2}^{N};\mathsf{z}_{1}^{N}\,\overset{}{=}\,\mathsf{z}_{2}^{N}}}%
%BeginExpansion
{\displaystyle\int\limits_{\mathsf{z}_{1}^{N}\,\overset{_{1}}{=}%
\,\mathsf{z}_{2}^{N};\mathsf{z}_{1}^{N}\,\overset{}{=}\,\mathsf{z}_{2}^{N}}}
%EndExpansion
k_{j}^{d}\left(  \mathsf{z}^{N},\mathsf{z}^{\prime N}\right)  ^{\ast}X\left(
\mathsf{z}^{N},\mathsf{z}^{\prime N}\right)  d\mathsf{z}_{1}^{N}%
d\mathsf{z}_{2}^{N}\\
X_{\mathfrak{x}_{j}}  &  =%
%TCIMACRO{\dint \limits_{\mathsf{z}_{1}^{N}\,\overset{_{1}}{=}\,\mathsf{z}%
%_{2}^{N};\mathsf{z}_{1}^{N}\,\overset{}{=}\,\mathsf{z}_{2}^{N}}}%
%BeginExpansion
{\displaystyle\int\limits_{\mathsf{z}_{1}^{N}\,\overset{_{1}}{=}%
\,\mathsf{z}_{2}^{N};\mathsf{z}_{1}^{N}\,\overset{}{=}\,\mathsf{z}_{2}^{N}}}
%EndExpansion
\mathfrak{x}_{j}^{d}\left(  \mathsf{z}^{N},\mathsf{z}^{\prime N}\right)
^{\ast}X\left(  \mathsf{z}^{N},\mathsf{z}^{\prime N}\right)  d\mathsf{z}%
_{1}^{N}d\mathsf{z}_{2}^{N}%
\end{align}
and
\end{subequations}
\begin{equation}
X\left(  \mathsf{z}^{N},\mathsf{z}^{\prime N}\right)  =%
%TCIMACRO{\dsum \limits_{1\leq j\leq\infty}}%
%BeginExpansion
{\displaystyle\sum\limits_{1\leq j\leq\infty}}
%EndExpansion
\left\{  X_{k_{j}}k_{j}\left(  \mathsf{z}^{N},\mathsf{z}^{\prime N}\right)
+X_{\mathfrak{x}_{j}}\mathfrak{x}_{j}\left(  \mathsf{z}^{N},\mathsf{z}^{\prime
N}\right)  \right\}  ;\qquad\mathsf{z}_{1}^{N}\,\overset{_{1}}{=}%
\,\mathsf{z}_{2}^{N};\mathsf{z}_{1}^{N}\,\overset{}{=}\,\mathsf{z}_{2}^{N}%
\end{equation}


\subsubsection{Non Canonical Basis of Complimentary Subspaces of $\ker
L_{N}^{1}.$}

We can utilize the maps $\tilde{\Gamma}$ introduced in eq. $\left(
\ref{Csubspace}\right)  $ to define a map $\Upsilon:\mathcal{B}_{1}\left(
\mathcal{H}^{N}\right)  \rightarrow\mathcal{B}_{1}\left(  \mathcal{H}%
^{N}\right)  $ to produce bases adapted to the decompositions eq. $\left(
\ref{nonorthog}\right)  $ as
\begin{align}
&  \left\{  \left\{  \left\vert \mathsf{z}^{N}\right\rangle \left\langle
\mathsf{z}^{\prime N}\right\vert \,\forall\mathsf{z}^{N}\overset{2}{\nsim
}\mathsf{z}^{\prime N}\right\}  ,\left\{  k_{\kappa}\,;1\leq\lambda\leq
\infty\right\}  ,\left\{  \tilde{\Gamma}\left(  \mathfrak{x}_{\lambda}\right)
\,;1\leq\lambda\leq\infty\right\}  \right\} \nonumber\\
&  \equiv\left\{  \left\{  \Upsilon\left(  \left\vert \mathsf{z}%
^{N}\right\rangle \left\langle \mathsf{z}^{\prime N}\right\vert \right)
\,\forall\mathsf{z}^{N}\overset{2}{\nsim}\mathsf{z}^{\prime N}\right\}
,\left\{  \Upsilon\left(  k_{\kappa}\right)  \,;1\leq\lambda\leq
\infty\right\}  ,\left\{  \Upsilon\left(  \mathfrak{x}_{\lambda}\right)
\,;1\leq\lambda\leq\infty\right\}  \right\}  \label{noncanbasis}%
\end{align}
where one should note that it is only the bases elements of $\left[  \ker
L_{N}^{1}\right]  ^{c}$ that have been changed. The\emph{\ }dual basis
elements $\left\{  \left(  k_{\kappa}^{d}\right\vert \right\}  $ are now no
longer biorthogonal to $\left\{  \left\vert \tilde{\Gamma}\left(
\mathfrak{x}_{\lambda}\right)  \right)  \right\}  $ as%
\begin{equation}
\tilde{\Gamma}\left(  \mathfrak{x}_{\lambda}\right)  =\Gamma\left(
\mathfrak{x}_{\lambda}\right)  +\mathfrak{x}_{\lambda}.
\end{equation}
and
\begin{equation}
\Gamma\left(  \mathfrak{x}_{\lambda}\right)  =\sum_{\kappa}\left\vert
k_{\kappa}\right)  \left[  \Gamma_{\mathcal{K}_{D}}\right]  _{\kappa\lambda}+%
%TCIMACRO{\dint _{\mathsf{z}^{N}\,\overset{2}{\neq\,}\mathsf{z}^{\prime N}}}%
%BeginExpansion
{\displaystyle\int_{\mathsf{z}^{N}\,\overset{2}{\neq\,}\mathsf{z}^{\prime N}}}
%EndExpansion
\left[  \Gamma_{\mathcal{K}_{OD}}\right]  _{\lambda}\left(  \mathsf{z}%
^{N};\mathsf{z}^{\prime N}\right)  \left\vert \mathsf{z}^{N}\right\rangle
\left\langle \mathsf{z}^{\prime N}\right\vert d\mathsf{z}^{N}d\mathsf{z}%
^{\prime N}%
\end{equation}
However one can construct a basis that is biorthogonal to eq. $\left(
\ref{noncanbasis}\right)  $ by considering
\begin{align}
&  \left(  \Upsilon_{\Gamma}^{d}\left(  \underline{\left(  \left\vert
\mathsf{z}^{N}\right\rangle _{2}\left\langle \mathsf{z}^{\prime N}\right\vert
\right\vert },\underline{\left(  k^{d}\right\vert },\left(  \underline
{\mathfrak{x}^{d}}\right\vert \right)  \right)  \left(  \Upsilon_{\Gamma
}\left(  \underline{\left\vert \left\vert \mathsf{z}^{N}\right\rangle
_{2}\left\langle \mathsf{z}^{\prime N}\right\vert \right)  },\underline
{\left\vert k\right)  },\underline{\left\vert \mathfrak{x}\right)  }\right)
\right)  ^{t}=\nonumber\\
&  \left(  \underline{\left(  \left\vert \mathsf{z}^{N}\right\rangle
_{2}\left\langle \mathsf{z}^{\prime N}\right\vert \right\vert },\underline
{\left(  k^{d}\right\vert },\left(  \underline{\mathfrak{x}^{d}}\right\vert
\right)  \left(
\begin{array}
[c]{ccc}%
\mathbb{I} & \mathbb{O} & -\mathbf{\Gamma}_{\mathcal{K}_{OD}}\\
\mathbb{O} & \mathbb{I} & -\mathbf{\Gamma}_{\mathcal{K}_{D}}\\
\mathbb{O} & \mathbb{O} & \mathbb{I}%
\end{array}
\right)  \left(
\begin{array}
[c]{ccc}%
\mathbb{I} & \mathbb{O} & \mathbf{\Gamma}_{\mathcal{K}_{OD}}\\
\mathbb{O} & \mathbb{I} & \mathbf{\Gamma}_{\mathcal{K}_{D}}\\
\mathbb{O} & \mathbb{O} & \mathbb{I}%
\end{array}
\right)  \left(
\begin{array}
[c]{c}%
\underline{\left\vert \left\vert \mathsf{z}^{N}\right\rangle _{2}\left\langle
\mathsf{z}^{\prime N}\right\vert \right)  }\\
\underline{\left\vert k\right)  }\\
\underline{\left\vert \mathfrak{x}\right)  }%
\end{array}
\right) \nonumber\\
&  \left.  =\right.  \left(  \underline{\left(  \left\vert \mathsf{z}%
^{N}\right\rangle _{2}\left\langle \mathsf{z}^{\prime N}\right\vert
\right\vert },\underline{\left(  k^{d}\right\vert },\left(  \underline
{\mathfrak{x}^{d}}\right\vert \right)  \left(
\begin{array}
[c]{c}%
\underline{\left\vert \left\vert \mathsf{z}^{N}\right\rangle _{2}\left\langle
\mathsf{z}^{\prime N}\right\vert \right)  }\\
\underline{\left\vert k\right)  }\\
\underline{\left\vert \mathfrak{x}\right)  }%
\end{array}
\right)  =\left(
\begin{array}
[c]{ccc}%
\mathbb{I}_{\mathcal{K}_{OD}} & \mathbb{O} & \mathbb{O}\\
\mathbb{O} & \mathbb{I}_{\mathcal{K}_{D}} & \mathbb{O}\\
\mathbb{O} & \mathbb{O} & \mathbb{I}_{\left[  \ker L_{N}^{1}\right]  ^{c}}%
\end{array}
\right)
\end{align}
which shows that the dual map $\Upsilon^{d}$ to $\Upsilon$, which is defined
by%
\begin{equation}
\left(  \Upsilon^{d}\left(  A\right)  |B\right)  =\left(  A|\Upsilon\left(
B\right)  \right)  \quad\forall B\in\mathcal{B}_{1}\left(  \mathcal{H}%
^{N}\right)  \quad\forall A\in\mathcal{B}\left(  \mathcal{H}^{N}\right)  ,
\end{equation}
is given by
\begin{align}
\Upsilon_{\Gamma}^{d}\left(  \underline{\left(  \left\vert \mathsf{z}%
^{N}\right\rangle _{2}\left\langle \mathsf{z}^{\prime N}\right\vert
\right\vert },\underline{\left(  k^{d}\right\vert },\left(  \underline
{\mathfrak{x}^{d}}\right\vert \right)   &  =\left(  \underline{\left(
\left\vert \mathsf{z}^{N}\right\rangle _{2}\left\langle \mathsf{z}^{\prime
N}\right\vert \right\vert },\underline{\left(  k^{d}\right\vert },\left(
\underline{\mathfrak{x}^{d}}\right\vert \right)  \mathbf{\Upsilon}_{\Gamma
}^{d}\nonumber\\
&  =\left(  \underline{\left(  \left\vert \mathsf{z}^{N}\right\rangle
_{2}\left\langle \mathsf{z}^{\prime N}\right\vert \right\vert },\underline
{\left(  k^{d}\right\vert },\left(  \underline{\mathfrak{x}^{d}}\right\vert
\right)  \mathbf{\Upsilon}_{\Gamma}^{-1}%
\end{align}
where
\begin{equation}
\Upsilon_{\Gamma}\left(  \underline{\left\vert \left\vert \mathsf{z}%
^{N}\right\rangle _{2}\left\langle \mathsf{z}^{\prime N}\right\vert \right)
},\underline{\left\vert k\right)  },\underline{\left\vert \mathfrak{x}\right)
}\right)  =\left(  \underline{\left\vert \left\vert \mathsf{z}^{N}%
\right\rangle _{2}\left\langle \mathsf{z}^{\prime N}\right\vert \right)
},\underline{\left\vert k\right)  },\underline{\left\vert \mathfrak{x}\right)
}\right)  \mathbf{\Upsilon}_{\Gamma}%
\end{equation}
and
\begin{subequations}
\begin{align}
\mathbf{\Upsilon}_{\Gamma}  &  =\left(
\begin{array}
[c]{ccc}%
\mathbb{I}_{\mathcal{K}_{OD}} & \mathbb{O} & \mathbf{\Gamma}_{\mathcal{K}%
_{OD}}\\
\mathbb{O} & \mathbb{I}_{\mathcal{K}_{D}} & \mathbf{\Gamma}_{\mathcal{K}_{D}%
}\\
\mathbb{O} & \mathbb{O} & \mathbb{I}_{\left[  \ker L_{N}^{1}\right]  ^{c}}%
\end{array}
\right) \\
\mathbf{\Upsilon}_{\Gamma}^{-1}  &  \mathbb{=}\left(
\begin{array}
[c]{ccc}%
\mathbb{I}_{\mathcal{K}_{OD}} & \mathbb{O} & -\mathbf{\Gamma}_{\mathcal{K}%
_{OD}}\\
\mathbb{O} & \mathbb{I}_{\mathcal{K}_{D}} & -\mathbf{\Gamma}_{\mathcal{K}_{D}%
}\\
\mathbb{O} & \mathbb{O} & \mathbb{I}_{\left[  \ker L_{N}^{1}\right]  ^{c}}%
\end{array}
\right)
\end{align}


The non canonical representations, $\tilde{\Gamma}\left(  \mathcal{V}\right)
$, of the base of the vector bundle $\mathfrak{B}_{L_{N}^{1}}$ are thus given
by
\end{subequations}
\begin{align}
\mathcal{V}_{\Gamma}  &  =\left\{  \left(  \underline{\left\vert \left\vert
\mathsf{z}^{N}\right\rangle _{2}\left\langle \mathsf{z}^{\prime N}\right\vert
\right)  },\underline{\left\vert k\right)  },\underline{\left\vert
\mathfrak{x}\right)  }\right)  \left(
\begin{array}
[c]{ccc}%
\mathbb{I}_{\mathcal{K}_{OD}} & \mathbb{O} & \mathbf{\Gamma}_{\mathcal{K}%
_{OD}}\\
\mathbb{O} & \mathbb{I}_{\mathcal{K}_{D}} & \mathbf{\Gamma}_{\mathcal{K}_{D}%
}\\
\mathbb{O} & \mathbb{O} & \mathbb{I}_{\left[  \ker L_{N}^{1}\right]  ^{c}}%
\end{array}
\right)  \left(
\begin{array}
[c]{c}%
\mathbf{0}\\
\mathbf{0}\\
\mathfrak{y}%
\end{array}
\right)  ;\mathfrak{y\in}\left[  \ker L_{N}^{1}\right]  ^{c}\right\}
\nonumber\\
&  =\left\{  \left(  \underline{\left\vert \left\vert \mathsf{z}%
^{N}\right\rangle _{2}\left\langle \mathsf{z}^{\prime N}\right\vert \right)
},\underline{\left\vert k\right)  },\underline{\left\vert \mathfrak{x}\right)
}\right)  \left(
\begin{array}
[c]{c}%
\mathbf{\Gamma}_{\mathcal{K}_{OD}}\mathfrak{y}\\
\mathbf{\Gamma}_{\mathcal{K}_{D}}\mathfrak{y}\\
\mathfrak{y}%
\end{array}
\right)  ;\mathfrak{y\in}\left[  \ker L_{N}^{1}\right]  ^{c}\right\}
\nonumber\\
&  =\left\{  \left(  \mathfrak{b}_{0},\cdots,\mathfrak{b}_{N}\right)  \left(
\begin{array}
[c]{c}%
\mathfrak{b}_{0}\\
\vdots\\
\mathfrak{b}_{N}%
\end{array}
\right)  \left(  \underline{\left\vert \left\vert \mathsf{z}^{N}\right\rangle
_{2}\left\langle \mathsf{z}^{\prime N}\right\vert \right)  },\underline
{\left\vert k\right)  },\underline{\left\vert \mathfrak{x}\right)  }\right)
\left(
\begin{array}
[c]{c}%
\mathbf{\Gamma}_{\mathcal{K}_{OD}}\mathfrak{y}\\
\mathbf{\Gamma}_{\mathcal{K}_{D}}\mathfrak{y}\\
\mathfrak{y}%
\end{array}
\right)  ;\mathfrak{y\in}\left[  \ker L_{N}^{1}\right]  ^{c}\right\}
\nonumber\\
&  =\left\{  \left(  \mathfrak{b}_{0},\cdots,\mathfrak{b}_{N}\right)  \left(
\begin{array}
[c]{ccc}%
\mathbb{I} & \mathbb{O} & \mathbb{O}\\
\mathbb{O} & \left(  \mathfrak{b}_{N-1}|\left\vert \underline{k}\right)
\right)  & \left(  \mathfrak{b}_{N}|\left\vert \underline{k}\right)  \right)
\\
\mathbb{O} & \left(  \mathfrak{b}_{N-1}|\underline{\left\vert \mathfrak{x}%
\right)  }\right)  & \left(  \mathfrak{b}_{N}|\underline{\left\vert
\mathfrak{x}\right)  }\right)
\end{array}
\right)  \left(
\begin{array}
[c]{c}%
\mathbf{\Gamma}_{\mathcal{K}_{OD}}\mathfrak{y}\\
\mathbf{\Gamma}_{\mathcal{K}_{D}}\mathfrak{y}\\
\mathfrak{y}%
\end{array}
\right)  ;\mathfrak{y\in}\left[  \ker L_{N}^{1}\right]  ^{c}\right\}
\end{align}


\subsubsection{Non Canonical Bases that Contain States \label{noncanstate}}

We concentrate on those representations of the base that are generated by the
set $\mathfrak{P}$ of maps $\left\{  \Gamma\right\}  $ that induce base
representatives and contains a maximal number of states
\begin{equation}
\mathfrak{P}=\left\{  \Gamma|\tilde{\Gamma}\left(  \mathfrak{d}_{1}\right)
\in\mathcal{B}_{1+}\left(  \mathcal{H}^{N}\right)  \,\forall\,\mathfrak{d}%
_{1}\in\ \mathfrak{C}_{1}\right\}
\end{equation}
where $\mathfrak{C}_{1}$ is the set of $\left[  \ker L_{N}^{1}\right]  ^{c}%
$-components of unnormalized states i.e.
\begin{equation}
\mathfrak{C}_{1}=\left\{  \mathfrak{d}_{1}|\mathfrak{d}_{1}\in\left[  \ker
L_{N}^{1}\right]  ^{c}\,\text{such that }\exists\mathfrak{d\in}\left[  \ker
L_{N}^{1}\right]  \text{ so that }\mathfrak{d+d}_{1}\in\mathcal{B}_{1+}\left(
\mathcal{H}^{N}\right)  \right\}
\end{equation}
Thus the cone
\begin{equation}
\mathcal{V}_{\Gamma+}=\mathcal{V}_{\Gamma}\cap\mathcal{B}_{1+}\left(
\mathcal{H}^{N}\right)  \equiv\left\{  \tilde{\Gamma}\left(  \mathfrak{d}%
_{1}\right)  |\mathfrak{d}_{1}\in\ \mathfrak{C}_{1};\tilde{\Gamma}\left(
\mathfrak{d}_{1}\right)  \in\mathcal{B}_{1+}\left(  \mathcal{H}^{N}\right)
\right\}
\end{equation}
is given by
\begin{equation}
\mathcal{V}_{\Gamma+}=\left\{  \left(  \underline{\left\vert \left\vert
\mathsf{z}^{N}\right\rangle _{2}\left\langle \mathsf{z}^{\prime N}\right\vert
\right)  },\underline{\left\vert k\right)  },\underline{\left\vert
\mathfrak{x}\right)  }\right)  \left(
\begin{array}
[c]{ccc}%
\mathbb{I}_{\mathcal{K}_{OD}} & \mathbb{O} & \mathbf{\Gamma}_{\mathcal{K}%
_{OD}}\\
\mathbb{O} & \mathbb{I}_{\mathcal{K}_{D}} & \mathbf{\Gamma}_{\mathcal{K}_{D}%
}\\
\mathbb{O} & \mathbb{O} & \mathbb{I}_{\left[  \ker L_{N}^{1}\right]  ^{c}}%
\end{array}
\right)  \left(
\begin{array}
[c]{c}%
\mathbf{0}\\
\mathbf{0}\\
\mathfrak{d}_{1}%
\end{array}
\right)  ;\mathfrak{d}_{1}\mathfrak{\in C}_{1}\right\}
\end{equation}
We denote a typical element of $\mathcal{V}_{\Gamma+}$ as $D_{0_{F}}^{N}$ thus
for $\mathfrak{d}_{1}\mathfrak{\in C}_{1}$
\begin{align}
D_{0_{F}}^{N}  &  =\Upsilon_{\Gamma}\left(  \mathfrak{d}_{1}\right)
=\tilde{\Gamma}\left(  \mathfrak{d}_{1}\right)  =\left(  \underline{\left\vert
\left\vert \mathsf{z}^{N}\right\rangle _{2}\left\langle \mathsf{z}^{\prime
N}\right\vert \right)  },\underline{\left\vert k\right)  },\underline
{\left\vert \mathfrak{x}\right)  }\right)  \left(
\begin{array}
[c]{c}%
\mathbf{\Gamma}_{\mathcal{K}_{OD}}\mathfrak{d}_{1}\\
\mathbf{\Gamma}_{\mathcal{K}_{D}}\mathfrak{d}_{1}\\
\mathfrak{d}_{1}%
\end{array}
\right) \nonumber\\
&  =\left(  \mathfrak{b}_{0},\cdots,\mathfrak{b}_{N}\right)  \left(
\begin{array}
[c]{ccc}%
\mathbb{I} & \mathbb{O} & \mathbb{O}\\
\mathbb{O} & \left(  \mathfrak{b}_{N-1}|\left\vert \underline{k}\right)
\right)  & \left(  \mathfrak{b}_{N}|\left\vert \underline{k}\right)  \right)
\\
\mathbb{O} & \left(  \mathfrak{b}_{N-1}|\underline{\left\vert \mathfrak{x}%
\right)  }\right)  & \left(  \mathfrak{b}_{N}|\underline{\left\vert
\mathfrak{x}\right)  }\right)
\end{array}
\right)  \left(
\begin{array}
[c]{c}%
\mathbf{\Gamma}_{\mathcal{K}_{OD}}\mathfrak{d}_{1}\\
\mathbf{\Gamma}_{\mathcal{K}_{D}}\mathfrak{d}_{1}\\
\mathfrak{d}_{1}%
\end{array}
\right) \nonumber\\
&  =\left(  \mathfrak{b}_{0},\cdots,\mathfrak{b}_{N}\right)  \left(
\begin{array}
[c]{c}%
D_{0_{F}\mathfrak{B}_{0}}^{N}\\
\vdots\\
D_{0_{F}\mathfrak{B}_{N}}^{N}%
\end{array}
\right)
\end{align}
recall that the domain of $\Upsilon_{\Gamma}$ is all of $\mathcal{B}%
_{1}\left(  \mathcal{H}^{N}\right)  $, while that of $\tilde{\Gamma}$ is
$\left[  \ker L_{N}^{1}\right]  ^{c}$ and
\begin{equation}
D_{0_{F}\mathfrak{B}_{m}}^{N}=%
%TCIMACRO{\dint \limits_{\mathsf{z}^{\prime N}\,\overset{m}{=}\,\mathsf{z}^{N}%
%}}%
%BeginExpansion
{\displaystyle\int\limits_{\mathsf{z}^{\prime N}\,\overset{m}{=}%
\,\mathsf{z}^{N}}}
%EndExpansion
D_{0_{F}}^{N}\left(  \mathsf{z}^{\prime N},\mathsf{z}^{N}\right)  \left\vert
\mathsf{z}^{\prime N}\right\rangle \left\langle \mathsf{z}^{N}\right\vert
d\mathsf{z}^{N}d\mathsf{z}^{\prime N}%
\end{equation}
so in particular%
\begin{align}
&  D_{0_{F}\mathfrak{B}_{N}}^{N}=%
%TCIMACRO{\dint }%
%BeginExpansion
{\displaystyle\int}
%EndExpansion
D_{0_{F}}^{N}\left(  \mathsf{z}^{N},\mathsf{z}^{N}\right)  \left\vert
\mathsf{z}^{N}\right\rangle \left\langle \mathsf{z}^{N}\right\vert
d\mathsf{z}^{N}\nonumber\\
&  \qquad\quad\left.  =\right.  \left(  \mathfrak{b}_{N}|\underline{\left\vert
\mathfrak{x}\right)  }\right)  \mathbf{\Gamma}_{\mathcal{K}_{D}}%
\mathfrak{d}_{1}+\left(  \mathfrak{b}_{N}|\underline{\left\vert \mathfrak{x}%
\right)  }\right)  \mathfrak{d}_{1}\nonumber\\
&  \qquad\quad\left.  =\right.  \left(  \mathfrak{b}_{N}|\underline{\left\vert
k\right)  }\right)  \left(  \underline{\left(  k^{d}\right\vert }%
|\Gamma\left(  \mathfrak{d}_{1}\right)  \right)  +\left(  \mathfrak{b}%
_{N}|\mathfrak{d}_{1}\right) \nonumber\\
&  \left.  =\right.
%TCIMACRO{\dsum \limits_{1\leq\nu\leq\infty}}%
%BeginExpansion
{\displaystyle\sum\limits_{1\leq\nu\leq\infty}}
%EndExpansion%
%TCIMACRO{\dint }%
%BeginExpansion
{\displaystyle\int}
%EndExpansion
k_{\nu}\left(  \mathsf{z}^{N},\mathsf{z}^{N}\right)  \Gamma\left[
\mathfrak{d}_{1}\right]  \left(  \mathsf{z}^{N},\mathsf{z}^{N}\right)
\left\vert \mathsf{z}^{N}\right\rangle \left\langle \mathsf{z}^{N}\right\vert
d\mathsf{z}^{N}+\int\mathfrak{d}_{1}\left(  \mathsf{z}^{N},\mathsf{z}%
^{N}\right)  \left\vert \mathsf{z}^{N}\right\rangle \left\langle
\mathsf{z}^{N}\right\vert d\mathsf{z}^{N}%
\end{align}
so that
\begin{equation}
D_{0_{F}}^{N}\left(  \mathsf{z}^{N},\mathsf{z}^{N}\right)  =D_{0_{F}%
\mathfrak{B}_{N}}^{N}\left(  \mathsf{z}^{N},\mathsf{z}^{N}\right)  =%
%TCIMACRO{\dsum \limits_{1\leq\nu\leq\infty}}%
%BeginExpansion
{\displaystyle\sum\limits_{1\leq\nu\leq\infty}}
%EndExpansion
k_{\nu}\left(  \mathsf{z}^{N},\mathsf{z}^{N}\right)  \Gamma\left[
\mathfrak{d}_{1}\right]  \left(  \mathsf{z}^{N},\mathsf{z}^{N}\right)
+\mathfrak{d}_{1}\left(  \mathsf{z}^{N},\mathsf{z}^{N}\right)
\end{equation}
while
\begin{align}
D_{0_{F}\mathfrak{B}_{N-1}}^{N}  &  =%
%TCIMACRO{\dint \limits_{\mathsf{z}^{\prime N}\,\overset{N-1}{=}\,\mathsf{z}%
%^{N}}}%
%BeginExpansion
{\displaystyle\int\limits_{\mathsf{z}^{\prime N}\,\overset{N-1}{=}%
\,\mathsf{z}^{N}}}
%EndExpansion
D_{0_{F}}^{N}\left(  \mathsf{z}^{\prime N},\mathsf{z}^{N}\right)  \left\vert
\mathsf{z}^{\prime N}\right\rangle \left\langle \mathsf{z}^{N}\right\vert
d\mathsf{z}^{N}d\mathsf{z}^{\prime N}\nonumber\\
&  =\left(  \mathfrak{b}_{N-1}|\underline{\left\vert \mathfrak{x}\right)
}\right)  \mathbf{\Gamma}_{\mathcal{K}_{D}}\mathfrak{d}_{1}+\left(
\mathfrak{b}_{N-1}|\underline{\left\vert \mathfrak{x}\right)  }\right)
\mathfrak{d}_{1}\nonumber\\
&  =\left(  \mathfrak{b}_{N-1}|\underline{\left\vert k\right)  }\right)
\left(  \underline{\left(  k^{d}\right\vert }|\Gamma\left(  \mathfrak{d}%
_{1}\right)  \right)  +\left(  \mathfrak{b}_{N-1}|\mathfrak{d}_{1}\right)
\nonumber\\
&  =%
%TCIMACRO{\dsum \limits_{1\leq\nu\leq\infty}}%
%BeginExpansion
{\displaystyle\sum\limits_{1\leq\nu\leq\infty}}
%EndExpansion%
%TCIMACRO{\dint \limits_{\mathsf{z}^{\prime N}\,\overset{N-1}{=}\,\mathsf{z}%
%^{N}}}%
%BeginExpansion
{\displaystyle\int\limits_{\mathsf{z}^{\prime N}\,\overset{N-1}{=}%
\,\mathsf{z}^{N}}}
%EndExpansion
k_{\nu}\left(  \mathsf{z}^{N},\mathsf{z}^{\prime N}\right)  \Gamma\left[
\mathfrak{d}_{1}\right]  \left(  \mathsf{z}^{N},\mathsf{z}^{\prime N}\right)
\left\vert \mathsf{z}^{N}\right\rangle \left\langle \mathsf{z}^{\prime
N}\right\vert d\mathsf{z}^{N}\nonumber\\
&  \qquad\qquad\qquad\qquad\qquad\qquad+%
%TCIMACRO{\dint \limits_{\mathsf{z}^{\prime N}\,\overset{N-1}{=}\,\mathsf{z}%
%^{N}}}%
%BeginExpansion
{\displaystyle\int\limits_{\mathsf{z}^{\prime N}\,\overset{N-1}{=}%
\,\mathsf{z}^{N}}}
%EndExpansion
\mathfrak{d}_{1}\left(  \mathsf{z}^{N},\mathsf{z}^{\prime N}\right)
\left\vert \mathsf{z}^{N}\right\rangle \left\langle \mathsf{z}^{\prime
N}\right\vert d\mathsf{z}^{N}d\mathsf{z}^{\prime N}%
\end{align}
so that
\begin{equation}
D_{0_{F}}^{N}\left(  \mathsf{z}^{\prime N},\mathsf{z}^{N}\right)  =%
%TCIMACRO{\dsum \limits_{1\leq\nu\leq\infty}}%
%BeginExpansion
{\displaystyle\sum\limits_{1\leq\nu\leq\infty}}
%EndExpansion
k_{\nu}\left(  \mathsf{z}^{N},\mathsf{z}^{\prime N}\right)  \Gamma\left[
\mathfrak{d}_{1}\right]  \left(  \mathsf{z}^{N},\mathsf{z}^{\prime N}\right)
+\mathfrak{d}_{1}\left(  \mathsf{z}^{N},\mathsf{z}^{\prime N}\right)
;\quad\mathsf{z}^{\prime N}\overset{N-1}{=}\mathsf{z}^{N}%
\end{equation}
where we have used the notation $\Gamma\left[  \mathfrak{d}_{1}\right]
\left(  \mathsf{z}^{N},\mathsf{z}^{\prime N}\right)  $ to distinguish between
arguments that are representation dependent and those that are not.

\section{Transformation of the Kinetic Energy Operator\label{KinEnTrans}}

\subsection{Transition Amplitudes}

We consider the commutation relationships between the operators $\mathbf{P,}%
e^{i\omega}$ and $\nu$ and observe that

\begin{enumerate}
\item
\begin{align}
&  \mathbf{P}e^{i\omega}\left\vert \Psi\right\rangle =\left(  -i\nabla\right)
e^{i\omega}\left\vert \Psi\right\rangle =e^{i\omega}\left(  -i\nabla\right)
\left\vert \Psi\right\rangle +e^{i\omega}\left(  \nabla\omega\right)
\left\vert \Psi\right\rangle \nonumber\\
&  \qquad\qquad\qquad\left.  =\right.  e^{i\omega}\left(  -i\nabla
+\nabla\omega\right)  \left\vert \Psi\right\rangle =e^{i\omega}\left\{
\mathbf{P}+\nabla\omega\right\}  \left\vert \Psi\right\rangle
\end{align}


\item
\begin{align}
&  \mathbf{P}\nu\left\vert \Psi\right\rangle =\left(  -i\nabla\right)
\nu\left\vert \Psi\right\rangle \nonumber\\
&  =\nu\left(  -i\nabla\right)  \left\vert \Psi\right\rangle +\left(
-i\nabla\nu\right)  \left\vert \Psi\right\rangle =\left\{  \nu\mathbf{P}%
-i\nabla\nu\right\}  \left\vert \Psi\right\rangle
\end{align}


\item
\begin{equation}
\nu e^{i\omega}\left\vert \Psi\right\rangle =e^{i\omega}\nu\left\vert
\Psi\right\rangle
\end{equation}

\end{enumerate}

Which can be summarized in the commutation relationship
\begin{align}
\left[  \mathbf{P,}e^{i\omega}\right]   &  =e^{i\omega}\nabla\omega\nonumber\\
\left[  \mathbf{P,}\nu\right]   &  =i\nabla\nu\nonumber\\
\left[  \nu,e^{i\omega}\right]   &  =0
\end{align}
that produce
\begin{align}
\mathbf{P}\nu e^{i\omega}\left\vert \Psi\right\rangle  &  =\left\{
\nu\mathbf{P}-i\nabla\nu\right\}  e^{i\omega}\left\vert \Psi\right\rangle
=\left\{  \nu e^{i\omega}\left\{  \mathbf{P}+\nabla\omega\right\}
-ie^{i\omega}\nabla\nu\right\}  \left\vert \Psi\right\rangle \nonumber\\
&  =e^{i\omega}\left\{  \nu\mathbf{P+}\nu\nabla\omega-i\nabla\nu\right\}
\left\vert \Psi\right\rangle
\end{align}
This allows us to express the transition amplitude%
\begin{equation}
\left\langle T\left(  \nu_{a},\omega_{a}\right)  \Psi|\left(  \mathbf{P\bullet
P}\right)  T\left(  \nu_{b},\omega_{b}\right)  \Psi\right\rangle
\end{equation}
as
\begin{equation}
\left\langle e^{i\omega_{a}}\left\{  \nu_{a}\mathbf{P+}\nu_{a}\nabla\omega
_{a}-i\nabla\nu_{a}\right\}  \Psi|e^{i\omega_{b}}\left\{  \nu_{b}%
\mathbf{P+}\nu_{b}\nabla\omega_{b}-i\nabla\nu_{b}\right\}  \Psi\right\rangle
\end{equation}
which gives

\begin{description}
\item[A]
\begin{equation}
\left\langle \nu_{a}e^{i\omega_{a}}\mathbf{P}\Psi|\nu_{b}e^{i\omega_{b}%
}\mathbf{P}\Psi\right\rangle =-\left\langle \Psi|\left\{  \nabla\bullet\left(
\nu_{a}\nu_{b}e^{i\left(  \omega_{b}-\omega_{a}\right)  }\right)
\mathbb{I}\bullet\nabla\right\}  \Psi\right\rangle \label{exA}%
\end{equation}
where
\begin{equation}
\nabla\bullet\left(  \nu_{a}\nu_{b}e^{i\left(  \omega_{b}-\omega_{a}\right)
}\right)  \mathbb{I}\bullet\nabla=\left(
\begin{array}
[c]{ccc}%
\nabla_{1} & \cdots & \nabla_{m}%
\end{array}
\right)  \left(
\begin{array}
[c]{ccc}%
\nu_{a}\nu_{b}e^{i\left(  \omega_{b}-\omega_{a}\right)  } & \cdots & 0\\
\vdots & \ddots & \vdots\\
0 & \cdots & \nu_{a}\nu_{b}e^{i\left(  \omega_{b}-\omega_{a}\right)  }%
\end{array}
\right)  \left(
\begin{array}
[c]{c}%
\nabla_{1}\\
\vdots\\
\nabla_{m}%
\end{array}
\right)
\end{equation}
The transition amplitude eq. $\left(  \ref{exA}\right)  $ is given explicitly
by
\begin{align}
&  -%
%TCIMACRO{\dint }%
%BeginExpansion
{\displaystyle\int}
%EndExpansion
\left\{  \nu_{a}\left(  \mathsf{z}\right)  e^{-i\omega_{a}\left(
\mathsf{z}\right)  }\nabla\Psi\left(  \mathsf{z}\right)  ^{\ast}\nu_{b}\left(
\mathsf{z}\right)  e^{i\omega_{b}\left(  \mathsf{z}\right)  }\nabla\Psi\left(
\mathsf{z}\right)  \right\}  d\mathsf{z}\nonumber\\
&  \left.  =\right.  -%
%TCIMACRO{\dint }%
%BeginExpansion
{\displaystyle\int}
%EndExpansion
\left\{  \nabla\Psi\left(  \mathsf{z}\right)  ^{\dagger}\bullet\nabla
\Psi\left(  \mathsf{z}\right)  e^{i\left(  \omega_{b}\left(  \mathsf{z}%
\right)  -\omega_{a}\left(  \mathsf{z}\right)  \right)  }\nu_{a}\left(
\mathsf{z}\right)  \nu_{b}\left(  \mathsf{z}\right)  \right\}  d\mathsf{z}%
\end{align}


\item[B]
\begin{equation}
\left\langle \nu_{a}e^{i\omega_{a}}\mathbf{P}\Psi|\left\{  \nu_{b}\nabla
\omega_{b}-i\nabla\nu_{b}\right\}  e^{-i\omega_{b}}\Psi\right\rangle
=-i\left\langle \Psi|\left\{  \nabla\bullet\nu_{a}e^{i\left(  \omega
_{b}-\omega_{a}\right)  }\mathbb{I}\bullet\left(  \nu_{b}\nabla\omega
_{b}-i\nabla\nu_{b}\right)  \right\}  \Psi\right\rangle
\end{equation}


\item[C]
\begin{equation}
\left\langle \left(  \nu_{a}\nabla\omega_{a}-i\nabla\nu_{a}\right)  \Psi
|v_{b}e^{i\omega_{b}}\mathbf{P}\Psi\right\rangle =-i\left\langle \Psi|\left\{
\left(  \nu_{a}\nabla\omega_{a}-i\nabla\nu_{a}\right)  \bullet v_{b}%
e^{i\left(  \omega_{b}-\omega_{a}\right)  }\mathbb{I}\bullet\nabla\right\}
\Psi\right\rangle
\end{equation}


\item[D]
\begin{align}
&  \left\langle \left(  \nu_{a}\nabla\omega_{a}-i\nabla\nu_{a}\right)
\Psi|\left(  \nu_{b}\nabla\omega_{b}-i\nabla\nu_{b}\right)  \Psi\right\rangle
\nonumber\\
&  \left.  =\right.  \left\langle \Psi|\left\{  \left(  \nu_{a}\nabla
\omega_{a}-i\nabla\nu_{a}\right)  \bullet e^{i\left(  \omega_{b}-\omega
_{a}\right)  }\mathbb{I}\bullet\left(  \nu_{b}\nabla\omega_{b}-i\nabla\nu
_{b}\right)  \right\}  \Psi\right\rangle
\end{align}
collecting together the terms A,B,C and D we obtain%
\begin{align}
&  \left\langle T\left(  \nu_{a},\omega_{a}\right)  \Psi|\left(
\mathbf{P\bullet P}\right)  T\left(  \nu_{b},\omega_{b}\right)  \Psi
\right\rangle =\nonumber\\
&  \qquad\left\langle \Psi\right.  |\left\{  -\nabla\bullet\left(  \nu_{a}%
\nu_{b}e^{i\left(  \omega_{b}-\omega_{a}\right)  }\right)  \mathbb{I}%
\bullet\nabla\right. \nonumber\\
&  \qquad-\nabla\bullet\nu_{a}e^{i\left(  \omega_{b}-\omega_{a}\right)
}\mathbb{I}\bullet\left(  \nu_{b}\nabla\omega_{b}-i\nabla\nu_{b}\right)
\nonumber\\
&  \qquad-\left(  \nu_{a}\nabla\omega_{a}-i\nabla\nu_{a}\right)  \bullet
v_{b}e^{i\left(  \omega_{b}-\omega_{a}\right)  }\mathbb{I}\bullet
\nabla\nonumber\\
&  \quad\left.  +\left(  \nu_{a}\nabla\omega_{a}-i\nabla\nu_{a}\right)
\bullet e^{i\left(  \omega_{b}-\omega_{a}\right)  }\mathbb{I}\bullet\left(
\nu_{b}\nabla\omega_{b}-i\nabla\nu_{b}\right)  \right\}  \left.
\Psi\right\rangle
\end{align}
We can also consider this as the expectation value wrt the pure state
$\left\vert \Psi\right\rangle $ of the effective operator $K\left(  \nu
_{a},\omega_{a},\nu_{b},\omega_{b}\right)  $ where
\begin{align}
&  K\left(  \nu_{a},\omega_{a},\nu_{b},\omega_{b}\right)  =-\nabla
\bullet\left(  \nu_{a}\nu_{b}e^{i\left(  \omega_{b}-\omega_{a}\right)
}\right)  \mathbb{I}\bullet\nabla-\nabla\bullet\nu_{a}e^{i\left(  \omega
_{b}-\omega_{a}\right)  }\mathbb{I}\bullet\left(  \nu_{b}\nabla\omega
_{b}-i\nabla\nu_{b}\right) \nonumber\\
&  -\left(  \nu_{a}\nabla\omega_{a}-i\nabla\nu_{a}\right)  \bullet
v_{b}e^{i\left(  \omega_{b}-\omega_{a}\right)  }\mathbb{I}\bullet
\nabla+\left(  \nu_{a}\nabla\omega_{a}-i\nabla\nu_{a}\right)  \bullet
e^{i\left(  \omega_{b}-\omega_{a}\right)  }\mathbb{I}\bullet\left(  \nu
_{b}\nabla\omega_{b}-i\nabla\nu_{b}\right)
\end{align}

\end{description}

\subsection{Expectation Values}

If we set $\nu_{a}=\nu_{b}=\nu$ and $\omega_{a}=\omega_{b}=\omega$ we obtain
the expectation value of the Kinetic Energy operator
\begin{equation}
\left\langle T\left(  \nu,\omega\right)  \Psi|\left(  \mathbf{P\bullet
P}\right)  T\left(  \nu,\omega\right)  \Psi\right\rangle
\end{equation}
which can be developed to produce
\begin{align}
&  \left\langle T\left(  \nu,\omega\right)  \Psi|\left(  \mathbf{P\bullet
P}\right)  T\left(  \nu,\omega\right)  \Psi\right\rangle =\nonumber\\
&  \qquad\left\langle \Psi\right.  |\left\{  -\nabla\bullet\left(  \nu
^{2}\right)  \mathbb{I}\bullet\nabla\right. \nonumber\\
&  \qquad-\nabla\bullet\left(  \nu^{2}\nabla\omega-i\nu\nabla\nu\right)
\nonumber\\
&  \qquad-\left(  \nu^{2}\nabla\omega-i\nu\nabla\nu\right)  \bullet
\nabla\nonumber\\
&  \quad\left.  \left(  \nu\nabla\omega-i\nabla\nu\right)  \bullet
\mathbb{I}\bullet\left(  \nu\nabla\omega-i\nabla\nu\right)  \right\}  \left.
\Psi\right\rangle
\end{align}
leading to the definition of an effective Kinetic Energy operator
\begin{align}
K\left(  \nu,\omega\right)   &  =-\nabla\bullet\left(  \nu^{2}\right)
\mathbb{I}\bullet\nabla-\nabla\bullet\left(  \nu^{2}\nabla\omega-i\nu\nabla
\nu\right) \nonumber\\
&  -\left(  \nu^{2}\nabla\omega-i\nu\nabla\nu\right)  \bullet\nabla+\left(
\nu\nabla\omega-i\nabla\nu\right)  \bullet\mathbb{I}\bullet\left(  \nu
\nabla\omega-i\nabla\nu\right)
\end{align}


\section{Semigroups and Representations\label{Groups}}

The semigroups we consider reflect the algebraic and vector space properties
of $\mathcal{B}_{1}\left(  \mathcal{H}^{N}\right)  $ inherited from
$\mathcal{B}\left(  \mathcal{H}^{N}\right)  $ and the contraction map
$L_{N}^{1}$. In the following "$\mathfrak{S}$" denotes non abelian semigroups
and $\mathfrak{A}$ abelian semigroups. We consider representations,
$\mathcal{\hat{T}}_{\mathcal{W}}\left(  \mathfrak{S}\right)  $, of
$\mathfrak{S}$ as bounded superoperators that are defined on \emph{all }of
$\mathcal{B}_{1}\left(  \mathcal{H}^{N}\right)  $ that map subspaces
$\mathcal{W}$ of $\mathcal{B}_{1}\left(  \mathcal{H}^{N}\right)  $ to
themselves i.e.
\begin{equation}
\mathcal{\hat{T}}_{\mathcal{W}}:\mathfrak{S}\rightarrow\mathcal{B}_{1}\left(
\mathcal{H}^{N}\right)
\end{equation}
and
\begin{equation}
\mathcal{\hat{T}}_{\mathcal{W}}\left(  \mathfrak{S}\right)  :\mathcal{W}%
\rightarrow\mathcal{W};\qquad s\in\mathfrak{S};\quad\mathcal{W\subseteq B}%
_{1}\left(  \mathcal{H}^{N}\right)
\end{equation}
that are defined by matrix, (or generalized matrix - in the proceeding the
term matrix will refer to either), representations $\mathcal{R}_{\mathcal{W}%
}^{B}\left(  \mathfrak{S}\right)  $ induced by a specific bases $\left\{
\underline{B}\right\}  $ of $\mathcal{W}$ in the following manner
\begin{equation}
\mathcal{\hat{T}}_{\mathcal{W}}\left(  \mathfrak{S}\right)  \left(
\underline{B}\right)  _{R}=\left(  \underline{B}\right)  _{R}\mathcal{R}%
_{\mathcal{W}}^{B}\left(  \mathfrak{S}\right)
\end{equation}
If the basis $\left\{  \underline{B}\right\}  $ is defined to be the reference
basis for the subspace $\mathcal{W}$ we omit the superscript "$B$" and set
$\mathcal{R}_{\mathcal{W}}\equiv\mathcal{R}_{\mathcal{W}}^{B}$ and if
$\mathcal{W}=\mathcal{B}_{1}\left(  \mathcal{H}^{N}\right)  $ we set
$\mathcal{\hat{T}}\left(  \mathfrak{S}\right)  \equiv\mathcal{\hat{T}%
}_{\mathcal{B}_{1}\left(  \mathcal{H}^{N}\right)  }\left(  \mathfrak{S}%
\right)  $.

The representations $\mathcal{\hat{T}}\left(  \mathfrak{S}\right)  $ define
superoperators acting in $\mathcal{W}$, through expansions defined by the
representations $\mathcal{R}_{\mathcal{W}}^{B}\left(  \mathfrak{S}\right)  $
wrt the basis $\left\{  \underline{B}\right\}  $. Obviously these
superoperators can also be expanded wrt to any other basis, $\left\{
\underline{B^{\prime}}\right\}  $, of $\mathcal{W}$ producing the
representation $\mathcal{R}_{\mathcal{W}}^{B^{\prime}}\left(  \mathfrak{S}%
\right)  $. We will utilize this property by first characterizing a semigroup
of super operators in one basis and then expanding elements of the semigroup
in a basis that is more amenable for carrying out calculations. In particular
we consider the following matrix representations:

\begin{enumerate}
\item $\mathcal{R}_{\mathcal{B}_{1}\left(  \mathcal{H}^{N}\right)  }$ is
defined by the basis $\left\{  \mathfrak{b}_{0},\cdots,\mathfrak{b}%
_{N}\right\}  $ of $\mathcal{B}_{1}\left(  \mathcal{H}^{N}\right)  $

\item $\mathcal{R}_{\mathcal{B}_{1}\left(  \mathcal{H}^{N}\right)  }^{\Gamma}$
is defined by the basis $\left\{  \underline{\left\vert \left\vert
\mathsf{z}^{N}\right\rangle _{2}\left\langle \mathsf{z}^{\prime N}\right\vert
\right)  },\underline{\left\vert k\right)  },\underline{\left\vert
\tilde{\Gamma}\left(  \mathfrak{x}\right)  \right)  }\right\}  $ of
$\mathcal{B}_{1}\left(  \mathcal{H}^{N}\right)  $

\item $\mathcal{R}_{\ker L_{N}^{1}}$ is defined by the basis $\left\{
\underline{\left\vert \left\vert \mathsf{z}^{N}\right\rangle _{2}\left\langle
\mathsf{z}^{\prime N}\right\vert \right)  },\underline{\left\vert k\right)
}\right\}  $ of $\ker L_{N}^{1}$

\item $\mathcal{R}_{\mathcal{V}_{\Gamma}}$ is defined by the basis $\left\{
\left\vert \tilde{\Gamma}\left(  \mathfrak{x}\right)  \right)  \right\}  $ of
$\mathcal{V}_{\Gamma}$
\end{enumerate}

Other representations are also used and introduced in context.

\subsection{The semigroups $\mathfrak{S}_{\mathcal{B}_{1}\left(
\mathcal{H}^{N}\right)  },\,\mathfrak{S}_{\mathcal{B}_{1+}\left(
\mathcal{H}^{N}\right)  } $}

The semigroup $\mathfrak{S}_{\mathcal{B}_{1}\left(  \mathcal{H}^{N}\right)  }$
is defined by the representation $\hat{T}\left(  \mathfrak{S}_{\mathcal{B}%
_{1}\left(  \mathcal{H}^{N}\right)  }\right)  $ as the set of bounded linear
maps $\mathcal{B}_{1}\left(  \mathcal{H}^{N}\right)  \rightarrow
\mathcal{B}_{1}\left(  \mathcal{H}^{N}\right)  $, this contains a sub
semigroup $\mathfrak{S}_{\mathcal{B}_{1+}\left(  \mathcal{H}^{N}\right)  }$ of
positive maps, $\mathcal{\hat{T}}\left(  \mathfrak{S}_{\mathcal{B}_{1+}\left(
\mathcal{H}^{N}\right)  }\right)  $, that map the cone $\mathcal{B}%
_{1+}\left(  \mathcal{H}^{N}\right)  $ to itself. Equivalent representations
of these groups can be generated by any basis or generalized basis of
$\mathcal{B}_{1}\left(  \mathcal{H}^{N}\right)  $. The form of the super
operators belonging to $\mathcal{\hat{T}}\left(  \mathfrak{S}_{\mathcal{B}%
_{1}\left(  \mathcal{H}^{N}\right)  }\right)  $ are defined through the
"matrix" representation $\mathcal{R}_{\mathcal{B}_{1}\left(  \mathcal{H}%
^{N}\right)  }$ by%
\begin{equation}
\hat{S}=%
%TCIMACRO{\dint }%
%BeginExpansion
{\displaystyle\int}
%EndExpansion
S\left(  \mathsf{z}_{1}^{N}\mathsf{z}_{2}^{N}\mathsf{z}_{3}^{N}\mathsf{z}%
_{4}^{N}\right)  \left\vert \left\vert \mathsf{z}_{1}^{N}\right\rangle
\left\langle \mathsf{z}_{2}^{N}\right\vert \right)  \left(  \left\vert
\mathsf{z}_{3}^{N}\right\rangle \left\langle \mathsf{z}_{4}^{N}\right\vert
\right\vert d\mathsf{z}_{1}^{N}d\mathsf{z}_{2}^{N}d\mathsf{z}_{3}%
^{N}d\mathsf{z}_{4}^{N}%
\end{equation}
with $S\left(  \mathsf{z}_{1}^{N}\mathsf{z}_{2}^{N}\mathsf{z}_{3}%
^{N}\mathsf{z}_{4}^{N}\right)  $ only restricted by the condition that
$\hat{S}:\mathcal{B}_{1}\left(  \mathcal{H}^{N}\right)  \rightarrow
\mathcal{B}_{1}\left(  \mathcal{H}^{N}\right)  $, while those belonging to
$\mathcal{\hat{T}}\left(  \mathfrak{S}_{\mathcal{B}_{1+}\left(  \mathcal{H}%
^{N}\right)  }\right)  $ must also map positive operators to positive operators.

\subsection{The semigroup $\mathfrak{A}_{pos}\subset\mathfrak{S}%
_{\mathcal{B}_{1+}\left(  \mathcal{H}^{N}\right)  }$}

\label{uker}

The abelian semigroup $\mathfrak{A}_{pos}$ $\subset\mathfrak{S}_{\mathcal{B}%
_{1+}\left(  \mathcal{H}^{N}\right)  }$ has a defining representation,
$\mathcal{R}_{\mathcal{B}_{1}\left(  \mathcal{H}^{N}\right)  }\left(
\mathfrak{A}_{pos}\right)  $, in the space-spin coordinate representation,
through the operators defined in eq. $\left(  \ref{TDiag}\right)  $ acting in
$\mathcal{H}^{N}$. If we use other generalized basis $\left\{  \left\vert
\mathsf{x}_{1}\cdots\mathsf{x}_{N}\right\rangle \right\}  \neq\left\{
\left\vert \mathsf{z}_{1}\cdots\mathsf{z}_{N}\right\rangle \right\}  $ of
$\mathcal{H}^{N}$ in eq. $\left(  \ref{TDiag}\right)  $ we obtain isomorphic
representations. The representation $\mathcal{R}_{\mathcal{B}_{1}\left(
\mathcal{H}^{N}\right)  }\left(  \mathfrak{A}_{pos}\right)  $ can be expressed
as a direct integral of one dimensional representations, $\mathcal{R}%
_{\left\vert \mathsf{z}^{N}\right\rangle \left\langle \mathsf{z}^{\prime
N}\right\vert }\left(  \mathfrak{A}_{pos}\right)  $ i.e.%
\begin{equation}
\mathcal{R}_{\mathcal{B}_{1}\left(  \mathcal{H}^{N}\right)  }\left(
\mathfrak{A}_{pos}\right)  =\int_{\oplus}\mathcal{R}_{\left\vert
\mathsf{z}^{N}\right\rangle \left\langle \mathsf{z}^{\prime N}\right\vert
}\left(  \mathfrak{A}_{pos}\right)  d\mathsf{z}^{N}d\mathsf{z}^{\prime N}.
\end{equation}
which is easy to see as
\begin{equation}
\hat{T}\left(  \left\vert \mathsf{z}^{N}\right\rangle \left\langle
\mathsf{z}^{N\prime}\right\vert \right)  =T\left(  \mathsf{z}^{N}\right)
T\left(  \mathsf{z}^{\prime N}\right)  ^{\ast}\left\vert \mathsf{z}%
^{N}\right\rangle \left\langle \mathsf{z}^{\prime N}\right\vert \quad
\forall\mathsf{z}^{N},\mathsf{z}^{\prime N};\quad\hat{T}\in\hat{T}%
_{\mathcal{B}_{1}\left(  \mathcal{H}^{N}\right)  }\left(  \mathfrak{A}%
_{pos}\right)
\end{equation}
and corresponds to the expansion
\begin{equation}
\hat{T}=%
%TCIMACRO{\dint }%
%BeginExpansion
{\displaystyle\int}
%EndExpansion
\left(  \left\vert \mathsf{z}^{N}\right\rangle \left\langle \mathsf{z}%
^{N\prime}\right\vert \left\vert \hat{T}\right.  \left\vert \mathsf{z}%
^{N}\right\rangle \left\langle \mathsf{z}^{N\prime}\right\vert \right)
\left\vert \left\vert \mathsf{z}^{N}\right\rangle \left\langle \mathsf{z}%
^{N\prime}\right\vert \right)  \left(  \left\vert \mathsf{z}^{N}\right\rangle
\left\langle \mathsf{z}^{N\prime}\right\vert \right\vert d\mathsf{z}%
^{N}d\mathsf{z}^{\prime N} \label{upos}%
\end{equation}
and hence one can view the representation $\mathcal{R}_{\mathcal{B}_{1}\left(
\mathcal{H}^{N}\right)  }\left(  \mathfrak{A}_{pos}\right)  $ as being
constituted by the generalized diagonal matrices $\left\{  \left(  \left\vert
\mathsf{z}^{N}\right\rangle \left\langle \mathsf{z}^{N\prime}\right\vert
\left\vert \hat{T}\right.  \left\vert \mathsf{z}^{N}\right\rangle \left\langle
\mathsf{z}^{N\prime}\right\vert \right)  \right\}  \equiv\left\{  T\left(
\mathsf{z}^{N}\right)  T\left(  \mathsf{z}^{\prime N}\right)  ^{\ast}\right\}
$. The expansion eq. $\left(  \ref{upos}\right)  $ clearly shows that one can
decompose $\hat{T}$ as
\begin{equation}
\hat{T}=%
%TCIMACRO{\dsum \limits_{0\leq m\leq N}}%
%BeginExpansion
{\displaystyle\sum\limits_{0\leq m\leq N}}
%EndExpansion
\hat{T}_{m}%
\end{equation}
where $\hat{T}_{m}:\mathfrak{B}_{m}\rightarrow\mathfrak{B}_{m}$i.e. $\hat
{T}_{m}\equiv\hat{T}_{\mathfrak{B}_{m}}$.

The elements of $\hat{T}\left(  \mathfrak{A}_{pos}\right)  $ can also be
expanded wrt the basis adapted to the canonical decomposition eq. $\left(
\ref{orthog}\right)  $ as (note the restrictions $\mathsf{z}^{N}\overset
{2}{\nsim}\mathsf{z}^{\prime N}$ on the integrals)
\begin{align}
\hat{T}  &  =%
%TCIMACRO{\dint \limits_{\mathsf{z}^{N}\overset{2}{\,\nsim\,}\mathsf{z}^{\prime
%N}}}%
%BeginExpansion
{\displaystyle\int\limits_{\mathsf{z}^{N}\overset{2}{\,\nsim\,}\mathsf{z}%
^{\prime N}}}
%EndExpansion
\left(  \left\vert \mathsf{z}^{N}\right\rangle \left\langle \mathsf{z}%
^{N\prime}\right\vert \left\vert \hat{T}\right.  \left\vert \mathsf{z}%
^{N}\right\rangle \left\langle \mathsf{z}^{N\prime}\right\vert \right)
\left\vert \overset{}{}\left\vert \mathsf{z}^{N}\right\rangle \left\langle
\mathsf{z}^{N\prime}\right\vert \right)  \left(  \overset{}{}\left\vert
\mathsf{z}^{N}\right\rangle \left\langle \mathsf{z}^{N\prime}\right\vert
\right\vert d\mathsf{z}^{N}d\mathsf{z}^{\prime N}\nonumber\\
&  +%
%TCIMACRO{\dsum \limits_{1\leq\kappa,\lambda\leq\infty}}%
%BeginExpansion
{\displaystyle\sum\limits_{1\leq\kappa,\lambda\leq\infty}}
%EndExpansion
\left(  k_{\kappa}^{d}\left\vert \hat{T}\right.  k_{\lambda}\right)
\left\vert k_{\lambda}\right)  \left(  k_{\kappa}^{d}\right\vert +%
%TCIMACRO{\dsum \limits_{1\leq\kappa,\lambda\leq\infty}}%
%BeginExpansion
{\displaystyle\sum\limits_{1\leq\kappa,\lambda\leq\infty}}
%EndExpansion
\left(  \mathfrak{x}_{\kappa}^{d}\left\vert \hat{T}\right.  \mathfrak{x}%
_{\lambda}\right)  \left\vert \mathfrak{x}_{\kappa}\right)  \left(
\mathfrak{x}_{\lambda}^{d}\right\vert \nonumber\\
&  +%
%TCIMACRO{\dsum \limits_{1\leq\kappa,\lambda\leq\infty}}%
%BeginExpansion
{\displaystyle\sum\limits_{1\leq\kappa,\lambda\leq\infty}}
%EndExpansion
\left(  k_{\kappa}^{d}\left\vert \hat{T}\right.  \mathfrak{x}_{\lambda
}\right)  \left\vert k_{\kappa}\right)  \left(  \mathfrak{x}_{\lambda}%
^{d}\right\vert +%
%TCIMACRO{\dsum \limits_{1\leq\kappa,\lambda\leq\infty}}%
%BeginExpansion
{\displaystyle\sum\limits_{1\leq\kappa,\lambda\leq\infty}}
%EndExpansion
\left(  \mathfrak{x}_{\kappa}^{d}\left\vert \hat{T}\right.  k_{\lambda
}\right)  \left\vert \mathfrak{x}_{\kappa}\right)  \left(  k_{\lambda}%
^{d}\right\vert
\end{align}
where $\left\{  \left(  \mathfrak{x}_{\kappa}\right\vert ,\left\vert
\mathfrak{x}_{\kappa}\right)  \right\}  $ and $\left\{  \left(  k_{\kappa}%
^{d}\right\vert ,\left\vert k_{\kappa}\right)  \right\}  $ are biorthogonal
bases of the dual spaces $\left\{  \left[  \ker L_{N}^{1}\right]
^{cd},\left[  \ker L_{N}^{1}\right]  ^{c}\right\}  $ and $\left\{
\ \mathcal{K}_{D}^{d},\mathcal{K}_{D}\right\}  $ respectively.

The action of the super operators $\hat{T}\in\hat{T}\left(  \mathfrak{A}%
_{pos}\right)  $ on these basis elements is given by%
\begin{equation}
\qquad\qquad\qquad\qquad\hat{T}\left(  \left\vert \mathsf{z}^{N}\right\rangle
\left\langle \mathsf{z}^{\prime N}\right\vert \right)  =T\left(
\mathsf{z}^{N}\right)  T\left(  \mathsf{z}^{\prime N}\right)  ^{\ast
}\left\vert \mathsf{z}^{N}\right\rangle \left\langle \mathsf{z}^{\prime
N}\right\vert ;\quad\mathsf{z}^{N}\overset{2}{\nsim}\mathsf{z}^{\prime
N}\qquad\qquad
\end{equation}%
\begin{align}
&  \hat{T}\left(
%TCIMACRO{\dint }%
%BeginExpansion
{\displaystyle\int}
%EndExpansion
k_{\kappa}\left(  \mathsf{z}_{1}\mathsf{z}^{N-1};\mathsf{z}_{1}^{\prime
}\mathsf{z}^{N-1}\right)  \left\vert \mathsf{z}_{1}\mathsf{z}^{N-1}%
\right\rangle \left\langle \mathsf{z}_{1}^{\prime}\mathsf{z}^{N-1}\right\vert
d\mathsf{z}_{1}d\mathsf{z}_{1}^{\prime}d\mathsf{z}^{N-1}\right)  =\nonumber\\
&
%TCIMACRO{\dint }%
%BeginExpansion
{\displaystyle\int}
%EndExpansion
k_{\kappa}\left(  \mathsf{z}_{1}\mathsf{z}^{N-1};\mathsf{z}_{1}^{\prime
}\mathsf{z}^{N-1}\right)  T\left(  \mathsf{z}_{1}\mathsf{z}^{N-1}\right)
T\left(  \mathsf{z}_{1}^{\prime}\mathsf{z}^{N-1}\right)  ^{\ast}\left\vert
\mathsf{z}_{1}\mathsf{z}^{N-1}\right\rangle \left\langle \mathsf{z}%
_{1}^{\prime}\mathsf{z}^{N-1}\right\vert d\mathsf{z}_{1}d\mathsf{z}%
_{1}^{\prime}d\mathsf{z}^{N-1}%
\end{align}%
\begin{align}
&  \hat{T}\left(
%TCIMACRO{\dint }%
%BeginExpansion
{\displaystyle\int}
%EndExpansion
\mathfrak{x}_{\lambda}\left(  \mathsf{z}_{1}\mathsf{z}^{N-1};\mathsf{z}%
_{1}^{\prime}\mathsf{z}^{N-1}\right)  \left\vert \mathsf{z}_{1}\mathsf{z}%
^{N-1}\right\rangle \left\langle \mathsf{z}_{1}^{\prime}\mathsf{z}%
^{N-1}\right\vert d\mathsf{z}_{1}d\mathsf{z}_{1}^{\prime}d\mathsf{z}%
^{N-1}\right)  =\nonumber\\
&
%TCIMACRO{\dint }%
%BeginExpansion
{\displaystyle\int}
%EndExpansion
\mathfrak{x}_{\lambda}\left(  \mathsf{z}_{1}\mathsf{z}^{N-1};\mathsf{z}%
_{1}^{\prime}\mathsf{z}^{N-1}\right)  T\left(  \mathsf{z}_{1}\mathsf{z}%
^{N-1}\right)  T\left(  \mathsf{z}_{1}^{\prime}\mathsf{z}^{N-1}\right)
^{\ast}\left\vert \mathsf{z}_{1}\mathsf{z}^{N-1}\right\rangle \left\langle
\mathsf{z}_{1}^{\prime}\mathsf{z}^{N-1}\right\vert d\mathsf{z}_{1}%
d\mathsf{z}_{1}^{\prime}d\mathsf{z}^{N-1}%
\end{align}
which shows that $\mathcal{R}_{\mathcal{B}_{1}\left(  \mathcal{H}^{N}\right)
}\left(  \mathfrak{A}_{pos}\right)  $ splits into the direct sum
\begin{equation}
\mathcal{R}_{\mathcal{B}_{1}\left(  \mathcal{H}^{N}\right)  }\left(
\mathfrak{A}_{pos}\right)  =\mathcal{R}_{\mathcal{K}_{OD}}\left(
\mathfrak{A}_{pos}\right)  \oplus\mathcal{R}_{\mathcal{D}}\left(
\mathfrak{A}_{pos}\right)
\end{equation}
where
\begin{equation}
\mathcal{R}_{\mathcal{K}_{OD}}\left(  \mathfrak{A}_{pos}\right)  =\int
_{\oplus;\mathsf{z}^{N}\overset{2}{\nsim}\mathsf{z}^{\prime N}}\mathcal{R}%
_{\left\vert \mathsf{z}^{N}\right\rangle \left\langle \mathsf{z}^{\prime
N}\right\vert }\left(  \mathfrak{A}_{pos}\right)  d\mathsf{z}^{N}%
d\mathsf{z}^{\prime N}%
\end{equation}
and
\begin{equation}
\mathcal{R}_{\mathcal{D}}\left(  \mathfrak{A}_{pos}\right)  =\mathcal{R}%
_{\mathcal{K}_{D}\oplus\left[  \ker L_{N}^{1}\right]  ^{c}}\left(
\mathfrak{A}_{pos}\right)
\end{equation}
where
\begin{equation}
\mathcal{D}=\mathfrak{B}_{N-1}\oplus\mathfrak{B}_{N}%
\end{equation}
(One should note that $\mathcal{R}_{\mathcal{K}_{OD}}\left(  \mathfrak{A}%
_{pos}\right)  $ is not a faithful representation even though $\mathcal{R}%
_{\mathcal{B}_{1}\left(  \mathcal{H}^{N}\right)  }\left(  \mathfrak{A}%
_{pos}\right)  $ is). The matrix representations $\mathcal{R}_{\mathcal{K}%
_{OD}}\left(  \mathfrak{A}_{pos}\right)  $ induces super operator
representations $\mathcal{\hat{T}}_{\mathcal{K}_{OD}}\left(  \mathfrak{A}%
_{pos}\right)  $ and $\mathcal{\hat{T}}_{\mathcal{D}}\left(  \mathfrak{A}%
_{pos}\right)  $ given by
\begin{equation}
\mathcal{\hat{T}}_{\mathcal{K}_{OD}}\left(  \mathfrak{A}_{pos}\right)
=\left(  \underline{\left\vert \left\vert \mathsf{z}^{N}\right\rangle
_{2}\left\langle \mathsf{z}^{\prime N}\right\vert \right)  }\right)
_{R}\mathcal{R}_{\mathcal{K}_{OD}}\left(  \mathfrak{A}_{pos}\right)
\end{equation}
and
\begin{equation}
\mathcal{\hat{T}}_{\mathcal{D}}\left(  \mathfrak{A}_{pos}\right)  =\left(
\underline{\left\vert \left\vert \mathsf{z}^{N}\right\rangle _{2}\left\langle
\mathsf{z}^{\prime N}\right\vert \right)  }\right)  _{R}\mathcal{R}%
_{\mathcal{D}}\left(  \mathfrak{A}_{pos}\right)
\end{equation}
which are constituted by super operators of the form
\begin{equation}
\mathcal{\hat{T}}_{\mathcal{W}}\left(  s\right)  =\int T_{\mathcal{W}}\left[
s\right]  \left(  \mathsf{z}^{N}\right)  T_{\mathcal{W}}\left[  s\right]
\left(  \mathsf{z}^{\prime N}\right)  ^{\ast}\left\vert \mathsf{z}^{\prime
N}\right\rangle \left\langle \mathsf{z}^{N}\right\vert d\mathsf{z}%
^{N}d\mathsf{z}^{\prime N};\quad\mathsf{z}^{N}\overset{2}{\nsim}%
\mathsf{z}^{\prime N};s\in\mathfrak{A}_{pos}%
\end{equation}
where $\mathcal{W}=$ $\mathcal{K}_{OD}$ and $\mathcal{D}$. The representation
$\mathcal{R}_{\mathcal{D}}\left(  \mathfrak{A}_{pos}\right)  $ splits into the
representation $\mathcal{R}_{\mathfrak{B}_{N}}\left(  \mathfrak{A}%
_{pos}\right)  $ given by
\begin{equation}
\mathcal{R}_{\mathfrak{B}_{N}}\left(  s\right)  =\int\left\vert T_{0}\left[
s\right]  \left(  \mathsf{z}^{N}\right)  \right\vert ^{2}\left\vert
\mathsf{z}^{N}\right\rangle \left\langle \mathsf{z}^{N}\right\vert
d\mathsf{z}^{N}\equiv\int\left\vert T\left[  s\right]  \left(  \mathsf{z}%
^{N}\right)  \right\vert ^{2}\left\vert \mathsf{z}^{N}\right\rangle
\left\langle \mathsf{z}^{N}\right\vert d\mathsf{z}^{N};s\in\mathfrak{A}_{pos}%
\end{equation}
and $\mathcal{R}_{\mathfrak{B}_{N-1}}\left(  \mathfrak{A}_{pos}\right)  $
given by
\begin{align}
\mathcal{R}_{\mathfrak{B}_{N-1}}\left(  s\right)   &  =\int T_{1}\left[
s\right]  \left(  \mathsf{z}^{N}\right)  T_{1}\left[  s\right]  \left(
\mathsf{z}^{\prime N}\right)  ^{\ast}\left\vert \mathsf{z}^{\prime
N}\right\rangle \left\langle \mathsf{z}^{N}\right\vert d\mathsf{z}%
^{N}d\mathsf{z}^{\prime N}\nonumber\\
&  \equiv\int T\left[  s\right]  \left(  \mathsf{z}^{N}\right)  T\left[
s\right]  \left(  \mathsf{z}^{\prime N}\right)  ^{\ast}\left\vert
\mathsf{z}^{\prime N}\right\rangle \left\langle \mathsf{z}^{N}\right\vert
d\mathsf{z}^{N}d\mathsf{z}^{\prime N};\quad\mathsf{z}^{N}\overset{N-1}%
{=}\mathsf{z}^{\prime N};s\in\mathfrak{A}_{pos}%
\end{align}
i.e.%
\begin{equation}
\mathcal{R}_{\mathcal{D}}\left(  \mathfrak{A}_{pos}\right)  =\mathcal{R}%
_{\mathfrak{B}_{N-1}}\left(  \mathfrak{A}_{pos}\right)  \oplus\mathcal{R}%
_{\mathfrak{B}_{N}}\left(  \mathfrak{A}_{pos}\right)
\end{equation}
and
\begin{equation}
\mathfrak{B}_{N-1}\oplus\mathfrak{B}_{N}=\mathcal{K}_{D}\oplus\left[  \ker
L_{N}^{1}\right]  ^{c}%
\end{equation}


\emph{It is important to observe that} the representation $\mathcal{R}%
_{\mathfrak{B}_{N}}\left(  \mathfrak{A}_{pos}\right)  $ clearly determines the
magnitude of the integral kernels $T\left[  s\right]  \left(  \mathsf{z}%
^{N}\right)  $ and thus \emph{determines} the integral kernels of the super
operators constituting the representations $\left\{  \mathcal{R}%
_{\mathfrak{B}_{m}}\left(  \mathfrak{A}_{pos}\right)  |0\leq m\leq
N-1\right\}  $ \emph{up to phase factors. }This can also be viewed in terms of
induced representations i.e. as constraints on the class of representations
$\mathcal{\hat{T}}$ induced by a specific representation $\mathcal{\hat{T}%
}_{0}$.

To summarize the representations we have discussed in this section have been
generated by the subspaces
\begin{subequations}
\begin{align}
\mathcal{B}_{1}\left(  \mathcal{H}^{N}\right)   &  =%
%TCIMACRO{\dbigoplus \limits_{\mathsf{z}^{N},\mathsf{z}^{\prime N}}}%
%BeginExpansion
{\displaystyle\bigoplus\limits_{\mathsf{z}^{N},\mathsf{z}^{\prime N}}}
%EndExpansion
\left\{  \left\vert \mathsf{z}^{N}\right\rangle \left\langle \mathsf{z}%
^{\prime N}\right\vert \right\} \\
\mathcal{K}_{OD}  &  =\mathfrak{V}_{2}=%
%TCIMACRO{\dbigoplus \limits_{\mathsf{z}^{N}\,\overset{2}{\nsim}\,\mathsf{z}%
%^{\prime N}}}%
%BeginExpansion
{\displaystyle\bigoplus\limits_{\mathsf{z}^{N}\,\overset{2}{\nsim}%
\,\mathsf{z}^{\prime N}}}
%EndExpansion
\left\{  \left\vert \mathsf{z}^{N}\right\rangle \left\langle \mathsf{z}%
^{\prime N}\right\vert \right\}  =%
%TCIMACRO{\dbigoplus \limits_{0\leq m\leq N-2}}%
%BeginExpansion
{\displaystyle\bigoplus\limits_{0\leq m\leq N-2}}
%EndExpansion
\mathfrak{B}_{m}\\
\mathcal{K}_{D}  &  =\ker L_{N}^{1}\ominus\mathcal{K}_{OD}\\
\mathcal{D}  &  =\mathfrak{B}_{N}\oplus\mathfrak{B}_{N-1}=\mathcal{K}%
_{D}\oplus\left[  \ker L_{N}^{1}\right]  ^{c}=\mathcal{B}_{1}\left(
\mathcal{H}^{N}\right)  \ominus\mathcal{K}_{OD}\\
\left[  \ker L_{N}^{1}\right]  ^{c}  &  =\mathcal{B}_{1}\left(  \mathcal{H}%
^{N}\right)  \ominus\ker L_{N}^{1}%
\end{align}


\subsection{The semigroup $\mathfrak{A}_{\ker}\subset\mathfrak{A}_{pos}$}

Using the representation $\mathcal{R}_{\mathcal{B}_{1}\left(  \mathcal{H}%
^{N}\right)  }\left(  \mathfrak{A}_{pos}\right)  $ of $\mathfrak{A}_{pos}$ we
define the sub semigroup $\mathfrak{A}_{\ker}$, which preserves the vector
bundle structure of $\mathfrak{B}_{\ker L_{N}^{1}}$, as
\end{subequations}
\begin{equation}
\mathfrak{A}_{\ker}=\left\{  s|\hat{T}\left(  s\right)  :\ker L_{N}%
^{1}\rightarrow\ker L_{N}^{1};s\in\mathfrak{A}_{pos}\right\}  .
\end{equation}
The representation $\mathcal{R}_{\mathcal{B}_{1}\left(  \mathcal{H}%
^{N}\right)  }\left(  \mathfrak{A}_{\ker}\right)  $ of $\mathfrak{A}_{\ker}$
in terms of the generalized basis $\left\{  \mathfrak{b}_{0},\cdots
,\mathfrak{b}_{N}\right\}  $ of $\mathcal{B}_{1}\left(  \mathcal{H}%
^{N}\right)  $ is given by
\begin{equation}
\mathcal{\hat{T}}\left(  \mathfrak{A}_{\ker}\right)  \left(  \mathfrak{b}%
_{0},\cdots,\mathfrak{b}_{N}\right)  =\left(  \mathfrak{b}_{0},\cdots
,\mathfrak{b}_{N}\right)  \mathcal{R}_{\mathcal{B}_{1}\left(  \mathcal{H}%
^{N}\right)  }\left(  \mathfrak{A}_{\ker}\right)
\end{equation}
and the representation, $\mathcal{R}_{\mathcal{B}_{1}\left(  \mathcal{H}%
^{N}\right)  }^{0}\left(  \mathfrak{A}_{\ker}\right)  $, of $\mathfrak{A}%
_{\ker}$ which is equivalent to $\mathcal{R}_{\mathcal{B}_{1}\left(
\mathcal{H}^{N}\right)  }\left(  \mathfrak{A}_{\ker}\right)  $ but is adapted
specifically to the decomposition
\begin{equation}
\mathcal{B}_{1}\left(  \mathcal{H}^{N}\right)  =\ker L_{N}^{1}\oplus\left[
\ker L_{N}^{1}\right]  ^{c}%
\end{equation}
has the form%
\begin{equation}
\mathcal{\hat{T}}\left(  \mathfrak{A}_{\ker}\right)  \left(  \underline
{\left\vert \left\vert \mathsf{z}^{N}\right\rangle _{2}\left\langle
\mathsf{z}^{\prime N}\right\vert \right)  },\underline{\left\vert k\right)
},\underline{\left\vert \mathfrak{x}\right)  }\right)  =\left(  \underline
{\left\vert \left\vert \mathsf{z}^{N}\right\rangle _{2}\left\langle
\mathsf{z}^{\prime N}\right\vert \right)  },\underline{\left\vert k\right)
},\underline{\left\vert \mathfrak{x}\right)  }\right)  \mathcal{R}%
_{\mathcal{B}_{1}\left(  \mathcal{H}^{N}\right)  }^{0}\left(  \mathfrak{A}%
_{\ker}\right)
\end{equation}
where
\begin{equation}
\mathcal{R}_{\mathcal{B}_{1}\left(  \mathcal{H}^{N}\right)  }^{0}\left(
\mathfrak{A}_{\ker}\right)  =\left(
\begin{array}
[c]{cc}%
\mathcal{R}_{\ker L_{N}^{1}}\left(  \mathfrak{A}_{\ker}\right)  & \ast\\
0 & \mathcal{R}_{\left[  \ker L_{N}^{1}\right]  ^{c}}\left(  \mathfrak{A}%
_{\ker}\right)
\end{array}
\right)  \label{K0}%
\end{equation}
producing representations $\mathcal{R}_{\ker L_{N}^{1}}$ and $\mathcal{R}%
_{\left[  \ker L_{N}^{1}\right]  ^{c}}$ of $\mathfrak{A}_{\ker}$ explicitly
formed by%
\begin{align}
\mathcal{\hat{T}}\left(  \mathfrak{A}_{\ker}\right)  \left(  \underline
{\left\vert \left\vert \mathsf{z}^{N}\right\rangle _{2}\left\langle
\mathsf{z}^{\prime N}\right\vert \right)  },\underline{\left\vert k\right)
}\right)   &  =\left(  \underline{\left\vert \left\vert \mathsf{z}%
^{N}\right\rangle _{2}\left\langle \mathsf{z}^{\prime N}\right\vert \right)
},\underline{\left\vert k\right)  }\right)  \mathcal{R}_{\ker L_{N}^{1}%
}\left(  \mathfrak{A}_{\ker}\right) \nonumber\\
\mathcal{\hat{T}}\left(  \mathfrak{A}_{\ker}\right)  \left(  \underline
{\left\vert \mathfrak{x}\right)  }\right)  _{R}  &  =\left(  \underline
{\left\vert \mathfrak{x}\right)  }\right)  _{R}\mathcal{R}_{\left[  \ker
L_{N}^{1}\right]  ^{c}}\left(  \mathfrak{A}_{\ker}\right)  .
\end{align}
It should be noted that $\mathcal{R}_{\mathcal{B}_{1}\left(  \mathcal{H}%
^{N}\right)  }^{0}\left(  \mathfrak{A}_{\ker}\right)  $ is \emph{not fully}
reducible over the direct sum $\ker L_{N}^{1}\oplus\left[  \ker L_{N}%
^{1}\right]  ^{c}$ i.e. it \emph{properly contains} the direct sum of
$\mathcal{R}_{\ker L_{N}^{1}}$ and $\mathcal{R}_{\left[  \ker L_{N}%
^{1}\right]  ^{c}}$
\begin{equation}
\mathcal{R}_{\mathcal{B}_{1}\left(  \mathcal{H}^{N}\right)  }^{0}\left(
\mathfrak{A}_{\ker}\right)  \supset\mathcal{R}_{\ker L_{N}^{1}}\left(
\mathfrak{A}_{\ker}\right)  \oplus\mathcal{R}_{\left[  \ker L_{N}^{1}\right]
^{c}}\left(  \mathfrak{A}_{\ker}\right)
\end{equation}
The lack of full reducibility of $\mathcal{R}_{\mathcal{B}_{1}\left(
\mathcal{H}^{N}\right)  }^{0}\left(  \mathfrak{A}_{\ker}\right)  $ is a
consequence of the fact that
\begin{equation}
\hat{T}:\ker L_{N}^{1}\rightarrow\ker L_{N}^{1}\nRightarrow\hat{T}:\left[
\ker L_{N}^{1}\right]  ^{c}\rightarrow\left[  \ker L_{N}^{1}\right]
^{c}\text{ for }\hat{T}\in\mathcal{R}_{\mathcal{B}_{1}\left(  \mathcal{H}%
^{N}\right)  }\left(  \mathfrak{A}_{\ker}\right)
\end{equation}
as $\mathcal{B}_{1}\left(  \mathcal{H}^{N}\right)  $ is not a Hilbert space
and $\mathcal{R}_{\mathcal{B}_{1}\left(  \mathcal{H}^{N}\right)  }$ is not a
unitary representation.

The representation $\mathcal{R}_{\ker L_{N}^{1}}\left(  \mathfrak{A}_{\ker
}\right)  $ splits into the direct sum of two representations $\mathcal{R}%
_{\mathcal{K}_{OD}}\left(  \mathfrak{A}_{\ker}\right)  $ and $\mathcal{R}%
_{\mathcal{K}_{D}}\left(  \mathfrak{A}_{\ker}\right)  $ i.e.
\begin{equation}
\mathcal{R}_{\ker L_{N}^{1}}\left(  \mathfrak{A}_{\ker}\right)  =\mathcal{R}%
_{\mathcal{K}_{OD}}\left(  \mathfrak{A}_{\ker}\right)  \oplus\mathcal{R}%
_{\mathcal{K}_{D}}\left(  \mathfrak{A}_{\ker}\right)
\end{equation}
where
\begin{equation}
\mathcal{R}_{\mathcal{K}_{D}}\left(  \mathfrak{A}_{\ker}\right)  =\left\{
\mathbb{\hat{T}}|\,\hat{T}\in\mathcal{R}_{\mathcal{B}_{1}\left(
\mathcal{H}^{N}\right)  }\left(  \mathfrak{A}_{pos}\right)  \text{ and }%
\hat{T}:\mathcal{K}_{D}\rightarrow\ \mathcal{K}_{D}\right\}
\end{equation}
and the representation $\mathcal{R}_{\mathcal{K}_{OD}}\left(  \mathfrak{A}%
_{\ker}\right)  $ is given in terms of the direct integral%
\begin{equation}
\mathcal{R}_{\mathcal{K}_{OD}}\left(  \mathfrak{A}_{\ker}\right)
=\int_{\oplus}\mathcal{R}_{\left\vert \mathsf{z}^{N}\right\rangle \left\langle
\mathsf{z}^{\prime N}\right\vert }\left(  \mathfrak{A}_{\ker}\right)
d\mathsf{z}^{N}d\mathsf{z}^{\prime N};\quad\mathsf{z}^{N}\overset{2}{\nsim
}\mathsf{z}^{\prime N}%
\end{equation}


\subsection{Restrictions on the Integral Kernel $T\left(  \mathsf{z}%
^{N}\right)  $}

Eq. $\left(  \ref{K0}\right)  $ places the following restrictions on the
integral kernel $T\left(  \mathsf{z}^{N}\right)  $ of elements of
$\mathcal{R}_{\mathcal{B}_{1}\left(  \mathcal{H}^{N}\right)  }^{0}\left(
\mathfrak{A}_{\ker}\right)  $
\begin{align}
&  \left(  \mathfrak{x}_{\kappa}^{d}\left\vert \hat{T}\right.  k_{\lambda
}\right)  =\nonumber\\
&
%TCIMACRO{\dint }%
%BeginExpansion
{\displaystyle\int}
%EndExpansion
\mathfrak{x}_{\lambda}^{d}\left(  \mathsf{z}_{1}\mathsf{z}^{N-1}%
;\mathsf{z}_{1}^{\prime}\mathsf{z}^{N-1}\right)  ^{\ast}T\left(
\mathsf{z}_{1}\mathsf{z}^{N-1}\right)  T\left(  \mathsf{z}_{1}^{\prime
}\mathsf{z}^{N-1}\right)  ^{\ast}k_{\kappa}\left(  \mathsf{z}_{1}%
\mathsf{z}^{N-1};\mathsf{z}_{1}^{\prime}\mathsf{z}^{N-1}\right)
d\mathsf{z}_{1}d\mathsf{z}_{1}^{\prime}d\mathsf{z}^{N-1}\nonumber\\
&  \left.  {}\right.  \qquad\qquad\qquad\qquad\qquad\qquad\qquad\qquad
\qquad\qquad\qquad\qquad=0\quad1\leq\kappa,\lambda\leq\infty
\end{align}
The conditions
\begin{align}
&  \left(  k_{\kappa}^{d}\left\vert \hat{T}\right.  \left\vert \mathsf{z}%
^{N}\right\rangle \left\langle \mathsf{z}^{N\prime}\right\vert \right)
=0\nonumber\\
&  \left(  \mathfrak{x}_{\lambda}^{d}\left\vert \hat{T}\right.  \left\vert
\mathsf{z}^{N}\right\rangle \left\langle \mathsf{z}^{N\prime}\right\vert
\right)  =0
\end{align}
where $\mathsf{z}^{N}\overset{2}{\nsim}\mathsf{z}^{\prime N}$ are
automatically satisfied. These conditions enforce $\hat{T}$ to be a vector
bundle map, but do not reflect any special property of $\ker L_{N}^{1}$ the
subspace defining the vector bundle structure. The conditions induced by the
specific nature of $\hat{T}$ i.e. $\hat{T}:\ker L_{N}^{1}\rightarrow\ker
L_{N}^{1}$ are the \textit{generalized strong orthogonality} constraints
\begin{equation}
\int T\left(  \mathsf{z}_{1}\mathsf{z}^{N-1}\right)  T\left(  \mathsf{z}%
_{1}^{\prime}\mathsf{z}^{N-1}\right)  ^{\ast}k_{\kappa}\left(  \mathsf{z}%
_{1}\mathsf{z}^{N-1};\mathsf{z}_{1}^{\prime}\mathsf{z}^{N-1}\right)
d\mathsf{z}^{N-1}=0\quad1\leq\kappa\leq\infty
\end{equation}


\subsection{Non Canonical Representations \label{noncanrep}}

Representations, $\mathcal{R}_{\mathcal{B}_{1}\left(  \mathcal{H}^{N}\right)
}^{\Gamma}\left(  \mathfrak{A}_{\ker}\right)  $, over the non canonical
decomposition eq. $\left(  \ref{nonorthog}\right)  $ are obtained from
$\mathcal{R}_{\mathcal{B}_{1}\left(  \mathcal{H}^{N}\right)  }^{0}\left(
\mathfrak{A}_{\ker}\right)  $ by considering the isomorphisms $\Upsilon
_{\Gamma}:\mathcal{B}_{1}\left(  \mathcal{H}^{N}\right)  \rightarrow
\mathcal{B}_{1}\left(  \mathcal{H}^{N}\right)  $ from the canonical direct sum
eq. $\left(  \ref{orthog}\right)  $ to the non canonical direct sum eq.
$\left(  \ref{nonorthog}\right)  $ in terms of the bases introduced in
Appendix \ref{bases} these isomorphisms are
\begin{align}
&  \Upsilon_{\Gamma}\left(  \left(  \underline{\left\vert \left\vert
\mathsf{z}^{N}\right\rangle _{2}\left\langle \mathsf{z}^{\prime N}\right\vert
\right)  }\right)  ,\left(  \underline{\left\vert k\right)  }\right)  ,\left(
\underline{\left\vert \mathfrak{x}\right)  }\right)  \right)  =\nonumber\\
&  \qquad\qquad\qquad\left(  \left(  \underline{\left\vert \left\vert
\mathsf{z}^{N}\right\rangle _{2}\left\langle \mathsf{z}^{\prime N}\right\vert
\right)  }\right)  ,\left(  \underline{\left\vert k\right)  }\right)  ,\left(
\underline{\left\vert \mathfrak{x}\right)  }\right)  \right)  \left(
\begin{array}
[c]{ccc}%
\mathbb{I}_{\mathcal{K}_{OD}} & \mathbb{O} & \mathbf{\Gamma}_{\mathcal{K}%
_{OD}}\\
\mathbb{O} & \mathbb{I}_{\mathcal{K}_{D}} & \mathbf{\Gamma}_{\mathcal{K}_{D}%
}\\
\mathbb{O} & \mathbb{O} & \mathbb{I}_{\left[  \ker L_{N}^{1}\right]  ^{c}}%
\end{array}
\right)
\end{align}
with inverses
\begin{align}
&  \Upsilon_{\Gamma}^{-1}\left(  \left(  \underline{\left\vert \left\vert
\mathsf{z}^{N}\right\rangle _{2}\left\langle \mathsf{z}^{\prime N}\right\vert
\right)  }\right)  ,\left(  \underline{\left\vert k\right)  }\right)  ,\left(
\underline{\left\vert \mathfrak{x}\right)  }\right)  \right)  =\nonumber\\
&  \qquad\qquad\qquad\left(  \left(  \underline{\left\vert \left\vert
\mathsf{z}^{N}\right\rangle _{2}\left\langle \mathsf{z}^{\prime N}\right\vert
\right)  }\right)  ,\left(  \underline{\left\vert k\right)  }\right)  ,\left(
\underline{\left\vert \mathfrak{x}\right)  }\right)  \right)  \left(
\begin{array}
[c]{ccc}%
\mathbb{I}_{\mathcal{K}_{OD}} & \mathbb{O} & -\mathbf{\Gamma}_{\mathcal{K}%
_{OD}}\\
\mathbb{O} & \mathbb{I}_{\mathcal{K}_{D}} & -\mathbf{\Gamma}_{\mathcal{K}_{D}%
}\\
\mathbb{O} & \mathbb{O} & \mathbb{I}_{\left[  \ker L_{N}^{1}\right]  ^{c}}%
\end{array}
\right)
\end{align}
One can define the similar representations $\mathcal{R}_{\mathcal{B}%
_{1}\left(  \mathcal{H}^{N}\right)  }^{\Gamma}\left(  \mathfrak{A}_{\ker
}\right)  $ to $\mathcal{R}_{\mathcal{B}_{1}\left(  \mathcal{H}^{N}\right)
}^{0}\left(  \mathfrak{A}_{\ker}\right)  $ by
\begin{equation}
\mathcal{R}_{\mathcal{B}_{1}\left(  \mathcal{H}^{N}\right)  }^{\Gamma}\left(
\mathfrak{A}_{\ker}\right)  =\mathbf{\Upsilon}_{\Gamma}\mathcal{R}%
_{\mathcal{B}_{1}\left(  \mathcal{H}^{N}\right)  }^{0}\left(  \mathfrak{A}%
_{\ker}\right)  \mathbf{\Upsilon}_{\Gamma}^{-1}%
\end{equation}
these representations correspond to the general decompositions eq. $\left(
\ref{nonorthog}\right)  $ of $\mathcal{B}_{1}\left(  \mathcal{H}^{N}\right)  $
and
\begin{equation}
\mathcal{R}_{\mathcal{B}_{1}\left(  \mathcal{H}^{N}\right)  }^{\Gamma}\left(
\mathfrak{A}_{\ker}\right)  \supset\mathcal{R}_{\ker L_{N}^{1}}\left(
\mathfrak{A}_{\ker}\right)  \oplus\mathcal{R}_{\mathcal{V}_{\Gamma}}\left(
\mathfrak{A}_{\ker}\right)
\end{equation}


\subsection{Fiber Representations of $\mathfrak{A}_{\ker}\label{fiberrep}$}

Consider the semigroup of intra fiber transformations $\mathcal{R}%
_{F_{\mathfrak{x}_{1}}}\left(  \mathfrak{A}_{\ker}\right)  \subset
\mathcal{R}_{\mathcal{B}_{1}\left(  \mathcal{H}^{N}\right)  }\left(
\mathfrak{A}_{pos}\right)  $ given by the representation $\mathcal{R}%
_{F_{\mathfrak{x}_{1}}}$ of $\mathfrak{A}_{\ker}$ where the super operators
defining this generalized matrix semigroup are characterized by
\begin{align}
\mathcal{\hat{T}}_{F_{\mathfrak{x}_{1}}}\left(  \mathfrak{A}_{\ker}\right)
&  =\left(  \underline{\left\vert \left\vert \mathsf{z}^{N}\right\rangle
_{2}\left\langle \mathsf{z}^{\prime N}\right\vert \right)  },\underline
{\left\vert k\right)  },\underline{\left\vert \mathfrak{x}\right)  }\right)
\mathcal{R}_{F_{\mathfrak{x}_{1}}}\left(  \mathfrak{A}_{\ker}\right)
\nonumber\\
&  =\left\{  \hat{T}|\hat{T}\in\mathcal{\hat{T}}\left(  \mathfrak{A}_{\ker
}\right)  ;\,X\in\left[  \mathfrak{x}_{1}\right]  \Rightarrow\hat{T}\left(
X\right)  \in\left[  \mathfrak{x}_{1}\right]  \right\}
\end{align}


\begin{theorem}
$\hat{T}\in\mathcal{\hat{T}}_{F_{\mathfrak{x}_{1}}}\left(  \mathfrak{A}_{\ker
}\right)  \Leftrightarrow\hat{T}:\ker L_{N}^{1}\rightarrow\ker L_{N}^{1}$ and
$\hat{T}\left(  \mathfrak{x}_{1}\right)  =\mathfrak{x}_{1}$

\begin{proof}
"$\Rightarrow$"Let $X,X^{\prime}\in\ker L_{N}^{1}$ then $\hat{T}\left(
X-X^{\prime}\right)  =\hat{T}\left(  X_{\ker}-X_{\ker}^{\prime}\right)
\in\ker L_{N}^{1}$, hence $\hat{T}:\ker L_{N}^{1}\rightarrow\ker L_{N}^{1}$
then $\hat{T}\left(  X_{\ker}+\mathfrak{x}_{1}\right)  =\hat{T}\left(
X_{\ker}\right)  +\hat{T}\left(  \mathfrak{x}_{1}\right)  =\tilde{X}_{\ker
}+\hat{T}\left(  \mathfrak{x}_{1}\right)  $ where $\tilde{X}_{\ker}\in\ker
L_{N}^{1}$, hence $\hat{T}\left(  \mathfrak{x}_{1}\right)  =\mathfrak{x}_{1}.$

"$\Leftarrow"$The implication that the right hand conditions $\Rightarrow$
$\mathbb{\hat{T}}\in\mathcal{R}_{F_{\mathfrak{x}_{1}}}\left(  \mathfrak{A}%
_{\ker}\right)  $ is obvious.
\end{proof}
\end{theorem}

As a result of this theorem one can construct a representation $\mathcal{\hat
{T}}_{F_{\mathfrak{x}_{1}}}\left(  \mathfrak{A}_{\ker}\right)  $ of
$\mathfrak{A}_{\ker}$ by superoperators of the form
\begin{equation}
\hat{T}=\hat{T}_{F}+\hat{T}_{FB}+\hat{T}\left[  \mathfrak{x}_{1}\right]  _{B}%
\end{equation}
where the components mapping $\left[  \ker L_{N}^{1}\right]  ^{c}%
\rightarrow\left[  \ker L_{N}^{1}\right]  ^{c}$ are given by%
\begin{align}
\hat{T}\left[  \mathfrak{x}_{1}\right]  _{B}  &  =\hat{t}_{\left\vert
\mathfrak{x}_{1}\right)  }+\hat{t}_{\left\vert \mathfrak{x}_{1}\right)  ^{c}%
}=\left(  \underline{\left\vert \mathfrak{x}\right)  }\right)  _{R}%
\mathbb{T}\left[  \mathfrak{x}_{1}\right]  _{B}\left(  \underline{\left(
\mathfrak{x}^{d}\right\vert }\right)  _{R}^{t}\nonumber\\
\hat{t}_{\left\vert \mathfrak{x}_{1}\right)  }  &  =\eta\left\vert
\mathfrak{x}_{1}\right)  \left(  \mathfrak{x}_{1}^{d}\right\vert ;\quad
\eta\neq0\in\mathbb{C}\nonumber\\
\hat{t}_{\left\vert \mathfrak{x}_{1}\right)  ^{c}}  &  =%
%TCIMACRO{\dsum \limits_{2\leq\kappa,\lambda\leq\infty}}%
%BeginExpansion
{\displaystyle\sum\limits_{2\leq\kappa,\lambda\leq\infty}}
%EndExpansion
\hat{T}_{B\kappa\lambda}\left\vert \mathfrak{x}_{\kappa}\right)  \left(
\mathfrak{x}_{\lambda}^{d}\right\vert =\left(  \underline{\left\vert
\mathfrak{x}\right)  }\right)  _{R}\mathbb{T}\left[  \mathfrak{x}_{1}%
^{c}\right]  _{B}\left(  \underline{\left(  \mathfrak{x}^{d}\right\vert
}\right)  _{R}^{t}%
\end{align}
with expansion coefficients given by
\begin{align}
&  \hat{T}_{B\kappa\lambda}=\left(  \mathfrak{x}_{\kappa}^{d}|\hat{T}\left[
\mathfrak{x}_{1}\right]  _{B}\mathfrak{x}_{\lambda}\right) \nonumber\\
&  =%
%TCIMACRO{\dint }%
%BeginExpansion
{\displaystyle\int}
%EndExpansion
Tr\left\{  \mathfrak{x}_{\lambda}^{d}\left(  \mathsf{z}_{1}\mathsf{z}%
^{N-1};\mathsf{z}_{1}^{\prime}\mathsf{z}^{N-1}\right)  ^{\ast}\left\vert
\mathsf{z}_{1}^{\prime}\mathsf{z}^{N-1}\right\rangle \left\langle
\mathsf{z}_{1}\mathsf{z}^{N-1}\right\vert T\left(  \mathsf{y}_{1}%
\mathsf{y}^{N-1}\right)  T\left(  \mathsf{y}_{1}^{\prime}\mathsf{y}%
^{N-1}\right)  ^{\ast}\left\vert \mathsf{y}_{1}^{\prime}\mathsf{y}%
^{N-1}\right\rangle \left\langle \mathsf{y}_{1}\mathsf{y}^{N-1}\right\vert
\right. \nonumber\\
&  \qquad\left.  \mathfrak{x}_{\kappa}\left(  \mathsf{x}_{1}\mathsf{x}%
^{N-1};\mathsf{x}_{1}^{\prime}\mathsf{x}^{N-1}\right)  \left\vert
\mathsf{x}_{1}\mathsf{x}^{N-1}\right\rangle \left\langle \mathsf{x}%
_{1}^{\prime}\mathsf{x}^{N-1}\right\vert \right\}  d\mathsf{z}_{1}%
d\mathsf{z}_{1}^{\prime}d\mathsf{z}^{N-1}d\mathsf{y}_{1}d\mathsf{y}%
_{1}^{\prime}d\mathsf{y}^{N-1}d\mathsf{x}_{1}d\mathsf{x}_{1}^{\prime
}d\mathsf{x}^{N-1}\nonumber\\
& \nonumber\\
&  =%
%TCIMACRO{\dint }%
%BeginExpansion
{\displaystyle\int}
%EndExpansion
\mathfrak{x}_{\kappa}^{d}\left(  \mathsf{z}_{1}\mathsf{z}^{N-1};\mathsf{z}%
_{1}^{\prime}\mathsf{z}^{N-1}\right)  ^{\ast}T\left(  \mathsf{z}_{1}%
\mathsf{z}^{N-1}\right)  T\left(  \mathsf{z}_{1}^{\prime}\mathsf{z}%
^{N-1}\right)  ^{\ast}\mathfrak{x}_{\kappa}\left(  \mathsf{z}_{1}%
\mathsf{z}^{N-1};\mathsf{z}_{1}^{\prime}\mathsf{z}^{N-1}\right)
d\mathsf{z}_{1}d\mathsf{z}_{1}^{\prime}d\mathsf{z}^{N-1}%
\end{align}
The components mixing $\left[  \ker L_{N}^{1}\right]  ^{c}$ with $\left[  \ker
L_{N}^{1}\right]  $ by
\begin{equation}
\hat{T}_{FB}=%
%TCIMACRO{\dsum \limits_{1\leq\kappa,\lambda\leq\infty}}%
%BeginExpansion
{\displaystyle\sum\limits_{1\leq\kappa,\lambda\leq\infty}}
%EndExpansion
\hat{T}_{\mathcal{K}_{D}B\kappa}\left\vert k_{\kappa}\right)  \left(
\mathfrak{x}_{\lambda}^{d}\right\vert =\left(  \underline{\left\vert k\right)
}\right)  _{R}\mathbb{T}_{\mathcal{K}_{D}B}\left(  \underline{\left(
\mathfrak{x}^{d}\right\vert }\right)  _{R}^{t}%
\end{equation}
and the expansion coefficients by
\begin{align}
&  \hat{T}_{FB\kappa\lambda}=\left(  k_{\kappa}^{d}|\hat{T}\mathfrak{x}%
_{\lambda}\right)  \qquad\qquad\qquad\qquad\qquad\qquad\qquad\qquad
\qquad\qquad\qquad\qquad\qquad\nonumber\\
&  =%
%TCIMACRO{\dint }%
%BeginExpansion
{\displaystyle\int}
%EndExpansion
k_{\kappa}^{d}\left(  \mathsf{z}_{1}\mathsf{z}^{N-1};\mathsf{z}_{1}^{\prime
}\mathsf{z}^{N-1}\right)  ^{\ast}T\left(  \mathsf{z}_{1}\mathsf{z}%
^{N-1}\right)  T\left(  \mathsf{z}_{1}^{\prime}\mathsf{z}^{N-1}\right)
^{\ast}\mathfrak{x}_{\lambda}\left(  \mathsf{z}_{1}\mathsf{z}^{N-1}%
;\mathsf{z}_{1}^{\prime}\mathsf{z}^{N-1}\right)  d\mathsf{z}_{1}%
d\mathsf{z}_{1}^{\prime}d\mathsf{z}^{N-1}%
\end{align}
The components mapping $\left[  \ker L_{N}^{1}\right]  \rightarrow\left[  \ker
L_{N}^{1}\right]  $ are given by
\begin{align}
\hat{T}_{F}  &  =%
%TCIMACRO{\dint \limits_{\mathsf{z}^{N}\overset{2}{\neq}\mathsf{z}^{\prime N}%
%}}%
%BeginExpansion
{\displaystyle\int\limits_{\mathsf{z}^{N}\overset{2}{\neq}\mathsf{z}^{\prime
N}}}
%EndExpansion
T\left(  \mathsf{z}^{N}\right)  T\left(  \mathsf{z}^{N\prime}\right)
\left\vert \left\vert \mathsf{z}^{N}\right\rangle \left\langle \mathsf{z}%
^{N\prime}\right\vert \right)  \left(  \left\vert \mathsf{z}^{N}\right\rangle
\left\langle \mathsf{z}^{N\prime}\right\vert \right\vert d\mathsf{z}%
^{N}d\mathsf{z}^{\prime N}\nonumber\\
&  \qquad\qquad\qquad\qquad\qquad\qquad\qquad\qquad+%
%TCIMACRO{\dsum \limits_{1\leq\kappa,\lambda\leq\infty}}%
%BeginExpansion
{\displaystyle\sum\limits_{1\leq\kappa,\lambda\leq\infty}}
%EndExpansion
\hat{T}_{F\kappa\lambda}\left\vert k_{\kappa}\right)  \left(  k_{\lambda}%
^{d}\right\vert
\end{align}
and the expansion coefficients constituting the matrix $\mathbb{T}_{F}$ by%
\begin{align}
&  \hat{T}_{F\kappa\lambda}=\left(  k_{\kappa}^{d}|\hat{T}k_{\lambda}\right)
=\qquad\qquad\qquad\qquad\qquad\qquad\qquad\qquad\qquad\qquad\qquad
\qquad\qquad\nonumber\\
&
%TCIMACRO{\dint }%
%BeginExpansion
{\displaystyle\int}
%EndExpansion
k_{\kappa}^{d}\left(  \mathsf{z}_{1}\mathsf{z}^{N-1};\mathsf{z}_{1}^{\prime
}\mathsf{z}^{N-1}\right)  ^{\ast}T\left(  \mathsf{z}_{1}\mathsf{z}%
^{N-1}\right)  T\left(  \mathsf{z}_{1}^{\prime}\mathsf{z}^{N-1}\right)
^{\ast}k_{\lambda}\left(  \mathsf{z}_{1}\mathsf{z}^{N-1};\mathsf{z}%
_{1}^{\prime}\mathsf{z}^{N-1}\right)  d\mathsf{z}_{1}d\mathsf{z}_{1}^{\prime
}d\mathsf{z}^{N-1}%
\end{align}
producing
\begin{equation}
\hat{T}_{F}=\left(  \underline{\left\vert \left\vert \mathsf{z}^{N}%
\right\rangle _{2}\left\langle \mathsf{z}^{N\prime}\right\vert \right)
}\right)  _{R}\mathbb{T}_{F\mathcal{K}_{OD}}\left(  \underline{\left(
\left\vert \mathsf{z}^{N}\right\rangle _{2}\left\langle \mathsf{z}^{N\prime
}\right\vert \right\vert }\right)  _{R}^{t}+\left(  \underline{\left\vert
k\right)  }\right)  _{R}\mathbb{T}_{F\mathcal{K}_{D}}\left(  \underline
{\left(  k^{d}\right\vert }\right)  _{R}^{t}%
\end{equation}
The magnitude of the integral kernel $T\left(  \mathsf{z}^{N}\right)  $ is
thus determined by
\begin{align}
\left\vert T\left(  \mathsf{z}^{N}\right)  \right\vert ^{2}  &  =\left(
\left\vert \mathsf{z}^{N}\right\rangle \left\langle \mathsf{z}^{N}\right\vert
\left\vert \hat{T}\right.  \left\vert \mathsf{z}^{N}\right\rangle \left\langle
\mathsf{z}^{N}\right\vert \right) \nonumber\\
&  =\left(  \left\vert \mathsf{z}^{N}\right\rangle \left\langle \mathsf{z}%
^{N}\right\vert \left\vert \left\{  \hat{T}_{F}+\hat{T}_{FB}+\hat{T}\left[
\mathfrak{x}_{1}\right]  _{B}\right\}  \right.  \left\vert \mathsf{z}%
^{N}\right\rangle \left\langle \mathsf{z}^{N}\right\vert \right) \nonumber\\
&  =\left(  \left\vert \mathsf{z}^{N}\right\rangle \left\langle \mathsf{z}%
^{N}\right\vert \left\vert \left\{  \hat{T}_{F\mathcal{K}_{OD}}+\hat
{T}_{F\mathcal{K}_{D}}+\hat{T}_{FB}+\hat{T}\left[  \mathfrak{x}_{1}\right]
_{B}\right\}  \right.  \left\vert \mathsf{z}^{N}\right\rangle \left\langle
\mathsf{z}^{N}\right\vert \right)
\end{align}
leading to
\begin{align}
\left\vert T\left(  \mathsf{z}^{N}\right)  \right\vert ^{2}  &  =\left(
\left\vert \mathsf{z}^{N}\right\rangle \left\langle \mathsf{z}^{N}\right\vert
\left\vert \underline{\left\vert k\right)  }\right.  \right)  \mathbb{T}%
_{F\mathcal{K}_{D}}\left(  \left.  \underline{\left(  k^{d}\right\vert
}\right\vert \left\vert \mathsf{z}^{N}\right\rangle \left\langle
\mathsf{z}^{N}\right\vert \right)  ^{t}\nonumber\\
&  +\left(  \left\vert \mathsf{z}^{N}\right\rangle \left\langle \mathsf{z}%
^{N}\right\vert \left\vert \underline{\left\vert k\right)  }\right.  \right)
\mathbb{T}_{\mathcal{K}_{D}B}\left(  \left.  \underline{\left(  \mathfrak{x}%
^{d}\right\vert }\right\vert \left\vert \mathsf{z}^{N}\right\rangle
\left\langle \mathsf{z}^{N}\right\vert \right)  ^{t}\nonumber\\
&  +\left(  \left\vert \mathsf{z}^{N}\right\rangle \left\langle \mathsf{z}%
^{N}\right\vert \left\vert \left\vert \mathfrak{x}_{1}\right)  \left(
\mathfrak{x}_{1}^{d}\right\vert \right.  \left\vert \mathsf{z}^{N}%
\right\rangle \left\langle \mathsf{z}^{N}\right\vert \right)
\end{align}
and finally to
\begin{align}
&  \left\vert T\left(  \mathsf{z}^{N}\right)  \right\vert ^{2}=\nonumber\\
&  \sum_{2\leq\kappa,\lambda\leq\infty}\left\{  k_{\kappa}\left(
\mathsf{z}^{N};\mathsf{z}^{N}\right)  \mathbb{T}_{F\mathcal{K}_{D\,}%
\kappa\lambda}k_{\lambda}^{d}\left(  \mathsf{z}^{N};\mathsf{z}^{N}\right)
^{\ast}+k_{\kappa}\left(  \mathsf{z}^{N};\mathsf{z}^{N}\right)  \mathbb{T}%
_{\mathcal{K}_{D}B}\,_{\kappa\lambda}\mathfrak{x}_{\lambda}^{d}\left(
\mathsf{z}^{N};\mathsf{z}^{N}\right)  ^{\ast}\right\} \nonumber\\
&  \qquad\qquad\qquad\qquad\qquad\qquad\qquad\qquad\qquad\qquad+\eta
\mathfrak{x}_{1}\left(  \mathsf{z}^{N};\mathsf{z}^{N}\right)  \mathfrak{x}%
_{1}^{d}\left(  \mathsf{z}^{N};\mathsf{z}^{N}\right)  ^{\ast}%
\end{align}
One should observe that the magnitudes $\left\{  \left\vert T\left(
\mathsf{z}^{N}\right)  \right\vert \right\}  $ constitute the representation
$\mathcal{R}_{\mathfrak{B}_{N}}\left(  \mathfrak{A}_{\ker}\right)  $ contained
in $\mathcal{R}_{F_{\mathfrak{x}_{1}}}\left(  \mathfrak{A}_{\ker}\right)  $
and \emph{determine} the representation $\mathcal{R}_{F_{\mathfrak{x}_{1}}%
}\left(  \mathfrak{A}_{\ker}\right)  $ up to the phase of integral kernels
$\left\{  T\left(  \mathsf{z}^{N}\right)  \right\}  $.

\subsubsection{Non Redundant Parameterization of States}

It is not possible to homogeneously parameterize a general fiber $\left[
\mathfrak{x}_{1}\right]  $ by elements of $\mathcal{\hat{T}}_{F_{\mathfrak{x}%
_{1}}}\left(  \mathfrak{A}_{\ker}\right)  $ \cite{weiner2003} \emph{but it is
}possible to parameterize all the states $\left\{  \rho_{1}\right\}  $ in the
fiber $\left[  \rho_{1}\right]  $ where $\rho_{1}\in\mathfrak{C}_{1}$ by
elements of $\mathcal{\hat{T}}_{F_{\rho_{1}}}\left(  \mathfrak{A}_{\ker
}\right)  $. However in order to construct non redundant parameterization of
$\left\{  \rho_{1}\right\}  $ one needs to put restrictions on which elements
of $\mathcal{\hat{T}}_{F_{\rho_{1}}}\left(  \mathfrak{A}_{\ker}\right)  $ one
uses so that the set of states $\left\{  \rho_{1}\right\}  $ in such a fiber
are given by $\left\{  \hat{T}_{F\rho_{1}}\left(  \rho_{1}\right)  \right\}  $
in a unique fashion. This is accomplished by noting that the representation
$\mathcal{\hat{T}}_{F\rho_{1}}$ contains the sub representation
\begin{align}
\mathcal{\hat{T}}_{\ker L_{N}^{1}\oplus\left\vert \mathfrak{\rho}_{1}\right)
}\left(  \mathfrak{A}_{\ker}\right)   &  =\left(  \underline{\left\vert
\left\vert \mathsf{z}^{N}\right\rangle _{2}\left\langle \mathsf{z}^{\prime
N}\right\vert \right)  },\underline{\left\vert k\right)  },\underline
{\left\vert \mathfrak{x}\right)  }\right)  \mathcal{R}_{F_{\ker L_{N}%
^{1}\oplus\left\vert \mathfrak{\rho}_{1}\right)  }}\left(  \mathfrak{A}_{\ker
}\right) \nonumber\\
&  =\left\{  \hat{T}|\hat{T}\in\mathcal{\hat{T}}_{F\rho_{1}}\left(
\mathfrak{A}_{\ker}\right)  ;\hat{T}:\ker L_{N}^{1}\oplus\left\vert
\mathfrak{\rho}_{1}\right)  \rightarrow\ker L_{N}^{1}\oplus\left\vert
\mathfrak{\rho}_{1}\right)  \right\}
\end{align}
i.e.%
\begin{equation}
\mathcal{\hat{T}}_{F\rho_{1}}\supset\mathcal{\hat{T}}_{\ker L_{N}^{1}%
\oplus\left\vert \mathfrak{\rho}_{1}\right)  }\left(  \mathfrak{A}_{\ker
}\right)
\end{equation}
and only considering transformations belonging to the coset of super operator
transformations, $\mathcal{C}_{F_{\rho_{1}}}$, defined by
\begin{equation}
\mathcal{C}_{F_{\rho_{1}}}=\mathcal{\hat{T}}_{\ker L_{N}^{1}\oplus\left\vert
\mathfrak{\rho}_{1}\right)  }\left(  \mathfrak{A}_{\ker}\right)  /\left(
\mathcal{\hat{T}}_{\ker L_{N}^{1}}\left(  \mathfrak{A}_{\ker}\right)
\times\mathcal{\hat{T}}_{\left\vert \mathfrak{\rho}_{1}\right)  }\left(
\mathfrak{A}_{\ker}\right)  \right)  .
\end{equation}
which provides a representation of a coset $\mathfrak{Co}_{1}$ of
$\mathfrak{A}_{\ker}$. Explicitly this can be described as limiting oneself to
by super operators of the form%
\begin{equation}
\hat{T}_{F\rho_{1}}=\hat{T}\left[  \rho_{1}\right]  _{F}+\hat{T}\left[
\rho_{1}\right]  _{FB}+\hat{t}_{\left\vert \rho_{1}\right)  }%
\end{equation}
where%
\begin{equation}
\hat{T}\left[  \rho_{1}\right]  _{FB}=\hat{T}\left[  \rho_{1}\right]
_{\mathcal{K}_{D}B}=%
%TCIMACRO{\dsum \limits_{1\leq\kappa\leq\infty}}%
%BeginExpansion
{\displaystyle\sum\limits_{1\leq\kappa\leq\infty}}
%EndExpansion
\hat{T}\left[  \rho_{1}\right]  _{\mathcal{K}_{D}B\kappa}\left\vert k_{\kappa
}\right)  \left(  \mathfrak{\rho}_{1}^{d}\right\vert =\left(  \underline
{\left\vert k\right)  }\right)  _{R}\mathbb{T}\left[  \rho_{1}\right]
_{\mathcal{K}_{D}B}\left(  \underline{\left(  \mathfrak{x}^{d}\right\vert
}\right)  _{R}^{t}%
\end{equation}
as $\hat{T}\left[  \rho_{1}\right]  _{\mathcal{K}_{OD}B}=0$. In terms of the
generalized basis $\left\{  \mathfrak{b}_{0},\cdots,\mathfrak{b}_{N}\right\}
$ this term can be expanded as
\begin{equation}
\hat{T}\left[  \rho_{1}\right]  _{\mathcal{K}_{D}B}=\hat{T}\left[  \rho
_{1}\right]  _{FB}{}_{\mathfrak{B}_{N}}+\hat{T}\left[  \rho_{1}\right]
_{FB\mathfrak{B}_{N-1}}%
\end{equation}
The fiber component, $\hat{T}\left[  \rho_{1}\right]  _{F}$, of $\hat
{T}\left[  \rho_{1}\right]  $ is given by%
\begin{equation}
\hat{T}\left[  \rho_{1}\right]  _{F}=\hat{T}\left[  \rho_{1}\right]
_{F\mathcal{K}_{OD}}%
\end{equation}
as $\hat{T}\left[  \rho_{1}\right]  _{F\mathcal{K}_{D}}=0$, hence
\begin{align}
&  \hat{T}\left[  \rho_{1}\right]  _{F}=\hat{T}\left[  \rho_{1}\right]
_{F\mathcal{K}_{OD}}=%
%TCIMACRO{\dint \limits_{\mathsf{z}^{N}\overset{2}{\,\nsim\,}\mathsf{z}^{\prime
%N}}}%
%BeginExpansion
{\displaystyle\int\limits_{\mathsf{z}^{N}\overset{2}{\,\nsim\,}\mathsf{z}%
^{\prime N}}}
%EndExpansion
\left(  \left\vert \mathsf{z}^{N}\right\rangle \left\langle \mathsf{z}%
^{N\prime}\right\vert \left\vert \hat{T}\left[  \rho_{1}\right]  \right.
\left\vert \mathsf{z}^{N}\right\rangle \left\langle \mathsf{z}^{N\prime
}\right\vert \right)  \left\vert \left\vert \mathsf{z}^{N}\right\rangle
\left\langle \mathsf{z}^{N\prime}\right\vert \right)  \left(  \left\vert
\mathsf{z}^{N}\right\rangle \left\langle \mathsf{z}^{N\prime}\right\vert
\right\vert d\mathsf{z}^{N}d\mathsf{z}^{\prime N}\nonumber\\
&  \qquad\qquad\qquad\qquad\quad\left.  =\right.
%TCIMACRO{\dint \limits_{\mathsf{z}^{N}\overset{2}{\,\nsim\,}\mathsf{z}^{\prime
%N}}}%
%BeginExpansion
{\displaystyle\int\limits_{\mathsf{z}^{N}\overset{2}{\,\nsim\,}\mathsf{z}%
^{\prime N}}}
%EndExpansion
T\left(  \mathsf{z}^{N}\right)  T\left(  \mathsf{z}^{\prime N}\right)  ^{\ast
}\left\vert \mathsf{z}^{\prime N}\right\rangle \left\langle \mathsf{z}%
^{N}\right\vert d\mathsf{z}^{N}d\mathsf{z}^{\prime N}\in\mathfrak{B}_{2}%
\end{align}
The base term $\hat{t}_{\left\vert \rho_{1}\right)  }$ can be expanded as
\begin{equation}
\hat{t}_{\left\vert \rho_{1}\right)  }=\hat{t}_{\left\vert \rho_{1}\right)
\mathfrak{B}_{N}}+\hat{t}_{\left\vert \rho_{1}\right)  \mathfrak{B}_{N-1}}%
\end{equation}
The terms $\hat{t}_{\left\vert \rho_{1}\right)  \mathfrak{B}_{N-1}}$ and
$\hat{T}\left[  \rho_{1}\right]  _{FB\mathfrak{B}_{N-1}}$ are determined up to
phase factors by $\hat{t}_{\left\vert \rho_{1}\right)  \,\mathfrak{B}_{N}}$
$+$ $\hat{T}\left[  \mathfrak{x}_{1}\right]  _{FB\,\mathfrak{B}_{N}}$, which
can be expanded as
\begin{align}
&  \hat{t}_{\left\vert \mathfrak{x}_{1}\right)  \,\mathfrak{B}_{N}}+\hat
{T}\left[  \mathfrak{x}_{1}\right]  _{FB\,\mathfrak{B}_{N}}\nonumber\\
&  =%
%TCIMACRO{\dint \limits_{\mathsf{z}^{N}\overset{N-1}{\,=\,}\mathsf{z}^{\prime
%N}}}%
%BeginExpansion
{\displaystyle\int\limits_{\mathsf{z}^{N}\overset{N-1}{\,=\,}\mathsf{z}%
^{\prime N}}}
%EndExpansion
\mathfrak{\rho}_{1}^{d}\left(  \mathsf{z}^{N}\right)  ^{\ast}\mathfrak{\rho
}_{1}\left(  \mathsf{z}^{\prime N}\right)  \left\vert \mathsf{z}^{\prime
N}\right\rangle \left\langle \mathsf{z}^{N}\right\vert d\mathsf{z}%
^{N}d\mathsf{z}^{\prime N}+\eta%
%TCIMACRO{\dint }%
%BeginExpansion
{\displaystyle\int}
%EndExpansion
\mathfrak{\rho}_{1}^{d}\left(  \mathsf{z}^{N}\right)  ^{\ast}\mathfrak{\rho
}_{1}\left(  \mathsf{z}^{N}\right)  \left\vert \mathsf{z}^{N}\right\rangle
\left\langle \mathsf{z}^{N}\right\vert d\mathsf{z}^{N}%
\end{align}
In matrix form these transformations are%
\begin{equation}
\mathcal{\hat{T}}_{F\rho_{1}}\left(  s\right)  =\left(  \left(  \underline
{\left\vert \left\vert \mathsf{z}^{N}\right\rangle _{2}\left\langle
\mathsf{z}^{\prime N}\right\vert \right)  }\right)  ,\left(  \underline
{\left\vert k\right)  }\right)  ,\left(  \underline{\left\vert \mathfrak{x}%
\right)  }\right)  \right)  \mathcal{R}_{F\rho_{1}}^{0}\left(  s\right)
;\,\mathcal{\hat{T}}_{F\rho_{1}}\left(  s\right)  \in\mathcal{C}_{F_{\rho_{1}%
}};s\in\mathfrak{Co}_{1}%
\end{equation}
where
\begin{equation}
\mathcal{R}_{F\rho_{1}}^{0}\left(  s\right)  =\left(
\begin{array}
[c]{cccc}%
\mathbb{\hat{T}}\left[  \rho_{1}\right]  _{F\mathcal{K}_{OD}} & \mathbb{O} &
\mathbb{O} & \mathbb{O}\\
\mathbb{O} & \mathbb{O}_{\mathcal{K}_{D}} & \mathbb{O} & \mathbb{T}\left[
\rho_{1}\right]  _{\mathcal{K}_{D}B}\\
\mathbb{O} & \mathbb{O} & \mathbb{O}_{\left\vert \mathfrak{\rho}_{1}\right)
^{c}} & \mathbb{O}\\
\mathbb{O} & \mathbb{O} & \mathbb{O} & \eta
\end{array}
\right)
\end{equation}
and the matrix $\mathbb{\hat{T}}\left[  \rho_{1}\right]  _{F\mathcal{K}_{OD}}$
and the column vector $\mathbb{T}\left[  \rho_{1}\right]  _{\mathcal{K}_{D}B}$
depend on $s\in\mathfrak{Co}_{1}$.

The non canonical matrix representations are given by
\begin{equation}
\mathcal{R}_{F_{\tilde{\Gamma}\left(  \rho_{1}\right)  }}\left(
\mathfrak{A}_{_{\ker}}\right)  =\mathbf{\Upsilon}_{\Gamma}\mathcal{R}%
_{F\rho_{1}}\left(  \mathfrak{A}_{_{\ker}}\right)  \mathbf{\Upsilon}_{\Gamma
}^{-1}%
\end{equation}
and the coset of super operators $\mathcal{C}_{F_{\tilde{\Gamma}\left(
\rho_{1}\right)  }}$ by
\begin{equation}
\mathcal{C}_{F_{\tilde{\Gamma}\left(  \rho_{1}\right)  }}=\Upsilon_{\Gamma
}\mathcal{C}_{F\rho_{1}}\left(  \mathfrak{A}_{_{\ker}}\right)  \Upsilon
_{\Gamma}^{-1}%
\end{equation}
The super operators, $\hat{T}\in\mathcal{C}_{F_{\tilde{\Gamma}\left(  \rho
_{1}\right)  }}$, have in particular the property
\begin{equation}
\left(  \mathcal{D}_{0_{F}}^{N}\left(  \rho_{1}\right)  ^{d}|\hat
{T}\mathcal{D}_{0_{F}}^{N}\left(  \rho_{1}\right)  \right)  =\eta\neq0
\end{equation}
where $\mathcal{D}_{0_{F}}^{N}\left(  \rho_{1}\right)  =\tilde{\Gamma}\left(
\rho_{1}\right)  $.

The coset $\mathcal{C}_{F_{\tilde{\Gamma}\left(  \rho_{1}\right)  }}$ provides
a non redundant parameterization of the states $\left\{  \mathcal{D}%
^{N}\left(  \rho_{1}\right)  \right\}  $ in the fiber $\left[  \mathcal{D}%
_{0_{F}}^{N}\left(  \rho_{1}\right)  \right]  $ as $\left\{  \left.  \hat
{T}\left(  \mathcal{D}_{0_{F}}^{N}\left(  \rho_{1}\right)  \right)
\right\vert \hat{T}\in\mathcal{C}_{F_{\tilde{\Gamma}\left(  \rho_{1}\right)
}}\right\}  $ and the non canonical representative matrices of the
transformations parameterizing the fiber $\left[  \mathfrak{\rho}_{1}\right]
$ are
\begin{align}
&  \mathbb{\hat{T}}_{F_{\tilde{\Gamma}\left(  \rho_{1}\right)  }%
}=\mathbf{\Upsilon}_{\Gamma}\mathbb{\hat{T}}_{F_{\rho_{1}}}\mathbf{\Upsilon
}_{\Gamma}^{-1}=\nonumber\\
&  \left(
\begin{array}
[c]{ccc}%
\mathbb{I}_{\mathcal{K}_{OD}} & \mathbb{O} & \mathbf{\Gamma}_{\mathcal{K}%
_{OD}}\\
\mathbb{O} & \mathbb{I}_{\mathcal{K}_{D}} & \mathbf{\Gamma}_{\mathcal{K}_{D}%
}\\
\mathbb{O} & \mathbb{O} & \mathbb{I}_{\left[  \ker L_{N}^{1}\right]  ^{c}}%
\end{array}
\right)  \left(
\begin{array}
[c]{cccc}%
\mathbb{\hat{T}}\left[  \rho_{1}\right]  _{F\mathcal{K}_{OD}} & \mathbb{O} &
\mathbb{O} & \mathbb{O}\\
\mathbb{O} & \mathbb{O}_{\mathcal{K}_{D}} & \mathbb{O} & \mathbb{T}\left[
\rho_{1}\right]  _{\mathcal{K}_{D}B}\\
\mathbb{O} & \mathbb{O} & \mathbb{O}_{\left\vert \mathfrak{\rho}_{1}\right)
^{c}} & \mathbb{O}\\
\mathbb{O} & \mathbb{O} & \mathbb{O} & 1
\end{array}
\right)  \times\\
&  \qquad\qquad\qquad\qquad\qquad\qquad\qquad\qquad\qquad\times\left(
\begin{array}
[c]{ccc}%
\mathbb{I}_{\mathcal{K}_{OD}} & \mathbb{O} & -\mathbf{\Gamma}_{\mathcal{K}%
_{OD}}\\
\mathbb{O} & \mathbb{I}_{\mathcal{K}_{D}} & -\mathbf{\Gamma}_{\mathcal{K}_{D}%
}\\
\mathbb{O} & \mathbb{O} & \mathbb{I}_{\left[  \ker L_{N}^{1}\right]  ^{c}}%
\end{array}
\right)
\end{align}
so that
\begin{align}
\left\{  \rho_{1}\right\}   &  =\left\{  \left.  \hat{T}\left\vert \overset
{}{\Upsilon_{\Gamma}\left(  \rho_{1}\right)  }\right)  \right\vert \text{fixed
}\Upsilon_{\Gamma}\in\mathfrak{P};\hat{T}\in\mathcal{C}_{F_{\tilde{\Gamma
}\left(  \rho_{1}\right)  }}\right\} \nonumber\\
&  \equiv\left\{  \left(  \left(  \underline{\left\vert \left\vert
\mathsf{z}^{N}\right\rangle _{2}\left\langle \mathsf{z}^{\prime N}\right\vert
\right)  }\right)  ,\left(  \underline{\left\vert k\right)  }\right)  ,\left(
\underline{\left\vert \mathfrak{x}\right)  }\right)  \right)  \mathbb{\hat{T}%
}_{F_{\tilde{\Gamma}\left(  \rho_{1}\right)  }}\left(
\begin{array}
[c]{c}%
\mathbf{\Gamma}_{\mathcal{K}_{OD}}\rho_{1}\\
\mathbf{\Gamma}_{\mathcal{K}_{D}}\rho_{1}\\
\rho_{1}%
\end{array}
\right)  \right. \nonumber\\
&  \qquad\qquad\qquad\qquad\qquad\qquad\qquad\qquad\left.  \text{fixed
}\Upsilon_{\Gamma}\in\mathfrak{P};s\in\mathfrak{Co}_{1}\right\}
\end{align}
It is especially important to note that the super operator $\hat{T}_{F\rho
_{1}}$ can also be expanded in the form
\begin{equation}
\hat{T}_{F\rho_{1}}=\int T_{F\rho_{1}}\left(  \mathsf{z}^{N}\right)
T_{F\rho_{1}}\left(  \mathsf{z}^{\prime N}\right)  ^{\ast}\left\vert
\left\vert \mathsf{z}^{N}\right\rangle \left\langle \mathsf{z}^{N\prime
}\right\vert \right)  \left(  \left\vert \mathsf{z}^{N}\right\rangle
\left\langle \mathsf{z}^{N\prime}\right\vert \right\vert d\mathsf{z}%
^{N}d\mathsf{z}^{\prime N}%
\end{equation}


The magnitude of the integral kernels $T\left(  \mathsf{z}^{N}\right)  $ are
given in this case by
\begin{align}
&  \left\vert \hat{T}_{F\rho_{1}}\left(  \mathsf{z}^{N}\right)  \right\vert
^{2}=\nonumber\\
&  \sum_{2\leq\kappa,\lambda\leq\infty}\left\{  k_{\kappa}\left(
\mathsf{z}^{N};\mathsf{z}^{N}\right)  \mathbb{T}_{F\mathcal{K}_{D}%
\kappa\lambda}k_{\lambda}^{d}\left(  \mathsf{z}^{N};\mathsf{z}^{N}\right)
^{\ast}+k_{\kappa}\left(  \mathsf{z}^{N};\mathsf{z}^{N}\right)  \mathbb{T}%
_{\mathcal{K}_{D}B}{}_{\kappa\lambda}\mathfrak{x}_{\lambda}^{d}\left(
\mathsf{z}^{N};\mathsf{z}^{N}\right)  ^{\ast}\right\} \nonumber\\
&  \qquad\qquad\qquad\qquad\qquad\qquad\qquad\qquad+\eta\mathcal{D}_{0_{F}%
}^{N}\left[  \rho_{1}\right]  \left(  \mathsf{z}^{N};\mathsf{z}^{N}\right)
\mathcal{D}_{0_{F}}^{N}\left[  \rho_{1}\right]  ^{d}\left(  \mathsf{z}%
^{N};\mathsf{z}^{N}\right)  ^{\ast}%
\end{align}


\subsection{Base Representations of $\mathfrak{A}_{\ker}\label{baserep}$}

In this case we wish to start from a reference state $\mathcal{D}_{0_{B}}%
^{N}=\Upsilon\left(  \mathfrak{\rho}_{r}\right)  $ and generate reference
states $\left\{  \mathcal{D}_{0_{F}}^{N}\left(  \rho_{1}\right)  \right\}  $
in a non redundant manner for all of the fibers that contain states. In order
to accomplish this one uses superoperators, associated with transformations
that belong to $\hat{T}_{B_{\mathfrak{\rho}_{r}}}\left(  \mathfrak{A}_{\ker
}\right)  $, of the form $\hat{T}_{B_{\tilde{\Gamma}\left(  \mathfrak{\rho
}_{r}\right)  }}=\Upsilon_{\Gamma}\hat{T}_{B_{\mathfrak{\rho}_{r}}}%
\Upsilon_{\Gamma}^{-1}$ where
\begin{align}
\hat{T}_{B_{\mathfrak{\rho}_{r}}}  &  =%
%TCIMACRO{\dint \limits_{\mathsf{z}^{N}\overset{2}{\nsim}\mathsf{z}^{\prime N}%
%}}%
%BeginExpansion
{\displaystyle\int\limits_{\mathsf{z}^{N}\overset{2}{\nsim}\mathsf{z}^{\prime
N}}}
%EndExpansion
T_{B_{\mathfrak{\rho}_{r}}}\left(  \mathsf{z}^{N}\right)  T_{B_{\mathfrak{\rho
}_{r}}}\left(  \mathsf{z}^{N\prime}\right)  \left\vert \left\vert
\mathsf{z}^{N}\right\rangle \left\langle \mathsf{z}^{N\prime}\right\vert
\right)  \left(  \left\vert \mathsf{z}^{N}\right\rangle \left\langle
\mathsf{z}^{N\prime}\right\vert \right\vert d\mathsf{z}^{N}d\mathsf{z}^{\prime
N}\nonumber\\
&  \qquad\qquad\qquad\qquad+%
%TCIMACRO{\dsum \limits_{1\leq\kappa\leq\infty}}%
%BeginExpansion
{\displaystyle\sum\limits_{1\leq\kappa\leq\infty}}
%EndExpansion
\left(  \mathfrak{x}_{\kappa}^{d}\left\vert \hat{T}\right.  \mathfrak{\rho
}_{r}\right)  \left\vert \mathfrak{x}_{\kappa}\right)  \left(  \mathfrak{\rho
}_{r}^{d}\right\vert
\end{align}
where $\mathfrak{x}_{1}=\mathfrak{\rho}_{r}$ and $\hat{T}_{B_{\mathfrak{\rho
}_{r}}}$ have matrix representatives wrt the canonical basis
\begin{equation}
\mathbb{\hat{T}}_{B_{\mathfrak{\rho}_{r}}}=\left(
\begin{array}
[c]{cccc}%
\mathbb{\hat{T}}\left[  \mathfrak{\rho}_{r}\right]  _{F\mathcal{K}_{OD}} &
\mathbb{O} & \mathbb{O} & \mathbb{O}\\
\mathbb{O} & \mathbb{O}_{\mathcal{K}_{D}} & \mathbb{O} & \mathbb{T}\left[
\mathfrak{\rho}_{r}\right]  _{\mathcal{K}_{D}B}\\
\mathbb{O} & \mathbb{O} & \mathbb{O}_{\left\vert \mathfrak{\rho}_{1}\right)
^{c}} & \left(
\begin{array}
[c]{c}%
\vdots\\
\left[  \mathfrak{\rho}_{r}\right]  _{\kappa}\\
\vdots
\end{array}
\right) \\
\mathbb{O} & \mathbb{O} & \mathbb{O} & \left[  \mathfrak{\rho}_{r}\right]
_{1}%
\end{array}
\right)  .
\end{equation}
The matrix representatives of $\hat{T}_{B_{\tilde{\Gamma}\left(
\mathfrak{\rho}_{r}\right)  }}$ are hence
\begin{align}
&  \mathbb{\hat{T}}_{B}=\nonumber\\
&  \left(
\begin{array}
[c]{ccc}%
\mathbb{I}_{\mathcal{K}_{OD}} & \mathbb{O} & \mathbf{\Gamma}_{\mathcal{K}%
_{OD}}\\
\mathbb{O} & \mathbb{I}_{\mathcal{K}_{D}} & \mathbf{\Gamma}_{\mathcal{K}_{D}%
}\\
\mathbb{O} & \mathbb{O} & \mathbb{I}_{\left[  \ker L_{N}^{1}\right]  ^{c}}%
\end{array}
\right)  \left(
\begin{array}
[c]{cccc}%
\mathbb{\hat{T}}\left[  \mathfrak{\rho}_{r}\right]  _{F\mathcal{K}_{OD}} &
\mathbb{O} & \mathbb{O} & \mathbb{O}\\
\mathbb{O} & \mathbb{O}_{\mathcal{K}_{D}} & \mathbb{O} & \mathbb{T}\left[
\mathfrak{\rho}_{r}\right]  _{\mathcal{K}_{D}B}\\
\mathbb{O} & \mathbb{O} & \mathbb{O}_{\left\vert \mathfrak{\rho}_{r}\right)
^{c}} & \left(
\begin{array}
[c]{c}%
\vdots\\
\left[  \mathfrak{\rho}_{B1}\right]  _{\kappa}\\
\vdots
\end{array}
\right) \\
\mathbb{O} & \mathbb{O} & \mathbb{O} & \left[  \mathfrak{\rho}_{r}\right]
_{1}%
\end{array}
\right)  \times\nonumber\\
&  \qquad\qquad\qquad\qquad\qquad\qquad\qquad\qquad\qquad\qquad\qquad
\times\left(
\begin{array}
[c]{ccc}%
\mathbb{I}_{\mathcal{K}_{OD}} & \mathbb{O} & -\mathbf{\Gamma}_{\mathcal{K}%
_{OD}}\\
\mathbb{O} & \mathbb{I}_{\mathcal{K}_{D}} & -\mathbf{\Gamma}_{\mathcal{K}_{D}%
}\\
\mathbb{O} & \mathbb{O} & \mathbb{I}_{\left[  \ker L_{N}^{1}\right]  ^{c}}%
\end{array}
\right)
\end{align}
As was the case for intra fiber transformation it is again especially
important to note that the inter fiber super operator $\hat{T}_{B_{\tilde
{\Gamma}\left(  \mathfrak{\rho}_{r}\right)  }}$ can also be expanded in the
diagonal form
\begin{equation}
\hat{T}_{B_{\tilde{\Gamma}\left(  \mathfrak{\rho}_{r}\right)  }}=\int
T_{B}\left(  \mathsf{z}^{N}\right)  T_{B}\left(  \mathsf{z}^{\prime N}\right)
^{\ast}\left\vert \left\vert \mathsf{z}^{N}\right\rangle \left\langle
\mathsf{z}^{N\prime}\right\vert \right)  \left(  \left\vert \mathsf{z}%
^{N}\right\rangle \left\langle \mathsf{z}^{N\prime}\right\vert \right\vert
d\mathsf{z}^{N}d\mathsf{z}^{\prime N}%
\end{equation}


Any $N$-particle state can thus be obtained by the action of $\hat
{T}_{B_{\tilde{\Gamma}\left(  \mathfrak{\rho}_{r}\right)  }}$ and $\hat
{T}_{F_{\tilde{\Gamma}\left(  \rho_{1}\right)  }}$ in a two step process
\begin{equation}
D^{N}\left(  \rho_{1}\right)  =\hat{T}_{F_{\tilde{\Gamma}\left(  \rho
_{1}\right)  }}\circ\hat{T}_{B_{\tilde{\Gamma}\left(  \mathfrak{\rho}%
_{r}\right)  }}\left(  \mathcal{D}_{0_{B}}^{N}\right)  ;\quad\hat
{T}_{F_{\tilde{\Gamma}\left(  \rho_{1}\right)  }}\in\mathcal{C}_{F_{\tilde
{\Gamma}\left(  \rho_{1}\right)  }},\hat{T}_{B_{\tilde{\Gamma}\left(
\mathfrak{\rho}_{r}\right)  }}\in\mathcal{C}_{F_{\tilde{\Gamma}\left(
\rho_{1}\right)  }}%
\end{equation}
the transformation traversing states in the base in a unique fashion and the
second transformation states in a particular fiber associated with a common FORDO.

\section{Properties of the Integral Kernels of the Representations of the
Semigroup $\mathfrak{A}_{\ker}\label{sgkernel}$}

It is convenient to express the integral kernels $\left\{  T\left(
\mathsf{z}^{N}\right)  \right\}  $ defining the super operators belonging
$\left\{  \mathcal{\hat{T}}\left(  \mathfrak{A}_{\ker}\right)  \right\}  $ to
as
\begin{equation}
T\left(  \mathsf{z}^{N}\right)  =\delta\left(  \mathsf{z}^{N}\right)
-\Phi\left(  \mathsf{z}^{N}\right)
\end{equation}
where
\begin{equation}
\delta\left(  \mathsf{z}^{N}\right)  =%
%TCIMACRO{\dprod \limits_{1\leq j\leq N}}%
%BeginExpansion
{\displaystyle\prod\limits_{1\leq j\leq N}}
%EndExpansion
\delta\left(  \mathsf{z}_{j}\right)
\end{equation}
leading to
\begin{equation}
T\left(  \mathsf{z}^{N}\right)  T\left(  \mathsf{z}^{\prime N}\right)
=\delta\left(  \mathsf{z}^{N}-\mathsf{z}^{\prime N}\right)  +\Lambda\left(
\mathsf{z}^{N},\mathsf{z}^{\prime N}\right)
\end{equation}%
\begin{equation}
\Lambda\left(  \mathsf{z}^{N},\mathsf{z}^{\prime N}\right)  =\Phi\left(
\mathsf{z}^{N}\right)  \Phi\left(  \mathsf{z}_{1}^{\prime},\mathsf{z}^{\prime
N}\right)  ^{\ast}-\Phi\left(  \mathsf{z}^{\prime N}\right)  -\Phi\left(
\mathsf{z}^{N}\right)  ^{\ast}.
\end{equation}
We can then prove the following theorems:

\begin{theorem}
Incremental Strong Orthogonality to the Kernel
\begin{align}
\hat{T}  &  \in\mathcal{\hat{T}}\left(  \mathfrak{A}_{\ker}\right)
\Leftrightarrow%
%TCIMACRO{\dint }%
%BeginExpansion
{\displaystyle\int}
%EndExpansion
\Lambda\left(  \mathsf{z}_{1},\mathsf{z}^{N-1},\mathsf{z}_{1}^{\prime
},\mathsf{z}^{N-1}\right)  X\left(  \mathsf{z}_{1},\mathsf{z}^{N-1}%
,\mathsf{z}_{1}^{\prime},\mathsf{z}^{N-1}\right)  d\mathsf{z}^{N-1}%
\overset{a.e}{=}0\nonumber\\
&  \qquad\qquad\qquad\qquad\qquad\qquad\qquad\qquad\qquad\qquad\qquad
\qquad\forall X\in\ker L_{N}^{1}%
\end{align}


\begin{proof}
$"\Rightarrow"$%
\begin{align}
&
%TCIMACRO{\dint }%
%BeginExpansion
{\displaystyle\int}
%EndExpansion
\left\{  \delta\left(  \mathsf{z}_{1}-\mathsf{z}_{1}^{\prime}\right)
\delta\left(  \mathsf{z}^{N-1}-\mathsf{z}^{\prime N-1}\right)  +\Lambda\left(
\mathsf{z}_{1},\mathsf{z}^{N-1},\mathsf{z}_{1}^{\prime},\mathsf{z}%
^{N-1}\right)  X\left(  \mathsf{z}_{1},\mathsf{z}^{N-1},\mathsf{z}_{1}%
^{\prime},\mathsf{z}^{N-1}\right)  \right. \nonumber\\
&  \qquad\qquad\qquad\qquad\qquad\qquad\qquad\qquad\qquad\qquad\qquad
\qquad\qquad\left.  d\mathsf{z}^{N-1}\overset{a.e}{=}0\right\} \nonumber\\
&  \Rightarrow%
%TCIMACRO{\dint }%
%BeginExpansion
{\displaystyle\int}
%EndExpansion
\left\{  \Lambda\left(  \mathsf{z}_{1},\mathsf{z}^{N-1},\mathsf{z}_{1}%
^{\prime},\mathsf{z}^{N-1}\right)  X\left(  \mathsf{z}_{1},\mathsf{z}%
^{N-1},\mathsf{z}_{1}^{\prime},\mathsf{z}^{N-1}\right)  \right\}
d\mathsf{z}^{N-1}\overset{a.e}{=}0\forall X\in\ker L_{N}^{1}%
\end{align}


$"\Leftarrow"$%

\begin{align}
&
%TCIMACRO{\dint }%
%BeginExpansion
{\displaystyle\int}
%EndExpansion
\left\{  \Lambda\left(  \mathsf{z}_{1},\mathsf{z}^{N-1},\mathsf{z}_{1}%
^{\prime},\mathsf{z}^{N-1}\right)  X\left(  \mathsf{z}_{1},\mathsf{z}%
^{N-1},\mathsf{z}_{1}^{\prime},\mathsf{z}^{N-1}\right)  \right\}
d\mathsf{z}^{N-1}\overset{a.e}{=}0\forall X\in\ker L_{N}^{1}\nonumber\\
&  \Rightarrow%
%TCIMACRO{\dint }%
%BeginExpansion
{\displaystyle\int}
%EndExpansion
\left\{  \delta\left(  \mathsf{z}_{1}-\mathsf{z}_{1}^{\prime}\right)
\delta\left(  \mathsf{z}^{N-1}-\mathsf{z}^{\prime N-1}\right)  +\Lambda\left(
\mathsf{z}_{1},\mathsf{z}^{N-1},\mathsf{z}_{1}^{\prime},\mathsf{z}%
^{N-1}\right)  X\left(  \mathsf{z}_{1},\mathsf{z}^{N-1},\mathsf{z}_{1}%
^{\prime},\mathsf{z}^{N-1}\right)  \right. \nonumber\\
&  \qquad\qquad\qquad\qquad\qquad\qquad\qquad\qquad\qquad\qquad\left.
d\mathsf{z}^{N-1}\overset{a.e}{=}0\right\}  \forall X\in\ker L_{N}^{1}%
\end{align}

\end{proof}
\end{theorem}

\begin{theorem}
Incremental Strong Orthogonality to an Equivalence Class
\begin{align}
\hat{T}  &  \in\mathcal{\hat{T}}_{F_{\left[  L_{N}^{1}\left(  X\right)
\right]  _{N}}}\left(  \mathfrak{A}_{\ker}\right)  \Leftrightarrow%
%TCIMACRO{\dint }%
%BeginExpansion
{\displaystyle\int}
%EndExpansion
\Lambda\left(  \mathsf{z}_{1},\mathsf{z}^{N-1},\mathsf{z}_{1}^{\prime
},\mathsf{z}^{N-1}\right)  X\left(  \mathsf{z}_{1},\mathsf{z}^{N-1}%
,\mathsf{z}_{1}^{\prime},\mathsf{z}^{N-1}\right)  d\mathsf{z}^{N-1}%
\overset{a.e}{=}0\nonumber\\
&  \quad\qquad\qquad\qquad\qquad\qquad\qquad\qquad\qquad\qquad\qquad
\qquad\qquad\qquad\forall X\in\left[  L_{N}^{1}\left(  X\right)  \right]  _{N}%
\end{align}

\end{theorem}

\begin{proof}
$"\Rightarrow"$

If $\hat{T}\in$ $\mathcal{\hat{T}}_{F_{\left[  L_{N}^{1}\left(  X\right)
\right]  _{N}}}\left(  \mathfrak{A}_{\ker}\right)  $ then
\begin{equation}
\hat{T}\left(  X\right)  \in\left[  L_{N}^{1}\left(  X\right)  \right]
_{N}\,\forall X\in\left[  L_{N}^{1}\left(  X\right)  \right]  _{N}%
\end{equation}
thus
\begin{equation}%
%TCIMACRO{\dint }%
%BeginExpansion
{\displaystyle\int}
%EndExpansion
T\left(  \mathsf{z}^{N}\right)  T\left(  \mathsf{z}^{\prime N}\right)
X\left(  \mathsf{z}_{1},\mathsf{z}^{N-1},\mathsf{z}_{1}^{\prime}%
,\mathsf{z}^{N-1}\right)  d\mathsf{z}^{N-1}=L_{N}^{1}\left(  X\right)  \left(
\mathsf{z}_{1},\mathsf{z}_{1}^{\prime}\right)  \,\forall X\in\left[  L_{N}%
^{1}\left(  X\right)  \right]  _{N}%
\end{equation}


so that
\begin{equation}%
%TCIMACRO{\dint }%
%BeginExpansion
{\displaystyle\int}
%EndExpansion
\Lambda\left(  \mathsf{z}_{1},\mathsf{z}^{N-1},\mathsf{z}_{1}^{\prime
},\mathsf{z}^{N-1}\right)  X\left(  \mathsf{z}_{1},\mathsf{z}^{N-1}%
,\mathsf{z}_{1}^{\prime},\mathsf{z}^{N-1}\right)  d\mathsf{z}^{N-1}%
\overset{a.e}{=}0\quad\forall X\in\left[  L_{N}^{1}\left(  X\right)  \right]
_{N}%
\end{equation}


$"\Leftarrow"$%
\begin{align}
&
%TCIMACRO{\dint }%
%BeginExpansion
{\displaystyle\int}
%EndExpansion
T\left(  \mathsf{z}^{N}\right)  T\left(  \mathsf{z}^{\prime N}\right)
X\left(  \mathsf{z}_{1},\mathsf{z}^{N-1},\mathsf{z}_{1}^{\prime}%
,\mathsf{z}^{N-1}\right)  d\mathsf{z}^{N-1}\nonumber\\
&  \left.  =\right.
%TCIMACRO{\dint }%
%BeginExpansion
{\displaystyle\int}
%EndExpansion
\left\{  \delta\left(  \mathsf{z}^{N}-\mathsf{z}^{\prime N}\right)
+\Lambda\left(  \mathsf{z}^{N},\mathsf{z}^{\prime N}\right)  \right\}
X\left(  \mathsf{z}_{1},\mathsf{z}^{N-1},\mathsf{z}_{1}^{\prime}%
,\mathsf{z}^{N-1}\right)  d\mathsf{z}^{N-1}\nonumber\\
&  \left.  =\right.
%TCIMACRO{\dint }%
%BeginExpansion
{\displaystyle\int}
%EndExpansion
X\left(  \mathsf{z}_{1},\mathsf{z}^{N-1},\mathsf{z}_{1}^{\prime}%
,\mathsf{z}^{N-1}\right)  d\mathsf{z}^{N-1}\left.  =\right.  L_{N}^{1}\left(
X\right)  \left(  \mathsf{z}_{1},\mathsf{z}_{1}^{\prime}\right)  \,\forall
X\in\left[  L_{N}^{1}\left(  X\right)  \right]  _{N}%
\end{align}


hence $\hat{T}\in$ $\mathcal{\hat{T}}_{F_{\left[  L_{N}^{1}\left(  X\right)
\right]  _{N}}}\left(  \mathfrak{A}_{\ker}\right)  $
\end{proof}

\begin{theorem}
Incremental Non Orthogonality to the Base
\begin{align}
\hat{T}  &  \in\mathcal{\hat{T}}_{\mathcal{V}_{\Gamma}}\left(  \mathfrak{A}%
_{\ker}\right)  \subset\mathcal{\hat{T}}\left(  \mathfrak{A}_{\ker}\right)
\text{ for some }\mathcal{V}_{\Gamma}\Leftrightarrow\nonumber\\
&
%TCIMACRO{\dint }%
%BeginExpansion
{\displaystyle\int}
%EndExpansion
\Lambda\left(  \mathsf{z}^{N},\mathsf{z}^{\prime N}\right)  X\left(
\mathsf{z}_{1},\mathsf{z}^{N-1},\mathsf{z}_{1}^{\prime},\mathsf{z}%
^{N-1}\right)  d\mathsf{z}^{N-1}\overset{a.e}{\neq}0\quad\text{for some }%
X\in\mathcal{V}_{\Gamma}%
\end{align}

\end{theorem}

\begin{proof}
$"\Rightarrow"$

Clearly if
\begin{equation}%
%TCIMACRO{\dint }%
%BeginExpansion
{\displaystyle\int}
%EndExpansion
\Lambda\left(  \mathsf{z}^{N},\mathsf{z}^{\prime N}\right)  X\left(
\mathsf{z}_{1},\mathsf{z}^{N-1},\mathsf{z}_{1}^{\prime},\mathsf{z}%
^{N-1}\right)  d\mathsf{z}^{N-1}\overset{a.e}{\neq}0
\end{equation}
then $\hat{T}\left(  X\right)  \notin\left[  L_{N}^{1}\left(  X\right)
\right]  _{N}$ and belongs to $\mathcal{V}_{\Gamma}$ for some $\Gamma$. One
can see this by letting
\begin{equation}
\hat{T}\left(  X\right)  \notin\left[  L_{N}^{1}\left(  X\right)  \right]
_{N}%
\end{equation}
then
\begin{align}
\hat{T}\left(  X\left(  \mathfrak{x}_{1}\right)  \right)   &  =\hat{T}\left(
k\left(  \mathfrak{x}_{1}\right)  +X_{ref}\left(  \mathfrak{x}_{1}\right)
\right)  =\hat{T}\left(  k\left(  \mathfrak{x}_{1}\right)  \right)  +\hat
{T}\left(  X_{ref}\left(  \mathfrak{x}_{1}\right)  \right) \nonumber\\
&  =k^{\prime}\left(  \mathfrak{x}_{1}\right)  +X_{ref}\left(  \mathfrak{x}%
_{1}^{\prime}\right)
\end{align}


$"\Leftarrow"$

If $\hat{T}\in$ $\mathcal{\hat{T}}_{\mathcal{V}_{\Gamma}}\left(
\mathfrak{A}_{\ker}\right)  $ then
\begin{equation}%
%TCIMACRO{\dint }%
%BeginExpansion
{\displaystyle\int}
%EndExpansion
\Lambda\left(  \mathsf{z}^{N},\mathsf{z}^{\prime N}\right)  X\left(
\mathsf{z}_{1},\mathsf{z}^{N-1},\mathsf{z}_{1}^{\prime},\mathsf{z}%
^{N-1}\right)  d\mathsf{z}^{N-1}\overset{a.e}{\neq}0\text{for some }%
X\in\mathcal{V}_{\Gamma}%
\end{equation}

\end{proof}

\section{Transition Amplitudes and Expectation Values wrt IPSs\label{Transamp}%
}

We will need to calculate expectation values and transition amplitudes between
IPS\ composed of non orthogonal 1-particle states, we will follow the
presentation found in "On the Evaluation of Nonorthogonal Matrix
Elements\textit{" }by Jacob Verbeek and J.H. van Lenthe THEOCHEM 1990.
Consider two IPSs $\left\vert \Psi\right\rangle $ and $\left\vert
\Phi\right\rangle $ composed of nonorthogonal general spin orbitals $\left\{
\psi_{j}|1\leq j\leq N\right\}  $ and $\left\{  \varphi_{j}|1\leq j\leq
N\right\}  $ and the transition amplitude
\begin{equation}
\left\langle \Psi|H\Phi\right\rangle
\end{equation}
where the Hamiltonian is a sum of one and two particle operators
\begin{equation}
H=%
%TCIMACRO{\dsum \limits_{1\leq j\leq N}}%
%BeginExpansion
{\displaystyle\sum\limits_{1\leq j\leq N}}
%EndExpansion
h\left(  j\right)  +\frac{1}{2}%
%TCIMACRO{\dsum \limits_{1\leq j\neq k\leq N}}%
%BeginExpansion
{\displaystyle\sum\limits_{1\leq j\neq k\leq N}}
%EndExpansion
g\left(  j,k\right)
\end{equation}


The transition amplitude of the $1$-particle part is given by
\begin{equation}
\left\langle \Psi\left\vert
%TCIMACRO{\dsum \limits_{1\leq j\leq N}}%
%BeginExpansion
{\displaystyle\sum\limits_{1\leq j\leq N}}
%EndExpansion
h\left(  j\right)  \right.  \Phi\right\rangle =Tr\left\{  \mathbf{h}%
\operatorname{adj}\left(  \mathbf{S}\right)  \right\}
\end{equation}
where the matrix $\mathbf{h}$ has matrix elements $\left\{  h_{jk}|1\leq
j,k\leq N\right\}  $ given by
\begin{equation}
h_{jk}=\left\langle \psi_{j}|h\left(  1\right)  \varphi_{k}\right\rangle =%
%TCIMACRO{\dint }%
%BeginExpansion
{\displaystyle\int}
%EndExpansion
\psi_{j}\left(  \mathsf{z}\right)  ^{\ast}h\left(  \mathsf{z,z}^{\prime
}\right)  \varphi_{k}\left(  \mathsf{z}^{\prime}\right)  d\mathsf{z}%
d\mathsf{z}^{\prime}%
\end{equation}
and the classical $N\times N$ adjugate matrix $\operatorname{adj}\left(
\mathbf{S}\right)  $ of the overlap matrix $\mathbf{S\equiv}\left\{
\left\langle \psi_{j}|\varphi_{k}\right\rangle |1\leq j,k\leq N\right\}  $ is
defined by
\begin{equation}
\left(  \operatorname{adj}\left(  \mathbf{S}\right)  \right)  _{jk}=\left(
-1\right)  ^{j+k}\det\left(  \mathbf{S}\left[  k|j\right]  \right)
\end{equation}
where $\mathbf{S}\left[  k|j\right]  $ is the $\left(  N-1\right)  \times$
$\left(  N-1\right)  $ matrix obtained from $\mathbf{S}$ by deleting the row
$k$ and the column $j$.

The transition amplitude of the $2$-particle part is given by
\begin{equation}
\left\langle \Psi\left\vert \frac{1}{2}%
%TCIMACRO{\dsum \limits_{1\leq j\neq k\leq N}}%
%BeginExpansion
{\displaystyle\sum\limits_{1\leq j\neq k\leq N}}
%EndExpansion
g\left(  j,k\right)  \right.  \Phi\right\rangle =Tr\left\{  \mathbf{G}%
\operatorname{adj}^{\left(  2\right)  }\left(  \mathbf{S}\right)  \right\}
\end{equation}
where the matrix $\mathbf{G}$ has matrix elements $\left\{  g_{jklm}|1\leq
j<k\leq N;1\leq l<m\leq N\right\}  $ given by
\begin{align}
g_{jklm}  &  =\left\langle \psi_{j}\psi_{k}|g\left(  1,2\right)  \varphi
_{l}\varphi_{m}\right\rangle \nonumber\\
&  =%
%TCIMACRO{\dint }%
%BeginExpansion
{\displaystyle\int}
%EndExpansion
\left\{  \psi_{j}\left(  \mathsf{z}_{1}\right)  ^{\ast}\psi_{k}\left(
\mathsf{z}_{2}\right)  ^{\ast}g\left(  \mathsf{z}_{1}\mathsf{,\mathsf{z}%
_{2},\mathsf{z}_{1}^{\prime},z}_{2}^{\prime}\right)  \varphi_{l}\left(
\mathsf{z}_{1}^{\prime}\right)  \varphi_{m}\left(  \mathsf{z}_{2}^{\prime
}\right)  \right. \nonumber\\
&  \left.  -\psi_{j}\left(  \mathsf{z}_{1}\right)  ^{\ast}\psi_{k}\left(
\mathsf{z}_{2}\right)  ^{\ast}g\left(  \mathsf{z}_{1}\mathsf{,\mathsf{z}%
_{2},\mathsf{z}_{1}^{\prime},z}_{2}^{\prime}\right)  \varphi_{m}\left(
\mathsf{z}_{1}^{\prime}\right)  \varphi_{l}\left(  \mathsf{z}_{2}^{\prime
}\right)  \right\}  d\mathsf{z}^{2}d\mathsf{z}^{\prime2}%
\end{align}
and the second order $\binom{N}{2}\times$ $\binom{N}{2}$ adjugate matrix
$\operatorname{adj}^{\left(  2\right)  }\left(  \mathbf{S}\right)  $ is
defined by
\begin{equation}
\left(  \operatorname{adj}^{\left(  2\right)  }\left(  \mathbf{S}\right)
\right)  _{jklm}=\left(  -1\right)  ^{j+k}\det\left(  \mathbf{S}\left[
lm|jk\right]  \right)
\end{equation}
where $\mathbf{S}\left[  lm|jk\right]  $ is the $\left(  N-2\right)
\times\left(  N-2\right)  $ matrix obtained from $\mathbf{S}$ by deleting the
rows $l,m$ and the column $j,k$.

In general the transition amplitude of the $p$-particle part is given by%
\begin{equation}
\left\langle \Psi\left\vert \frac{1}{p!}%
%TCIMACRO{\dsum \limits_{1\leq j_{1}\neq\cdots\neq j_{p}\leq N}}%
%BeginExpansion
{\displaystyle\sum\limits_{1\leq j_{1}\neq\cdots\neq j_{p}\leq N}}
%EndExpansion
O_{p}\left(  j_{1},\cdots,j_{p}\right)  \right.  \Phi\right\rangle =Tr\left\{
\mathbf{O}_{p}\operatorname{adj}^{\left(  p\right)  }\left(  \mathbf{S}%
\right)  \right\}
\end{equation}
where the matrix $\mathbf{O}_{p}$ has matrix elements
\begin{equation}
\left\{  \left[  \mathbf{O}_{p}\right]  _{j_{1},\cdots,j_{p};k_{1}%
,\cdots,k_{p}}|1\leq j_{1}<\cdots<j_{p}\leq N;1\leq k_{1}<\cdots<k_{p}\leq
N\right\}
\end{equation}
given by
\begin{align}
&  \left[  \mathbf{O}_{p}\right]  _{j_{1},\cdots,j_{p};k_{1},\cdots,k_{p}%
}=\left\langle \psi_{j_{1}}\cdots\psi_{j_{p}}|g\left(  j,k\right)
\varphi_{k_{1}}\cdots\varphi_{k_{p}}\right\rangle \nonumber\\
&  =%
%TCIMACRO{\dint }%
%BeginExpansion
{\displaystyle\int}
%EndExpansion
\left\{  \psi_{j_{1}}\left(  \mathsf{z}_{1}\right)  ^{\ast}\cdots\psi_{j_{p}%
}\left(  \mathsf{z}_{p}\right)  ^{\ast}O_{p}\left(  \mathsf{z}_{1\cdots
}\mathsf{z}_{p}\mathsf{,z}_{1\cdots}^{\prime}\mathsf{z}_{p}^{\prime}\right)
%TCIMACRO{\dsum \limits_{\sigma\in S_{p}}}%
%BeginExpansion
{\displaystyle\sum\limits_{\sigma\in S_{p}}}
%EndExpansion
\left(  -1\right)  ^{\pi\left(  \sigma\right)  }\varphi_{k_{\sigma_{\left(
1\right)  }}}\left(  \mathsf{z}_{1}^{\prime}\right)  \cdots\varphi
_{k_{\sigma\left(  p\right)  }}\left(  \mathsf{z}_{p}^{\prime}\right)
\right\}  d\mathsf{z}^{p}d\mathsf{z}^{\prime p}%
\end{align}
and the $p^{th}$ order $\binom{N}{p}\times$ $\binom{N}{p}$ adjugate matrix
$\operatorname{adj}^{\left(  p\right)  }\left(  \mathbf{S}\right)  $ is
defined by
\begin{equation}
\left(  \operatorname{adj}^{\left(  p\right)  }\left(  \mathbf{S}\right)
\right)  _{j_{1},\cdots,j_{p};k_{1},\cdots,k_{p}}=\left(  -1\right)  ^{%
%TCIMACRO{\tsum \limits_{1\leq j,k\leq p}}%
%BeginExpansion
{\textstyle\sum\limits_{1\leq j,k\leq p}}
%EndExpansion
\left\{  j_{l}+k_{l}\right\}  }\det\left(  \mathbf{S}\left[  j_{1}%
,\cdots,j_{p}|k_{1},\cdots,k_{p}\right]  \right)
\end{equation}
where $\mathbf{S}\left[  j_{1},\cdots,j_{p}|k_{1},\cdots,k_{p}\right]  $ is
the $\left(  N-p\right)  \times\left(  N-p\right)  $ matrix obtained from
$\mathbf{S}$ by deleting the rows $j_{1},\cdots,j_{p}$ and the column
$k_{1},\cdots,k_{p}$.

\bibliographystyle{apsrev}
\bibliography{acompat,DMF}

\end{document}